\pdfoutput=1                      
% to make the paper available to visually impaired people

\documentclass[10 pt,a4paper]{article}


\usepackage[accsupp]{axessibility}    
% to make the paper available to visually impaired people




\usepackage{mathrsfs,amsfonts,amsmath,comment,verbatim,amsthm,amssymb,graphicx,fancyhdr,makeidx,amscd} 
\usepackage{color}


%,natbib}
%\makeindex
%
% inserire \index{espo@$\exp$} subito prima di un simbolo matematico (esempio, $\exp$) per identificare la pagina dove appare.
%\printindex dove voglio che compaia l'indice
%
%
% \usepackage[all]{xy}

%\usepackage[notref,notcite]{showkeys}
%\setcounter{section}{-1}
%\numberwithin{equation}{section}


\newtheorem{definition}{Definition}
\newtheorem{proposition}[definition]{Proposition}
\newtheorem{lemma}[definition]{Lemma}
\newtheorem{theorem}[definition]{Theorem}
\newtheorem{corollary}[definition]{Corollary}
\newtheorem{example}[definition]{Example}
\newtheorem{conjecture}[definition]{Conjecture}
\newtheorem{question}{Question}


%\newtheorem{definition}{Definition}[section]
%\newtheorem{proposition}{Proposition}[section]
%\newtheorem{lemma}{Lemma}[section]
%\newtheorem{theorem}{Theorem}[section]
%\newtheorem{corollary}{Corollary}[section]
%\newtheorem{example}{Example}[section]
%\newtheorem{conjecture}{Conjecture}[section]
%\newtheorem{question}{Question}
%\newtheorem{remark}{Remark}[section]



\theoremstyle{definition}
\newtheorem{remark}[definition]{Remark}


\newtheorem*{notation}{Notational convenience}
\newtheorem*{teononum}{Theorem}

\setcounter{secnumdepth}{1}
%\pagestyle{headings}

\addtolength{\textwidth}{2cm}


\usepackage{Commands_Giada}




\begin{document}
%\pagestyle{fancy}
%% i comandi seguenti impediscono la scrittura in maiuscolo
%% dei nomi dei capitoli e dei paragrafi nelle intestazioni
%\renewcommand{\chaptermark}[1]{\markboth{#1}{}}
%\renewcommand{\sectionmark}[1]{\markright{\thesection\ #1}}
%\fancyhf{} % rimuove l'attuale contenuto dell'intestazione
%% e del pi\`e di pagina
%\fancyhead[LE,RO]{\bfseries\thepage}
%\fancyhead[LO]{\bfseries\rightmark}
%\fancyhead[RE]{\bfseries\leftmark}
%\renewcommand{\headrulewidth}{0.5pt}
%\renewcommand{\footrulewidth}{0pt}
%\addtolength{\headheight}{0.5pt} % riserva spazio per la linea
%\fancypagestyle{plain}{%
%\fancyhead{} % ignora, nello stile plain, le intestazioni
%\renewcommand{\headrulewidth}{0pt} % e la linea
%}


\newcommand{\riem}{(M^n, \langle \, , \, \rangle)}
\newcommand{\Hess}{\mathrm{Hess}\, }
\newcommand{\hess}{\mathrm{hess}\, }
\newcommand{\cut}{\mathrm{cut}}
\newcommand{\ind}{\mathrm{ind}}

\newcommand{\ess}{\mathrm{ess}}

\newcommand{\longra}{\longrightarrow}

\newcommand{\eps}{\varepsilon}

\newcommand{\ra}{\rightarrow}

\newcommand{\vol}{\mathrm{vol}}

\newcommand{\di}{\mathrm{d}}

\newcommand{\R}{\mathbb R}

\newcommand{\C}{\mathbb C}

\newcommand{\Z}{\mathbb Z}

\newcommand{\N}{\mathbb N}

\newcommand{\HH}{\mathbb H}

\newcommand{\esse}{\mathbb S}

\newcommand{\bull}{\rule{2.5mm}{2.5mm}\vskip 0.5 truecm}

\newcommand{\binomio}[2]{\genfrac{}{}{0pt}{}{#1}{#2}} %identico ad \atop

\newcommand{\metric}{\langle \, , \, \rangle}

\newcommand{\lip}{\mathrm{Lip}}

\newcommand{\loc}{\mathrm{loc}}

\newcommand{\diver}{\mathrm{div}}

\newcommand{\disp}{\displaystyle}

\newcommand{\rad}{\mathrm{rad}}

\newcommand{\inj}{\mathrm{inj}}

\newcommand{\Ric}{\mathrm{Ric}}
\newcommand{\Sec}{\mathrm{Sec}}

\newcommand{\LL}{\mathscr{L}}

\newcommand{\MAX}{\mathrm{MAX}} 

\newcommand{\mmetric}{\langle\langle \, , \, \rangle\rangle}

\newcommand{\hol}{\mathrm{H\ddot{o}l}}

\newcommand{\capac}{\mathrm{cap}}

\newcommand{\bmo}{\{b <0\}}

\newcommand{\bmuo}{\{b \le 0\}}

\newcommand{\dist}{\mathrm{dist}}

\newcommand{\MM}{\mathscr{M}}
\newcommand{\GG}{\mathscr{G}}


\newcommand{\vp}{\varphi}

\renewcommand{\div}[1]{{\mathop{\mathrm div}}\left(#1\right)}

\newcommand{\divphi}[1]{{\mathop{\mathrm div}}\bigl(\vert \nabla #1

\vert^{-1} \varphi(\vert \nabla #1 \vert)\nabla #1   \bigr)}

\newcommand{\nablaphi}[1]{\vert \nabla #1\vert^{-1}

\varphi(\vert \nabla #1 \vert)\nabla #1}

\newcommand{\modnabla}[1]{\vert \nabla #1\vert }

\newcommand{\modnablaphi}[1]{\varphi\bigl(\vert \nabla #1 \vert\bigr)

\vert \nabla #1\vert }

\newcommand{\ds}{\displaystyle}

\newcommand{\cL}{\mathcal{L}}

%

%









%\newcommand{\esse}{\mathds{S}}

\newcommand{\essem}{\mathds{S}^n}

\newcommand{\erre}{\mathds{R}}

\newcommand{\errem}{\mathds{R}^n}

\newcommand{\enne}{\mathds{N}}

\newcommand{\acca}{\mathds{H}}



\newcommand{\cvett}{\Gamma(TM)}

\newcommand{\cinf}{C^{\infty}(M)}

\newcommand{\sptg}[1]{T_{#1}M}

\newcommand{\partder}[1]{\frac{\partial}{\partial {#1}}}

\newcommand{\partderf}[2]{\frac{\partial {#1}}{\partial {#2}}}

\newcommand{\ctloc}{(\mathcal{U}, \varphi)}

\newcommand{\fcoord}{x^1, \ldots, x^n}

\newcommand{\ddk}[2]{\delta_{#2}^{#1}}

\newcommand{\christ}{\Gamma_{ij}^k}

%\newcommand{\Ric}{\operatorname{Ricc}}

\newcommand{\supp}{\operatorname{supp}}

\newcommand{\sgn}{\operatorname{sgn}}

\newcommand{\rg}{\operatorname{rg}}

\newcommand{\inv}[1]{{#1}^{-1}}

\newcommand{\id}{\operatorname{id}}

\newcommand{\jacobi}[3]{\sq{\sq{#1,#2},#3}+\sq{\sq{#2,#3},#1}+\sq{\sq{#3,#1},#2}=0}

\newcommand{\lie}{\mathfrak{g}}

\newcommand{\wedgedot}{\wedge\cdots\wedge}

\newcommand{\rp}{\erre\mathds{P}}

\newcommand{\II}{\operatorname{II}}

\newcommand{\gradh}[1]{\nabla_{H^n}{#1}}

\newcommand{\absh}[1]{{\left|#1\right|_{H^n}}}

\newcommand{\mob}{\mathrm{M\ddot{o}b}}

\newcommand{\mab}{\mathfrak{m\ddot{o}b}}

%\newcommand{\cut}{\mathrm{cut}}

\newcommand{\foc}{\mathrm{foc}}

%\newcommand{\C}{\mathcal{C}}

\newcommand{\F}{\mathcal{F}}

\newcommand{\Cf}{\mathcal{C}_f}

\newcommand{\cutf}{\mathrm{cut}_{f}}

\newcommand{\Cn}{\mathcal{C}_n}

\newcommand{\cutn}{\mathrm{cut}_{n}}

\newcommand{\Ca}{\mathcal{C}_a}

\newcommand{\cuta}{\mathrm{cut}_{a}}

\newcommand{\cutc}{\mathrm{cut}_c}

\newcommand{\cutcf}{\mathrm{cut}_{cf}}

%\newcommand{\ind}{\mathrm{ind}}

\newcommand{\rk}{\mathrm{rk}}

%\newcommand{\vol}{\mathrm{vol}}

\newcommand{\crit}{\mathrm{crit}}

\newcommand{\diam}{\mathrm{diam}}

\newcommand{\haus}{\mathscr{H}}

\newcommand{\po}{\mathrm{po}}

\newcommand{\gr}{\mathcal{G}}

\newcommand{\sn}{\mathrm{sn}_H}

\newcommand{\cn}{\mathrm{cn}_H}

\newcommand{\Tr}{\mathrm{Tr}}

\newcommand{\bh}{\mathbb{B}}
\newcommand{\Gr}{\mathrm{Gr}}
\newcommand{\Pa}{\mathcal{P}}

\newcommand{\aaa}{a(|\nabla u|)}
\newcommand{\aap}{a'(|\nabla u|)}
\newcommand{\inte}{\mathrm{Int}}

\newcommand{\tcr}{\textcolor{red}}
\newcommand{\tcb}{\textcolor{blue}}
\newcommand{\tco}{\textcolor{red}}
\newcommand{\critu}{{\rm crit}(u)}













\author{Giulio Colombo${}^{(1)}$ 
	\and 
	Alberto Farina${}^{(2)}$
\and Marco Magliaro${}^{(3)}$
\and Luciano Mari${}^{(1)}$
\and Marco Rigoli${}^{(1)}$ 
}
%\title{\textbf{On the classification of capillary graphs in $\R^2$, $\R^3$ and more general Riemannian manifolds}}
\title{\textbf{On the classification of capillary graphs in Euclidean and non-Euclidean spaces}}
\date{}
\maketitle
\scriptsize \begin{center} (1) -- Dipartimento di Matematica ``Federigo Enriques"\\
Universit\`a degli Studi di Milano\\ 
Via C. Saldini 50, 20133 Milano (Italy).
\end{center}    
\scriptsize \begin{center} (2) -- Laboratoire 
Ami\'enois de Math\'ematique Fondamentale et Appliqu\'ee\\
UMR CNRS 7352, Universit\'e Picardie ``Jules Verne'' \\33 Rue St Leu, 80039 Amiens (France).\\ 
\end{center}
\scriptsize \begin{center} (3) -- Dipartimento di Scienza e Alta Tecnologia\\
Universit\`a degli Studi dell'Insubria\\
Via Valleggio 11, 22100 Como (Italy)
\end{center}    

\bigskip

\noindent \textbf{E-mail addresses:}  giulio.colombo@unimi.it, alberto.farina@u-picardie.fr, marco.magliaro@uninsubria.it,\\
luciano.mari@unimi.it, marco.rigoli55@gmail.com

\vspace{0.2cm}

\noindent\textbf{Keywords:} overdetermined problem, capillarity, CMC graph, splitting, minimal hypersurface.

\vspace{0.2cm}

\noindent \textbf{Mathematics subject classification:} Primary 53C21, 53C42; Secondary 53C24, 58J32, %31C12, 31B05.\\
35B06, 35B35


\normalsize
%
%
%

\begin{abstract}
	We prove some rigidity and classification results for graphs with prescribed mean curvature and locally constant Dirichlet and Neumann data, for instance as they appear in capillarity problems. We consider domains in Riemannian manifolds, with emphasis on $\R^2$ and $\R^3$. We classify both the underlying domain and the resulting solution, providing general splitting theorems in this setting. 	
\end{abstract}	

\tableofcontents

\section{Introduction}

If a cylindrical tube in vertical position $\Omega \times (a,\infty) \subseteq \R^{n+1} = \R^n \times \R$, with $a<<-1$, is dipped into a large reservoir filled with liquid modelled by the half-space $\R^n \times (-\infty,0]$, the combination of gravity and surface tension demands that the height function $u : \overline\Omega \to \R$ reached by the liquid in the tube satisfies the capillarity equation
\[
	\diver \left( \dfrac{\nabla u}{\sqrt{1+|\nabla u|^2}} \right) - \kappa u = 0 \qquad \text{in } \, \Omega,
\]
where $\kappa > 0$ is a physical constant, $\nabla, \diver$ are the gradient and divergence in $\R^n$ and 
\[
	\MM[u] : = \diver \left( \dfrac{\nabla u}{\sqrt{1+|\nabla u|^2}} \right)
\]
is the mean curvature of the graph of the liquid profile
\[
	\GG_u \ : \ \overline\Omega \to \R^{n+1}, \qquad \GG_u(x) \doteq \big(x, u(x)\big)
\]
with respect to the upward pointing unit normal. Furthermore, $\GG_u$ must intersect the boundary of the tube with constant angle. For a detailed account on the capillarity problem, we refer the interested reader to \cite{finn}.

A classical question, motivating the present paper, asks for which domains $\Omega$ the surface of the liquid also attains a (locally) constant height on $\partial \Omega$. In this case, the angle condition rewrites as constant Neumann data for $u$. Henceforth we will consider mean curvature prescriptions more general than $\kappa u$, so given $f \in C^1(\R)$ we call a domain $\Omega$ a \emph{capillary domain} if it supports a non-constant solution $u$ to 
\begin{equation}\label{eq_PMC}
	\left\{ \begin{array}{ll}
		\MM[u]  + f(u) = 0 & \text{in } \, \Omega, \\[0.2cm]
		u, \partial_\eta u \, \text{ locally constant} & \text{on $\partial \Omega$,} 
	\end{array}\right.
\end{equation}
where $\eta$ is the interior unit normal to $\partial \Omega$. The associated graph  $\GG_u(\Omega)$ will then be called a \emph{capillary graph}. If $f(u)$ is constant  the equation describes CMC graphs, in particular minimal ones if $f(u) \equiv 0$, while the capillarity equation corresponds to the choice $f(u) = -\kappa u$. The goal of this paper is to investigate the shape of capillary graphs over $\Omega \subseteq \R^n$ or, more generally, over domains in a complete Riemannian manifold.

Hereafter, 
\begin{itemize}
\item $\{\partial_j\Omega\}$, $j \le \infty$ denotes  the collection of connected components of $\partial \Omega$;
\item $\partial_\star \Omega = \big\{ x \in \partial \Omega \ : \ \partial_\eta u(x) \neq 0\big\}$ 
\end{itemize}

In his pioneering paper on the moving plane method, Serrin \cite[Theorem 2]{serrin} treated the case where $\Omega \subseteq \R^n$ is a bounded $C^2$ domain and $u \in C^2(\overline\Omega)$ solves
\begin{equation}\label{eq_PMC_serrin}
	\left\{ \begin{array}{ll}
		\MM[u]  + f(u) = 0 & \text{in } \, \Omega, \\[0.2cm]
		u > 0 & \text{in $\Omega$.} \\[0.2cm]
		u=0, \ \partial_\eta u = {\rm const} & \text{on $\partial \Omega$}
	\end{array}\right.
\end{equation}
for some $f \in C^1(\R)$, proving that $\Omega$ must be a ball and $u$ is radially symmetric about its center. The method was then refined by Reichel \cite{reichel_1} and Sirakov \cite{sirakov} to be applied to exterior domains. The most general result to our knowledge, \cite[Theorem 2]{sirakov}, states that a $C^2$ domain $\Omega$ with bounded complement supporting a solution $u \in C^2 (\overline{\Omega})$ to
\begin{equation}\label{eq_PMC_reichel}
	\left\{ \begin{array}{ll}
		\MM[u] + f(u) = 0 & \text{in } \, \Omega, \\[0.2cm]
		u \ge 0 & \text{in } \, \Omega, \\[0.2cm]
		u(x) \to 0 & \text{as $x \to \infty$} \\[0.2cm]
		u=b_j > 0 & \text{on $\partial_j \Omega$,} \\[0.2cm]
		\partial_\eta u = c_j \le 0 & \text{on $\partial_j \Omega$,}
	\end{array}\right.
\end{equation}
for some constants $b_j,c_j$ and for $f \in \lip_\loc(\R)$ non-increasing in a neighbourhood of zero, must be the exterior of a ball, and $u$ must be radially symmetric. The paper \cite{sirakov} also contains a similar characterization for $\Omega$ a $C^2$ bounded domain with holes.
 
If $\partial \Omega$ is unbounded, or if the word ``constant'' in \eqref{eq_PMC_serrin} is replaced by ``locally constant'', the problem is widely open even for CMC graphs. In fact, even though the results in \cite{serrin,reichel_1,sirakov} hold for a large class of quasilinear operators, some of which have been extensively studied in the literature, relatively few results have appeared so far for the mean curvature operator. 
In particular, an analogue of the famous Berestycki-Caffarelli-Nirenberg conjecture \cite{bcn} for capillary graphs is still completely open. More generally, one may ask:
 
\begin{question}\label{Q1}
	Can we classify all smooth enough domains $\Omega \subseteq \R^n$ supporting a solution to \eqref{eq_PMC}? 
	%Are they just balls, cylinders, half-spaces and %their complements?
\end{question}
 
Among the non-constant solutions to \eqref{eq_PMC} with  $\partial \Omega$ unbounded we single out those that are one-dimensional, namely, those for which 
\[
	\Omega = (0,T) \times \R^{n-1} \qquad \text{for some } \, T \le \infty 
\]
and $u$ is 1D. More precisely, there exists an affine hyperplane $\pi$ with unit normal $w$ such that $\Omega = \{ y + tw : y \in \pi, \ t \in (0,T)\}$, and in the coordinates $(t,y)$ the function $u$ only depends on $t$. For instance, when  $f(u)$ is constant, analyzing the resulting ODE for $u$ leads to the following classification of 1D solutions:



\begin{itemize}
	\item If $f(u) \equiv 0$, then either $T = \infty$ and $\GG_u$ is a half-hyperplane, or $T< \infty$ and $\GG_u$ is a piece of a half-hyperplane.
	\item If $f(u) \equiv H \neq 0$, then $T< \infty$ and $\GG_u$ is a piece of a cylinder with circular cross-section, i.e. the graph of $u(t)$ is a piece of $\mathbb{S}^1$ of suitable radius. Hereafter, we refer to this case as ``a piece of cylinder'', leaving implicit that the cross-section is circular. 	 
\end{itemize}

Question \ref{Q1} was first raised in a paper by three of the authors \cite{cmr}, where some classification results were obtained for domains in manifolds which are globally Lipschitz epigraphs or globally Lipschitz slabs, see below. Among them, we highlight \cite[Theorem 1.3]{cmr}, the first result on Question 1 when $\partial \Omega$ is unbounded: if $\Omega \subseteq \R^2$ is a globally Lipschitz epigraph supporting a solution to 
\[
	\left\{ \begin{array}{ll}
		\MM[u] + H = 0 & \text{in } \, \Omega, \\[0.2cm]
		\inf_\Omega u > -\infty, & \text{in } \, \Omega, \\[0.2cm]
		u=0, \, \partial_\eta u = c \neq 0 & \text{on $\partial \Omega$} \\[0.2cm]
	\end{array}\right.
\]
for some $H,c \in \R$, then $\Omega$ is a half-plane, $H=0$ and $\GG_u$ is a half-plane. Below, we shall significantly improve on this result. Very recently, Lian and Sicbaldi treated the case of any $C^1$ domain $\Omega \subseteq \R^2$ with $\partial \Omega$ unbounded and connected. By \cite[Theorem 1]{lian_sicbaldi}, if $u$ solves   
\begin{equation}\label{eq_PMC_sic}
	\left\{ \begin{array}{ll}
		\MM[u] + f(u) = 0 & \text{in } \, \Omega, \\[0.2cm]
		0 < u \le b & \text{in } \, \Omega, \\[0.2cm]
		u= 0, \, \partial_\eta u = c \neq 0 & \text{on $\partial \Omega$} \\[0.2cm]
		|\nabla u| \in L^\infty(\Omega) 		
	\end{array}\right.
\end{equation}
for some $b,c \in \R$, and if $f \in C^1(\R)$ admits a primitive $F$ satisfying 
\begin{equation}\label{eq_condiprimi}
	F \le 0 \quad \text{on } \, \R, \qquad F(0) \ge \frac{1}{\sqrt{1+c^2}} - 1, 
\end{equation}
then $\Omega$ is a half-plane and $u$ is 1D. Moreover, if $f' \le 0$ on $[0,b]$ the gradient request in \eqref{eq_PMC_sic} can be dropped. They also observed that condition \eqref{eq_condiprimi} is necessary for the existence of a 1D-solution in a half-plane.

We call a 1D solution \emph{monotone} if $u(t)$ is monotone in $t$. In particular, for a monotone solution $\partial_\eta u$ has different signs in the two components of $\partial \Omega$, allowing for $\partial_\eta u \equiv 0$ in one of them.
 
It is known that monotone solutions are a particular class of \emph{stable} solutions to the equation $\MM[u] + f(u) = 0$ (see \cite {FSV}), i.e. those for which the linearized operator is non-negative in the spectral sense on $C^\infty_c(\Omega)$. Variationally, for each $\Omega' \Subset \Omega$ they correspond to stationary points of the energy
\[
	\int_{\Omega'} \sqrt{1+|\nabla \psi|^2} - \int_{\Omega'} F(\psi), \qquad F(t) \doteq \int_0^t f(s) \di s 
\]
whose second variation is non-negative. 
 
\begin{remark}
	Since the linearized mean curvature operator is elliptic, if $f'(u) \le 0$ then every solution to $\MM[u] + f(u) = 0$ is automatically stable. This applies both to the capillary case $f(u) = -\kappa u$, $\kappa >0$ and to the CMC case. 
\end{remark}

%Differently to the case where the driving operator is the Laplacian, solutions to 
%\begin{equation}\label{eq_rigid}
%\MM[u] + f(u) = 0 \qquad \text{with} \qquad f'(u) \le 0 
%\end{equation}
%exhibit much stronger rigidity, see Finn \cite{} and Tkachev \cite{tkachev, farina} and also \cite{bmpr,acr}. For instance, by \cite{acr}, a weak solution $u \in \lip_\loc(\overline{\Omega})$ to \eqref{eq_rigid} must satisfy
%\[
%\min\left\{ 0,\inf_{\partial \Omega} f(u) \right\} \le f(u) \le \max\left\{ 0,\sup_{\partial \Omega} f(u) \right\}
%\]
%%In particular, if $u$ is constant on $\partial \Omega$ (say, $u=0$), by applying the result to $u$ and to $- u$ we deduce that $\GG_u$ has constant mean curvature $H = f(b)$. 
%%Consequently, for instance, any solution to 
%%\[
%%\left\{ \begin{array}{ll} 
%%	\MM[u] - \kappa u = 0 & \quad \text{in } \, \Omega \subseteq \R^n, \ \ \ \kappa \in \R^+ \\[0.2cm]
%%	u = b & \quad \text{on } \, \partial \Omega
%%\end{array}\right.
%%\]
%%must be constant even without any overdetermined, sign or growth condition on $u$. This was observed by Finn in $\R^2$ (\tcr{citations})
 
In this paper, we will obtain some  characterization results for stable solutions to \eqref{eq_PMC} in the direction of Question 1, new even in the case of constant $f$. We first introduce the class of domains we will be interested in. 

\begin{definition}[\textbf{Mildly $\mathbf{v}$-transverse boundary}]\label{def_mildlytransv}
	Let $\Omega \subseteq \R^n$ be a $C^1$ domain, and let $\partial^\dagger\Omega \subseteq \partial \Omega$ be the union of some connected components of $\partial \Omega$ (possibly, a single component or the entire $\partial \Omega$). Then, $\partial^\dagger \Omega$ is said to be \emph{mildly $v$-transverse} for some $v \in \mathbb{S}^{n-1}$ if $\langle \eta,v \rangle$ does not change sign on any component of $\partial^\dagger\Omega$.
%	and $\langle \eta,v \rangle \not \equiv 0$ on at least one component of $\partial^\dagger\Omega$.
%	; the whole $\partial \Omega$ is called \emph{mildly $v$-transverse} if $\langle \eta,v\rangle$ does not change sign on any connected component of $\partial \Omega$, and $\langle \eta,v \rangle \not \equiv 0$ on $\partial \Omega$.
	%does not change sign on any component of $\partial \Omega$.  	
\end{definition}

Note that, if $\partial^\dagger \Omega$ is mildly $v$-transverse, the product $\langle \eta,v \rangle$ is allowed to take different signs on different components of $\partial^\dagger \Omega$. Also, a mildly $v$-transverse portion of $\partial\Omega$ may have any number of connected components, even countably many. Letting $e_j = \partial_{x^j}$, and writing points $x \in \R^n$ as $x = (x^1,x')$ with $x' \in \R^{n-1}$, the following two examples have mildly $e_1$-transverse boundaries:
 
\begin{itemize}
	\item An \emph{epigraph}: a set of the type $\{x^1 > \phi(x')\}$ for some function $\phi \in C^1(\R^n)$.
	\item A \emph{slab}: a set of the type $\{\phi_1(x') < x^1 < \phi_2(x')\}$ for some functions $\phi_j \in C^1(\R^{n-1})$ with $\phi_1(x') < \phi_2(x')$ for each $x' \in \R^{n-1}$.
\end{itemize}
In what follows, an epigraph (respectively, slab) is said to be bounded, or globally Lipschitz, if so is its defining function $\phi$ (resp. $\phi_1,\phi_2$). 

\vspace{0.3cm} 

Our first main result regards capillary CMC graphs in $\R^2$. Recently, Cui in \cite{cui} made a significant advance in the minimal case by proving the following theorem: if $\Omega$ is a $C^2$ domain with connected boundary, then the graph of a solution to 
\begin{equation}\label{eq_PMC_cui}
	\left\{ \begin{array}{ll}
		\MM[u] = 0 & \text{in } \, \Omega \subseteq \R^2, \\[0.2cm]
		u > 0 & \text{in $\Omega$} \\[0.2cm]
		u=0, \ \partial_\eta u = {\rm const} & \text{on $\partial \Omega$}
	\end{array}\right.
\end{equation}
must either be a half-plane or part of a vertical catenoid (thus, $\Omega$ is either a half-space or the complement of a ball). The argument relies on a reflection technique developed by Choe \cite{choe} and deep classification results for complete, boundaryless  minimal surfaces. As such, these techniques are specific to the minimal setting and to three-dimensional ambient spaces.



%\begin{remark}
%	The case $\partial_\eta u \equiv +\infty$ in \eqref{eq_PMC_cui}, not considered in \cite{cui}, corresponds to free boundary minimal graphs, that is, to minimal graphs meeting the horizontal hyperplane $\{x^3=0 \}$ orthogonally. Traizet \cite{traizet} classified such solutions for $\Omega \subseteq \R^2$.
%\end{remark}
%
%\tcr{Il Remark 3 mi pare poco rilevante. Suggerisco di toglierlo.}


%We henceforth denote by $\{\partial_j\Omega\}$, $j \le \infty$ the connected components of $\partial \Omega$.

We obtain: 

\begin{theorem}\label{teo_CMC_intro_R2}
	Let $\Omega \subseteq \R^2$ be a $C^2$ domain supporting a non-constant solution $u \in C^2(\overline\Omega)$ to 
	\[
		\left\{ 
		\begin{array}{ll}
			\MM[u]  + H = 0 & \text{on } \Omega, \\[0.2cm]
			u = b_j, \ \partial_\eta u = c_j  & \text{on $\partial_j \Omega$} \\[0.2cm]
			\inf_\Omega u > -\infty,
		\end{array} \right.
	\]
	for some $H,b_j,c_j \in \R$.
	\begin{itemize}
		\item[(i)] If $\partial \Omega$ is connected and mildly $v$-transverse for some $v$, then $\Omega$ is a half-plane and $\GG_u$ is a half-plane;
		\item[(ii)] Assume that $\partial \Omega$ is disconnected, and  that $\partial_\star \Omega$ is connected and mildly $v$-transverse for some $v$. Then, $\Omega = (0,T) \times \R$ and $\GG_u$ is a monotone piece of a cylinder.
		\item[(iii)] Assume that $\partial_\star \Omega = \partial_1\Omega \cup \partial_2\Omega$ is disconnected and mildly $v$-transverse for some $v$, with $\langle \eta,v \rangle \ge 0$ on $\partial_1\Omega$ and $\langle \eta,v \rangle \le 0$ on $\partial_2\Omega$. If $c_1c_2<0$, then $\Omega = (0,T) \times \R$ and $u$ is a monotone piece of either a cylinder or a plane.
%		\item[(iii)] Assume that $\partial \Omega$ is disconnected, that $\partial_1\Omega \cup \partial_2\Omega$ is mildly $v$-transverse for some $v$ with $\langle \eta,v \rangle \ge 0$ on $\partial_1\Omega$ and $\langle \eta,v \rangle \le 0$ on $\partial_2\Omega$, and that $c_j = 0$ for $j \ge 3$. If $c_1c_2<0$, then necessarily $b_1 \neq b_2$. Moreover, if 
%		\begin{equation}\label{eq_areagrowth_intro}
%		\text{for some $j \in \{1,2\}$,} \qquad \mathscr{H}^{1}(\partial_j \Omega \cap B_R) = o (R \log R) \qquad \text{as } \, R \to \infty,
%		\end{equation}
%		then $\Omega = (0,T) \times \R$ and $\GG_u$ is a monotone piece of either a cylinder or a plane.	
		\item[(iv)] Assume that $\partial_\star \Omega = \partial_1\Omega \cup \partial_2\Omega$ is disconnected and that $\partial_j \Omega$ is mildly $v_j$-transverse for $j \in \{1,2\}$ and some $v_j$ (possibly with $v_1 \neq v_2$). If $\inf_\Omega u$ is attained in a set $\Gamma$ disconnecting $\Omega$, which is the case for instance if  
		\[
		\inf_\Omega u < b_j \qquad \forall \, j \ge 3 
		\]
		and $\Gamma$ has an accumulation point in $\overline{\Omega}$, then $\Omega=(0,T) \times \R$ and $\GG_u$ is a piece of cylinder.
	\end{itemize}
\end{theorem}

\noindent The figures below show potentially capillary domains to which Theorem \ref{teo_CMC_intro_R2} applies, thereby excluding them. 

\begin{figure}[ht]\label{fig_1}
\caption{Cases $(i)$ and $(ii)$}
\begin{minipage}[b]{0.45\linewidth}
\centering
\includegraphics[scale=0.052]{figura1.png}
%\vspace{3cm}
\end{minipage}
\hspace{0.1cm}
\begin{minipage}[b]{0.45\linewidth}
\centering
\includegraphics[scale=0.051]{figura2.png}
\end{minipage}
\end{figure}

\begin{figure}[ht]\label{fig_2}
\caption{Cases $(iii)$ and $(iv)$}
\begin{minipage}[b]{0.45\linewidth}
\centering
\includegraphics[scale=0.05]{figura3.png}
%\vspace{3cm}
\end{minipage}
\hspace{0.1cm}
\begin{minipage}[b]{0.45\linewidth}
\centering
\includegraphics[scale=0.054]{figura4.png}
\end{minipage}
\end{figure}

\begin{remark}
	Some remarks are in order:
	\begin{itemize}
		\item[-] Differently from most of the literature quoted above, we make no assumption on the sign of $u$ compared to that of its boundary values. In particular, if $H\equiv 0$ the conclusion in case $(i)$ does not follow from Cui's result \cite{cui}. Moreover, $u$ may possibly be unbounded above. 
		\item[-] In $(ii),(iii)$ and $(iv)$ no restriction is made a-priori on the number of connected components of $\partial \Omega$, which may possibly be infinite. 
		\item[-] Cases $(i)$ and $(iii)$ improve on \cite[Thm. 1.3]{cmr}, where the conclusion was obtained under the stronger assumptions that $\Omega$ is a globally Lipschitz epigraph (for $(i)$) or slab (for $(iii)$). 
		\item[-] Case $(iv)$ is related to the recent \cite{abm}. There, Agostiniani, Borghini and Mazzieri proved that a solution to 
		\begin{equation}\label{pb_abm}
			\left\{
			\begin{array}{ll}
				\Delta u - 1 = 0 & \text{in an annular domain }  \Omega \subseteq \R^2,  \\[0.2cm]
				u < 0 & \text{in } \Omega, \\[0.2cm]
				u=0, \ \partial_\eta u = c_1 \in \R & \text{on } \, \partial_1\Omega, \\[0.2cm]
				u=0, \ \partial_\eta u = c_2 \in \R  & \text{on } \, \partial_2\Omega
			\end{array}\right.
		\end{equation}
		must be radially symmetric provided that the set of minima of $u$ has an accumulation point. They also exhibit a non-radial annular domain $\Omega$ supporting a solution to \eqref{pb_abm} with a finite number of minimum points. This may suggest that our assumption on $\Gamma$ in $(iv)$ is necessary. We mention that an analogue of the main theorem in \cite{abm} for solutions to $\MM[u] - 1 = 0$ is yet to be established. The higher dimensional case of \eqref{pb_abm} was recently tackled in \cite{abbm}. 		
	\end{itemize}
\end{remark}

For domains $\Omega \subseteq \R^3$, we have a similar statement. 

\begin{theorem}\label{teo_CMC_intro_R3}
	Let $\Omega \subseteq \R^3$ be a $C^2$ domain satisfying
	\begin{equation}\label{eq_quadraticgrowth}
		|\Omega \cap B_R| = o(R^2 \log R) \qquad \text{as } \, R \to \infty,
	\end{equation}
	and supporting a non-constant solution $u \in C^2(\overline\Omega)$ to 
	\[
		\left\{ 
		\begin{array}{ll}
			\MM[u]  + H = 0 & \text{in } \Omega, \\[0.2cm]
			u = b_j, \ \partial_\eta u = c_j  & \text{on $\partial_j \Omega$} \\[0.2cm]
			\inf_\Omega u > -\infty,
		\end{array} \right.
	\]
	for some $H,b_j,c_j \in \R$. Then,   
	\begin{itemize}
		\item[(i)] If $\partial_1 \Omega$ is mildly $v$-transverse for some $v$ and $\langle v,\eta \rangle \not\equiv 0$ on $\partial_1 \Omega$, then $\partial \Omega$ must be disconnected;
		\item[(ii)] Assume that $\partial \Omega$ is disconnected, that $\partial_\star \Omega$ is connected and mildly $v$-transverse for some $v$, and that $\langle v,\eta \rangle \not\equiv 0$ on $\partial_\star \Omega$. Then, $\Omega = (0,T) \times \R^2$ and $\GG_u$ is a monotone piece of a cylinder.
		\item[(iii)] Assume that $\partial_\star \Omega = \partial_1\Omega \cup \partial_2\Omega$ is disconnected and mildly $v$-transverse for some $v$, with $\langle \eta,v \rangle \ge 0$ on $\partial_1\Omega$, $\langle \eta,v \rangle \le 0$ on $\partial_2\Omega$ and $\langle \eta, v \rangle \not\equiv 0$ on $\partial_\star \Omega$. If $c_1c_2<0$, then $\Omega = (0,T) \times \R^2$ and $u$ is a monotone piece of either a cylinder or a hyperplane.
		\item[(iv)] Assume that $\partial_\star \Omega = \partial_1\Omega \cup \partial_2\Omega$ is disconnected and that $\partial_j \Omega$ is mildly $v_j$-transverse for $j \in \{1,2\}$ and some $v_j$ (possibly with $v_1 \neq v_2$), with $\langle \eta, v_j \rangle \not \equiv 0$ on $\partial_j\Omega$. If $\inf_\Omega u$ is attained in a set $\Gamma$ disconnecting $\Omega$, which is the case for instance if  
		\[
		\inf_\Omega u < b_j \qquad \forall \, j \ge 3 
		\]
		and $\Gamma$ has Hausdorff measure $\haus^{2}(\Gamma)>0$, then $\Omega=(0,T) \times \R^2$ and $\GG_u$ is a piece of cylinder.
	\end{itemize}
\end{theorem}

%\begin{remark}
%	An example of domain satisfying the assumptions in Theorem \ref{teo_CMC_intro_R3} is a slab whose defining functions $\phi_1,\phi_2$ saytisfy
%	\[
%	|\phi_j(x')| = o\Big( \log |x'|\Big) \qquad \text{as } \, |x'| \to \infty.
%	\]
%\end{remark}

\begin{remark}
	While not including the example of half-spaces in $\R^3$, condition \eqref{eq_quadraticgrowth} is still quite general and satisfied, for instance, if $\Omega$ is a slab  whose defining functions $\phi_1,\phi_2$ grow as follows:
	\[
		|\phi_j(x')| = o\Big( \log |x'|\Big) \qquad \text{as } \, |x'| \to \infty
	\]
	(note that no $L^\infty$-bound is imposed on the gradient of $\phi_j$). Also, mildly $e_1$-transverse regions like, for instance,
	\[
		\left\{ (x^1,x')\in \R^3 \ : \ |x'|\le 1 \ \text{or } \, |x^1| \le \frac{|x'|^2 + 1}{|x'|^2 -1} + \sqrt{\log|x'|} \right\}
	\]
	satisfy \eqref{eq_quadraticgrowth}. 
\end{remark}

\begin{remark}
	Case $(ii)$ improves on \cite[Thm 1.3]{cmr}, where the authors obtained the same conclusion by assuming that $\Omega$ is a globally Lipschitz, globally bounded slab.
\end{remark}


%
%
%
%
%Interestingly, Theorem \ref{teo_CMC_intro_R3} does not apply to the setting in $(iii)$ of Theorem \ref{teo_CMC_intro_R2}. In fact, to treat the problem in $\R^n$ our methods require to replace  \eqref{eq_areagrowth_intro} by 
%\[
%\text{for some $j \in \{1,2\}$,} \qquad \mathscr{H}^{n-1}(\partial_j \Omega \cap B_R) = o (R \log R) \qquad \text{as } \, R \to \infty,
%\]
%which fails for pieces of cylinder or hyperplanes if $n \ge 3$.
%%and \ref{teo_CMC_intro_R3} does not include the possibility that $c_1>0$, $c_2<0$, which occurs for instance if $\GG_u$ is a strictly monotone piece of cylinder. In fact, the first condition in \eqref{condi_Hcj} fails in this case. 
%However, this case was considered in \cite{cmr} for $C^3$ domains, and in particular we have: 
%%\tcr{[sottolineare che, nel teorema di \cite{cmr} riportato qui, l'ipotesi di regolarità sul bordo di $\Omega$ era $C^3$, ma può essere indebolita a $C^2$ sfruttando i risultati di questo lavoro]}
%
%\begin{theorem}\cite[Theorem 1.3]{cmr}
%	Let $\Omega \subseteq \R^3$ be a globally Lipschitz, globally bounded $C^3$ slab supporting a solution to 
%	\begin{equation}\label{eq_PMC_cmc1}
%	\left\{ 
%	\begin{array}{ll}
%		\MM[u]  + H = 0 & \text{in } \Omega, \\[0.2cm]
%		u = 0, \ \partial_\eta u = c_1 > 0  & \text{on $\partial_1 \Omega$} \\[0.2cm]
%		u = b > 0, \ \partial_\eta u = c_2 < 0  & \text{on $\partial_2 \Omega$} \\[0.2cm]
%		\inf_\Omega u > -\infty
%	\end{array} \right.	
%	\end{equation}
%\tcb{for some $b,c_1,c_2 \in \R$.} Then, $\Omega=(0,T) \times \R^2$ and $\GG_u$ is a monotone piece of either a cylinder or a hyperplane. 
%\end{theorem}
%



Theorems \ref{teo_CMC_intro_R2} and \ref{teo_CMC_intro_R3} follow from a general splitting result, which holds in a manifold setting without dimensional restriction and for general sources $f$. Hereafter, $(M^n,\metric)$ will be a complete Riemannian manifold of dimension $n \ge 2$ with Ricci curvature $\Ric$ and sectional curvature $\Sec$. Having fixed an origin $o \in M$, we denote by $B_r$ the geodesic ball of radius $r$ centered at $o$. Hereafter, a vector field $X$ of class $C^1$ in $\overline{\Omega}$ will be said to be Killing in $\overline{\Omega}$ if the Killing condition $\mathscr{L}_X g = 0$ is satisfied pointwise there. 


\begin{theorem}\label{teo_main_MC_manifolds}
	Let $\Omega \subseteq M^n$ be a $C^2$ domain satisfying
	\[
		\Ric \ge 0 \quad \text{in } \, \Omega, \qquad |\Omega \cap B_R| = o(R^2 \log R) \qquad \text{as } \, R \to \infty
	\]
	and supporting a non-constant, stable solution $u \in C^3(\Omega) \cap C^2(\overline\Omega)$ of 
	\begin{equation}\label{eq_PMC_cmc2}
		\left\{ 
		\begin{array}{ll}
			\MM[u]  + f(u) = 0 & \quad\text{in } \Omega, \\[0.2cm]
			u = b_j, \ \partial_\eta u = c_j  & \quad \text{on $\partial_j \Omega$} \\[0.2cm]
			\inf_\Omega u > -\infty
		\end{array} \right.
	\end{equation}
	for some $b_j,c_j \in \R$ and $f \in C^1(\R)$. Assume that one of the following sets of assumptions is met for some constant $C_0$:
	\begin{itemize}
		\item[(A)] $- C_0 \le f(u) \le 0$ in $\Omega$ and $u$ is constant on $\partial_\star \Omega$.
%		\item[(B)] $f(u) \le 0$ in $\Omega$, $u$ is a constant $b$ on $\partial_\star \Omega$ and $\inf_{\Omega} u = b$ \tcr{(NON MOLTO RILEVANTE, LA TOGLIEREI MAGARI A BENEFICIO DI QUESTA SOTTO)}.
		\item[(B)] $- C_0 \le f(u) \le 0$ in $\Omega$, $u$ is bounded on $\partial_\star \Omega$ and 
		\[
		\haus^{n-1}(\partial_\star^{b}\Omega \cap B_R)= o \big(R^2 \log R\big) \qquad \text{as } \, R \to \infty
		\]
		for some $b \in \R$, where $\partial_\star^b\Omega=\{ x \in \partial_\star \Omega : u(x) \neq b\}$.
		\item[(C)] $u$ is constant on $\partial_\star \Omega$ and bounded in $\Omega$.
		\item[(D)] $|f(u)| \leq C_0$ and $f'(u) \le 0$ in $\Omega$, $\sup_j |c_j| < \infty$ and $\Sec \ge - \kappa$ in $M$ for some $\kappa \in \R^+$. 
	\end{itemize}			
%		
%	
%	
%	
%	
%	either
%	\begin{equation} \label{condi_fu1}
%		f(u) \leq 0 \quad \text{in } \, \Omega , \qquad u \geq 0 \quad \text{in } \, \Omega , \qquad u \equiv 0 \quad \text{on } \, \partial_\star \Omega
%	\end{equation}
%	or
%	\begin{equation} \label{condi_fu2}
%		f'(u) \leq 0 \quad \text{in } \, \Omega , \qquad \inf_\Omega u > -\infty \, , \qquad \Sec > -\kappa^2 \quad \text{in } \, M
%	\end{equation}
%	for some $\kappa\in\R$, and that
%	\begin{equation}\label{condi_fu}
%		\begin{array}{l}
%			\text{$f(u)$ does not change sign in $\Omega$;} \\[0.2cm] 
%			\text{$u$ is constant on } \,  \partial_\star \Omega \doteq \bigcup_{j : c_j \neq 0} \partial_j\Omega,
%		\end{array} 
%	\end{equation}
	If there exists a Killing field $X$ on $\overline\Omega$ with $|X| \in L^\infty(\Omega)$ and  
	\begin{equation}\label{condi_Hcj}
		\begin{array}{l}
			%			\text{$u$ is constant on } \,  \bigcup_{j : c_j \neq 0} \partial_j\Omega, \\[0.3cm] 
			%			H c_j \le 0, \quad 
			\disp c_j \langle \eta, X \rangle \ge 0 \qquad \text{for each $j$,} \\[0.2cm]
			\disp c_j \langle \eta, X \rangle \not \equiv 0 \qquad \text{for some $j$,}
		\end{array} 
	\end{equation}
	Then:
	\begin{itemize}
		\item[i)] $\Omega = (0,T) \times N$ with the product metric, for some $T \le \infty$ and some complete, totally geodesic hypersurface $N \subseteq M$. 
		\item[ii)] $u= u(t)$ is 1D and strictly monotone in the arclength $t \in (0,T)$, 
		\item[iii)] $\langle \partial_t, X \rangle$ is constant in $\Omega$.
	\end{itemize}
%	$\Omega = (0,T) \times N$ with the product metric, for some $T \le \infty$ and some complete, totally geodesic hypersurface $N \subseteq M$. Moreover, $u$ is 1D and strictly monotone in the arclength $t \in (0,T)$.
\end{theorem}

%\tcr{[Precisare che cosa intendiamo per $X$ campo di Killing in $\overline{\Omega}$. Per la regolarità credo che basti assumere $X$ definito in $\overline{\Omega}$ e di classe $C^1(\overline{\Omega})$. Possiamo dare come definizione di campo di Killing la condizione puntuale di antisimmetria su $\nabla X$, precisando che sugli aperti ciò equivale al fatto che il flusso locale di $X$ sia costituito da isometrie locali.]}


%
%Here is our first theorem. Denote by  $\{\partial_j\Omega\}$, $j \le \infty$ the connected components of $\partial \Omega$, and by $B_R$ balls centered at the origin of $\R^n$.
%
%\begin{theorem}\label{teo_mainR2R3}
%	Let $\Omega \subseteq \R^n$, $n \in \{2,3\}$ be a $C^3$-domain supporting a non-constant stable solution $u \in C^3(\Omega) \cap C^2(\overline\Omega)$ to 
%	\begin{equation}\label{eq_PMC_cmc2}
%		\left\{ 
%		\begin{array}{ll}
%			\MM[u]  + f(u) = 0 & \text{on } \Omega, \\[0.2cm]
%			u = b_j, \ \partial_\eta u = c_j  & \text{on $\partial_j \Omega$} 	
%		\end{array} \right.
%	\end{equation}
%	for some $b_j,c_j \in \R$ and $f \in C^1(\R)$. Assume that either $n=2$ or 
%	\begin{equation}\label{eq_volumegrowth}
%		n=3, \qquad |\Omega \cap B_R| = o(R^2 \log R) \qquad \text{as } \, R \to \infty.
%	\end{equation}
%	If
%	\begin{equation}\label{condi_fu}
%		\begin{array}{l}
%			\text{$f(u)$ does not change sign in $\Omega$;} \\[0.2cm] 
%			\text{$u$ is constant on } \,  \disp\bigcup_{j : c_j \neq 0} \partial_j\Omega,
%		\end{array} 
%	\end{equation}
%	and, for some $v \in \mathbb{S}^{n-1}$, 
%	\begin{equation}\label{condi_Hcj}
%		\begin{array}{l}
%			%			\text{$u$ is constant on } \,  \bigcup_{j : c_j \neq 0} \partial_j\Omega, \\[0.3cm] 
%			%			H c_j \le 0, \quad 
%			\disp c_j \langle \eta, v \rangle \ge 0 \qquad \text{for each $j$,} \\[0.2cm]
%			\disp c_j \langle \eta, v \rangle \not \equiv 0 \qquad \text{for some $j$,}
%		\end{array} 
%	\end{equation}
%	Then $\Omega = (0,T) \times \R^{n-1}$ for some $T \le \infty$ and $u$ is 1D . Moreover, $u$ is strictly monotone in the coordinate $t$ on $(0,T)$.
%	%		\item If $\partial \Omega$ has two connected components, $u$ attains a global maximum in $\Omega$ and the set of maxima is not discrete, then		
%\end{theorem}

\begin{remark}
	\item[-] Again, no growth from above is imposed on $u$, and no restriction is made on the number of components of $\partial \Omega$. 
%	\item[-] If the second in \eqref{condi_fu} fails, i.e. if $u$ is not globally constant in the union of all components where $\partial_\eta u \neq 0$, then one achieves the conclusion in Theorem \ref{teo_main_MC_manifolds} by requiring that 
%	\[
%	\haus^{n-1}(\partial_\star \Omega \cap B_R) = o(R \log R) \qquad \text{as } \, R \to \infty.
%	\]

%$$
%|\nabla u|\Delta|\nabla u|\ge|\nabla \di u|^2-|\nabla|\nabla u|^2+\Ric(\nabla u,\nabla u)-f'(u)|\nabla u|^2
%$$
%$$
%\int_{\partial\Omega}+\int\phi^2|\nabla|\nabla u||^2
%$$
%
%	\item[-] While not including the example of half-spaces in $\R^3$, condition \eqref{eq_volumegrowth} is still quite general and satisfied, for instance, if $\Omega$ is a slab  whose defining functions $\phi_1,\phi_2$ saytisfy
%	\[
%	|\phi_j(x')| = o\Big( \log |x'|\Big) \qquad \text{as } \, |x'| \to \infty
%	\]
%	(note that no $L^\infty$-bound is imposed on the gradient of $\phi_j$). Also, mildly $e_1$-transverse regions like, for instance,
%	\[
%	\left\{ (x^1,x')\in \R^3 \ : \ |x'|\le 1 \ \text{or } \, |x^1| \le \frac{|x'|^2 + 1}{|x'|^2 -1} + \sqrt{ \log|x'|} \right\}
%	\]
%	satisfy \eqref{eq_volumegrowth}. 
	%
	\item[-] The $t$-direction in the splitting $\Omega = (0,T) \times \R^{n-1}$ may not coincide with $X$. %However, as a byproduct of our proof, $\langle X, \partial_t \rangle$ is constant in $\Omega$.
	\item[-] Condition \eqref{condi_Hcj}, automatically satisfied by monotone solutions with $X = \pm \partial_t$, implies the request that $\partial_\star \Omega$ be mildly $X$-transverse. We stress that \eqref{condi_Hcj} only depends on $X$ and on the boundary conditions imposed on $u$, and is therefore easily checkable in explicit examples.
	\item[-] The actual existence of monotone solutions depends on the properties of $f$, and should be examined case by case (for solutions on $\R$, see \cite{FSV}). Since our methods do not require the a priori knowledge of such solutions, we have not performed the 1D analysis here. Observe that, in case no monotone 1D solutions exist, Theorems \ref{teo_main_MC_manifolds} and \ref{teo_splitting_intro} can be seen as non-existence results. 	
	%	
	%	, and that the $c_j$'s have signs compatible with the possible existence of a monotone solution. 
\end{remark}

\begin{remark}
	Conditions $(A)$ to $(D)$ are instrumental to prove that $u$ has \emph{moderate energy growth}, a crucial property to obtain Theorem \ref{teo_main_MC_manifolds}, see Section \ref{sec_energy} for more details. Indeed, further conditions are shown to imply moderate energy growth, and we refer to Theorem \ref{teo_goodcutoff_MC} for a more extensive list. We have  selected $(A)-(D)$ for the sake of simplicity. 	
\end{remark}

\begin{remark}
	A closely related result was obtained in \cite[Theorem 1.6]{cmr} for constant $f$. In fact, the assumption that $f$ be constant is essential for the methods therein to work. However, even in this case Theorem \ref{teo_main_MC_manifolds} improves on \cite{cmr} in some aspects, including the fact that the set $\{c_j\}$ is not a priori required to be bounded in cases $(A)$, $(B)$ and $(C)$.
\end{remark}

In view of Theorem \ref{teo_main_MC_manifolds}, Theorems \ref{teo_CMC_intro_R2} and \ref{teo_CMC_intro_R3} allow for extension to  stable solutions to $\MM[u] + f(u) = 0$. We only state the result in $\R^2$.

\begin{theorem}\label{teo_CMC_intro_R2_gen}
	Let $\Omega \subseteq \R^2$ be a $C^2$ domain supporting a non-constant stable solution $u \in C^3(\Omega) \cap C^2(\overline\Omega)$ to 
	\[
		\left\{ 
		\begin{array}{ll}
			\MM[u]  + f(u) = 0 & \text{on } \Omega, \\[0.2cm]
			u = b_j, \ \partial_\eta u = c_j  & \text{on $\partial_j \Omega$} \\[0.2cm]
			\inf_\Omega u > -\infty,
		\end{array} \right.
	\]
	for some $b_j,c_j \in \R$ and $f \in C^1(\R)$. Consider the following assumptions for some constant $C_0$:
	\begin{itemize}
		\item[(A)] $- C_0 \le f(u) \le 0$ in $\Omega$; %, and $u$ is constant on $\partial_\star \Omega$.
%		\item[(B)] $f(u) \le 0$ in $\Omega$, $u$ is a constant $b$ on $\partial_\star \Omega$ and $\inf_{\Omega} u = b$ \tcr{(NON MOLTO RILEVANTE, LA TOGLIEREI)}.
		\item[(B)] $- C_0 \le f(u) \le 0$ in $\Omega$ and $\haus^{1}(\partial_1\Omega \cap B_R)= o \big(R^2 \log R\big)$ as $R \to \infty$;
		\item[(C)] $u$ is bounded in $\Omega$; %constant on $\partial_\star \Omega$ and 
		\item[(D)] $|f(u)| \leq C_0$ and $f'(u) \le 0$ in $\Omega$.
	\end{itemize}			
	Then, 
	\begin{itemize}
		\item[(i)] Assume one among $(A),(C)$ and $(D)$.\\
		If $\partial \Omega$ is connected and mildly $v$-transverse for some $v$, then $\Omega$ is a half-plane and $u$ is a $1D$, monotone solution;
		\item[(ii)] Assume one among $(A),(C)$ and $(D)$.\\
		If $\partial \Omega$ is disconnected, and $\partial_\star \Omega$ is connected and mildly $v$-transverse for some $v$, then $\Omega = (0,T) \times \R$ and $u$ is a $1D$, monotone solution.
		\item[(iii)] Assume one between $(B)$ and $(D)$.\\
		Assume that $\partial_\star \Omega = \partial_1\Omega \cup \partial_2\Omega$ is disconnected and mildly $v$-transverse for some $v$, with $\langle \eta,v \rangle \ge 0$ on $\partial_1\Omega$ and $\langle \eta,v \rangle \le 0$ on $\partial_2\Omega$. If $c_1c_2<0$, then $\Omega = (0,T) \times \R$ and $u$ is a $1D$, monotone solution.	
		\item[(iv)] Assume one among $(A),(C)$ and $(D)$.\\
		Assume that $\partial_\star \Omega = \partial_1\Omega \cup \partial_2\Omega$ is disconnected, and that $\partial_j \Omega$ is mildly $v_j$-transverse for $j \in \{1,2\}$ and some $v_j$ (possibly with $v_1 \neq v_2$). If $\inf_\Omega u$ is attained in a set $\Gamma$ disconnecting $\Omega$, which is the case for instance if all of the following are satisfied: 
		\begin{itemize}
			\item $\inf_\Omega u < b_j$ for each $j \ge 3$;
			\item $f$ is analytic in a neighbourhood of $\inf_\Omega u$ and $f(\inf_\Omega u) \neq 0$, 
			\item $\Gamma$ has an accumulation point in $\overline\Omega$,
		\end{itemize}	
		then $\Omega=(0,T) \times \R$ and $u$ is a $1D$ solution.
	\end{itemize}
\end{theorem}

The domains depicted in Figures \ref{fig_1} and 2 above provide examples to which Theorem \ref{teo_CMC_intro_R2_gen} applies.

\begin{remark}
	Observe that the capillarity case $f(u) = - \kappa u$ is included by $(D)$. Moreover, because of $(C)$ Theorem \ref{teo_CMC_intro_R2_gen} allows to treat the pairs $(u,f)$ considered by Lian and Sicbaldi \cite{lian_sicbaldi}, providing complementary classification results when $\partial \Omega$ is disconnected.
\end{remark}

%\begin{theorem}\label{teo_CMC_intro_R2}
%	Let $\Omega \subseteq \R^2$ be a $C^3$ domain supporting a non-constant \tcb{weak solution $u \in C^1(\overline\Omega)$} to 
%	\begin{equation}\label{eq_PMC_cmc1}
%		\left\{ 
%		\begin{array}{ll}
%			\MM[u]  + H = 0 & \text{on } \Omega, \\[0.2cm]
%			u = b_j, \ \partial_\eta u = c_j  & \text{on $\partial_j \Omega$} 	
%		\end{array} \right.
%	\end{equation}
%	for some $H,b_j,c_j \in \R$. 
%	\begin{itemize}
%		\item[(i)] If $\partial \Omega$ is connected and mildly $v$-transverse for some $v$, then $\Omega$ is a half-plane and $\GG_u$ is a half-plane;
%		\item[(ii)] Assume that $\partial \Omega$ is disconnected, that $\partial_1\Omega$ is mildly $v$-transverse for some $v$ and that $c_j = 0$ for $j \ge 2$. Then, $\Omega = (0,T) \times \R$ and $\GG_u$ is a monotone piece of cylinder.
%		\item[(iii)] Assume that $\partial \Omega$ has two connected components and is mildly $v$-transverse for some $v$, with $\langle \eta,v \rangle \ge 0$ on $\partial_1\Omega$ and $\langle \eta,v \rangle \le 0$ on $\partial_2\Omega$. If $c_1c_2<0$, then necessarily $b_1 \neq b_2$. Moreover, if 
%		\begin{equation}\label{eq_areagrowth_intro}
%			\text{for some $j \in \{1,2\}$,} \qquad \mathscr{H}^{1}(\partial_j \Omega \cap B_R) = o (R \log R) \qquad \text{as } \, R \to \infty,
%		\end{equation}
%		then $\Omega = (0,T) \times \R$ and $\GG_u$ is a monotone piece of either a cylinder or a plane.	
%		\item[(iv)] Assume that $\partial \Omega$ has two connected components, that $\partial_j\Omega$ is mildly $v_j$-transverse for some $v_j$ (possibly with $v_1 \neq v_2$) and that $c_1c_2 \neq 0$. If $\sup_\Omega u > \max\{b_1,b_2\}$ and it is attained in a set $\Gamma$ with an accumulation point, then $\Omega=(0,T) \times \R$ and $\GG_u$ is a piece of cylinder.
%	\end{itemize}
%\end{theorem}








We conclude this introduction by describing the main novelties of our work. In the generality of Theorem \ref{teo_main_MC_manifolds} the manifold does not possess enough symmetries to exploit the moving plane method as in \cite{serrin,reichel_1,sirakov,lian_sicbaldi}; indeed, even in our applications to $\R^2$ and $\R^3$ the method is not used. The strategy is instead based on some integral formulas relating $u$ to the function $\langle\nabla u,X\rangle$, which is a particular solution of the linearized problem. This approach follows \cite{cmr}, itself inspired by the integral inequalities obtained in \cite{Arma, Arma2} for the semilinear equation $\Delta u+f(u)=0$ in $\R^n$ and extended to manifolds in \cite{FarMarVal}. In the latter, the authors proved that if $u$ is monotone in the direction of a bounded Killing field $X$ in $\Omega$, then
\begin{equation}\label{eq_poica_lapla}
\int_\Omega\left\{|\nabla u|^2 |\II|^2 +\left|\nabla^\top|\nabla u|\right|^2+\Ric(\nabla u,\nabla u)\right\}\varphi^2\le\int_\Omega|\nabla\varphi|^2|\nabla u|^2 
\end{equation}
holds for each $\varphi\in\lip_c(\overline\Omega)$, where $\nabla^\top$ and $\II$ are, respectively, the tangential gradient and the second fundamental form of the level sets of $u$. In particular, test functions $\varphi$ whose support intersects $\partial\Omega$ are allowed. In fact, by a remarkable cancellation, the overdetermined boundary conditions force the boundary term appearing in the computations leading to \eqref{eq_poica_lapla} to vanish identically \emph{provided that} $\langle \nabla u, X \rangle$ has a sign in $\Omega$, say it is positive there. Testing \eqref{eq_poica_lapla} with suitable cut-off functions $\varphi$ and assuming $\Ric \ge 0$, one is able to deduce that 
\begin{equation}\label{eq_split}
|\nabla u|^2 |\II|^2 +\left|\nabla^\top|\nabla u|\right|^2 \equiv 0 \qquad \text{in } \, \Omega
\end{equation}
whenever the energy of $u$ satisfies 
\begin{equation}\label{eq_bounden_u}
\int_{B_R}|\nabla u|^2 = O \left( R^2 \log R \right) \qquad \text{as } \, R \to \infty 
\end{equation}
for balls $B_R$ centered at a fixed origin. Property \eqref{eq_split} implies that $\Omega$ splits and $u$ is 1D. In the present paper, we first obtain a new integral inequality for monotone solutions in the direction of $X$, extending \eqref{eq_poica_lapla} and valid for  all quasilinear equations of the form
\begin{equation}\label{eq_alapla_gen}
\diver \big( a(|\nabla u|)\nabla u \big) +f(u) = 0
\end{equation}
under mild regularity assumptions on $a$ (essentially, those guaranteeing locally uniform ellipticity even at critical points of $u$), see Theorem \ref{teo_Poincare_loc}. Clearly, to get rigidity, condition \eqref{eq_bounden_u} shall be replaced by a growth tailored to $a$: for operators of mean curvature type and under suitable conditions on $f$ and $u_{|\partial_\star \Omega}$, by a calibration argument,  such a growth is granted by a mild volume growth condition on $\Omega$ and (in some cases) $\partial_\star \Omega$, see Theorem \ref{teo_goodcutoff_MC}.

A key point is therefore to show that $w \doteq \langle \nabla u, X \rangle > 0$ in $\Omega$. To the best of our knowledge, proofs of the monotonicity of $u$ in the direction of a Killing field (usually, a parallel field in $\R^n$) are currently available only for certain classes of nonlinearities $f$ and non-compact domains $\Omega$, and typically rely on the moving plane and/or the sliding method (see for instance \cite{bcn,FarMarVal} and the references therein). These techniques often require  \emph{global} gradient or H\"older estimates on $u$ to be applied. For the mean curvature operator such estimates turn out to be subtle, especially when $f$ is not monotone or under just lower bounds on the Ricci curvature of $M$, see \cite{cmmr,cmr}. In \cite{cmr}, we proposed a new approach to get monotonicity for CMC capillary graphs, by combining pointwise gradient estimates and parabolicity arguments. We here complement \cite{cmr} with a new technique depending on energy estimates.
%The underlying observation is that the energy growth on $u$ that allows to get rigidity from \eqref{eq_poica_lapla} is the same request that allows to get rigidity for positive supersolutions of the linearized problem
%\[
%\phi \longmapsto \MM'(u)[\phi] + f'(u) \phi.
%\]
The method is quite versatile, and allows to consider much more general classes of nonlinearities $f$. In particular, we do not  need $f' \le 0$. Observe that request \eqref{condi_Hcj} rephrases as $w \ge 0$ on $\partial \Omega$, a necessary condition for monotonicity, and $w \not\equiv 0$ on $\partial \Omega$, the latter just to avoid a pathological case (which can be handled in the case of $\R^2$, see Section \ref{sec_critical}).

The solutions in Theorem \ref{teo_main_MC_manifolds} are monotone, and so are those in all items in Theorems \ref{teo_CMC_intro_R2}, \ref{teo_CMC_intro_R3} and \ref{teo_CMC_intro_R2_gen} but $(iv)$. To obtain $(iv)$, where 1D solutions are not monotone, one needs to separately treat each connected component $\Omega_0$ of $\Omega \backslash \Gamma$, prove rigidity for $\Omega_0$ and then glue the resulting 1D solutions. This requires a careful analysis, carried out in Proposition \ref{prop_splitting_alt}, Lemma \ref{lem_Du0} and Theorem \ref{teo_twopieces}, of independent interest.

A further feature to point out is that our techniques are not specific to the mean curvature operator, but allow to treat a large class of quasilinear equations under mild structural conditions, as described in the next section. Besides allowing for a greater generality, such an extension simplifies the resulting integral formulas making the role of the linearized operator of \eqref{eq_alapla_gen} and its eigenvalues more transparent.
%
%
%
%\begin{theorem}
%	Let $\Omega \subseteq \R^n$, $n \in \{2,3\}$ be a mildly $v$-transverse domain supporting a non-constant solution to 
%	\begin{equation}\label{eq_PMC_f}
%	\left\{ 
%	\begin{array}{ll}
%		\MM[u]  + f(u) = 0 & \text{on } \Omega, \\[0.2cm]
%		u = b_j, \ \partial_\eta u = c_j  & \text{on $\partial_j \Omega$} 	
%	\end{array} \right.
%	\end{equation}
%	for some $b_j,c_j \in \R$. Assume that either $n=2$ or 
%	\[
%	n=3, \qquad |\Omega \cap B_R| = o(R^2 \log R) \qquad \text{as } \, R \to \infty.
%	\]
%	\tcr{check the asusmptions on $f$ compared to the sign of $\partial_\eta u$, there must be some assumption more!!!}
%	\begin{itemize}
%	\item If 
%	\[
%	\begin{array}{l}
%	\disp c_j \langle \eta, v \rangle \ge 0 \qquad \text{for each $j$,} \\[0.2cm]
%	\disp c_j \langle \eta, v \rangle \not \equiv 0 \qquad \text{for some $j$,}
%	\end{array} 
%	\] 
%	Then $\Omega = (0,T) \times \R$ and $u$ is 1D  and monotone in the coordinate $t$ on $\R$.
%	\item If $\Omega$ has two connected components, $f$ is analytic and the set of maxima of $f$ consists of isolated points
%	
%	\end{itemize}
%\end{theorem}
%


\subsection{The general setting}

Let $(M^n,\metric)$ be a complete Riemannian manifold of dimension $n \ge 2$, $f \in C^1(\R)$ and consider the quasilinear equation 
\begin{equation}\label{equazu}
	\Delta_a u + f(u)=0 \qquad \text{on } \Omega \subseteq M,
\end{equation}
where
\begin{equation}\label{def_T}
	\Delta_a u = \diver\big(a(|\nabla u|)\nabla u\big) 
\end{equation}
and $a$ satisfies
\begin{equation}\label{assu_A}
	\begin{array}{rcl}
		 & \quad a(t) \in C^1(\R^+) \cap C(\R^+_0). \\[0.1cm]
		& \quad a(t)>0 \ \  \text{ for } t \in \R^+, \qquad \big(ta(t) \big)' >0 \ \ \text{ for } t \in \R^+,
	\end{array}
\end{equation}
For instance: 
\begin{itemize}
	\item[(a)] the mean curvature operator is obtained for the choice $a(t) = \frac{1}{\sqrt{1+t^2}}$, 
	\item[(b)] the $p$-Laplacian $\Delta_p$ for $p \in [2,\infty)$ is obtained with $a(t) = t^{p-2}$. More generally, for $2 \le p < q < \infty$ the $(p,q)$-Laplacian, introduced in  \cite{zhikov,marcellini}, is obtained with the choice $a(t) = t^{p-2} + t^{q-2}$. 
	\item[(c)] the operator of exponentially harmonic functions, first introduced in \cite{duc_eels, eels_lemaire}, is obtained for $a(t)= e^{t^2}$.
\end{itemize}

More examples can be found in \cite{serrin2}. In most of our results we shall restrict to operators satisfying
\begin{equation}\label{assu_A_strong}
a \in C^1(\R^+_0), \qquad a(t), \big( ta(t)\big)' > 0 \quad \text{for } \, t \in \R^+_0
\end{equation}
%
or the stronger
\begin{equation}\label{assu_A_superstrong}
	a\in C^{1,1}_\loc(\R_0^+), \qquad a(t), \big( ta(t)\big)' > 0 \quad \text{for } \, t \in \R^+_0.
\end{equation}
In particular, referring to the above examples, (a) and (c) satisfy \eqref{assu_A_superstrong}, as well as the standard $2$-Laplacian, while $\Delta_{2,q}$ in (b) satisfies \eqref{assu_A} for each $q>2$, \eqref{assu_A_strong} for $q \ge 3$ and \eqref{assu_A_superstrong} for $q \in \{3\} \cup [4,\infty)$. Further interesting operators treatable with the techniques below include that of the polytropic flow from gas dynamics, see \cite{sibner_sibner}, for which   
\[
a(t) = \left( 1 - \frac{\gamma-1}{2}t^2 \right)^\frac{1}{\gamma-1}, \qquad \gamma >1.
\]
Here, $a$ represents the density of a fluid whose velocity is $\nabla u$, and \eqref{assu_A_strong} is met if $t$ restricts to the interval 
\[
t \in (0, \upsilon^*), \qquad \upsilon^* \doteq \sqrt{ \frac{2}{\gamma+1}},
\]
which identifies the subsonic regime. Hence, our results are applicable provided that $\|\nabla u\|_\infty < \upsilon^*$.

 
The formal linearization of $\Delta_a$ at $u$ is
\begin{equation}\label{def_linearized}
%	\Delta_a'(u)[\varphi] = 
\left.\frac{\di}{\di s} \right|_{s=0} \Delta_a(u+s\varphi) =  \diver\big(A(\nabla u) \nabla \varphi\big),
\end{equation}
where $A : \Gamma(TM\backslash \{0\}) \ra \Gamma(S^{(1,1)}(M))$ is the field of endomorphisms of $TM$ given by 
\begin{equation}\label{def_A}
	A(X) = \frac{a'(|X|)}{|X|} X_\flat \otimes X + a(|X|) {\rm Id},
\end{equation}
and $\flat$ is the musical isomorphism induced by $\metric$. Note that the eigenvalues of $A(X)$ are $\{\lambda_j(|X|)\}_{j=1}^n$ with
\begin{equation}\label{eigenvalues}
	\lambda_1(t) = a(t) + ta'(t) = \big(t a(t)\big)' \qquad \lambda_2(t) = \ldots = \lambda_n(t) = a(t),
\end{equation}
hence the second in \eqref{assu_A_strong} rephrases as $\lambda_1(t),\lambda_2(t) > 0$ on $\R^+_0$, equivalently $A(X)$ is elliptic, nonsingular and non-degenerate. Henceforth, we write
\[
\lambda_{\max}(t) = \max \big\{ \lambda_1(t), \lambda_2(t) \big\}.
\]
For instance, for the mean curvature operator, 
\[
\lambda_1(t) = (1+t^2)^{-3/2}, \qquad \lambda_2(t) = (1+t^2)^{-1/2} = \lambda_{\max}(t).
\]
%

\begin{definition}
	Let $a$ satisfy \eqref{assu_A_strong}. A function $u \in C^1(\Omega)$ weakly solving \eqref{equazu} is said to be \emph{stable} if
	\begin{equation}\label{hardy}
		\int_{M} f'(u) \varphi^2\di x \le \int_M \langle A(\nabla u)\nabla\varphi, \nabla \varphi\rangle \di x \qquad \forall \, \varphi \in \lip_c(\Omega).
	\end{equation}
\end{definition}
%

%\begin{remark}
%	\tcr{SHOULD WE KEEP IT?}\\
%	\tcr{Our regularity restrictions on $u$ and $a$ are imposed to guarantee, among other things,  the validity of some key tools for non-negative solutions to 
%	\[
%	\diver\left( A(\nabla u) \nabla w\right) + f'(u) w \le 0, 
%	\]
%	including the Hopf Lemma and the Harnack inequality, see Section \ref{sec_stability}. It might be possible that the techniques in the present paper allow to treat certain classes of singular or degenerate operators.
%	} 
%\end{remark}


%\begin{remark}
%	\emph{The integral in the RHS is not automatically finite under the sole assumptions $(a1),(a2)$. Indeed, if the eigenvalue $\lambda_1(|X|)$ of $A$ diverges as $|X| \ra 0^+$ (as in the case of the $p$-Laplacian for $p \in (1,2)$), then the RHS might diverge when $\nabla u$ is zero on a large set. However, this unpleasant situation can be avoided by assuming that
%		$$
%		(a3) \qquad t a(t) \in \lip_\loc(\R^+_0), 
%		$$
%		or in other words that $\lambda_1(t)$ be bounded in a neighbourhood of zero.
%	}
%\end{remark}

%We shall study overdetermined problems for $\Delta_a u + f(u) = 0$ with $f \in C^1(\R)$ and $a$ satisfying the structural condition

%
%
% In this paper, we shall investigate solutions to overdetermined problems of the type
%\begin{equation}\label{overdetermined}
%	\left\{ \begin{array}{ll}
%		%
%		\Delta_a u + f(u) = 0 & \quad \text{on } \, \Omega \subseteq M,\\[0.2cm]
%		u=0 & \quad \text{on } \, \partial \Omega \\[0.2cm]
%		\partial_{\eta} u & \quad \text{locally constant on } \, \partial \Omega,
%	\end{array}\right.
%\end{equation}
%
%\tcr{IMPORTANT: try to consider the case of an annulus... can we prove that if the annulus is star-shaped and the solution is monotone, then the annulus is a ball.}

%where $\Omega$ is an open domain wih smooth boundary and exterior normal $\eta$, $f \in C^1(\R)$ and 
%%
%This class encompasses relevant examples such as 
%
%is formally characterized, on each open subset $U \Subset \Omega$, as a stationary point of the energy functional $E_U : \lip_c(U) \ra\R$ given by
%\begin{equation}\label{energia}
%	E_U(w) =  \int_{U} \Big\{\tcr{\Theta(|\nabla w|)} - F(w)\Big\} \di x, 
%\end{equation}
%where
%\begin{equation}\label{def_FG}
%	F(t) = \int_0^t f(s) \di s, \qquad \Theta(t) = \int_0^t sa(s)\di s,
%\end{equation}
%with respect to compactly supported, smooth variations in $U$. The solution is called \emph{stable} if $E_U'' \ge 0$, that is, if  
%\begin{equation}
%\int_\Omega \Big[ \langle A(\nabla u) \nabla \phi, \nabla \phi \rangle - f'(u)\phi^2\Big] \qquad \forall \, \phi \in \lip_c(\Omega), 
%\end{equation}
%where $A(X)$ is the endomorphism of $TM$ associated to the symmetric bilinear form 
%\begin{equation}\label{def_A_intro}
%\frac{a'(|X|)}{|X|} X_\flat \otimes X_\flat + a(|X|) \metric
%\end{equation}
%and $\flat$ is the musical isomorphism induced by $\metric$. 

Here is our first main theorem for general operators: 

\begin{theorem} \label{teo_splitting_intro}
	Let $(M^n, \metric)$ be a complete Riemannian manifold. Let $\Omega$ be a $C^2$ domain with interior normal $\eta$, and assume that $\Ric\geq 0$ in $\Omega$. Let $a$ satisfy \eqref{assu_A_superstrong}, $f \in C^1(\R)$ and let $u \in C^3(\Omega)\cap C^2(\overline{\Omega})$ be a non-constant, stable solution to
	\begin{equation}\label{equazu_conbordo}
	\left\{
	\begin{array}{l@{\qquad}l}
		\Delta_a u + f(u) = 0 & \text{in } \Omega \\[0.3cm]
		u = b_j, \ \partial_\eta u = c_j  & \text{on $\partial_j \Omega$}. 
	\end{array}
	\right.
	\end{equation}
	%\tcr{and such that $\overline{u(\partial\Omega)}$ is not an interval}.
	If there exists a Killing field $X$ on $\overline\Omega$ with $|X| \in L^\infty(\Omega)$ and  
	\begin{equation}\label{condi_Hcj_con_a}
		\begin{array}{l}
			\disp c_j \langle \eta, X \rangle \ge 0 \qquad \text{for each $j$,} \\[0.2cm]
			\disp c_j \langle \eta, X \rangle \not \equiv 0 \qquad \text{for some $j$,}
		\end{array} 
	\end{equation}
	and 
	\begin{equation}\label{eq_lowenergy}
			\int_{\Omega \cap B_R} |\nabla u|^2 \lambda_{\max}(|\nabla u|) = O(R^2 \log R) \qquad \text{as } \, R \to \infty \, ,
	\end{equation}
	then: 
	\begin{itemize} 
		\item[i)] $\Omega = (0,T) \times N$ with the product metric, for some $T \le \infty$ and some complete, totally geodesic hypersurface $N \subseteq M$. 
		\item[ii)] $u= u(t)$ is 1D and strictly monotone in the arclength $t \in (0,T)$, 
		\item[iii)] $\langle \partial_t, X \rangle$ is constant in $\Omega$.
	\end{itemize}
	%
	%
	%     QUESTA È LA RISCRITTURA PIÙ DETTAGLIATA 
%	\begin{itemize}
%		\item [i)] there exists an isometry $\Psi : I \times N \to \Omega$, where $I\subseteq\R$ is an open interval and $N$ is a complete, totally geodesic hypersurface of $M$ with $\Ric_N \geq 0$
%		\item [ii)] $u(\Psi(t,x)) = u_0(t)$ for all $(t,x)\in I\times N$, for a suitable function $u_0 : I \to \R$ with $u_0'>0$
%		\item [iii)] $\langle \nu,X \rangle$ is a positive constant in $\Omega$, where $\nu = \frac{\nabla u}{|\nabla u|}$ \emph{(}in particular, $\langle\nabla u,X\rangle > 0$ in $\Omega$\emph{)}.
%	\end{itemize}
\end{theorem}

For operators of mean curvature type, that is, satisfying
\[
t^2 \lambda_2(t) \le C_1 \lambda_1(t) \le \frac{C_2}{t} \qquad \forall \, t \in \R^+
\]
for some positive constants $C_1,C_2$, condition \eqref{eq_lowenergy} can be replaced by 
\[
|\Omega \cap B_R| = o\big( R^2 \log R \big) \qquad \text{as } \, R \to \infty, 
\]
provided that any of assumptions $(i)$ to $(vi)$ in Theorem \ref{teo_goodcutoff_MC} are satisfied in $\Omega_0 = \Omega$.

\vspace{0.5cm}

The paper is structured as follows: in Section \ref{sec_energy} we introduce the moderate energy growth assumption needed for our main theorems, and obtain some sufficient conditions for its validity. Among them, we stress the bound on the growth of $|\Omega \cap B_R|$ for operators of mean curvature type in Theorem \ref{teo_goodcutoff_MC}. In Section \ref{sec_stability} we study the stability condition and its consequences for solutions with moderate energy growth. The main result is Theorem \ref{teo_monotonicity}, which guarantees the monotonicity of $u$ in the direction of a Killing field $X$ under mild conditions. The goal of Section \ref{sec_poinc} is to prove the key integral formula for monotone solutions to \eqref{equazu_conbordo}, see Theorem \ref{teo_Poincare_loc} and the related splitting results including Theorems \ref{teo_main_MC_manifolds} and \ref{teo_splitting_intro}. In Section \ref{sec_critical} we will focus on Euclidean space, analysing the set of critical points of $u$, and finally prove Theorems \ref{teo_CMC_intro_R2}, \ref{teo_CMC_intro_R3} and \ref{teo_CMC_intro_R2_gen}.
 
\vspace{0.5cm}

\noindent \textbf{Acknowledgements:} A.F. was partially supported by ANR EINSTEIN CONSTRAINTS:
past, present, and future. Research project ANR-23-CE40-0010-02. L.M. and M.R. were supported by the PRIN project no. 20225J97H5 ``Differential-geometric aspects of manifolds via Global Analysis''. 

\vspace{0.2cm}

\noindent \textbf{Conflict of Interests.} The authors have no conflict of interest.
%
\vspace{0.2cm}

\noindent \textbf{Data availability statement.} No data was generated by this research.


%\subsection{General setting and main results}
%
%Hereafter, we will consider a smooth Riemannian manifold $(M_o,g)$ of dimension $m \ge 2$, without boundary. For notational convenience, we will also denote the metric with $\metric$. We let $\Ric$ be its Ricci tensor. Having fixed an origin $o$, we set $r(x)= \mathrm{dist}(x,o)$, and we write $B_R$ for geodesic balls centered at $o$. Throughout the paper, with the symbol $\{\Omega_j\}\uparrow M_o$ we mean a family $\{\Omega_j\}$, $j \in \N$, of relatively compact, open sets with smooth boundary and satisfying
%$$
%\Omega_j \Subset \Omega_{j+1} \Subset M, \qquad M = \bigcup_{j=0}^{\infty} \Omega_j,
%$$
%where $A \Subset B$ means $\overline A \subseteq B$. Such a family will be called an exhaustion of $M$. \\[0.3cm]
%%
%%Let $a: \R^+_0 \ra \R$ be a function satisfying the structural conditions
%%\begin{equation}\label{assu_A}
%%\begin{array}{rcl}
%%(a1) & \quad a(t) \in C^1(\R^+) \cap C(\R^+_0). \\[0.1cm]
%%(a2) & \quad a(t)>0 \ \  \text{ for } t \in \R^+, \qquad \big(ta(t) \big)' >0 \ \ \text{ for } t \in \R^+,
%%\end{array}
%%\end{equation}
%%%
%%and let $f \in C^1(\R)$. Hereafter, we will consider weak solutions $u \in C^1(M)$ of the differential equation
%%%
%%Where $\Delta_a$ is the quasilinear operator defined as follows:
%%\begin{equation}\label{def_T}
%%\Delta_a u = \diver\big(a(|\nabla u|)\nabla u\big). 
%%\end{equation}
%%%
%%This class encompasses relevant examples such as 
%%\begin{itemize}
%%\item[-] the mean curvature operator, for which $a(t) = \frac{1}{\sqrt{1+t^2}}$, 
%%\item[-] the $p$-Laplacian $\Delta_p$, where $a(t) = t^{p-2}$,
%%\item[-] the operator of exponentially harmonic functions, for which $a(t)= e^{t^2}$, considered for instance in \cite{CITAZIONE}.
%%\end{itemize}
%
%Our aim is to prove the following splitting and half-space theorem:
%
%
%\begin{theorem}\label{teo_main}
%	Let $(M^n,\metric)$ be a complete Riemannian manifold, let $\Omega \subseteq M$ be a smooth open subset with $\partial \Omega$ connected and $\Ric \ge 0$, supporting a non-constant solution to 
%	\begin{equation}\label{overdetermined}
%		\left\{ \begin{array}{ll}
%			%
%			\Delta_a u + f(u) = 0 & \quad \text{on } \, \Omega,\\[0.2cm]
%			u > 0 & \quad \text{on } \, \Omega \\[0.2cm]
%			u= 0 & \quad \text{on } \, \partial \Omega \\[0.2cm]
%			\partial_{\eta} u = c > 0& \quad \text{on }  \, \partial \Omega, 
%		\end{array}\right.
%	\end{equation}
%	where $\eta$ is the interior normal. Assume that there exists a Killing field $X$ with $\sup_\Omega|X| < \infty$ and satisfying $\langle X, \eta \rangle \ge 0$ on $\partial \Omega$. If
%	\[
%	\int_{\Omega \cap B_R} |\nabla u|^2 a(|\nabla u|) = o(R^2 \log R) \qquad \text{as } \, R \to \infty
%	\]
%	Then:
%	\begin{itemize}
%		\item $\Omega = [0,\infty) \times N$ for some complete $N$ with $\Ric_N \ge 0$;
%		\item $u$ only depends on the coordinate on the $\R$-direction. 
%	\end{itemize}
%\end{theorem}


%\begin{theorem}\label{teo_global}
%Let $M$ be a parabolic manifold with $\Ric \ge 0$, and let $u$ be an entire solution of the minimal surface equation. Then, $M = N \times \R$ for some parabolic manifold $N$ with $\Ric \ge 0$, and $u=u(t)$ is an affine function. In particular, if $u$ is bounded from above or below on $M$, then $u$ is constant.
%\end{theorem}

%\begin{conjecture}
%\tcr{This is not important now} Let $M$ be a parabolic manifold with $\Ric \ge 0$, and let $u \ge 0$ be a solution of the minimal surface equation on an unbounded subset $\Omega$ with smooth compact boundary. If $u_{\partial \Omega} = 0$, then $u \equiv 0$.
%\end{conjecture}




%The existence of stable solutions of \eqref{equazu} is tightly related to a global property of the %operator $\Delta $ on $M$, which is called parabolicity. 
%In Appendix 2, we recall some basic %definitions and properties, referring to \cite{grigoryan} for a beautiful, detailed account. 
%
%-ç;
%


%\section{The one-dimensional profile*****REMOVABLE}
%
%We want to prove that if a manifold $\Omega$ splits isometrically as $N\times 
%\R_+$ for a manifold $N$ without boundary, then there exists a bounded positive 
%solution $u$ of the equation 
%\[
%\begin{cases}
%\mc u + f(u) = 0 & \text{on $\Omega$}\\
%u = 0 & \text{on $\partial \Omega$} \point
%\end{cases}
%\]
%
%It is sufficient to show that there exists a bounded positive solution of the 
%equation
%\[
%\begin{cases}
%\frac{d}{d t} (a(u')u') + f(u) = 0 & \text{on $\R_+$}\\
%u(0) = 0 \point
%\end{cases}
%\]
%
%Define the functions $F$ and $\phi$ as
%\[
%F(t) \eqdef \int_0^t f(s) d s
%\quad
%\text{and}\quad
%\phi(t) = - \int_0^t (a'(s)s + a(s)) s d s \comma
%\]
%then the function $\phi(u') - F(u)$ is invariant, indeed
%\[
%\frac{d}{dt}(\phi(u') - F(u)) = \phi'(u') u'' - F'(u) u' = -(a'(u')u' + 
%a(u'))u'u'' - f(u) u' = 0 \point
%\]
%Thus in particular it holds that
%\begin{equation} \label{eq:IntegralePrimo}
%\phi(u') - F(u) = \phi(u'(0)) \point
%\end{equation}
%
%
%Moreover notice that $\phi$ is a decreasing function since $a'(s)s + a(s) = 
%(a(s)s)' > 0$ thanks to assumption $(a2)$.
%
%Now choose $u'(0) = m>0$ to be such that
%\[
%\phi(m) = \phi(0) - F(\lambda) \comma
%\]
%which is possible since $\phi$ is decreasing and we have the estimate on 
%$F(\lambda)$. 
%
%\begin{remark}
%\emph{In the case of the mean curvature operator 
%	\[
%	\phi(t) = \frac{1}{\sqrt{1+t^2}}-1 \in (-1,0),  
%	\]	
%Suppose that $f \ge 0$ on $(0,\lambda)$, in particular $F$ is increasing. Then, the ODE problem has an entire solution only if $F$ satisfies the consistency condition
%	\[
%	F(\lambda) < 1.
%	\]
%In this case, $\phi(u'(0)) = 1 - F(\lambda)$.
%}
%\end{remark}
%
%We have existence for short times, because $a'(t)t + a(t)$ is always different 
%from $0$ (in particular is always strictly positive).
%Now let us call $[0,T)$ the maximal interval of definition of $u$ with the 
%chosen initial conditions.
%
%We first notice that $u<\lambda$ in the whole interval of definition. Indeed if 
%$u(t) = \lambda$ for some $t$, then for \eqref{eq:IntegralePrimo} it 
%holds
%\[
%\phi(u'(t)) - F(\lambda) = \phi(u'(0)) = \phi(0) - F(\lambda) \implies 
%\phi(u'(t)) = \phi(0) 
%\]
%and thus $u'(t) = 0$. However $u$ can not touch $\lambda$ with derivative $0$ 
%because of uniqueness of solutions, since the constant $\lambda$ is still a 
%solution.
%
%Now observe that $u'$ is strictly positive in $[0,T)$ because $u'(0)>0$ and 
%$u'$ can not touch $0$ (with a reasoning completely analogous to the one for 
%proving that $u$ can not touch $\lambda$).
%
%Furthermore $u'$ is bounded from above thank to \eqref{eq:IntegralePrimo}, 
%indeed $\phi(u') \ge \phi(u') - F(u) \ge \phi(0) - F(u)$ and $\phi$ is 
%decreasing.
%Consequently we have existence of $u$ at all the times, $u$ must be strictly 
%increasing and bounded by $\lambda$. Finally $u$ must converge to $\lambda$, 
%because we have that $\lim_{t\to\infty}\phi(u'(t)) = 0$ and thus, always thanks 
%to \eqref{eq:IntegralePrimo} we must have that $\lim_{t\to\infty}F(u(t)) = 
%F(\lambda)$ and hence $\lim_{t\to\infty}u(t) = \lambda$.
%
%\begin{remark}
%Always using \eqref{eq:IntegralePrimo} is it possible to show that for $u'(0)< 
%m$ the solution $u$ is concave and touches $0$ in finite time, while for 
%$u'(0)>m$ the solution $u$ explodes to $\infty$ in finite time.
%\end{remark}



\section{Moderate energy growth}\label{sec_energy}


%\begin{lemma} \label{lem_phij0}
%	Let $(M^n,\metric)$ be a complete Riemannian manifold and $0\leq f\in L^1_\loc(M)$. If
%	\[
%	\int_{B_R} f = O(R^2\log R) \qquad \text{as } \, R \to +\infty
%	\]
%	then there exists a sequence $\{\varphi_j\} \subseteq \lip_c(M)$ satisfying
%	\[
%	\varphi_j \to 1 \quad \text{in} \quad W^{1,\infty}_\loc(M) , \qquad \int_M f |\nabla\varphi_j|^2 \to 0
%	\]
%	as $j\to\infty$.
%\end{lemma}
%
%\begin{proof}
%	Let $\{R_j\}$ be a sequence of real numbers satisfying $R_j > 1$ for each $j\in\N$ and $R_j \to \infty$ as $j\to\infty$. For each $j\in\N$ set
%	\[
%	\varphi_j(r) = \left\{
%	\begin{array}{l@{\quad}l}
%		1 & \text{if } \, r \leq \exp(\sqrt{\log R_j}) \\
%		2\left(1-\dfrac{\log(\log r)}{\log(\log R_j)}\right) & \text{if } \, \exp(\sqrt{\log R_j}) < r \leq R_j \\
%		0 & \text{if } \, r > R_j
%	\end{array}
%	\right.
%	\]
%	Let $H : \R^+ \to \R^+_0$ be defined by
%	\[
%	H(r) = \int_{B_r} f \, .
%	\]
%	Then, for each $j\in\N$ we have
%	\begin{align*}
%		\int_M f|\nabla\varphi_j|^2 & = \int_{e^{\sqrt{\log(R_j)}}}^{R_j} \left( \int_{\partial B_r} f|\nabla\varphi_j|^2 \right) \di r \\
%		& = \frac{4}{\log^2(\log R_j)} \int_{e^{\sqrt{\log(R_j)}}}^{R_j} \frac{1}{r^2\log^2 r} \left(\int_{\partial B_r} f\right) \di r \\
%		& = \frac{4}{\log^2(\log R_j)} \int_{e^{\sqrt{\log(R_j)}}}^{R_j} \frac{H'(r)}{r^2\log^2 r} \, \di r \\
%		& = \frac{4}{\log^2(\log R_j)} \left[ \left.\frac{H(r)}{r^2\log^2 r}\right|_{r=e^{\sqrt{\log(R_j)}}}^{r=R_j} + 2 \int_{e^{\sqrt{\log(R_j)}}}^{R_j} H(r)\frac{\log^2(r)+1}{r^3\log^4 r} \, \di r \right] .
%	\end{align*}
%	Using that $0 \leq H(r) \leq C r^2 \log r$ for all sufficiently large $r>0$, we have
%	\[
%	\left.\frac{H(r)}{r^2\log^2 r}\right|_{r=e^{\sqrt{\log(R_j)}}}^{r=R_j} \leq \frac{H(R_j)}{R_j^2 \log^2(R_j)} \leq \frac{C}{\log R_j}
%	\]
%	and
%	\begin{align*}
%		\int_{e^{\sqrt{\log(R_j)}}}^{R_j} H(r)\frac{\log^2(r)+1}{r^3\log^4 r} \, \di r & \leq \int_{e^{\sqrt{\log(R_j)}}}^{R_j} \frac{2H(r)\log^2 r}{r^3\log^4 r} \, \di r \\
%		& \leq \int_{e^{\sqrt{\log(R_j)}}}^{R_j} \frac{2C}{r\log r} \, \di r \\
%		& = 2C\log(\log r) \Big|_{r=e^{\sqrt{\log(R_j)}}}^{r=R_j} \\
%		& = C\log(\log R_j)
%	\end{align*}
%	for all sufficiently large $j\in\N$. This shows that
%	\[
%	\int_M f|\nabla\varphi_j|^2 \leq \frac{4C}{\log(\log R_j)} + \frac{4C}{\log^2(\log R_j)\log(R_j)}
%	\]
%	as $j\to\infty$, concluding the proof.
%\end{proof}

Let $\Omega \subseteq M$ be an open subset with $C^1$ boundary, and $u \in C^1(\overline\Omega)$ be a weak solution to 
\[
\Delta_a u + f(u) = 0,
\]
where $a$ satisfies \eqref{assu_A} and $f \in C(\R)$. 
One of the core steps to achieve our classification results is to prove that $u$ has moderate energy growth, according to the next

\begin{definition}
	Let $\Omega_0\subseteq\Omega$ be an open set. We say that $u$ has \emph{moderate energy growth in $\Omega_0$} if there exists a sequence of functions $\{\varphi_j\} \subseteq \lip_c(\overline{\Omega}_0)$ such that, as $j \to \infty$, 
	\[
		0 \le \varphi_j \to 1 \quad \text{in } \, C_\loc(\overline{\Omega}_0), \qquad \int_{\Omega_0} |\nabla u|^2 \langle A(\nabla u) \nabla \varphi_j, \nabla \varphi_j \rangle \to 0 .
	\]
	If the condition is satisfied for $\Omega_0 = \Omega$ we simply say that $u$ has moderate energy growth.
\end{definition}

In this section, we describe some sufficient conditions for this property to hold. Hereafter, we consider balls $B_R$ centered at a fixed origin $o \in M$ and we define $r(x) = {\rm dist}(x,o)$. We begin with the following extension of the well-known ``logarithmic cut-off trick'', which is related to an argument by Karp \cite{karp}.

\begin{lemma} \label{lem_phij}
	Let $(M^n,\metric)$ be a complete Riemannian manifold and $0\leq f\in L^1_\loc(M)$. If
	\[
	\int_{B_R} f \leq R^2 \, \theta(R) \qquad \text{for all } \, R > 0
	\]
	for some non-decreasing, $C^1$ function $\theta : [R_1,\infty) \to \R^+$ such that
	\[
	\int_{R_1}^\infty \frac{\di s}{s \theta(s)} = \infty
	\]
	then there exists a sequence $\{\varphi_j\} \subseteq \lip_c(M)$ satisfying
	\[
	\varphi_j \to 1 \quad \text{in} \quad W^{1,\infty}_\loc(M) , \qquad \int_M f |\nabla\varphi_j|^2 \to 0
	\]
	as $j\to\infty$.
\end{lemma}

\begin{remark}
	A typical example of function $\theta(R)$ satisfying the above assumptions, that we will hereafter use in the paper, is $\theta(R) = \log R$.	
\end{remark}

\begin{proof}
	We define $\Theta(s) = s\theta(s)$ and for each $r\geq 1$ we set
	\[
	\psi(t) = \int_{R_1}^t \frac{\di s}{\Theta(s)} \, .
	\]
	The function $\psi : [R_1,\infty) \to \R^+_0$ is $C^2$, strictly increasing, strictly positive on $(R_1,\infty)$ and such that $\psi(t) \to \infty$ as $t\to\infty$. Let $\{R_j\} \subseteq [R_1,\infty)$ be a diverging sequence. For each $j\in\N$ let $r_j > R_j$ be the real number such that $\psi(r_j) = 2\psi(R_j)$ and define
	\[
	\psi_j(t) = \begin{cases}
		1 & \text{if } \, 0 < t \leq R_j \\
		2\left(1-\dfrac{\psi(t)}{\psi(r_j)}\right) & \text{if } \, R_j < t \leq r_j \\
		0 & \text{if } \, t > r_j \, .
	\end{cases}
	\]
	Then for $R_j < t < r_j$ we have
	\[
	\psi_j'(t) = -\frac{2}{\psi(r_j)} \frac{1}{\Theta(t)} \, , \qquad \psi_j''(t) = \frac{2}{\psi(r_j)} \frac{\Theta'(t)}{\Theta(t)^2} \, .
	\]
	In particular, by our assumption on $\theta$ we have $\psi_j'\leq 0$ and $\psi_j''\geq 0$ on $(R_j,r_j)$. 
	%	\[
	%		\psi_j'(t) = \begin{cases}
		%			0 & \text{if } \, 0 < t < R_j \\
		%			-\dfrac{2}{\psi(r_j)} \dfrac{1}{\Theta(t)} & \text{if } \, R_j < t < r_j \\
		%			0 & \text{if } \, t > r_j \, .
		%		\end{cases}
	%	\]
	%	and
	%	\[
	%		\psi_j''(t) = \begin{cases}
		%			0 & \text{if } \, 0 < t < R_j \\
		%			\dfrac{2}{\psi(r_j)} \dfrac{G'(t)}{\Theta(t)^2} & \text{if } \, R_j < t < r_j \\
		%			0 & \text{if } \, t > r_j \, .
		%		\end{cases}
	%	\]
	For each $j\in\N$ define $\varphi_j = \psi_j\circ r \in \lip_c(M)$. 
	%where $r$ denotes the distance function in $(M,\metric)$ from a fixed origin $o\in M$. 
	%Each function $\varphi_j$ is Lipschitz continuous and compactly supported in $M$. 
	Since $\{\varphi_j = 1\} = \overline{B}_{R_j}$ and $R_j\to\infty$ as $j\to\infty$, we also have $\varphi_j \to 1$ in $W^{1,\infty}_\loc(M)$. We now prove that $f|\nabla\varphi_j|^2 \to 0$ in $L^1(M)$. The function $H : \R^+ \to \R^+_0$ defined by
	\[
	H(t) = \int_{B_t} f
	\]
	is absolutely continuous on compact intervals contained in $\R^+$. By the co-area formula
	\begin{align*}
		\int_M f|\nabla\varphi_j|^2 & = \int_0^\infty \left(\int_{\partial B_t} f|\nabla\varphi_j|^2 \right) \, \di t \\
		& = \int_0^\infty H'(t) \psi_j'(t)^2 \, \di t \\
		& = \int_{R_j}^{r_j} H'(t) \psi_j'(t)^2 \, \di t \, .
	\end{align*}
	Integrating by parts, using that $0 \leq H(t)\leq t \Theta(t)$, $\psi_j' \leq 0$ and $\psi_j'' \geq 0$ and also using the explicit expression for $\psi_j'$ we obtain
	\begin{align*}
		\int_{R_j}^{r_j} H'(t) \psi_j'(t)^2 \, \di t & = H(t) \psi_j'(t)^2 \Big|_{t=R_j}^{t=r_j} - 2 \int_{R_j}^{r_j} H(t) \psi_j'(t) \psi_j''(t) \, \di t \\
		& \leq H(r_j) \psi_j'(r_j)^2 - 2 \int_{R_j}^{r_j} H(t) \psi_j'(t) \psi_j''(t) \, \di t \\
		& \leq r_j \Theta(r_j) \psi_j'(r_j)^2 - 2 \int_{R_j}^{r_j} t \Theta(t) \psi_j'(t) \psi_j''(t) \, \di t \\
		& = \frac{4}{\psi(r_j)^2} \frac{r_j}{\Theta(r_j)} + \frac{4}{\psi(r_j)} \int_{R_j}^{r_j} t \psi_j''(t) \, \di t \, .
	\end{align*}
	Another integration by parts yields
	\begin{align*}
		\int_{R_j}^{r_j} t \psi_j''(t) \, \di t & = t\psi_j'(t) \Big|_{t=R_j}^{t=r_j} - \int_{R_j}^{r_j} \psi_j'(t) \, \di t \\
		& = -\frac{2}{\psi(r_j)}\frac{t}{\Theta(t)} \Big|_{t=R_j}^{t=r_j} - \psi_j(t) \Big|_{t=R_j}^{t=r_j} \\
		& = \frac{2}{\psi(r_j)} \left(\frac{1}{\theta(R_j)} - \frac{1}{\theta(r_j)}\right) + 1
	\end{align*}
	so we get
	\[
	\int_M f|\nabla\varphi_j|^2 \leq \frac{8}{\psi(r_j)^2}\frac{1}{\theta(R_j)} + \frac{4}{\psi(r_j)} \to 0 \qquad \text{as } \, j \to \infty
	\]
	as desired.
	%	So,
	%	\begin{align*}
		%		\int_{R_j}^{r_j} H'(s) \varphi'(s)^2 \, \di s & = H(r)\varphi'(r)^2 \Big|_{r=R_j}^{r=r_j} - 2\int_{R_j}^{r_j} H(s)\varphi'(s)\varphi''(s) \, \di s \\
		%		& = \frac{4}{\psi(r_j)^2} \frac{r^2}{H(r)} \Big|_{r=R_j}^{r=r_j} + \frac{4}{\psi(r_j)} \int_{R_j}^{r_j} s\varphi''(s) \, \di s \\
		%		& = \frac{4}{\psi(r_j)^2} \frac{r^2}{H(r)} \Big|_{r=R_j}^{r=r_j} + \frac{4}{\psi(r_j)} \left( r\varphi'(r) \Big|_{r=R_j}^{r=r_j} - \int_{R_j}^{r_j} \varphi'(s) \, \di s \right) \\
		%		& = \frac{4}{\psi(r_j)^2} \frac{r^2}{H(r)} \Big|_{r=R_j}^{r=r_j} - \frac{8}{\psi(r_j)^2} \frac{r^2}{H(r)} \Big|_{r=R_j}^{r=r_j} - \frac{4}{\psi(r_j)} \varphi(r) \Big|_{r=R_j}^{r=r_j} \\
		%		& = \frac{4}{\psi(r_j)^2} \left(\frac{R_j^2}{H(R_j)} - \frac{r_j^2}{H(r_j)^2} \right) + \frac{4}{\psi(r_j)}
		%	\end{align*}
\end{proof}

As a direct consequence, we have

\begin{corollary}\label{cor_lowenergy}
	If $a$ satisfies \eqref{assu_A}, $u$ has moderate energy growth in $\Omega_0$ whenever
	\[
	\int_{\Omega_0 \cap B_R} |\nabla u|^2\lambda_{\max}(|\nabla u|) = O(R^2 \log R) \qquad \text{as } \, R \to \infty. 
	\]
%	then there exists $\{\varphi_j\} \subseteq \lip_c(M)$ such that  
%	\[
%	\varphi_j \to 1 \quad \text{in } \, W^{1,\infty}_\loc(M), \qquad \int_\Omega |\nabla u|^2 \langle A(\nabla u) \nabla \varphi_j, \nabla \varphi_j \rangle \to 0.
%	\]
%	as $j \to \infty$.
\end{corollary}

\begin{proof}
	It is enough to apply Lemma \ref{lem_phij} with the choices 
	\[
	f = |\nabla u|^2\lambda_{\max}(|\nabla u|) \mathbf{1}_{\Omega_0}, \qquad \theta(R) = \log R 
	\]
	and observe that 
	\[
	0 \le |\nabla u|^2 \langle A(\nabla u)\nabla \varphi_j, \nabla \varphi_j \rangle \le |\nabla u|^2 \lambda_{\max}(|\nabla u|)|\nabla \varphi_j|^2.   
	\]
\end{proof}

We next examine operators of mean curvature type, i.e. those for which 
\begin{equation}\label{eq_meancurv_type}
t^2 \lambda_1(t) \le C_1 \lambda_2(t) \le \frac{C_2}{t} \qquad \text{on } \, \R^+
\end{equation}
for some constants $C_1,C_2>0$. In this case, under suitable assumptions on $f$ and $u$ we can guarantee that $u$ has moderate energy growth by controlling the measures 
\[
|\Omega \cap B_R| \qquad \text{and possibly} \qquad \mathscr{H}^{n-1}(\partial \Omega \cap B_R).
\]

We first observe the following

\begin{proposition}\label{prop_moderate_boundgradient}
	Assume that $a$ satisfies \eqref{assu_A} and is of mean curvature type. If a solution $u$ to \eqref{equazu} satisfies
	\[
	|\nabla u| \in L^\infty(\Omega),
	\]
	and $|\Omega \cap B_R| = o(R^2 \log R)$ as $R \to \infty$, then $u$ has moderate energy growth.
\end{proposition}

\begin{proof}
	The mean curvature type condition implies that the eigenvalues of $A$ can be bounded as follows: 
	\[
	\lambda_{\max}(t) \le \max_{t \in [0,1]}(\lambda_{\max}) + \frac{2C_2}{1+t} \le \frac{C_3}{1+t}
	\]
	whence
	\begin{equation}\label{eq_uppermean}
	\int_{\Omega \cap B_R} |\nabla u|^2 \lambda_{\max} (|\nabla u|) \le C_3 \int_{\Omega \cap B_R} \frac{|\nabla u|^2}{1+|\nabla u|} \le C_3 \|\nabla u\|_\infty |\Omega \cap B_R|,  
	\end{equation}
	and we conclude by Corollary \ref{cor_lowenergy}.
\end{proof}

For the mean curvature operator, condition $|\nabla u| \in L^\infty(\Omega)$ is satisfied under mild assumptions on $M$, $u$ and $f$ including $f'(u) \le 0$. Indeed, we have the following theorem which improves on \cite[Theorem 1.5]{lian_sicbaldi}. We postpone its proof to the appendix, see Theorem \ref{thm_prop_grad}. The case where $f(u)$ is constant was considered in \cite{cmmr,cmr}.

\begin{theorem} \label{prop_grad}
	Let $(M^n,\metric)$ be a complete Riemannian manifold and $\Omega\subseteq M$ a domain. Let $f\in C^1(\R)$ and let $u\in C^3(\Omega) \cap C^1(\overline{\Omega})$ be a solution to
	\[
		\MM[u] + f(u) = 0 \qquad \text{in } \, \Omega \, .
	\]
	Suppose that
	\[
		\sup_\Omega |f(u)| < \infty \, , \qquad f'(u) \leq 0 \, , \qquad \inf_\Omega u > -\infty
	\]
	and that
	\[
		\Ric \geq 0 \quad \text{in } \, \Omega , \qquad \Sec \geq - \kappa \quad \text{in } \, M
	\]
	for some $\kappa\in\R^+$. Then
	\[
		\sup_\Omega |\nabla u| \leq \max\left\{ \sqrt{n/2}, \, \sup_{\partial\Omega} |\nabla u|\right\} .
	\]
\end{theorem}


%\tcr{[The following example is removable, we can cut things short]}
%
%\begin{example}
%	More generally, arguing as for \eqref{eq_uppermean} one gets
%	\[
%	\int_{\Omega} |\nabla u|^2 \langle A(\nabla u) \nabla \varphi,\nabla \varphi \rangle \le C_3 \int_{\Omega} \frac{|\nabla u|^2}{1+|\nabla u|}|\nabla \varphi|^2
%	\]
%	The right-hand side can be interpreted in terms of the graph of $u$,  $\mathrm{graph}(u) \subseteq M\times\R$, as is apparent if $A$ is the mean curvature operator. If we equip $\mathrm{graph}(u)$ with the Riemannian metric induced from the ambient space, and $A$ is the mean curvature operator, then for each $\varphi\in\lip_c(\overline{\Omega})$ we have
%	\[
%		\int_{\Omega} |\nabla u|^2 \langle A(\nabla u)\nabla\varphi,\nabla\varphi\rangle \di x \leq \int_{\Omega} \|\nabla^g\varphi\|^2_g \, \di v_g
%	\]
%	where $v_g$ is the Riemannian volume measure induced by the graph metric, which we identify via the graph map $F : \Omega \to M\times\R$ with $(\Omega,g)$, $g = F^\ast(\metric+\di t^2)$. We deduce that $u$ has moderate energy growth whenever $\mathrm{graph}(u)$ is parabolic. In turn, a sufficient condition for parabolicity of a Riemannian manifold with boundary is a slow volume growth \tcr{CITAZIONE}.
%\end{example}

If we do not assume $f'(u) \le 0$ in $\Omega$, global gradient estimates are, to our knowledge, currently not available, not even for the mean curvature operator. However, if $f(u)$ satisfies suitable bounds in $\Omega$ and $u$ suitable conditions on $\partial_\star \Omega = \partial\Omega \cap \{\partial_\eta u \neq 0\}$ we can still guarantee that $u$ has moderate energy growth by controlling the measures 
\[
|\Omega \cap B_R| \qquad \text{and} \qquad \mathscr{H}^{n-1}(\partial_\star \Omega \cap B_R),
\]
without needing to control $|\nabla u|$. Moreover, in certain cases all requests on $\partial_\star \Omega \cap B_R$ can be dropped. The calibration argument below is inspired by \cite[p. 403]{gilbargtrudinger}, \cite[p. 24]{simon} and \cite{FSV}. We localize our estimates to subdomains $\Omega_0 \subseteq \Omega$ whose boundary satisfies
\[
\partial \Omega_0 \subseteq \partial \Omega \cup \critu,
\]
where
\[
\critu = \big\{ x \in \Omega \ : \ |\nabla u(x)| = 0 \big\}
\]
is the \emph{interior critical set of $u$}. 

Define
\[
	\rho = \sqrt{ r^2 + u^2}
\]
that is, $\rho$ is the restriction to the graph of $u$ of the distance from $(o,0)$ in $M \times \R$. Hereafter, positive constants will be denoted by $C,C_0,C_1$ and so forth.

%
%\begin{proposition}\label{prop_1_calib}   %[Lemma 9.1] 
%	Let $u \in C^2(\Omega) \cap C^1(\overline\Omega)$ be a weak solution to
%	\[
%		\Delta_a u + f(u) = 0 \qquad \text{in } \, \Omega
%	\]
%	where $0 \le a \in C(\R^+_0)$ satisfies $ta(t) \in L^\infty(\R^+) \cap \lip_\loc(\R^+_0)$ and $f \in C(\R)$. Let $\Omega_0\subseteq\Omega$ be an open set with locally finite perimeter in $\Omega$ such that $\partial\Omega_0 \subseteq \partial\Omega \cup \critu$.
%	%	\begin{equation}\label{eq_extra}
%		%	\begin{array}{l}
%			%	f(u) \ \ \text{ has a sign on } \, \Omega, \\[0.2cm] 
%			%	f(u) \partial_\eta u \le 0 \qquad \text{on } \, \partial \Omega,
%			%	\end{array}
%		%	\end{equation}
%	%	where $\eta$ is the interior normal on $\partial \Omega$. 
%	\begin{itemize}
%		\item [i)] Suppose that $f(u)$ does not change sign in $\Omega_0$. Then \begin{equation}\label{primointe}
%			\int_{\Omega_0 \cap\{\rho \le R\}} \aaa |\nabla u|^2 \le C' \Big( R \, \haus^{n-1}(\overline{\Omega}_0\cap \partial_\star \Omega \cap B_{2R}) + |\Omega_0 \cap B_{2R}| \Big)
%		\end{equation}
%	%
%		for some constant $C'>0$. Moreover, in case $u\equiv b$ on $\overline{\Omega}_0 \cap \partial_\star\Omega$ for some constant $b\in\R$, and
%	\[
%		b = \begin{cases}
%		\inf_{\Omega_0} u & \quad \text{if } \, f(u) \leq 0 \, \text{ in } \, \Omega \\[0.2cm]
%		\sup_{\Omega_0} u & \quad \text{if } \, f(u) \geq 0 \, \text{ in } \, \Omega
%	\end{cases}
%	%		u_\ast \quad \text{on } \overline{\Omega}_0 \cap \partial_\star\Omega & \qquad \text{if } \, f(u) \leq 0 \, \text{ in } \, \Omega, \\
%	%		u \equiv u^\ast \quad \text{on } \overline{\Omega}_0 \cap \partial_\star\Omega & \qquad \text{if } \, f(u) \geq 0 \, \text{ in } \, \Omega,
%	\]
%	then it holds 
%	\begin{equation}\label{primointe_better}
%		\int_{\Omega_0 \cap \{\rho \le R\}} \aaa |\nabla u|^2 \le C'\, |\Omega_0 \cap B_{4R}| \qquad \forall \, R \ge 2|b|
%	\end{equation}
%		\item [ii)] Suppose that $f(u)$ is bounded in $\Omega$ and let $b\in\R$. Then
%		\begin{equation}\label{secondointe}
%			\int_{\Omega_0 \cap\{\rho \le R\}} \aaa |\nabla u|^2 \le C'' \Big( R\haus^{n-1}(\overline{\Omega}_0\cap \partial^b_\star \Omega \cap B_{4R}) + R|\Omega_0 \cap B_{4R}| \Big)
%		\end{equation}
%		for some constant $C''>0$ and every $R \ge 2|b|$, where we set $\partial_\star^b\Omega = \{ x \in \partial_\star \Omega : u(x) \neq b\}$.
%	\end{itemize}
%	
%%	holds in any of the following cases:
%%	\begin{itemize}
%%		\item[(i)] $u$ is constant on $\overline{\Omega}_0\cap \partial_\star \Omega$ and;
%%		\item[(ii)] $\partial_\eta u \equiv 0$ on $\overline{\Omega}_0\cap \partial \Omega$;
%%		\item[(iii)] $f(u) \partial_\eta u \le 0$ on $\overline{\Omega}_0\cap\partial \Omega$. 
%%	\end{itemize}
%\end{proposition}

\begin{proposition} \label{prop_1_calib}
	Let $u \in C^2(\Omega) \cap C^1(\overline\Omega)$ be a weak solution to
	\[
		\Delta_a u + f(u) = 0 \qquad \text{in } \, \Omega,
	\]
	where $0 \le a \in C(\R^+_0)$ satisfies $ta(t) \in L^\infty(\R^+) \cap \lip_\loc(\R^+_0)$ and $f \in C(\R)$. Let $\Omega_0\subseteq\Omega$ be an open set with locally finite perimeter in $\Omega$ such that $\partial\Omega_0 \subseteq \partial\Omega \cup \critu$.
	\begin{itemize}
		\item [$(i)$] Suppose that
		\[
			u_{|\overline{\Omega}_0 \cap \partial_\star\Omega} \, \text{ is bounded,} \qquad -C_0 \leq f(u) \leq 0 \quad \text{in } \, \Omega_0, \qquad \inf_{\Omega_0} u > -\infty \, .
		\]
		and let $b\in\R$. Then there exists $C>0$ such that
		\begin{equation}
			\int_{\Omega_0\cap\{\rho\leq R\}} \aaa |\nabla u|^2 \leq C \left(|\Omega_0\cap B_{4R}| + \haus^{n-1}(\overline{\Omega}_0 \cap \partial^b_\star\Omega \cap B_{4R}) \right)
		\end{equation}
		for all sufficiently large $R$, where $\partial_\star^b\Omega=\{ x \in \partial_\star \Omega : u(x) \neq b\}$. %$R>\max\{-\inf_{\Omega_0}u,2b-2\inf_{\Omega_0}u\}$. 
		\item[$(ii)$] Suppose that
		\[
		u_{|\overline{\Omega}_0 \cap \partial_\star\Omega} \, \text{ is constant,} \qquad -C_0 \leq f(u) \leq 0 \quad \text{in } \, \Omega_0, \qquad \inf_{\Omega_0} u > -\infty \, ,
		\]
		then there exists $C>0$ such that
		\begin{equation}
			\int_{\Omega_0\cap\{\rho\leq R\}} \aaa |\nabla u|^2 \leq C |\Omega_0\cap B_{4R}|
		\end{equation}
		for all sufficiently large $R$. %$R>\max\{-\inf_{\Omega_0}u,2b-2\inf_{\Omega_0}u\}$, with $b$ the constant value of $u$ on $\overline{\Omega}_0 \cap \partial_\star\Omega$. 
		\item[$(iii)$] Suppose that 
		\[
		u_{|\overline{\Omega}_0 \cap \partial_\star\Omega} = b \, \text{ is constant,} \qquad f(u) \leq 0 \quad \text{in } \, \Omega_0, \qquad \inf_{\Omega_0} u = b ,
		\]
		then there exists $C>0$ such that
		\begin{equation}
			\int_{\Omega_0\cap\{\rho\leq R\}} \aaa |\nabla u|^2 \leq C |\Omega_0\cap B_{4R}|
		\end{equation}
		for all sufficiently large $R$.
		\item [$(iv)$] Suppose that
		\[
			u_{|\overline{\Omega}_0 \cap\partial_\star\Omega} \, \text{ is constant,} \qquad |f(u)| \leq C_0 \quad \text{in } \, \Omega_0 \, ,
		\]
		then there exists $C>0$ such that
		\begin{equation}
			\int_{\Omega_0\cap\{\rho\leq R\}} \aaa |\nabla u|^2 \leq (C + C_0 R) |\Omega_0\cap B_{8R}|
		\end{equation}
		for all sufficiently large $R$. %for all $R>$\tcr{determinare}.
		\item [$(v)$] Suppose that 
		\[
		u_{|\overline{\Omega}_0 \cap\partial_\star\Omega} \, \text{ is constant,}  \qquad u \in L^\infty(\Omega_0), 
		\]
		then there exists $C'>0$ such that
		\begin{equation}
			\int_{\Omega_0\cap\{\rho\leq R\}} \aaa |\nabla u|^2 \leq C' |\Omega_0\cap B_{8R}|
		\end{equation}
		for all sufficiently large $R$. %for all $R>$\tcr{determinare}.
		\item [$(vi)$] Suppose that
		\[
			f(u) \leq 0 \qquad \text{in } \, \Omega_0
		\]
		with no further assumptions on $u$. Then there exists $C>0$ such that
		\[
			\int_{\Omega_0\cap\{\rho\leq R\}} \aaa |\nabla u|^2 \leq C \left(|\Omega_0\cap B_{2R}| + R\haus^{n-1}(\overline{\Omega}_0 \cap \partial_\star\Omega \cap B_{2R}) \right)
		\]
		for all $R>1$.		
	\end{itemize}
\end{proposition}

\begin{remark}
	The constant $C$ depends on $C_0, \|a\|_{L^\infty(\R^+)}, \|t a\|_{L^\infty(\R^+)}$ and, according to the case, $b$ (in $(i)$), $\|u\|_{L^\infty(\overline{\Omega}_0 \cap \partial_\star \Omega)}$ (in $(i)-(v)$) and $\|u\|_{L^\infty(\Omega_0)}$ (in $(v)$). Likewise, the lower bounds for $R$ guaranteeing the inequalities in $(i)-(vi)$ depend on the aforementioned quantities. 
\end{remark}

\begin{proof}
	For fixed $R>0$, define $\psi\in\lip_c(M)$ by
	\[
	\psi(x) = \left\{
	\begin{array}{l@{\quad}l}
		1 & \text{if } \, r(x) \leq R \\[0.2cm]
		2 - \dfrac{r(x)}{R} & \text{if } \, R < r(x) < 2R \\[0.3cm]
		0 & \text{if } \, r(x) \geq 2R
	\end{array}
	\right.
	\]
	Let $b_0,b\in\R$ be given satisfying $0 \leq b - b_0 \leq R$ and define $\gamma\in\lip(\R)$ by
	\[
		\gamma(t) = \left\{
			\begin{array}{l@{\quad}l}
				\dfrac{b_0-b}{R} & \text{if } \, t \leq b_0 \\[0.3cm]
				\dfrac{t-b}{R} & \text{if } \, b_0 < t < b_0 + R \\[0.3cm]
				\dfrac{b_0-b}{R}+1 & \text{if } \, t \geq b_0 + R
			\end{array}
		\right.
	\]
	Note that by construction
	\[
		\gamma(b) = 0, \qquad 0\leq\gamma' = \frac{1}{R}\mathbf{1}_{(b_0,b_0+R)}, \qquad |\gamma| \leq 1 \quad \text{on } \, \R .
	\]
	Consider the vector field
	\[
		W = \psi\gamma(u)a(|\nabla u|)\nabla u
	\]
	which is continuous and compactly supported in $\overline{\Omega}$, locally Lipschitz\footnote{For $|X|,|Y|\leq R$ we have $|a(|X|)X-a(|Y|)Y| \leq C_R|X-Y|$, with $C_R = 2\|a\|_{L^\infty} + [ta]_{\lip([0,R])}$, since
			\begin{align*}
			a(|X|)X - a(|Y|)Y & = [a(|X|) - a(|Y|)]X + a(|Y|)(X-Y), \\
			[a(|X|) - a(|Y|)]|X| & = |X|a(|X|) - |Y|a(|Y|) + (|Y|-|X|)a(|Y|)
			\end{align*}
	so $a(|\nabla u|)\nabla u$ is locally Lipschitz in $\Omega$ as $u \in C^2(\Omega)$.} in $\Omega$ and satisfies
	\[
		\diver \, W = - \psi\gamma(u)f(u) + \gamma(u) a(|\nabla u|) \langle\nabla u,\nabla\psi\rangle + \psi \gamma'(u) a(|\nabla u|) |\nabla u|^2
	\]
	weakly in $\Omega$. We apply the divergence theorem (see Lemma \ref{lem_divergence} in the appendix) to the vector field $W$ on the domain $\Omega_0$, observing that $W = 0$ on $\{\nabla u = 0\}$ since $ta(t) \to 0$ as $t\to0$. We obtain
	\begin{align*}
		\int_{\Omega_0} \psi \gamma'(u) \aaa |\nabla u|^2 & =  
		\int_{\Omega_0} \psi \gamma(u)f(u) - \int_{\Omega_0} \gamma(u) \aaa \langle \nabla u, \nabla \psi \rangle \\
		& \phantom{=\;} - \int_{\overline{\Omega}_0\cap\partial \Omega} a(|\nabla u|) \gamma(u)\psi \partial_\eta u 
	\end{align*}
	where, we recall, $\eta$ is the inward normal. We next  estimate the various terms. First, by our choice of $\gamma$ and $\psi$ we have
	\[
		\int_{\Omega_0} \psi \gamma'(u) \aaa |\nabla u|^2 \geq \frac{1}{R} \int_{\Omega_0 \cap B_R \cap \{b_0 < u < b_0+R\}} \aaa |\nabla u|^2
	\]
	and, since $ta(t) \in L^\infty(\R^+)$,
	\[
		- \int_{\Omega_0} \gamma(u) \aaa \langle \nabla u, \nabla \psi \rangle \leq \frac{1}{R} \int_{\Omega_0 \cap B_{2R}} \aaa |\nabla u| \leq \frac{C}{R} |\Omega_0 \cap B_{2R}|
	\]
	Noting that $\gamma(u)\partial_\eta u = 0$ on $\partial\Omega \backslash \partial^b_\star\Omega$ we have
	\[
		- \int_{\overline{\Omega}_0\cap\partial \Omega} a(|\nabla u|) \gamma(u) \psi \partial_\eta u = - \int_{\overline{\Omega}_0\cap\partial^b_\star \Omega} a(|\nabla u|) \gamma(u) \psi \partial_\eta u \, .
	\]
	For notational convenience, let us set $F^b_\star = \overline{\Omega}_0\cap\partial^b_\star \Omega$. Then, by also using the expression of $\gamma$ we have
	\[
		- \int_{\overline{\Omega}_0\cap\partial \Omega} a(|\nabla u|) \gamma(u) \psi \partial_\eta u \leq C \cdot \min\left\{1,\dfrac{\sup_{F^b_\star \cap B_{2R}}|u-b|}{R}\right\} \cdot \haus^{n-1}(F^b_\star\cap B_{2R})
	\]
	Finally, splitting $f(u) = [f(u)]_+ - [f(u)]_-$ into its positive and negative part and using that $\gamma(u) < 0$ on $\{u<b\}$ and $\gamma(u)>0$ on $\{u>b\}$, we have
	\begin{align*}
		\int_{\Omega_0} \psi \gamma(u)f(u) & = \int_{\Omega_0 \cap \{u<b\}} \psi \gamma(u)f(u) + \int_{\Omega_0 \cap \{u>b\}} \psi \gamma(u)f(u) \\
		& \leq \int_{\Omega_0 \cap \{u<b\}} \psi |\gamma(u)| [f(u)]_- + \int_{\Omega_0 \cap \{u>b\}} \psi \gamma(u) [f(u)]_+ .
	\end{align*}
	By using the expression of $\gamma$ we further estimate
	\[
		\int_{\Omega_0 \cap \{u>b\}} \psi \gamma(u) [f(u)]_+ \leq \min\left\{\frac{b_0-b+R}{R},\frac{\sup_{\Omega_0\cap B_{2R}} (u-b)_+}{R} \right\} \cdot \sup_{\Omega_0\cap B_{2R}} [f(u)]_+ \cdot |\Omega_0 \cap B_{2R}|
	\]
	and similarly
	\[
		\int_{\Omega_0 \cap \{u<b\}} \psi |\gamma(u)| [f(u)]_- \leq \min\left\{\frac{b-b_0}{R},\frac{\sup_{\Omega_0\cap B_{2R}} (u-b)_-}{R} \right\} \cdot \sup_{\Omega_0\cap B_{2R}} [f(u)]_- \cdot |\Omega_0 \cap B_{2R}|
	\]
	Putting everything together, we obtain
	\begin{align*}
		\int_{\Omega_0 \cap B_R \cap \{b_0 < u < b_0+R\}} \aaa |\nabla u|^2 & \leq C |\Omega_0 \cap B_{2R}| \\
		& \phantom{=\;} + C \cdot \min\left\{R,\sup_{F^b_\star\cap B_{2R}}|u-b|\right\} \cdot \haus^{n-1}(F^b_\star\cap B_{2R}) \\
		& \phantom{=\;} + \min\left\{b-b_0+R,\sup_{\Omega_0\cap B_{2R}} (u-b)_+ \right\} \cdot \sup_{\Omega_0\cap B_{2R}} [f(u)]_+ \cdot |\Omega_0 \cap B_{2R}| \\
		& \phantom{=\;} + \min\left\{b-b_0,\sup_{\Omega_0\cap B_{2R}} (u-b)_- \right\} \cdot \sup_{\Omega_0\cap B_{2R}} [f(u)]_- \cdot |\Omega_0 \cap B_{2R}|
	\end{align*}
	
	Now, we consider several cases. \\[0.2cm]
	\textbf{Cases $(i)$ and $(ii)$.} Suppose that the assumptions in $(i)$ hold:
	\[
		u_{|\overline{\Omega}_0 \cap \partial_\star\Omega} \, \text{ is bounded,} \qquad -C_0 \leq f(u) \leq 0 \quad \text{in } \, \Omega_0, \qquad \inf_{\Omega_0} u > -\infty \, .
	\]
	Then we can choose $b_0 = \inf_{\Omega_0} u$ and fixing any $b\geq b_0$ we obtain
	\begin{align*}
		\int_{\Omega_0 \cap B_R \cap \{u < \inf_{\Omega_0} u+R\}} \aaa |\nabla u|^2 & \leq (C + C_0(b-\inf_{\Omega_0} u)) \cdot |\Omega_0 \cap B_{2R}| \\
		& \phantom{=\;} + C \haus^{n-1}(F^b_\star \cap B_{2R}) \qquad \qquad \text{for all } \, R > b-\inf_{\Omega_0} u
	\end{align*}
	where we use that $[f(u)]_+ \equiv 0$, $[f(u)]_- \leq C_0$ and $(u-b)_- \leq b-b_0$. 
	
	If we further assume that
	\[
		u \text{ is constant on } \overline{\Omega}_0 \cap \partial_\star\Omega
	\]
	i.e. that the assumptions in $(ii)$ are met, then choosing $b$ as the constant value of $u$ on that set we have $F^b_\ast = \emptyset$ and
	\[
		\int_{\Omega_0 \cap B_R \cap \{u < \inf_{\Omega_0} u+R\}} \aaa |\nabla u|^2 \leq (C + C_0(b-\inf_{\Omega_0} u)) \cdot |\Omega_0 \cap B_{2R}| \qquad \text{for all } \, R > b-\inf_{\Omega_0} u
	\]
	Observe that if $\inf_{\Omega_0}u > -R/2$, then
	\[
		\Omega_0 \cap \{ \rho<R/2 \} \subseteq \Omega_0 \cap B_{R/2} \cap \{u<R/2\} \subseteq \Omega_0 \cap B_R \cap \{u<\inf_{\Omega_0} u + R\}
	\]
	hence
	\[
		\int_{\Omega_0 \cap \{\rho\leq R/2\}} \aaa |\nabla u|^2 \leq \int_{\Omega_0 \cap B_R \cap \{u < \inf_{\Omega_0} u+R\}} \aaa |\nabla u|^2
	\]
	and we therefore obtain the desired conclusions in $(i)$ and $(ii)$ up to replacing $R$ by $2R$. \\[0.2cm]
	\textbf{Case $(iii)$.} Suppose that
	\[
		u_{|\overline{\Omega}_0 \cap\partial_\star\Omega} \, \text{ is constant,} \qquad f(u) \leq 0 \quad \text{in } \, \Omega_0, \qquad \text{and} \qquad u_{|\overline{\Omega}_0 \cap\partial_\star\Omega} \equiv \inf_{\Omega_0} u > -\infty \, .
	\]
	Then we can choose $b = b_0 = \inf_{\Omega_0} u$ to obtain
	\[
		\int_{\Omega_0 \cap B_R \cap \{u < \inf_{\Omega_0} u+R\}} \aaa |\nabla u|^2 \leq C |\Omega_0 \cap B_{2R}| \qquad \text{for all } \, R>0 \, 
	\]
	and we conclude as in the previous step.\\[0.2cm]
	\textbf{Cases $(iv)$ and $(v)$.} Suppose that the assumptions in $(iv)$ are met:
	\[
		u_{|\overline{\Omega}_0 \cap \partial_\star\Omega} \, \text{ is constant,} \qquad |f(u)| \leq C_0 \quad \text{in } \, \Omega_0 \, .
	\]
	Then we can choose $b$ the constant value of $u$ on $\overline{\Omega}_0 \cap \partial_\star\Omega$ and $b_0 = b - R/2$ to obtain
	\[
		\int_{\Omega_0 \cap B_R \cap\left\{b - \frac{R}{2} < u < b + \frac{R}{2}\right\}} \aaa |\nabla u|^2 \leq \left( C + 3 C_0 \min\left\{ \frac{R}{2}, \sup_{\Omega_0\cap B_{2R}} |u-b| \right\}\right) |\Omega_0 \cap B_{2R}| .
	\]
	In particular, for some $C'>0$ we have
	\[
		\int_{\Omega_0 \cap B_R \cap \left\{b - \frac{R}{2} < u < b + \frac{R}{2}\right\}} \aaa |\nabla u|^2 \leq C' R |\Omega_0 \cap B_{2R}| \qquad \text{for all } \, R\geq 1.
	\]
	If we assume the requests in $(v)$, condition $|f(u)| \le C_0$ is then automatic for suitable $C_0$, and for some $C''>0$ we have
	\[
		\int_{\Omega_0 \cap B_R \cap \left\{b - \frac{R}{2} < u < b + \frac{R}{2}\right\}} \aaa |\nabla u|^2 \leq C'' |\Omega_0 \cap B_{2R}| \qquad \text{for all } \, R\ge 1.
	\]
	The conclusions in $(iv)$ and $(v)$ follow by observing that 
	\[
	\{\rho \le R/4\} \subset B_R \cap \{|u| \le R/4\} \subset B_R \cap \left\{b - \frac{R}{2} < u < b + \frac{R}{2} \right\}
	\]
	for large enough $R$. \\[0.2cm]
	\textbf{Case $(vi)$.} Suppose that
	\[
		f(u) \leq 0 \qquad \text{in } \, \Omega_0
	\]
	with no further assumptions on $u$. Then choosing $b = b_0 = 0$ we get
	\begin{align*}
		\int_{\Omega_0 \cap B_R \cap \{0 < u < R\}} \aaa |\nabla u|^2 & \leq C|\Omega_0\cap B_{2R}| \\
		& \phantom{=\;} + C\cdot\min\{R,\sup_{F^0_\star \cap B_{2R}} |u|\} \cdot \haus^{n-1}(F^0_\star \cap B_{2R})
	\end{align*}
	while choosing $b = b_0 = -R$ we have
	\begin{align*}
		\int_{\Omega_0 \cap B_R \cap \{-R < u < 0\}} \aaa |\nabla u|^2 & \leq C|\Omega_0\cap B_{2R}| \\
		& \phantom{=\;} + C\cdot\min\{R,\sup_{F^{-R}_\star \cap B_{2R}} |u+R|\} \cdot \haus^{n-1}(F^{-R}_\star \cap B_{2R}) \, .
	\end{align*}
	Since $F^0_\star$ and $F^{-R}_\star$ are contained in $\overline{\Omega}_0 \cap \partial_\star\Omega$, and using that $\nabla u = 0$ almost everywhere on $\{u=0\}$, we get
	\begin{align*}
		\int_{\Omega_0 \cap B_R \cap \{|u| < R\}} \aaa |\nabla u|^2 & \leq C\left(|\Omega_0\cap B_{2R}| + R \cdot \haus^{n-1}(\overline{\Omega}_0 \cap \partial_\star\Omega \cap B_{2R})\right)
	\end{align*}
	for all $R\geq 1$. The desired conclusion follows. 
	\end{proof}
	
	
	%
	%
%	DIMOSTRAZIONE PRECEDENTE
%	
%	\tcr{[Dimostrazione precedente]}
%	%
%	Define $\psi\in\lip_c(M)$ by
%	\[
%		\psi(x) = \left\{
%			\begin{array}{l@{\quad}l}
%				1 & \text{if } \, r(x) \leq R \\[0.2cm]
%				2 - \dfrac{r(x)}{R} & \text{if } \, R < r(x) < 2R \\[0.3cm]
%				0 & \text{if } \, r(x) \geq 2R
%			\end{array}
%		\right.
%	%, \qquad \xi(x) = \min \left\{ \frac{r_{\partial \Omega}(x)}{T},1\right\}.  
%	\]
%	and $\gamma_1\in\lip(\R)$ by %$0 \le \gamma \in \lip(\R)$ be linear in $[0,R]$ and satisfying:
%	%		\[
%	%	\gamma \equiv 0 \quad \text{in } (-\infty, -R], \qquad \gamma=1 \qquad  \text{in } [R,\infty).
%	%	\]
%	\[
%		\gamma_1(t) = \left\{
%		\begin{array}{l@{\quad}l}
%			0 & \text{if } \, t \leq 0 \\[0.2cm]
%			t/R & \text{if } \, 0 < t < R \\[0.2cm]
%			1 & \text{if } \, t \geq R
%		\end{array}
%		\right.
%	\]
%	Now, choose $\gamma\in\lip(\R)$ as follows:
%	\[
%		\begin{cases}
%			\gamma(t) = \gamma_1(t) & \qquad \text{if } \, f(u) \leq 0 \text{ in } \Omega \\
%			\gamma(t) = \gamma_1(t)-1 & \qquad \text{if } \, f(u) \geq 0 \text{ in } \Omega \, .
%		\end{cases}
%	\]
%%	\[
%%	\begin{array}{lllll}
%%		\gamma(t) = \big(\min\left\{ t/R, 1\right\}\big)_+ & \quad \text{if } f(u) \le 0 \ \ \text{in } \, \Omega \\[0.2cm]
%%		\gamma(t) = \big(\min\left\{ t/R, 1\right\}\big)_+ -1 & \quad \text{if } f(u) \ge 0 \ \ \text{in } \, \Omega.
%%	\end{array}
%%	\]
%	Note that in both cases we have $\gamma(u)f(u) \leq 0$ in $\Omega$ and $0 \leq \gamma' = R^{-1}\mathbf{1}_{(0,R)}$.
%%	\begin{equation} \label{gamma_prop}
%%		\gamma(u)f(u)\leq 0 \quad \text{in } \Omega, \qquad \gamma(b) = 0, \qquad \gamma(u)\partial_\eta u = 0 \quad \text{in } \partial\Omega\backslash\partial^b_\star\Omega.
%%	\end{equation}
%	%observe that
%	%	\[
%	%	\begin{array}{ll}
%		%		\gamma(u)\partial_\eta u \ge 0 & \quad \text{on } \, \partial \Omega, \\[0.2cm]
%		%		\gamma(u)f(u) \le 0 & \quad \text{on } \, \Omega. 
%		%	\end{array}
%	%	\]
%	We consider the vector field
%	\[
%		W = \psi\gamma(u)a(|\nabla u|)\nabla u
%	\]
%	which is continuous and compactly supported in $\overline{\Omega}$, locally Lipschitz in $\Omega$ and satisfies
%	\[
%		\diver \, W = - \psi\gamma(u)f(u) + \gamma(u) a(|\nabla u|) \langle\nabla u,\nabla\psi\rangle + \psi \gamma'(u) a(|\nabla u|) |\nabla u|^2
%	\]
%	weakly in $\Omega$. 
%%	Since $\psi\geq 0$, $\gamma(u)f(u) \leq 0$ on $\Omega$ and
%%	\[
%%		\gamma(u) \langle\nabla u,\nabla\psi\rangle \leq |\gamma(u)||\nabla u||\nabla\psi| \leq \frac{|\nabla u|}{R} \mathbf{1}_{B_{2R}}
%%	\]
%%	on $\Omega$, we infer
%%	\[
%%		\diver\, V \geq - \frac{1}{R} a(|\nabla u|)|\nabla u| \mathbf{1}_{B_{2R}} + \psi \gamma'(u) a(|\nabla u|) |\nabla u|^2 \, .
%%	\]
%	We apply the divergence theorem (see Lemma \ref{lem_divergence} in the appendix) to the vector field $W$ on the domain $\Omega_0$, observing that $W = 0$ on $\{\nabla u = 0\}$ since $ta(t) \to 0$ as $t\to0$. Using that $\gamma(u)f(u) \leq 0$ we obtain
%%	\[
%%		\int_{\Omega_0} \psi \gamma'(u) a(|\nabla u|) |\nabla u|^2 \leq \int_{\overline{\Omega}_0\cap\partial\Omega} \langle V,\eta\rangle + \frac{1}{R} \int_{\Omega_0\cap B_{2R}} a(|\nabla u|)|\nabla u| \, .
%%	\]
%%	By our choice of $\gamma$ and $\psi$, we have
%%	\[
%%		\frac{1}{R} \int_{\Omega_0 \cap \{\rho\leq R\}\cap\{u\geq 0\}} a(|\nabla u|)|\nabla u|^2 \leq \int_{\Omega_0} \psi \gamma'(u) a(|\nabla u|) |\nabla u|^2
%%	\]
%%	and also $\langle V,\eta\rangle \leq |V| \leq a(|\nabla u|)|\nabla u|$ on $\partial\Omega$, hence
%	%
%	%We integrate $\diver(\aaa \nabla u)= - f(u)$ against the test function $\varphi = \gamma(u) \psi$ and use that $\gamma(u)f(u) \le 0$ on $\Omega$ to deduce 
%	%
%	%
%	%	
%	%	vanishes on $\partial \Omega$, hence it is an admissible test function in the weak definition of $\diver(\aaa \nabla u)= - f(u)$. Integrating by parts, we get
%	%
%	%	$$
%	%	\begin{array}{lcl}
%		%	\disp \int_\Omega \gamma'\psi \aaa |\nabla u|^2 & = & \disp 
%		%	-\int_{\partial \Omega} a(|\nabla u|) \gamma(u)\psi \partial_\eta u + 
%		%	\int_\Omega \psi \gamma(u)f(u) - \int_\Omega \gamma \aaa \langle \nabla u, \nabla \psi \rangle \\[0.5cm]
%		%	& \le & \disp \frac{C_4}{r} \int_{\Omega \cap C_{r,2r}} \aaa |\nabla u| \le \frac{C_5}{r}|\Omega \cap C_{r,2r}|.
%		%	\end{array}
%	%	$$
%	\[
%	\begin{array}{lcl}
%		\disp \int_{\Omega_0} \gamma'(u) \psi \aaa |\nabla u|^2 & \leq & \disp  
%		-\int_{\overline{\Omega}_0\cap\partial \Omega} a(|\nabla u|) \gamma(u)\psi \partial_\eta u - \int_{\Omega_0} \gamma(u) \aaa \langle \nabla u, \nabla \psi \rangle \\[0.5cm]
%		& \le & \disp -\int_{\overline{\Omega}_0\cap\partial_\star \Omega} a(|\nabla u|) \gamma(u)\psi \partial_\eta u + \frac{1}{R} \int_{\Omega_0 \cap B_{2R}} \aaa |\nabla u| \\[0.5cm]
%		& \le & \disp -\int_{\overline{\Omega}_0\cap\partial_\star \Omega} a(|\nabla u|) \gamma(u)\psi \partial_\eta u + C \frac{|\Omega_0 \cap B_{2R}|}{R}.
%	\end{array}
%	\]
%	By our choice of $\gamma$ and $\psi$,
%	\[
%		\int_{\Omega_0} \gamma'(u) \psi \aaa |\nabla u|^2 \geq \frac{1}{R} \int_{\Omega_0\cap\{\rho \le R\} \cap \{u \ge 0\}} \aaa |\nabla u|^2, 
%	\]
%	whence we deduce
%	\begin{equation}\label{eq_basic}
%		\int_{\Omega_0\cap\{\rho \le R\} \cap \{u \ge 0\}} \aaa |\nabla u|^2 \le - R\int_{\overline{\Omega}_0\cap\partial_\star \Omega} a(|\nabla u|) \gamma(u)\psi \partial_\eta u + C \, |\Omega \cap B_{2R}|
%	\end{equation}
%	From $|\partial_\eta u| \le |\nabla u|$, we can estimate the boundary term as follows:
%	\[
%		\left|-\int_{\overline{\Omega}_0\cap\partial_\star \Omega} a(|\nabla u|) \gamma(u)\psi \partial_\eta u \right| \le C \haus^{n-1}(\overline{\Omega}_0\cap\partial_\star \Omega \cap B_{2R}),
%	\]
%	whence 
%	\[
%		\int_{\Omega_0\cap\{\rho \le R\} \cap \{u\ge 0\}} \aaa |\nabla u|^2 \le C' \big(R \haus^{n-1}(\overline{\Omega}_0\cap\partial_\star \Omega \cap B_{2R}) + |\Omega_0 \cap B_{2R}| \big).
%	\]
%	To obtain an analogous inequality in $\Omega_0\cap \{\rho \le R\} \cap \{u\le 0\}$, and thus \eqref{primointe}, simply replace $u$ by $\bar u = -u$ and $f$ by $\bar f(t) = -f(-t)$.
%	
%	Suppose that $u \equiv b = \inf_{\Omega_0} u$ on $\overline{\Omega}_0\cap\partial_\star\Omega$ and that $f(u)\leq 0$ in $\Omega_0$. We can repeat the argument at the beginning of the proof with $u$ replaced by $\bar u = u-b$ up to inequality \eqref{eq_basic}, that is, 
%	\[
%		\int_{\Omega_0 \cap \{\bar\rho \le R\} \cap \{\bar u \ge 0\}} \aaa |\nabla u|^2 \le - R\int_{\overline{\Omega}_0\cap \partial \Omega} a(|\nabla u|) \gamma(\bar u)\psi \partial_\eta u + C \, |\Omega_0 \cap B_{2R}|
%	\]
%	where $\bar\rho = \sqrt{r^2 + \bar u^2}$. Since $\gamma(\bar u) = \gamma(0) = 0$ on $\overline{\Omega}_0 \cap \partial_\star \Omega$ and $\bar u \geq 0$ in $\Omega_0$, we obtain
%	\[
%		\int_{\Omega_0 \cap \{\bar\rho \le R\}} \aaa |\nabla u|^2 \le C \, |\Omega_0 \cap B_{2R}|.
%	\]
%	Taking into account that for $R \ge |b|$ the set $\{\rho \le R/2\}$ is enclosed in $\{\bar \rho \le R\}$, up to replacing $R$ with $2R$ we obtain  \eqref{primointe_better} . The case where $u \equiv \sup_{\Omega_0} u$ on $\overline{\Omega}_0\cap\partial_\star\Omega$ and $f(u)\geq 0$ in $\Omega_0$ can be deduced from the previous one by replacing $u$ with $-u$ and $f$ with $\bar f(t) = -f(-t)$.
%	
%%	We next show that \eqref{primointe_better} also holds in case $u$ is constant on $\overline{\Omega}_0 \cap \partial_\ast\Omega$ and satisfies \eqref{uastast}. If $f(u)\leq 0$, without loss of generality we can assume that $u_\star = 0$ by replacing $u$ by $\bar -u$
%%	
%%	We next show that $(i),(ii),(iii)$ allow to discard the boundary term in \eqref{eq_basic}. This is enough to conclude \eqref{primointe_better}, once we observe that any of $(i),(ii),(iii)$ is preserved by the substitution $u \mapsto \bar u$ and $f \mapsto \bar f$. Case $(ii)$ is obvious. If $(iii)$ holds, then by construction $\gamma(u) \partial_\eta u \ge 0$ on $\partial \Omega$ and again we can get rid of the boundary term \tcr{[Attenzione: nei punti di $\partial\Omega$ in cui $f(u)=0$, la condizione $f(u)\partial_\eta u \leq 0$ non dà vincoli sul segno di $\partial_\eta u$]}. Eventually, assume that $(i)$ holds and let $b$ be the constant value of $u$ on $\tco{\overline{\Omega}_0\cap\partial_\star\Omega}$. We can repeat the argument at the beginning of the proof with $u$ replaced by $\bar u = u-b$ up to inequality \eqref{eq_basic}, that is, 
%%	\[
%%		\int_{\tco{\Omega_0 \cap }\{\rho \le R\} \cap \{\bar u \ge 0\}} \aaa |\nabla u|^2 \le - R\int_{\tco{\overline{\Omega}_0\cap}\partial \Omega} a(|\nabla u|) \gamma(\bar u)\psi \partial_\eta u + C \, |\tco{\Omega_0} \cap B_{2R}|.
%%	\]
%%	Since $\gamma(\bar u) = 0$ \tco{on $\overline{\Omega}_0 \cap \partial_\star \Omega$ and $\partial_\eta u = 0$ everywhere else on $\overline{\Omega}_0 \cap \partial\Omega$}, we obtain \eqref{primointe_better}.
%	
%	\tcr{Proof of case ii)}
%\end{proof}

%\begin{remark}\label{rem_very_useful}
%	\tcr{[TO REMOVE]} Inequality \eqref{primointe} can be improved by restricting the request on $\mathscr H^{n-1}(\partial \Omega \cap B_R)$ to some components of $\partial \Omega$. We reason as follows: assume that $u = b_j$ and $\partial_\eta u = c_j$ on $\partial_j \Omega$, for some $b_j,c_j \in \R$. Fix a component where $c_j \neq 0$, say $\partial_1\Omega$ (we assume its existence, otherwise the stronger \eqref{primointe_better} holds by $(ii)$), and define  
%	\[
%	\partial^\dagger \Omega = \bigcup_{j \, : \, b_j \neq b_1, \, c_j \neq 0} \partial_j \Omega.
%	\]
%	Then, replacing $u$ by $\bar u = u-b_1$ and proceeding as in the proof of \eqref{primointe_better} for case $(i)$, the only possibly nonzero contribution to the boundary term of \eqref{eq_basic} occurs in  $\partial^\dagger \Omega$. Hence, \eqref{primointe} improves to
%	\begin{equation}\label{primointe_slightbetter}
%		\int_{\{\rho \le R\}} \aaa |\nabla u|^2 \le C \Big( R \, \haus^{n-1}(\partial^\dagger\Omega \cap B_{2R}) + |\Omega \cap B_{2R}| 
%		\Big).
%	\end{equation}
%\end{remark}



\begin{proposition}\label{prop_inteannuli}
	Let $0 \le a \in C(\R^+_0)$ and $u$ satisfy
	\begin{equation}\label{primointe_abstract}
		\int_{\Omega_0 \cap \{\rho \le R\}} \aaa |\nabla u|^2 \le V(R)
	\end{equation}
	for some $0 \le V\in C(\R^+)$. Then, there exists $C_3>1$ such that, for each $R>0$, 
	\[
		\int_{\Omega_0 \cap \{\sqrt{R} \le \rho \le R\}} \frac{\aaa |\nabla u|^2}{\rho^2} \le C_3\left( \frac{V(R)}{R^2} + \int_{\sqrt{R}}^R \frac{V(\sigma)}{\sigma^3} \di \sigma\right).
	\]
\end{proposition}
%
\begin{proof}
	Writing $\rho^{-2}-R^{-2}$ as an integral and using Fubini's theorem and \eqref{primointe_abstract}, we get
	\[
	\begin{array}{lcl}
		\disp \int_{\Omega_0 \cap \{\sqrt{R} \le \rho \le R\}} \aaa |\nabla u|^2 \big[\rho^{-2}-R^{-2}\big] & = & \disp 2\int_{\Omega_0 \cap \{\sqrt{R} \le \rho \le R\}} \aaa |\nabla u|^2 \left[\int^{R}_{\rho} \sigma^{-3}\di \sigma\right]\di x, \\[0.5cm]
		& = & \disp \disp 2\int_{\sqrt{R}}^R \sigma^{-3} \left[\int_{\Omega_0 \cap \{\sqrt{R} \le \rho \le \sigma\}} \aaa |\nabla u|^2 \di x \right] \di \sigma \\[0.5cm]
		& \le & \disp \disp 2\int_{\sqrt{R}}^R \sigma^{-3} V(\sigma) \di \sigma \\[0.5cm]
	\end{array}
	\]
	Again by \eqref{primointe_abstract},
	\[
	\begin{array}{lcl}
		\disp \int_{\Omega_0 \cap \{\sqrt{R} \le \rho \le R\}} \frac{\aaa |\nabla u|^2}{\rho^2} & = & \disp \frac{1}{R^2}\int_{\Omega_0 \cap \{\sqrt{R} \le \rho \le R\}} \aaa |\nabla u|^2 \\[0.5cm]
		& & \disp + \int_{\Omega_0 \cap \{\sqrt{R} \le \rho \le R\}} \aaa |\nabla u|^2 \big[\rho^{-2}-R^{-2}\big] \\[0.5cm] 
		& \le & \disp C\left( \frac{V(R)}{R^2} + \int_{\sqrt{R}}^R \frac{V(\sigma)}{\sigma^3} \di \sigma\right).
	\end{array}
	\]
\end{proof}

\begin{proposition}\label{prop_volume_MC}    %[Lemma 5.2]
	Let $a$ satisfy \eqref{assu_A} and 
	\begin{equation}\label{eq_ass_lambda}
	\begin{array}{ll}
	t^2 \lambda_{\max}(t) \to 0 & \quad \text{as } \, t \to 0, \\[0.3cm]
	t^2 \lambda_1(t) \le C \lambda_2(t) & \quad \text{on } \, (0, \infty).
	\end{array}
	\end{equation}
	Assume that $u$ satisfies \eqref{primointe_abstract} for some $0 \le V \in C(\R^+)$, and that
	\begin{equation}\label{crescivol}
		\lim_{R \ra \infty} \frac{V(R)}{R^2 \log R}=0. 
	\end{equation}
	Then, $u$ has moderate energy growth in $\Omega_0$.
%	there exists a sequence $\{\varphi_j\} \subseteq \lip_c(\overline\Omega)$, $\varphi_j \ra 1$ locally uniformly, such that 
%	$$
%	\int |\nabla u|^2\langle A(\nabla u)\nabla \varphi_j, \nabla \varphi_j \rangle \ra 0 \qquad \text{as } \, j \ra \infty.
%	$$
\end{proposition}

\begin{proof}
	Choose 
	$$
	\varphi_R = \left\{ \begin{array}{ll}
		1 & \quad \text{if } \, \rho \le \sqrt{R}, \\[0.1cm]
		\disp \frac{2\log (R/\rho)}{\log R} & \quad \text{if } \, \rho \in [\sqrt{R},R] \\[0.3cm]
		0 & \quad \text{if } \, \rho \ge R.
	\end{array}\right.
	$$
	Then, 
	$$
	\nabla \varphi_R = -2\frac{r\nabla r+ u\nabla u}{\rho^2 \log R} \qquad \text{for } \, \rho \in [\sqrt{R},R].
	$$
	We thus have on $\{|\nabla u|>0\}$
	$$
	\begin{array}{lcl}
		\disp \langle A(\nabla u) \nabla \varphi_R, \nabla \varphi_R \rangle & = & \disp \frac{4r^2}{\rho^4 \log^2 R}\langle A(\nabla u)\nabla r, \nabla r\rangle +  \frac{4u^2}{\rho^4 \log^2 R} \langle A(\nabla u) \nabla u, \nabla u\rangle \\[0.4cm]
		& & \disp + \frac{8r u}{\rho^4 \log^2R} \langle A(\nabla u)\nabla u, \nabla r\rangle.
	\end{array}
	$$
	Hence, by Young's inequality and using $\rho^2 = r^2 + u^2$, 
	$$
	\begin{array}{lcl}
		\disp \langle A(\nabla u) \nabla \varphi_R, \nabla \varphi_R \rangle & \le & \disp C \left(\frac{r^2}{\rho^4 \log^2 R}\langle A(\nabla u)\nabla r, \nabla r\rangle +  \frac{u^2}{\rho^4 \log^2 R} \langle A(\nabla u) \nabla u, \nabla u\rangle\right) \\[0.5cm]
		& \le & \disp \frac{C}{\rho^2 \log^2 R} \Big(\langle A(\nabla u)\nabla r, \nabla r\rangle + \langle A(\nabla u) \nabla u, \nabla u\rangle\Big).
	\end{array}
	$$
	In our assumptions, 
	$$
	\begin{array}{lcl}
		\disp \langle A(\nabla u)\nabla r, \nabla r\rangle \le \lambda_{\max}(|\nabla u|) \le C \aaa \\[0.2cm]
		\langle A(\nabla u) \nabla u, \nabla u\rangle = \lambda_1(|\nabla u|)|\nabla u|^2 \le C\aaa.
	\end{array}
	$$
	Hence 
	$$
	\begin{array}{lcl}
		\disp \langle A(\nabla u) \nabla \varphi_R, \nabla \varphi_R \rangle & \le & \disp C\frac{\aaa}{\rho^2 \log^2 R}.
	\end{array}
	$$
	holds in $\{|\nabla u|>0\}$. Integrating and using Proposition \ref{prop_inteannuli}, 
	$$
	\begin{array}{lcl}
		\disp \int_{\Omega_0 \cap \{|\nabla u|>0\}} |\nabla u|^2 \langle A(\nabla u) \nabla \varphi_R, \nabla \varphi_R \rangle & \le & \disp \frac{C}{\log^2 R} \int_{\Omega_0 \cap \{\sqrt{R} \le \rho \le R\}} \frac{\aaa |\nabla u|^2}{\rho^2} \\[0.5cm]
		& \le & \disp C\left( \frac{V(R)}{R^2\log^2R} + \frac{1}{\log^2 R}\int_{\sqrt{R}}^R \sigma^{-3} V(\sigma) \di \sigma\right).
	\end{array}
	$$
	By the first in \eqref{eq_ass_lambda}, $|\nabla u|^2 A(\nabla u) \to 0$ as a tensor if $|\nabla u| \to 0$, so the set $\{|\nabla u|=0\}$ can be added in the left-hand side keeping the validity of the inequality. Because of \eqref{crescivol}, for a suitable function $o_R(1)$ of $R$ vanishing as $R \to \infty$ we therefore have 
	$$
	\begin{array}{lcl}
		\disp \int_{\Omega_0} |\nabla u|^2 \langle A(\nabla u) \nabla \varphi_R, \nabla \varphi_R \rangle & \le & \disp o_R(1) \left( \frac{1}{\log R} + \frac{1}{\log^2 R}\int_{\sqrt{R}}^R \frac{\log \sigma}{\sigma}\di \sigma\right) \\[0.5cm]
		& \le & \disp o_R(1) \left( \frac{1}{\log R} + \frac{1}{\log R}\int_{\sqrt{R}}^R \frac{\di \sigma}{\sigma}\right) = o_R(1) \left(\frac{1}{\log R}+ 1\right)
	\end{array}
	$$
	and the thesis follows.
\end{proof}

Putting together Propositions \ref{prop_1_calib}, \ref{prop_inteannuli} and \ref{prop_volume_MC}, we readily deduce the following when $a$ satisfies the stronger \eqref{assu_A_strong}.


\begin{theorem}\label{teo_goodcutoff_MC}
	Assume that  $f \in C^1(\R)$, that $a$ satisfies \eqref{assu_A_strong} and that
	\[
			\disp t^2 \lambda_1(t) \le C_1 \lambda_2(t) \le \frac{C_2}{t} \qquad \text{on } \, (0, \infty).
	\]
	for some constants $C_1,C_2>0$. Let $u \in C^2(\Omega) \cap C^1(\overline\Omega)$ solve $\Delta_a u + f(u) = 0$. Let $\Omega_0\subseteq\Omega$ be an open set with locally finite perimeter in $\Omega$ such that $\partial\Omega_0 \subseteq \partial\Omega \cup \critu$ and
	\[
			|\Omega_0 \cap B_R| = o\big( R^2 \log R\big) \qquad \text{as } \, R \to \infty.
	\] 
	Then, $u$ has moderate energy growth in $\Omega_0$ in any of the following cases:
	\begin{itemize}
		\item[(i)] $- C_0 \le f(u) \le 0$ in $\Omega_0$, 
		\[
			u_{|\overline{\Omega}_0 \cap \partial_\star\Omega} \, \text{ is bounded,} \qquad \inf_{\Omega_0} u > -\infty \, .
		\]
		and for some $b \in \R$ 
		\[
			\haus^{n-1}(\overline{\Omega}_0 \cap \partial^b_\star\Omega \cap B_{R}) = o\big( R^2 \log R\big) \qquad \text{as } \, R \to \infty, 
		\]
		where $\partial_\star^b\Omega = \{ x \in \partial_\star \Omega : f(x) \neq b\}$.
		\item[(ii)] $- C_0 \le f(u) \le 0$ in $\Omega_0$, 
		\[
			u_{|\overline{\Omega}_0 \cap \partial_\star\Omega} \, \text{ is constant,} \qquad \inf_{\Omega_0} u > -\infty \, .
		\]
		\item[(iii)] $f(u) \le 0$ in $\Omega_0$, 
		\[
			u_{|\overline{\Omega}_0 \cap \partial_\star\Omega} \, \text{ is a constant $b$,} \qquad \inf_{\Omega_0} u = b\, 
		\]
		\item[(iv)] $|f(u)| \le C_0$ in $\Omega_0$, $u$ is constant on $\overline{\Omega}_0 \cap \partial_\star \Omega$ and 
		\[
			|\Omega_0 \cap B_R| = o\big(R \log R\big) \qquad \text{as } \, R \to \infty.
		\]
		\item[(v)] $u$ is constant on $\overline{\Omega}_0 \cap \partial_\star \Omega$ and bounded in $\Omega_0$.
		\item[(vi)] $f(u) \le 0$ in $\Omega_0$ and 
		\[
			\haus^{n-1}(\overline{\Omega}_0 \cap \partial_\star\Omega \cap B_{R}) = o\big(R \log R\big) \qquad \text{as } \, R \to \infty.
		\]
	\end{itemize}			
\end{theorem}

%Remark \ref{rem_very_useful} still applies to the context of Theorem \ref{teo_goodcutoff_MC}, allowing to localize \eqref{crescivol_2} to some subset of $\partial \Omega$ under suitable boundary conditions.



\section{Stability and monotonicity}\label{sec_stability}

Let $u \in C^2(\Omega)$ be a non-constant, stable solution to $\Delta_a u + f(u) =0$ with $f \in C^1(\R)$. In this section, we shall investigate some consequences of the stability property, so we assume that:

\begin{itemize}
	\item $a$ satisfies \eqref{assu_A_strong}. In particular, $A(\nabla u)$ is a $(1,1)$ tensor field on $\Omega$  whose eigenvalues are locally bounded away from $0$ and $\infty$ (we shortly say that $A(\nabla u)$ is locally uniformly elliptic);
	\item \eqref{hardy} holds for each test function $\varphi \in \lip_c(\Omega)$. 
\end{itemize}

Observe that 

\begin{itemize}
	\item[-] If $u \in C^2(\Omega)$ and $a$ satisfies \eqref{assu_A_strong}, then $A(\nabla u)$ has $L^\infty_\loc(\Omega)$ coefficients;
	\item[-] If $u \in C^2(\Omega)$ and $a$ satisfies \eqref{assu_A_superstrong}, then $A(\nabla u)$ has $\lip_\loc(\Omega)$ coefficients.
\end{itemize} 

If $A(\nabla u)$ has $\lip_\loc$ coefficients, it is well-known that stability is equivalent to the existence of $0 < v \in C^{1,\alpha}_\loc(\Omega)$ solving the linearized equation
\[
	\diver \left(A(\nabla u) \nabla v\right) + f'(u)v = 0,
\]
see \cite{mosspie} for a proof. An analogous result holds when $A(\nabla u)$ has only $L^\infty_\loc$ coefficients. This is probably well-known, but we include a brief proof for the sake of completeness.

\begin{lemma}\label{lem_mosspie}
	Let $A$ be a $(1,1)$ tensor field with $L^\infty_\loc(\Omega)$ coefficients which is locally uniformly elliptic, and let $V \in L^\infty_\loc(\Omega)$. Then, the following are equivalent:
	\begin{itemize}
		\item[(i)] For each $\varphi \in \lip_c(\Omega)$,
		\[
			I(\varphi) \doteq \int_M \langle A\nabla\varphi, \nabla \varphi\rangle \di x - \int_{M} V \varphi^2\di x \ge 0.
		\]
		\item[(ii)] There exists $0 < v \in C^{0,\alpha}_\loc(\Omega) \cap H^1_\loc(\Omega)$ solving $\diver(A \nabla v) + Vv = 0$ weakly in $\Omega$;
		\item[(iii)] There exists $0 < w \in H^1_\loc(\Omega)$ solving $\diver(A \nabla w) + Vw \le 0$ weakly in $\Omega$, where $w>0$ means that $w$ has locally positive essential infimum on $\Omega$.
	\end{itemize}
\end{lemma}

\begin{proof}
	$(ii) \Rightarrow (iii)$ is obvious, and $(iii) \Rightarrow (i)$ follows the standard path by integrating 
	\[
		\diver(A \nabla w) + Vw \le 0
	\]
	against $\varphi^2/w$ and using Young's inequality, see \cite[Lemma 3.10]{prs}. To prove $(i) \Rightarrow (ii)$, we follow the argument in \cite[Theorem 1]{fischercolbrieschoen}: we choose an increasing, smooth exhaustion $\Omega_j \uparrow M$ and associated solutions to
	\[
		\left\{ \begin{array}{ll}
			\diver \left( A\nabla v_j\right) + V v_j = 0	& \text{in } \Omega_j \\[0.2cm]
			v_j = 1 & \text{on } \, \partial \Omega_j
		\end{array} \right.
	\]
	The existence of $v_j$ follows from \cite[Theorems 8.6, 8.12, 8.29]{gilbargtrudinger} if we prove that 
	\[
		\lambda_1(\Omega_j) = \inf\Big\{ Q(\psi) \ : \ 0 \not \equiv \psi \in \lip_c(\Omega_j) \Big\} > 0, \qquad Q(\psi) = \frac{I(\psi)}{\int \psi^2}.
	\]
	Indeed, if $\lambda_1(\Omega_j) = 0$, pick a minimizer $\xi_j$ for $I$ in $H^1_0(\Omega_j)$. Since $I$ is homogeneous and $I(\psi) = I(|\psi|)$ we can suppose that $\xi_j \ge 0$, $\int \xi_j^2  = 1$. Extending $\xi_j$ with zero outside of $\Omega_j$ yields a function $\xi_j \in H^1_0(\Omega_{j+1})$. Now, from $0 = \lambda_1(\Omega_j) = I(\xi_j) \ge \lambda_1(\Omega_{j+1}) \ge 0$ we deduce that $\xi_j$  also minimizes $Q$ in $H^1_0(\Omega_{j+1})$, hence it is a non-negative solution to 
	\[
		\diver \left( A \nabla \xi_j\right) + V\xi_j = 0 \qquad \text{in } \, \Omega_{j+1}
	\]
	which vanishes in a neighbourhood of $\partial \Omega_{j+1}$, contradicting the Harnack inequality in \cite[Theorem 8.20 and Corollary 8.21]{gilbargtrudinger}. Once the sequence $\{v_j\}$ is produced, to obtain $v$ it is enough to rescale $v_j$ to $1$ at a fixed point $p \in \Omega_1$ and to pass to limits by using again the Harnack inequality together with the local uniform $C^{0,\alpha}$ and $H^1$ elliptic estimates.  
\end{proof}

\begin{remark}\label{rem_harnackHopf}
	If $a$ satisfies \eqref{assu_A_strong} and $u \in C^1(\Omega)$, non-negative weak solutions $0 \le w \in C^1(\Omega)$ of 
	\[
		\diver\big(A(\nabla u) \nabla w \big) + f'(u)w \le 0
	\]
	satisfy the (half)-Harnack inequality \cite[Thm. 7.1.2]{pucci_serrin}, so either  $w\equiv 0$ or $w>0$ everywhere in the domain $\Omega$. Moreover, if $a$ satisfies \eqref{assu_A_superstrong} and $u \in C^2(\overline\Omega)$, the Hopf Lemma holds for solutions $w \in C^2(\Omega) \cap C^1(\overline\Omega)$, see \cite[Theorem 2.8.3]{pucci_serrin}.
%	\cite[Theorems 2.8.3 and 2.1.1]{pucci_serrin}. 
\end{remark}

As is well-known, positive solutions to the linearized equation of $\Delta_a u + f(u) = 0$ can be produced if $u$ is monotone in the direction of a Killing vector field. We include the proof of the next result in our needed generality.

\begin{lemma}\label{lem_linearKilling}
	%
	Assume the validity of \eqref{assu_A_strong}, 
	%$(a1),(a2),(a3)$
 	let $\Omega \subseteq M$ be a domain and suppose that $u \in C^2(\Omega)$ solves $\Delta_a u + f(u)=0$. Let $X$ be a Killing field in $\overline\Omega$. Then, $w \doteq \langle \nabla u, X\rangle$ solves
	\begin{equation}\label{linearized_w} 
		\diver\big( A(\nabla u) \nabla w\big) + f'(u) w = 0 
	\end{equation}
	weakly in $\Omega$.
\end{lemma}

\begin{proof}
%	First of all we observe that, by $(a3)$, the map $X \ra a(|X|)X$ is locally Lipschitz from $\gamma(TM)$ to itself, and $a(|X|)X = 0$ if $X=0$. Therefore, by Stampacchia's theorem, $\nabla (a(|X|)X \equiv 0$ almost everywhere on $|X|=0$. 
Consider the flux $\Phi : \mathscr{D} \subseteq \Omega \times \R \to \Omega$ associated to $X$, defined on its maximal domain $\mathscr{D}$, set  $\Phi_t(x)=\Phi(t,x)$ and let $\mathscr{D}_t \subseteq \Omega$ be the maximal domain of $\Phi_t$. Set $u_t \doteq u \circ \Phi_t : \mathscr{D}_t \to \R$. The Killing property of $X$ implies the next identities for each $\psi \in C^2(\Omega)$ and $\psi_t = \psi \circ \Phi_t : \mathscr{D}_t \to \R$:
	\begin{equation}\label{propbase}
		\left\{\begin{array}{l}
			\partial_t \psi_t = \langle \nabla \psi, X\rangle \circ \Phi_t, \qquad \nabla \psi_t = \di \Phi_{-t}(\nabla \psi), \\[0.2cm]
			|\nabla \psi_t|^2 = %|(\di \Phi_t \circ \di g)^\sharp|^2 = 
			|\nabla \psi|^2 \circ \Phi_t, \\[0.2cm]
			\di (\partial_t \psi_t) = \di\big(\langle \nabla \psi, X\rangle \big) \circ \di \Phi_t = \nabla \di  \psi(X, \di \Phi_t) + \langle \nabla \psi, \nabla_{\di \Phi_t} X \rangle.
		\end{array}\right.
	\end{equation}
	%
	%Now, by the identity $|\nabla u_t|^2 = |\nabla u|^2 \circ \Phi_t$, we can choose $t_\eps$ small enough %that 
	%\begin{equation}
	%|\nabla u_t(x)|^2 \ge \frac{\eps}{2} \qquad \text{for all } x \in U_\eps \cap \supp(\varphi), \ \ t \in %(0, t_\eps).
	%\end{equation}
	In particular, from the second equation in \eqref{propbase} we deduce that, fixing $z \in U$ and considering the vector field $\nabla u_t$ along the constant curve $\gamma(t)=z$, 
	%
	\begin{equation}\label{derit}
		\left.\frac{\nabla}{\di t}\right|_{0}\big(\nabla u_t(z)\big) = \left.\frac{\nabla}{\di t}\right|_0 \Big(\di \Phi_{-t}\big[\nabla u(\Phi_t(z))\big]\Big) = (\mathscr{L}_X \nabla u)(z) = [X, \nabla u] = \nabla_X \nabla u - \nabla_{\nabla u}X.
	\end{equation}
	
	%\begin{equation}\label{simple}
	%\langle \nabla (u \circ \Phi_t), \nabla (w \circ \Phi_t) \rangle = \di (w \circ \Phi_t)(u \circ \Phi_t) = %\nabla \di u(X, \di \Phi_t(\nabla u)) + \langle \nabla u, \nabla_{\di \Phi_t(\nabla u)} X \rangle = %\nabla \di u(X, \nabla u_t).
	%
	%\end{equation}
	Since $\Phi_t$ is an isometry, $u_t$ solves weakly $\Delta_a u_t + f(u_t) = 0$ on $\mathscr{D}_t$. Fix $\varphi \in C^\infty_c(\Omega)$, and choose $t_0$ small enough so that, for $|t| \le t_0$, $\varphi \in C^\infty_c(\mathscr{D}_{t})$. Then,  
	\[
	\int_M \Big\{ a(|\nabla u_t(z)|) \langle \nabla u_t, \nabla \varphi\rangle_z - f(u_t(z))\varphi(z) \Big\} \di z = 0.
	\]
%	We plug $\varphi_{-t} = \varphi \circ \Phi_{-t}$ in the weak definition of $\Delta_a u + f(u)=0$ to obtain
%	$$
%	0 = \disp \int \Big\{a(|\nabla u(y)|) \langle \nabla u, \nabla \varphi_{-t} \rangle_y - f(u(y))\varphi(\Phi_{-t}(y)) \Big\}\di y.
%	$$
%	Changing variable according to $y = \Phi_t(z)$ and using \eqref{propbase} we deduce
%	$$
%	\begin{array}{lcl}
%		0 & = & \disp \disp \int_M \Big\{a(|\nabla u_t(z)|) \langle \di \Phi_t (\nabla u_t), \di \Phi_t(\nabla \varphi)\rangle_{\Phi_t(z)} - f(u_t(z))\varphi(z)\Big\} \di z. \\[0.4cm]
%		0 & = & \disp \disp \int_M \Big\{ a(|\nabla u_t(z)|) \langle \nabla u_t, \nabla \varphi\rangle_z - f(u_t(z))\varphi(z) \Big\} \di z.  
%	\end{array}
%	$$
	Next, by \eqref{assu_A_strong} and since $u \in C^2(\Omega)$ the integrand is a Lipschitz function of $(t, z) \in (-\eps, \eps) \times {\rm spt}(\varphi)$. Hence, by the differentiation theorem under the integral sign we can differentiate in $t$ to get 
	\begin{equation}\label{diffeinte}
		\begin{array}{lcl}         
			0 & = & \disp \disp \int \Big\{\frac{a'(|\nabla u_t|)}{|\nabla u_t|} \nabla \di  u_t(\nabla u_t,X) \langle \nabla u_t, \nabla \varphi\rangle + a(|\nabla u_t|) \frac{\di}{\di t}\langle \nabla u_t, \nabla \varphi \rangle - f'(u_t)(w \circ \Phi_t)\varphi\Big\}.
		\end{array}
	\end{equation}
	Evaluating at $t=0$ and noting that, by \eqref{derit} and since $X$ is Killing,
	$$
	\begin{array}{lcl}
		\disp \left.\frac{\di}{\di t}\right|_{t=0}\langle \nabla u_t, \nabla \varphi \rangle_z & = & \disp \langle \left.\frac{\nabla}{\di t}\right|_0(\nabla u_t), \nabla \varphi \rangle_z = \langle \nabla_X \nabla u, \nabla \varphi \rangle - \langle \nabla_{\nabla u} X, \nabla \varphi \rangle \\[0.3cm]
		& = & \disp \nabla \di  u(X, \nabla \varphi) - \langle \nabla_{\nabla u} X, \nabla \varphi \rangle = \disp \nabla \di  u(X, \nabla \varphi) + \langle \nabla_{\nabla \varphi} X, \nabla u \rangle
	\end{array}
	$$
	we obtain
	\begin{equation}\label{diffeinte_2}
		\begin{array}{lcl}
			0 & = & \disp \disp \disp \int \Big\{ \frac{a'(|\nabla u|)}{|\nabla u|} \nabla \di  u(\nabla u,X) \langle \nabla u, \nabla \varphi\rangle + a(|\nabla u|) \big[\nabla \di  u(X, \nabla \varphi) + \langle \nabla_{\nabla \varphi} X, \nabla u \rangle\big] - f'(u)w \varphi\Big\}.
			%& = & \disp \int \Big\{\frac{a'(|\nabla u|)}{|\nabla u|} \langle \nabla w, \nabla u \rangle \langle \nabla u, \nabla \varphi \rangle + a(|\nabla u|)\big[\nabla \di u(X, \nabla \varphi) + \langle \nabla_{\nabla \varphi} X, \nabla u \rangle\big] - f'(u)w \varphi\Big\}.
		\end{array}
	\end{equation}
	Using now that 
	$$
	\langle \nabla w, \nabla \varphi \rangle = \nabla \varphi\big( \langle \nabla u, X \rangle \big) = \nabla \di  u (\nabla \varphi, X) + \langle \nabla u, \nabla_{\nabla \varphi}X\rangle
	$$
	we get
	\begin{equation}\label{final}
		\begin{array}{lcl}
			0 & = & \disp \int \Big\{\frac{a'(|\nabla u|)}{|\nabla u|} \langle \nabla w, \nabla u \rangle \langle \nabla u, \nabla \varphi \rangle + a(|\nabla u|)\langle \nabla w, \nabla \varphi \rangle - f'(u)w \varphi\Big\} \\[0.4cm]
			& = & \disp \int \Big\{ \langle A(\nabla u)\nabla w, \nabla \varphi \rangle - f'(u)w \varphi \Big\},
		\end{array}
	\end{equation}
	that concludes the proof. 
\end{proof}

\begin{lemma} \label{Lemma_w-}
	Let $(M^n,\metric)$ be a Riemannian manifold and $\Omega \subseteq M$ a domain. Let $a$ satisfy \eqref{assu_A_strong}, $f \in C^1(\R)$ 
%	and let $u\in C^2(\Omega)$ be a stable solution to
%	\[
%		\Delta_a u + f(u) = 0 \qquad \text{on } \Omega. 
%	\]	
	and let $w\in C^1(\Omega)\cap C(\overline{\Omega})$ be a solution to
	\begin{equation} \label{w-eq}
	\left\{
	\begin{array}{ll}
		\diver\big( A(\nabla u) \nabla w\big) + f'(u) w = 0 & \qquad \text{weakly in } \Omega \\[0.2cm]
		w \geq 0 & \qquad \text{on } \, \partial\Omega \\
	\end{array}
	\right.
	\end{equation}
	for some $u \in C^1(\overline\Omega)$. Suppose that there exists $\{\varphi_j\} \subseteq \lip_c(\overline\Omega)$ satisfying 
	\begin{equation} \label{Hgrowth}
	\varphi_j \to 1 \quad \text{in } \, C_\loc(\overline\Omega) \qquad \text{and} \qquad \int_\Omega w_-^2 \langle A(\nabla u) \nabla \varphi_j, \nabla \varphi_j \rangle \to 0
\end{equation}
as $j \to \infty$, where $w_- = \max\{-w,0\}$ denotes the negative part of $w$. Then either
%	\begin{equation} \label{Hgrowth}
%		\int_{\Omega\cap B_R} w_-^2 \lambda_{\max}(|\nabla u|) = O(R^2\log R) \qquad \text{as } \, R \to \infty
%	\end{equation}
%	\tcr{and
%		\begin{equation} \label{degen}
%			\{ \nabla u = 0 \} \subseteq \{ w = 0 \}
%		\end{equation}
%	}
	\begin{equation} \label{w-alt3}
		w\equiv 0 \text{ in } \Omega \qquad \text{or} \qquad w<0 \text{ in } \Omega \qquad \text{or} \qquad w>0 \text{ in } \Omega.
	\end{equation}
%	$w\geq 0$ everywhere in $\Omega$ or $w<0$ everywhere in $\Omega$. Moreover,
%	\begin{itemize}
%		\item [i)] if $A(0)$ is nondegenerate then the assumption \eqref{degen} is not needed;
%		\item [ii)] if the strong maximum principle holds for $C^2$ solutions $\varphi$ to
%		\[
%		\diver(A(\nabla u)\nabla\varphi) + f'(u)\varphi = 0
%		\]
%		or, alternatively, if $w$ satisfies
%		\[
%		w_{|\partial\Omega} = 0 \qquad \text{and} \qquad \int_{\Omega\cap B_R} w^2 \lambda_{\max}(|\nabla u|) = O(R^2\log R) \quad \text{as } \, R \to \infty
%		\]
%		then we either have
%		\begin{equation} \label{w-alt3}
%			w\equiv 0 \text{ in } \Omega \qquad \text{or} \qquad w<0 \text{ in } \Omega \qquad \text{or} \qquad w>0 \text{ in } \Omega.
%		\end{equation}
%	\end{itemize}
\end{lemma}

\begin{remark}
		We emphasize that in Lemma \ref{Lemma_w-} no assumptions are made on the regularity of $\partial\Omega$, as we only require $\Omega$ to be open and connected.
\end{remark}

\begin{proof}
By Lemma \ref{lem_mosspie}, since $u$ is stable there exists a weak solution $0 < v \in C^{0,\alpha}_\loc(\Omega) \cap H^1_\loc(\Omega)$ to
\[
		\diver \left( A(\nabla u)\nabla v\right) + f'(u)v = 0.
\]

% Since $u$ is stable, there exists a weak solution $0 < v \in C^{0,\alpha}_\loc(\Omega) \cap H^1_\loc(\Omega)$ to
%	\[
%		\diver \left( A(\nabla u)\nabla v\right) + f'(u)v = 0.
%	\]
%	To show this fact first observe that, in our assumptions on $a$, a priori $A(\nabla u)$ only has $L^\infty_\loc$ coefficients. However, with few small changes we can follow the argument in \cite[Theorem 1]{fischercolbrieschoen}: we choose an increasing, smooth exhaustion $\Omega_j \uparrow M$ and associated solutions to
%	\[
%	\left\{ \begin{array}{ll}
%	\diver \left( A(\nabla u)\nabla v_j\right) + f'(u)v_j = 0	& \text{in } \Omega_j \\[0.2cm]
%	v_j = 1 & \text{on } \, \partial \Omega_j
%	\end{array} \right.
%	\]
%	The existence of $v_j$ follows from \cite[Theorems 8.6, 8.12, 8.29]{gilbargtrudinger} if we prove that 
%	\[
%	\lambda_1(\Omega_j) = \inf\left\{ \frac{I(\psi)}{\int \psi^2} \ : \ 0 \not \equiv \psi \in \lip_c(\Omega_j) \right\} > 0,
%	\]
%	where
%	\[
%	I(\psi) = \int \langle A(\nabla u)\nabla \psi, \nabla \psi\rangle  - f'(u)\psi^2.
%	\]
%	Indeed, if $\lambda_1(\Omega_j) = 0$, pick a minimizer $\xi_j$ for $I$ in $H^1_0(\Omega_j)$. Since $I$ is homogeneous and $I(\psi) = I(|\psi|)$ we can suppose that $\xi_j \ge 0$, $\int \xi_j^2  = 1$. Extending $\xi_j$ with zero outside of $\Omega_j$ yields a function $\xi_j \in H^1_0(\Omega_{j+1})$. Now, from $0 = \lambda_1(\Omega_j) = I(\xi_j) \ge \lambda_1(\Omega_{j+1}) \ge 0$ we deduce that $\xi_j$  also minimizes $I$ in $H^1_0(\Omega_{j+1})$, hence it is a non-negative solution to 
%	\[
%	\diver \left( A(\nabla u)\nabla \xi_j\right) + f'(u)\xi_j = 0 \qquad \text{in } \, \Omega_{j+1}
%	\]
%	which vanishes in a neighbourhood of $\partial \Omega_{j+1}$, contradicting the Harnack inequality in \cite[Theorem 8.20 and Corollary 8.21]{gilbargtrudinger}. Once the sequence $\{v_j\}$ is produced, to obtain $v$ it is enough to rescale $v_j$ to $1$ in a fixed point $p \in \Omega_1$ and to pass to limits by using again the Harnack inequality together with the local uniform $C^{0,\alpha}$ and $H^1$ elliptic estimates.  
%
Set $z = w_-/v$. We first show that
	\begin{equation} \label{A-L1}
		\int_\Omega \varphi^2 v^2 \langle A(\nabla u) \nabla z, \nabla z \rangle \leq 4 \int_\Omega w_-^2 \langle A(\nabla u)\nabla\varphi,\nabla\varphi \rangle %< \infty
		\qquad \text{for all } \, \varphi \in \lip_c(M) \, .
	\end{equation}
	To do so, for $\eps>0$ let us set $v_\eps \doteq v+\eps$ and $z_\eps = w_-/v_\eps$. The function $v_\eps$ solves
	\[
	\diver \left( A(\nabla u)\nabla v_\eps\right) + f'(u)v_\eps = \eps f'(u)
	\]
	weakly in  $\Omega$, so $z_\eps$ satisfies
	\begin{equation} \label{veps_weak}
		\diver \left( v_\eps^2 A(\nabla u)\nabla z_\eps \right) \ge -\eps w_-f'(u) \qquad \text{weakly in } \, \Omega.
	\end{equation}
	Let $\varphi \in \lip_c(\overline\Omega)$ be a cut-off function on $M$, $\delta>0$ and set $U_\delta \doteq \{z > \delta\}$ and $U_{\eps,\delta} \doteq \{z_\eps > \delta\}$.
	%Since we are assuming $z\not\equiv 0$, there exists $\delta>0$ such that the set $U_\delta \doteq \{z > \delta\}$ is non-empty. Let $\eps$ be small enough so that $U_{\eps,\delta} \doteq \{z_\eps > \delta\} \neq \emptyset$. By construction,
	By construction
	\[
	\overline{U}_{\eps,\delta} \subseteq \overline{\{w_- > \eps\delta\}} \subseteq \Omega
	\]
	so $(z_\eps -\delta)_+$ vanishes in a neighbourhood of $\partial \Omega$ and thus, testing with $\varphi^2 (z_\eps -\delta)_+ \in \lip_c(\Omega)$ we get
	\begin{align*}
		\disp - \eps \int_{\Omega} \varphi^2 f'(u) & w_-(z_\eps -\delta)_+ \le - \disp \int_\Omega \langle A(\nabla u) \nabla z_\eps, \nabla(\varphi^2(z_\eps -\delta)_+) \rangle v_\eps^2\\[0.5cm]
		& = \disp - \int_\Omega \varphi^2 \langle A(\nabla u) \nabla z_\eps, \nabla z_\eps \rangle v_\eps^2 \mathbf{1}_{U_{\eps,\delta}} + 2 \int_\Omega \varphi(z_\eps - \delta)_+ \langle A(\nabla u) \nabla z_\eps, \nabla \varphi \rangle v_\eps^2 \\[0.5cm]
		& \leq \disp - \frac{1}{2} \int_\Omega \varphi^2 v_\eps^2 \langle A(\nabla u) \nabla z_\eps, \nabla z_\eps \rangle \mathbf{1}_{U_{\eps,\delta}} + 2 \int_\Omega (z_\eps - \delta)_+^2v_\eps^2 \langle A(\nabla u) \nabla \varphi, \nabla \varphi \rangle.	    
	\end{align*}
	Since $(z_\eps-\delta)^2_+ v_\eps^2 \le w_-^2$ on $\Omega$, we have	
	\[
	\frac{1}{2} \int_\Omega \varphi^2 v_\eps^2 \langle A(\nabla u) \nabla z_\eps, \nabla z_\eps \rangle \mathbf{1}_{U_{\eps,\delta}} \le 2 \int_\Omega w_-^2 \langle A(\nabla u) \nabla \varphi, \nabla \varphi \rangle + \eps \int_\Omega \varphi^2 f'(u)w_-(z_\eps -\delta)_+ .
	\]
	We keep $\delta>0$ fixed and we let $\eps\to0^+$ in this inequality. Since
	\[
	|\eps\varphi^2 f'(u)w_-(z_\eps-\delta)_+| \leq \eps \varphi^2 |f'(u)| w_- z_\eps = \varphi^2 |f'(u)| w_-^2 \frac{\eps}{v_\eps} \leq \varphi^2 |f'(u)| w_-^2 \in L^1(\Omega)
	\]
	by dominated convergence we have
	\[
	\lim_{\eps\to0^+} \eps \int_\Omega \varphi^2 f'(u)w_-(z_\eps -\delta)_+ = 0 \, .
	\]
	On the other hand, expanding $\nabla z_\eps$ and using Fatou's lemma we have
	\[
	\liminf_{\eps\to0^+} \int_\Omega \varphi^2 v_\eps^2 \langle A(\nabla u) \nabla z_\eps, \nabla z_\eps \rangle \mathbf{1}_{U_{\eps,\delta}} \geq \int_{\Omega} \varphi^2 v^2 \langle A(\nabla u) \nabla z, \nabla z \rangle \mathbf{1}_{U_\delta}
	\]
	so we get
	\[
	\int_\Omega \varphi^2 v^2 \langle A(\nabla u) \nabla z, \nabla z \rangle \mathbf{1}_{U_\delta} \le 4 \int_\Omega w_-^2 \langle A(\nabla u) \nabla \varphi, \nabla \varphi \rangle
	\]
	and then letting $\delta\to0^+$ we obtain, by monotone convergence, \eqref{A-L1}. In our assumptions, choosing $\varphi = \varphi_j$ in \eqref{A-L1} and letting $j \to \infty$ we deduce
%	
%	From \eqref{A-L1} we further have
%	\[
%	\int_\Omega \varphi^2 v^2 \langle A(\nabla u)\nabla z,\nabla z\rangle \leq 4 \int_\Omega w_-^2 \lambda_{\max}(|\nabla u|) |\nabla\varphi|^2 \qquad \text{for all } \, \varphi \in \lip_c(M) \, .
%	\]
%	By assumption \eqref{Hgrowth} and Lemma \ref{lem_phij}, there exists a sequence $\{\varphi_j\} \subseteq \lip_c(M)$ satisfying $\varphi_j \to 1$ in $W^{1,\infty}_\loc(M)$ and
%	\[
%	\int_\Omega w_-^2 \lambda_{\max}(|\nabla u|) |\nabla\varphi_j|^2 \to 0
%	\]
%	as $j\to\infty$, so in the limit we get
	\[
	\int_\Omega v^2 \langle A(\nabla u)\nabla z,\nabla z\rangle = 0 \, .
	\]
	Since $v>0$ in $\Omega$ and $\langle A(\nabla u)\,\cdot\,,\,\cdot\,\rangle$ is positive definite we infer $\nabla z\equiv0$ almost everywhere in $\Omega$ and therefore $z$ is constant in $\Omega$, that is, there exists a constant $\lambda\geq0$ such that
	\[
		w_- = \lambda v \qquad \text{in } \, \Omega \, .
	\]
	Since $v>0$ in $\Omega$, this implies that either $w_->0$ everywhere in $\Omega$ or $w_- \equiv 0$ in $\Omega$, that is, either $w<0$ everywhere in $\Omega$ or $w\geq 0$ in $\Omega$. In the latter case, by the (half)-Harnack inequality (see Remark \ref{rem_harnackHopf}) we either have $w\equiv 0$ or $w>0$ everywhere in $\Omega$.
\end{proof}

\begin{remark} \label{rem-Lemma_w-}
	In the proof of Lemma \ref{Lemma_w-} we showed that if there exists a strictly positive solution $v\in C^{0,\alpha}_\loc(\Omega) \cap H^1_\loc(\Omega)$ to
	\[
		\diver(A(\nabla u)\nabla v) + f'(u)v = 0 \qquad \text{in } \, \Omega
	\]
	and $w\in C^1(\Omega)\cap C(\overline{\Omega})$ is a weak solution to
	\begin{equation} \label{weq-rem}
		\diver(A(\nabla u)\nabla w) + f'(u)w = 0 \qquad \text{in } \, \Omega
	\end{equation}
	satisfying $w\geq 0$ on $\partial\Omega$ and \eqref{Hgrowth}, then either $w > 0$ in $\Omega$ or $w = \lambda v$ in $\Omega$ for some constant $\lambda\leq 0$. 
%	The same argument shows that if $w\in C^2(\Omega)\cap C(\overline{\Omega})$ is a solution to \eqref{weq-rem} satisfying
%	\[
%		w \leq 0 \quad \text{on } \, \partial\Omega \qquad \text{and} \qquad \int_{\Omega\cap B_R} w_+^2 \lambda_{\max}(|\nabla u|) = O(R^2\log R)
%	\]
%	then either $w<0$ in $\Omega$ or $w = \lambda v$ in $\Omega$ for some constant $\lambda\geq 0$. 
	The same argument applied to both $w_+$ and $w_-$ shows that if $w\in C^1(\Omega)\cap C(\overline{\Omega})$ is a weak solution to \eqref{weq-rem} satisfying both $w=0$ on $\partial \Omega$ and \eqref{Hgrowth} with $w$ in place of $w_-$, then $w = \lambda v$ in $\Omega$ for some constant $\lambda\in\R$.
\end{remark}

%\begin{lemma} \label{Lemma_w-}
%	Let $(M^n,\metric)$ be a complete Riemannian manifold of dimension $m\geq 2$, let $\Omega$ be an open subset with $C^3$ boundary and let $u\in C^3(\overline{\Omega})$ be a non-constant, stable solution to
%	\[
%		\Delta_a u + f(u) = 0 \qquad \text{on } \Omega. 
%	\]	
%	Let $w\in C^2(\Omega)\cap C(\overline{\Omega})$ be a solution to
%	\[
%		\diver\big( A(\nabla u) \nabla w\big) + f'(u) w = 0 \qquad \text{on } \Omega \, .
%	\]
%	phiij $w\geq 0$ on $\partial\Omega$ and 
%	\begin{equation} \label{Hgrowth}
%		\limsup_{R\to+\infty} \frac{1}{R^2\log R} \int_{\Omega\cap B_R} w_-^2 \lambda_{\max}(|\nabla u|) < \infty
%	\end{equation}
%	where $\lambda_{\max}(t) = \max\{\lambda_1(t),\lambda_2(t)\}$, then $w\geq 0$ in $\Omega$.
%\end{lemma}
%
%\begin{proof}
%	Since $u$ is stable, there exists a solution $0 < v \in C^{2,\alpha}_\loc(\Omega)$ to
%	\[
%	\diver \left( A(\nabla u)\nabla v\right) + f'(u)v = 0.
%	\]
%	Set $z = w_-/v$. Note that the claim of the lemma is that $z\equiv 0$, so let us assume, by contradiction, that $z\not\equiv 0$. We first show that
%	\[
%		\int_\Omega \varphi^2 v^2 \langle A(\nabla u) \nabla z, \nabla z \rangle \tcr{\leq 4 \int_\Omega w_-^2 \langle A(\nabla u)\nabla\varphi,\nabla\varphi \rangle} < \infty
%	\]
%	for all $\varphi\in C^\infty_c(M)$. To do so, for $\eps>0$ let us set $v_\eps \doteq v+\eps$ and $z_\eps = w_-/v_\eps$. The function $v_\eps$ solves
%	\[
%		\diver \left( A(\nabla u)\nabla v_\eps\right) + f'(u)v_\eps = \eps f'(u)
%	\]
%	so $z_\eps$ satisfies
%	\begin{equation} \label{veps_weak}
%		\diver \left( v_\eps^2 A(\nabla u)\nabla z_\eps \right) \ge -\eps w_-f'(u) \qquad \text{weakly in } \, \Omega.
%	\end{equation}
%	Let $\varphi$ be a cut-off function on $M$, $\delta>0$ and set $U_\delta \doteq \{z > \delta\}$ and $U_{\eps,\delta} \doteq \{z_\eps > \delta\}$.
%	%Since we are assuming $z\not\equiv 0$, there exists $\delta>0$ such that the set $U_\delta \doteq \{z > \delta\}$ is non-empty. Let $\eps$ be small enough so that $U_{\eps,\delta} \doteq \{z_\eps > \delta\} \neq \emptyset$. By construction,
%	By construction
%	\[
%		\overline{U_{\eps,\delta}} \subseteq \overline{\{w_- > \eps\delta\}} \subseteq \Omega
%	\]
%	so $(z_\eps -\delta)_+$ vanishes in a neighbourhood of $\partial \Omega$ and thus, testing with $\varphi^2 (z_\eps -\delta)_+$ we get
%	\begin{align*}
%		\disp - \eps \int_{\Omega} \varphi^2 f'(u) & w_-(z_\eps -\delta)_+ \le - \disp \int_\Omega \langle A(\nabla u) \nabla z_\eps, \nabla(\varphi^2(z_\eps -\delta)_+) \rangle v_\eps^2\\[0.5cm]
%		& = \disp - \int_\Omega \varphi^2 \langle A(\nabla u) \nabla z_\eps, \nabla z_\eps \rangle v_\eps^2 \mathbf{1}_{U_{\eps,\delta}} + 2 \int_\Omega \varphi(z_\eps - \delta)_+ \langle A(\nabla u) \nabla z_\eps, \nabla \varphi \rangle v_\eps^2 \\[0.5cm]
%		& \leq \disp - \frac{1}{2} \int_\Omega \varphi^2 v_\eps^2 \langle A(\nabla u) \nabla z_\eps, \nabla z_\eps \rangle \mathbf{1}_{U_{\eps,\delta}} + 2 \int_\Omega (z_\eps - \delta)_+^2v_\eps^2 \langle A(\nabla u) \nabla \varphi, \nabla \varphi \rangle.	    
%	\end{align*}
%	Since $(z_\eps-\delta)^2_+ v_\eps^2 \le w_-^2$ on $\Omega$, we have	
%	\[
%		\frac{1}{2} \int_\Omega \varphi^2 v_\eps^2 \langle A(\nabla u) \nabla z_\eps, \nabla z_\eps \rangle \mathbf{1}_{U_{\eps,\delta}} \le 2 \int_\Omega w_-^2 \langle A(\nabla u) \nabla \varphi, \nabla \varphi \rangle + \eps \int_\Omega \varphi^2 f'(u)w_-(z_\eps -\delta)_+ .
%	\]
%	We keep $\delta>0$ fixed and we let $\eps\to0^+$ in this inequality. Since
%	\[
%		|\eps\varphi^2 f'(u)w_-(z_\eps-\delta)_+| \leq \eps \varphi^2 |f'(u)| w_- z_\eps = \varphi^2 |f'(u)| w_-^2 \frac{\eps}{v_\eps} \leq \varphi^2 |f'(u)| w_-^2 \in L^1(\Omega)
%	\]
%	by dominated convergence we have
%	\[
%		\lim_{\eps\to0^+} \eps \int_\Omega \varphi^2 f'(u)w_-(z_\eps -\delta)_+ = 0 \, .
%	\]
%	On the other hand, expanding $\nabla z_\eps$ and using Fatou's lemma we have
%	\[
%		\liminf_{\eps\to0^+} \int_\Omega \varphi^2 v_\eps^2 \langle A(\nabla u) \nabla z_\eps, \nabla z_\eps \rangle \mathbf{1}_{U_{\eps,\delta}} \geq \int_{\Omega} \varphi^2 v^2 \langle A(\nabla u) \nabla z, \nabla z \rangle \mathbf{1}_{U_\delta}
%	\]
%	so we get
%	\[
%		\int_\Omega \varphi^2 v^2 \langle A(\nabla u) \nabla z, \nabla z \rangle \mathbf{1}_{U_\delta} \le 4 \int_\Omega w_-^2 \langle A(\nabla u) \nabla \varphi, \nabla \varphi \rangle
%	\]
%	and then letting $\delta\to0^+$ we obtain, by monotone convergence,
%	\begin{equation} \label{A-L1}
%		\int_\Omega \varphi^2 v^2 \langle A(\nabla u) \nabla z, \nabla z \rangle \le 4 \int_\Omega w_-^2 \langle A(\nabla u) \nabla \varphi, \nabla \varphi \rangle < \infty
%	\end{equation}
%	as claimed.
%	
%	Now, we prove the claim of the lemma. For each $\eps>0$ and $\delta>0$ let $v_\eps$, $z_\eps$, $U_\delta$ and $U_{\eps,\delta}$ be as above. Let $\varphi\in C^\infty_c(M)$, $\varphi\geq0$. Testing \eqref{veps_weak} against $\varphi(z_\eps-\delta)_+$ we get
%	\begin{align*}
%		\disp - \eps \int_{\Omega} \varphi f'(u) & w_-(z_\eps -\delta)_+ \le - \disp \int_\Omega \langle A(\nabla u) \nabla z_\eps, \nabla(\varphi(z_\eps -\delta)_+) \rangle v_\eps^2\\
%		& = \disp - \int_\Omega \varphi \langle A(\nabla u) \nabla z_\eps, \nabla z_\eps \rangle v_\eps^2 \mathbf{1}_{U_{\eps,\delta}} + \int_\Omega (z_\eps - \delta)_+ \langle A(\nabla u) \nabla z_\eps, \nabla \varphi \rangle v_\eps^2 \, ,
%	\end{align*}
%	that is,
%	\[
%		\int_\Omega \varphi v_\eps^2 \langle A(\nabla u) \nabla z_\eps, \nabla z_\eps \rangle \mathbf{1}_{U_{\eps,\delta}} \le \int_\Omega (z_\eps-\delta)_+ \langle A(\nabla u) \nabla z_\eps, \nabla \varphi \rangle v_\eps^2 + \eps \int_\Omega \varphi f'(u)w_-(z_\eps -\delta)_+ .
%	\]
%	By Cauchy-Schwarz, and using that $(z_\eps-\delta)_+ v_\eps \leq w_- \mathbf{1}_{U_{\eps,\delta}}$,
%	\begin{equation} \label{CSweps}
%		\int_\Omega (z_\eps-\delta)_+ \langle A(\nabla u) \nabla z_\eps, \nabla \varphi \rangle v_\eps^2 \leq \int_\Omega w_- v_\eps \sqrt{\langle A(\nabla u) \nabla z_\eps, \nabla z_\eps \rangle} \sqrt{\langle A(\nabla u) \nabla \varphi, \nabla \varphi \rangle} \mathbf{1}_{U_{\eps,\delta}} \, .
%	\end{equation}
%	We now let $\eps\to0^+$ while keeping $\delta>0$ fixed. As above, we have
%	\begin{align*}
%		\lim_{\eps\to0^+} \eps \int_\Omega \varphi f'(u)w_-(z_\eps -\delta)_+ & = 0 \\
%		\liminf_{\eps\to0^+} \int_\Omega \varphi v_\eps^2 \langle A(\nabla u) \nabla z_\eps, \nabla z_\eps \rangle \mathbf{1}_{U_{\eps,\delta}} & \geq \int_\Omega \varphi v^2 \langle A(\nabla u) \nabla z, \nabla z \rangle \mathbf{1}_{U_\delta} \, .
%	\end{align*}
%	Furthermore, since
%	\begin{align*}
%		v_\eps^2 \nabla z_\eps & = v_\eps\nabla w_- - w_-\nabla v \\
%		& = \eps \nabla w_- + v\nabla w_- - w_-\nabla v \\
%		& = \eps \nabla w_- + v^2 \nabla z
%	\end{align*}
%	we have, by Young's inequality,
%	\begin{align*}
%		v_\eps^2 \langle A(\nabla u) \nabla z_\eps, \nabla z_\eps \rangle & = \frac{\langle A(\nabla u)(\eps\nabla w_- + v^2\nabla z),\eps\nabla w_- + v^2 \nabla z\rangle}{v_\eps^2} \\
%		& \leq 2\frac{\eps^2}{v_\eps^2} \langle A(\nabla u)\nabla w_-,\nabla w_-\rangle + 2\frac{v^4}{v_\eps^2} \langle A(\nabla u)\nabla z,\nabla z \rangle \\
%		& \leq 2 \langle A(\nabla u)\nabla w_-,\nabla w_-\rangle + 2 v^2 \langle A(\nabla u)\nabla z,\nabla z \rangle \, .
%	\end{align*}
%	By what we already observed, the last term in this chain of inequalities belongs to $L^1_\loc(\Omega)$, so by dominated convergence the RHS of \eqref{CSweps} converges to
%	\[
%		\int_\Omega w_- v \sqrt{\langle A(\nabla u) \nabla z, \nabla z \rangle} \sqrt{\langle A(\nabla u) \nabla \varphi, \nabla \varphi \rangle} \mathbf{1}_{U_\delta}
%	\]
%	as $\eps\to0^+$. Combining everything and recalling that $w_- = zv$ we obtain
%	\[
%		\int_\Omega \varphi v^2 \langle A(\nabla u) \nabla z, \nabla z \rangle \mathbf{1}_{U_\delta} \leq \int_\Omega w_- v \sqrt{\langle A(\nabla u) \nabla z, \nabla z \rangle} \sqrt{\langle A(\nabla u) \nabla \varphi, \nabla \varphi \rangle} \mathbf{1}_{U_\delta}
%	\]
%	and by monotone convergence, letting $\delta\to 0^+$, we have
%	\begin{align*}
%		\int_\Omega \varphi v^2 \langle A(\nabla u) \nabla z, \nabla z \rangle & \leq \int_\Omega w_- v \sqrt{\langle A(\nabla u) \nabla z, \nabla z \rangle} \sqrt{\langle A(\nabla u) \nabla \varphi, \nabla \varphi \rangle} \\
%		& \leq \int_\Omega w_- v \sqrt{\langle A(\nabla u) \nabla z, \nabla z \rangle} \sqrt{\lambda_{\max}(|\nabla u|)} |\nabla\varphi| \, .
%	\end{align*}
%	By H\"older's inequality,
%	\begin{equation} \label{almost_ode}
%		\left(\int_\Omega \varphi v^2 \langle A(\nabla u) \nabla z, \nabla z \rangle\right)^2 \leq \int_\Omega |\nabla\varphi| v^2 \langle A(\nabla u) \nabla z, \nabla z \rangle \cdot \int_\Omega |\nabla\varphi| w_-^2 \lambda_{\max}(|\nabla u|) .
%	\end{equation}
%	Let $G,H : \R^+ \to \R^+_0$ be defined by
%	\[
%		\Theta(r) = \int_{\Omega\cap B_r} v^2 \langle A(\nabla u) \nabla z, \nabla z \rangle \, , \qquad H(r) = \int_{\Omega\cap B_r} w_-^2 \lambda_{\max}(|\nabla u|) \, .
%	\]
%	By what we already observed, the functions $G$ and $H$ are absolutely continuous on compact sub-intervals of $\R^+$, so they are a.e.~differentiable on $\R^+$. For any $r\in\R^+$ such that $G$ and $H$ are both differentiable at $r$, and for each $\rho>0$ we choose $\varphi_\rho\in C^\infty_c(M)$, $\varphi_\rho\geq 0$, such that
%	\[
%		\left\{
%			\begin{array}{r@{\;}c@{\;}ll}
%				\varphi_\rho & = & 1 & \text{on } B_r \\[0.2cm]
%				\varphi_\rho & = & 0 & \text{on } M \backslash B_{r+\rho} \\[0.2cm]
%				|\nabla\varphi_\rho| & \leq & \dfrac{2}{\rho}
%			\end{array}
%		\right.
%	\]
%	Then, applying \eqref{almost_ode} for $\varphi = \varphi_\rho$ and letting $\rho\to0^+$ we obtain
%	\[
%		\Theta(r)^2 \leq G'(r) H'(r) \, .
%	\]
%	Since we are assuming $z\not\equiv 0$ in $\Omega$ but also $z\equiv 0$ on $\partial\Omega\neq\emptyset$, we have $\nabla z\not\equiv 0$ and therefore there exists $r_0 > 0$ such that $\Theta(r_0) > 0$. Therefore, $H'(r) > 0$ for a.e.~$r>r_0$ and we get
%	\[
%		\frac{1}{H'(r)} \leq \frac{G'(r)}{\Theta(r)^2} \qquad \text{for a.e } r > r_0 \, .
%	\]
%	Integrating this inequality, we get
%	\[
%		\int_{r_0}^r \frac{\di s}{H'(s)} \leq \int_{r_0}^r \frac{G'(s)}{\Theta(s)^2} \, \di s = \frac{1}{\Theta(r_0)} - \frac{1}{\Theta(r)} \leq \frac{1}{\Theta(r_0)}
%	\]
%	for all $r>r_0$, so in particular
%	\[
%		\int_{r_0}^{+\infty} \frac{\di s}{H'(s)} < \infty \, .
%	\]
%	By Proposition 1.3 in \cite{rigolisetti}, this further implies that
%	\[
%		\int_{r_0}^{+\infty} \frac{s}{H(s)} \, \di s < \infty \, .
%	\]
%	On the other hand, by assumption \eqref{Hgrowth} we have $H(s) = O(s^2\log s)$, so that
%	\[
%		\frac{s}{H(s)} \geq \frac{C}{s\log s} \qquad \text{for all sufficiently large } \, s > 0,
%	\]
%	contradiction. So, we conclude $z\equiv 0$, that is, $w \geq 0$ in $\Omega$.
%\end{proof}

\begin{theorem}\label{teo_monotonicity}
Let $(M, \metric)$ be a complete Riemannian manifold, let $\Omega$ be a $C^1$ domain and let $u \in C^2(\Omega) \cap C^1(\overline{\Omega})$ be a non-constant, stable solution of
\[
\Delta_a u + f(u) = 0 \qquad \text{in } \Omega, 
\]	
where $a$ satisfies \eqref{assu_A_strong} and $f \in C^1(\R)$. Let $\Omega_0 \subseteq \Omega$ be a subdomain such that $\partial \Omega_0 \subseteq \partial \Omega \cup \crit(u)$. Assume that $X$ is a bounded Killing vector field in $\Omega_0$, continuous in $\overline\Omega_0$ and satisfying 
\begin{equation}\label{eq_bound_Kil}
\langle \nabla u, X \rangle \ge 0, \quad \langle \nabla u, X \rangle \not \equiv 0 \qquad \text{on } \, \overline{\Omega}_0 \cap \partial \Omega. 
\end{equation}
If $u$ has moderate energy growth in $\Omega_0$, then $\langle \nabla u, X \rangle > 0$ in $\Omega_0$.
%there exists $\{\varphi_j\} \subseteq \lip_c(\overline\Omega)$ satisfying, as $j \to \infty$, 
%\[
%	\varphi_j \to 1 \quad \text{in } \, C_\loc(\overline\Omega), \qquad \text{and} \qquad \int_\Omega |\nabla u|^2 \langle A(\nabla u) \nabla \varphi_j, \nabla \varphi_j \rangle \to 0,
%\]
\end{theorem}

\begin{proof}
	By Lemma \ref{lem_linearKilling}, $w = \langle \nabla u,X \rangle$ solves the linearized equation 
	\[
	\diver \left( A(\nabla u)\nabla w\right) + f'(u)w = 0
	\]
	weakly in $\Omega_0$ and $w\geq 0$, $w \not \equiv 0$ on $\partial\Omega_0$. Since, in our assumptions, $w^2 \leq C|\nabla u|^2$ and $u$ has moderate energy growth in $\Omega_0$, we can apply Lemma \ref{Lemma_w-} in $\Omega_0$ to deduce that $w \geq 0$ there. Since $w\not\equiv 0$ on $\partial\Omega_0$, it must be $w>0$ somewhere in $\Omega_0$ and therefore, by the Harnack inequality in \cite[Thm. 7.1.2]{pucci_serrin}, $w>0$ everywhere in $\Omega_0$.
\end{proof}


%\tcr{[La Proposizione seguente viene in parte riassorbita nella prima parte della dimostrazione della Proposizione \ref{prop_Poincare_loc}, dove si prova \eqref{eq_picone} per la scelta particolare $\phi = \varphi|\nabla u|$. Quella prima parte di dimostrazione non è immediatamente sostituibile dalla Proposizione 32 perché richiede di avere una versione di \eqref{eq_picone} localizzata su un sottodominio $\Omega_0\subseteq\Omega$, e per ottenere tale versione della disuguaglianza si sfrutta la presenza di $|\nabla u|$ come fattore di $\phi$ per giustificare il passaggio al limite per $\delta,\tau\to0$ in prossimità delle componenti di $\partial\Omega$ adiacenti a $\Omega \cap \partial\Omega_0$.]}
%
%\medskip
%
%%
%We conclude this section with a stability inequality for test functions which may not vanish on $\partial \Omega$. 
%
%\begin{proposition}
%Let $\Omega \subseteq M$ be a $C^1$ domain with interior normal normal~$\eta$, let $f \in C^1(\R)$ and $a$ satisfy \eqref{assu_A_strong}. If $w \in C^1(\overline{\Omega})$ is a positive solution to 
%\begin{equation}\label{bonito_w}
%\diver\big( A(\nabla u) \nabla w\big) + f'(u) w \le 0 \qquad \text{weakly in $\Omega$.} 
%\end{equation}
%for some $u \in C^1(\overline{\Omega})$, then for each $\eps>0$ and each $\phi \in \lip_c(M)$,
%\begin{equation}\label{eq_picone}
%\begin{array}{l}
%\disp \int_\Omega \Big[ \langle A(\nabla u) \nabla \phi, \nabla \phi \rangle - f'(u) \frac{\phi^2w}{w+\eps}\Big] \\[0.5cm]
%\quad \ge \disp -\int_{\partial \Omega} \frac{\phi^2}{w+\eps} \langle A(\nabla u)\nabla w, \eta \rangle + \int_\Omega (w+\eps)^2 \langle A(\nabla u) \nabla\left(\frac{\phi}{w+\eps}\right),\nabla\left(\frac{\phi}{w+\eps}\right) \rangle. 
%\end{array}
%\end{equation}
%Furthermore, if either $\Omega = M$ or $w > 0$ on $\overline{\Omega}$, one can also take $\eps= 0$ inside the above inequality. The inequality is indeed an equality if $w$ solves \eqref{bonito_w} on $\Omega$ with the equality sign.
%\end{proposition}
%%
%
%%
%
%\begin{proof}
%Let $r(x) = \mathrm{dist}(x, \partial \Omega)$, and let $\tau_0$ be such that $r$ is $C^1$ in $\{r\in [0, \tau_0)\} \cap \overline\Omega$. For $\tau, \delta \in (0, \tau_0/2)$ Consider the cut-off function $\psi$ given by 
%$$
%\psi(x) = \left\{ \begin{array}{ll}
%0 & \quad \text{if } \, r(x) \in (0,\delta) \\[0.3cm]
%\frac{r(x)-\delta}{\tau} & \quad \text{if } \, r(x) \in [\delta, \delta+\tau] \\[0.3cm]
%1 & \quad \text{if } \, r(x) > \delta + \tau.
%\end{array}\right.
%$$
%We integrate $\diver\big( A(\nabla u) \nabla w\big) + f'(u) w \le 0$ against the test function $\psi\phi^2/(w+\eps) \in \lip_c(\Omega)$ to deduce
%\begin{equation}\label{step1}
%\begin{array}{lcl}
%0 & \le  & \disp - \int f'(u) \frac{\psi\phi^2w}{w+\eps} + \int \langle A(\nabla u)\nabla w, \nabla \left( \frac{\psi\phi^2}{w+\eps}\right)\rangle  \\[0.5cm]
%& = & \disp - \int f'(u) \frac{\psi\phi^2w}{w+\eps} + 2 \int \frac{\psi \phi}{w+\eps}\psi \langle A(\nabla u)\nabla w, \nabla \phi\rangle + \int \frac{\phi^2}{w+\eps} \langle A(\nabla u)\nabla w, \nabla \psi \rangle \\[0.5cm]
%& & \disp - \int \frac{\psi \phi^2}{(w+\eps)^2}\psi \langle A(\nabla u) \nabla w, \nabla w \rangle.
%\end{array}
%\end{equation}
%Inserting the identity
%\begin{equation}\label{picone}
%\begin{array}{l}
%\disp (w+\eps)^2 \langle A(\nabla u)\nabla \left(\frac{\phi}{w+\eps}\right),\nabla\left(\frac{\phi}{w+\eps}\right) \rangle \\[0.5cm]
%\disp = \langle A(\nabla u) \nabla \phi, \nabla \phi \rangle + \frac{\phi^2}{(w+\eps)^2}\langle A(\nabla u) \nabla w, \nabla w \rangle - 2 \frac{\phi}{w+\eps} \langle A(\nabla u)\nabla w,\nabla \phi\rangle 
%\end{array}
%\end{equation}
%into \eqref{step1} and using the definition of $\psi$ and the Coarea formula we obtain
%\begin{equation}\label{step2}
%\begin{array}{lcl}
%0 & \le  & \disp  - \int f'(u) \frac{\psi\phi^2w}{w+\eps} + \frac{1}{\tau} \int_{\delta}^{\delta+ \tau} \frac{\phi^2}{w+\eps} \langle A(\nabla u)\nabla w, \nabla r\rangle \\[0.5cm]
%& & \disp + \int \psi \langle A(\nabla u) \nabla \phi, \nabla \phi \rangle - \int \psi(w+\eps)^2 \langle A(\nabla u) \nabla\left(\frac{\phi}{w+\eps}\right),\nabla\left(\frac{\phi}{w+\eps}\right) \rangle. 
%\end{array}
%\end{equation}
%The desired result follows by letting $\delta \ra 0$ and then $\tau \ra 0$, taking into account that the vector field 
%$$
%\frac{\phi^2}{w+\eps} A(\nabla u)\nabla w
%$$
%is continuous in $\overline{\Omega}$ and that $\nabla r = \eta$ on $\partial \Omega$.
%\end{proof}


\section{The geometric Poincar\'e formula and splitting}\label{sec_poinc}
%
%
For $x \not \in \critu$, fix a local orthonormal frame $\{e_i\} = \{\nu, e_\alpha\}$, where $\nu=\nabla u/|\nabla u|$ and $1 \le \alpha \le n-1$. Given $\phi \in C^\infty(M)$, we denote by  $\nabla^\top  \phi$ the component of $\nabla \phi$ orthogonal to $\nu$, and by $\II(x)$ the second fundamental form of the level set $\{u=u(x)\}$ at $x$. The subscript $\nu$ will mean derivation in the direction of $\nu$, so that, for instance, $u_\nu = |\nabla u|$. The coefficients of (the $(0,2)$-version of) $A(X)$ will be denoted by $A^{ij}$.

\begin{lemma}\label{lem_basicomp}
	Let $A$ satisfy \eqref{assu_A}. Near a point where $|\nabla u| \neq 0$, the following identities hold:
	\begin{equation}
		\begin{array}{ll}\label{identities}
			1) & \quad \langle \nabla |\nabla u|, e_i\rangle = \nabla \di  u(\nu, e_i) = \frac{1}{u_\nu} u_{ki}u^k, \qquad \big| \nabla |\nabla u|\big|^2 = u_{\nu\nu}^2 + \big| \nabla^\top  |\nabla u|\big|^2 \\[0.3cm]
			2) & \quad |\nabla \di  u|^2 - \big| \nabla |\nabla u|\big|^2 = \big| \nabla^\top   |\nabla u|\big|^2 + u_\nu^2|\II|^2 \\[0.3cm]
			3) & \quad \langle A(\nabla u)\nabla |\nabla u|, \nabla |\nabla u|\rangle %= \frac{1}{u_\nu^2}A^{ij}u_{ik}u^ku_{jr}u^r 
			= \lambda_2|\nabla^\top  |\nabla u||^2 + \lambda_1 u_{\nu\nu}^2 \\[0.3cm]
			4) & \quad \nabla \di u^{(2)}(\nabla u, \nabla u) = u_\nu^2( u_{\nu\nu}^2 + |\nabla^\top  |\nabla u||^2) \\[0.3cm]
			5) & \quad A^{ij}u_{jk}u_{ir}A^{kr} = 2\lambda_1\lambda_2 \big|\nabla^\top  |\nabla u|\big|^2 + \lambda_1^2u_{\nu\nu}^2 + u_\nu^2\lambda_2^2|\II|^2 \\[0.3cm]
			6) & \quad \langle A(\nabla u) \nabla |\nabla u|, \nabla u \rangle 
			%= A^{ij}\left(\frac{u_{ik}u^k}{u_\nu}\right)u_j
			= u_\nu \lambda_1 u_{\nu\nu}. 
		\end{array}
	\end{equation}
	where $\nabla \di u^{(2)}$ is the tensor whose components are $u_{ik}u^k_j$ and $\lambda_j = \lambda_j(|\nabla u|)$.
\end{lemma}

\begin{remark}
	Item $2)$ was used by Sternberg and Zumbrun in \cite{sternbergzumbrun1}.
\end{remark}
%
%
\begin{proof}
	$1)$ and $2)$ are standard identities. As for $3)$, by $1)$
	$$
	A(\nabla u)\nabla\abs{\nabla u}=\nabla \di   u(\nu,\nu)\lambda_1\nu+\nabla \di   u(\nu,e_\alpha)\lambda_2e_\alpha=u_{\nu\nu}\lambda_1\nu+u_{\nu\alpha}\lambda_2e_\alpha,
	$$
	so
	$$
	\langle A(\nabla u)\nabla\abs{\nabla u},\nabla\abs{\nabla u}\rangle=\lambda_1u_{\nu\nu}^2+\lambda_2u_{\nu\alpha}u_{\nu\alpha}
	$$
	and, by $1)$, $u_{\nu\alpha}u_{\nu\alpha}=\abs{\nabla^\top  \abs{\nabla u}}^2$.\\
	To prove $4)$, observe that
	$$
	\nabla \di u^{(2)}(\nabla u,\nabla u)=u_{ik}u^k_ju^iu^j=u_\nu^2 u_{\nu k}u_{\nu k}=u_\nu^2\left(u_{\nu\nu}^2+\abs{\nabla^\top  \abs{\nabla u}}^2\right).
	$$
	For identity $5)$, we write
	\begin{align*}
		A^{ij}u_{jk}u_{ir}A^{kr}=&\frac{a'(\abs{\nabla u})^2}{\abs{\nabla u}^2}\left[\nabla \di  u(\nabla u,\nabla u)\right]^2+2\frac{a(\abs{\nabla u})a'(\abs{\nabla u})}{\abs{\nabla u}}\nabla \di u^{(2)}(\nabla u,\nabla u)+a(\abs{\nabla u})^2\abs{\nabla \di  u}^2\\
		=&\frac{a'(\abs{\nabla u})^2}{\abs{\nabla u}^2}u_{\nu\nu}^2u_\nu^4+2\frac{a(\abs{\nabla u})a'(\abs{\nabla u})}{\abs{\nabla u}}u_\nu^2\left(u_{\nu\nu}^2+\abs{\nabla^\top  \abs{\nabla u}}^2\right)+\\
		&+a(\abs{\nabla u})^2\left(u_{\nu\nu}^2+2\abs{\nabla^\top  \abs{\nabla u}}^2+u_\nu^2\abs{\operatorname{II}}^2\right)\\
		=&\left[\abs{\nabla u}a'(\abs{\nabla u})+a(\abs{\nabla u})\right]^2u_{\nu\nu}^2+\\
		&+2a(\abs{\nabla u})\left[a(\abs{\nabla u})+\abs{\nabla u}a'(\abs{\nabla u})\right]\abs{\nabla^\top  \abs{\nabla u}}^2+a(\abs{\nabla u})u_\nu^2\abs{\operatorname{II}}^2\\
		=&\lambda_1^2u_{\nu\nu}^2+2\lambda_1\lambda_2\abs{\nabla^\top  \abs{\nabla u}}^2+\lambda_2^2u_\nu^2\abs{\operatorname{II}}^2.
	\end{align*} 
	Finally, to prove $6)$, we just write
	$$
	\langle A(\nabla u)\nabla\abs{\nabla u},\nabla u\rangle=A^{ij}\frac1{u_\nu}u_{ik}u^ku_j=\lambda_1\frac{u^i}{u_\nu}u_{ik}u^k=\frac{\lambda_1}{u_\nu}\nabla \di  u(\nabla u,\nabla u)=\lambda_1u_{\nu\nu}u_\nu^2.
	$$
\end{proof}


We begin with the following Bochner type formulas:

\begin{proposition}
	Let $u \in C^3(\Omega)$, and let $a$ satisfy \eqref{assu_A_superstrong}. Then, at points of $\Omega \backslash \critu$ where $a(|\nabla u|)$ is twice differentiable, it holds
	\begin{equation}\label{bochner_better}
		\begin{array}{lcl}
			\disp \frac{1}{2} \diver\big( A(\nabla u) \nabla |\nabla u|^2 \big) & = & \lambda_1 u_{\nu\nu}^2 + (\lambda_1+\lambda_2) \big|\nabla^\top  |\nabla u|\big|^2 +\lambda_2u_\nu^2|\II|^2 \\[0.3cm]
			& & \disp + \lambda_2\Ric(\nabla u, \nabla u) + \langle \nabla \Delta_a u, \nabla u \rangle.
		\end{array} 
	\end{equation}
	and 
	\begin{equation}\label{bochner_betterfornablau}
		\begin{array}{lcl}
			\disp |\nabla u| \diver\big( A(\nabla u) \nabla |\nabla u| \big) & = & \lambda_1 \big|\nabla^\top  |\nabla u|\big|^2 +\lambda_2u_\nu^2|\II|^2 \\[0.3cm]
			& & \disp + \lambda_2\Ric(\nabla u, \nabla u) + \langle \nabla \Delta_a u, \nabla u \rangle,
		\end{array} 
	\end{equation}
	where $\lambda_j = \lambda_j(|\nabla u|)$. Moreover, if $\Delta_a u + f(u) = 0$ with $f \in C^1(\R)$ then \eqref{bochner_better} and \eqref{bochner_betterfornablau} hold, with the inequality sign $\geq$ in place of equality, weakly in the entire $\Omega$.
\end{proposition}

\begin{remark}\label{rem_stampa}
	The weak inequality requires to give a proper meaning to $\nabla^\top  |\nabla u|$, $u_{\nu\nu}$ and $u^2_\nu |\II|^2$ on $\critu$. By Stampacchia's inequality, $\nabla \di u \equiv 0$ a.e. on $\critu$. Since  
	\[
		u_{\nu\nu}^2 + 2|\nabla^\top |\nabla u||^2 + u_\nu^2|\II|^2 = |\nabla \di u|^2 \qquad \text{where } \, |\nabla u|> 0
	\]
	the terms $|\nabla^\top |\nabla u||^2$, $u_{\nu\nu}^2$ and $u_\nu^2|\II|^2$ extend continuously to the entire $\Omega$ by setting them to zero on $\critu$.
\end{remark}

\begin{proof}
	First, observe that $a(|\nabla u|) \in C^{1,1}_\loc(\{|\nabla u|>0\})$, hence the subset $E$ where it is twice differentiable has full measure in $\{|\nabla u|>0\}$. For $x \in E$ define
	\[
		(\star) = \diver\big( A(\nabla u) \nabla |\nabla u|^2 \big). 
	\]
	Using $1)$ in Lemma \ref{lem_basicomp}, we compute 
	$$
		\begin{array}{lcl}
			(\star) & = & \disp \diver \left( \frac{a'(|\nabla u|)}{|\nabla u|} \langle \nabla u , \nabla |\nabla u|^2 \rangle \nabla u + a(|\nabla u|) \nabla |\nabla u|^2\right) \\[0.4cm]
			& = & \disp \diver \left( 2\frac{a'(|\nabla u|)}{|\nabla u|} \nabla \di  u(\nabla u, \nabla u) \nabla u + a(|\nabla u|) \nabla |\nabla u|^2\right) \\[0.4cm]
			& = & \disp \left( 2\frac{a'(|\nabla u|)}{|\nabla u|}\right)' \frac{\big[\nabla \di  u(\nabla u, \nabla u)\big]^2}{|\nabla u|} + 2\frac{a'(|\nabla u|)}{|\nabla u|} \langle \nabla \big(\nabla \di  u(\nabla u, \nabla u)), \nabla u \rangle\\[0.4cm]
			& & \disp + 2\frac{a'(|\nabla u|)}{|\nabla u|} \nabla \di  u(\nabla u, \nabla u) \Delta u + \frac{a'(|\nabla u|)}{|\nabla u|} \nabla \di  u( \nabla u, \nabla |\nabla u|^2) + a(|\nabla u|) \Delta|\nabla u|^2\\[0.4cm]
			& = & \disp 2\left(\frac{a'(|\nabla u|)}{|\nabla u|}\right)' \frac{\big[\nabla \di  u(\nabla u, \nabla u)\big]^2}{|\nabla u|} + 2\frac{a'(|\nabla u|)}{|\nabla u|} \langle \nabla \big(\nabla \di  u(\nabla u, \nabla u)), \nabla u \rangle\\[0.4cm]
			& & \disp + 2\frac{a'(|\nabla u|)}{|\nabla u|} \nabla \di  u(\nabla u, \nabla u) \Delta u + 2\frac{a'(|\nabla u|)}{|\nabla u|} \nabla \di u^{(2)}(\nabla u, \nabla u) + \\[0.4cm]
			& & + \disp 2a(|\nabla u|) \big[ |\nabla \di  u|^2 + \langle \nabla \Delta u, \nabla u \rangle + \Ric(\nabla u, \nabla u)\big].
		\end{array}
	$$
	Next, 
	\begin{equation}\label{dos}
		\begin{array}{lcl}
			\disp a(|\nabla u|) \langle \nabla \Delta u, \nabla u \rangle & = & \disp \langle \nabla \Big(a(|\nabla u|)\Delta u\Big), \nabla u \rangle - \frac{a'(|\nabla u|)}{|\nabla u|}\nabla \di  u(\nabla u, \nabla u) \Delta u. \\[0.4cm]
		\end{array}
	\end{equation}
	Since 
	$$
		\Delta_a u = \frac{a'(|\nabla u|)}{|\nabla u|}\nabla \di  u(\nabla u, \nabla u) + a(|\nabla u|) \Delta u,
	$$
	inserting it into \eqref{dos} we get
	\begin{equation}\label{tres}
		\begin{array}{lcl}
			\disp a(|\nabla u|) \langle \nabla \Delta u, \nabla u \rangle & = & \disp \langle \nabla \left(-\frac{a'(|\nabla u|)}{|\nabla u|}\nabla \di  u(\nabla u, \nabla u) + \Delta_a u \right), \nabla u \rangle \\[0.4cm]
			& & \disp - \frac{a'(|\nabla u|)}{|\nabla u|}\nabla \di   u(\nabla u, \nabla u) \Delta u. \\[0.4cm]
			& = & \disp - \left(\frac{a'(|\nabla u|)}{|\nabla u|}\right)' \frac{\big[\nabla \di  u(\nabla u, \nabla u)\big]^2}{|\nabla u|} - \frac{a'(|\nabla u|)}{|\nabla u|} \langle \nabla \big( \nabla \di  u(\nabla u, \nabla u) \big), \nabla u \rangle \\[0.4cm]
			& &  \disp + \langle \nabla \left(\Delta_a u\right), \nabla u \rangle - \frac{a'(|\nabla u|)}{|\nabla u|}\nabla \di  u(\nabla u, \nabla u)\Delta u.
		\end{array}
	\end{equation}
	Inserting this in turn into $(\star)$ and simplifying, we obtain that
	$$
		\begin{array}{lcl}
			(\star) & = & \disp 2\frac{a'(|\nabla u|)}{|\nabla u|} \nabla \di u^{(2)}(\nabla u, \nabla u) + \disp 2a(|\nabla u|) \big[ |\nabla \di  u|^2 + \Ric(\nabla u, \nabla u)\big] \\[0.4cm]
			& & + 2\langle \nabla \Delta_a u, \nabla u \rangle.
		\end{array}
	$$
	Using then identities $2)$ and $4)$ in Lemma \ref{lem_basicomp},
	$$
		\begin{array}{lcl}
			(\star) & = & \disp \disp 2 a'(|\nabla u|)|\nabla u|\Big( u_{\nu\nu}^2 + |\nabla^\top   u_\nu|^2\Big) + 2\langle \nabla \Delta_a u, \nabla u \rangle \\[0.4cm]
			& & + \disp 2a(|\nabla u|) \big[u_{\nu\nu}^2 + 2|\nabla^\top   u_\nu|^2 + u_\nu^2|\II|^2 + \Ric(\nabla u, \nabla u)\big] \\[0.4cm]
			& = & \disp 2 \lambda_1 u_{\nu\nu}^2 + 2(\lambda_1+ \lambda_2) |\nabla^\top   u_\nu|^2 + 2\lambda_2 \Big[u_\nu^2|\II|^2 + \Ric(\nabla u, \nabla u)\Big] \\[0.4cm]
			& & \disp + 2\langle \nabla \Delta_a u, \nabla u \rangle,
		\end{array}
	$$
	proving \eqref{bochner_better}. To show \eqref{bochner_betterfornablau}, simply observe that expanding $(\star)$ and using $3)$ in Lemma \ref{lem_basicomp} we get
	$$
		(\star) = 2 |\nabla u| \diver \big( A(\nabla u) \nabla |\nabla u|\big) + 2\lambda_2 |\nabla^\top   u_\nu|^2 + 2 \lambda_1 u_{\nu\nu}^2.
	$$
	To prove that the inequalities hold weakly in $\Omega$, first observe that the vector fields 
	\[
		A(\nabla u)\nabla |\nabla u|^2, \qquad A(\nabla u) \nabla |\nabla u|
	\]
	are locally Lipschitz in $\{|\nabla u|>0\}$. Therefore, the pointwise equalities we just showed hold weakly against test functions supported there. Let now $0 \le \varphi \in \lip_c(\Omega)$. Fix a cutoff $0 \le \psi \in C^\infty_c(\R^+_0)$ supported in $[0,2)$, identically one in $[0,1]$ and with $\psi' \le 0$, and for $\eps > 0$ define the function
	\[
		\eta_\eps(t) = 1-\psi(t/\eps) \qquad \text{for } \, t \in [0,\infty).
	\]
	Let $(\star \star)$ be the right-hand side of \eqref{bochner_better}. We multiply \eqref{bochner_better} by $\eta_\eps(|\nabla u|^2)\varphi \in \lip_c(\{|\nabla u|>0\})$ and integrate by parts to get
	\[
		\begin{array}{lcl}
			\disp \int (\star\star)\eta_\eps(|\nabla u|^2)\varphi & = & \disp - \frac{1}{2} \int \langle A(\nabla u) \nabla |\nabla u|^2, \nabla \left( \eta_\eps(|\nabla u|^2)\varphi\right) \rangle \\[0.5cm]
			& \le & \disp - \frac{1}{2} \int \eta_\eps(|\nabla u|^2) \langle A(\nabla u) \nabla |\nabla u|^2, \nabla\varphi \rangle,  
		\end{array}
	\]
	where we used the ellipticity of $A$ and $\eta_\eps' \ge 0$. Letting $\eps \to 0$ we get
	\[
		\int_{\{|\nabla u|>0\}} (\star\star)\varphi \le \disp - \frac{1}{2} \int_{\{|\nabla u|>0\}} \langle A(\nabla u) \nabla |\nabla u|^2, \nabla\varphi \rangle.  
	\]
	We are left to prove that the domain of integration $\{|\nabla u|>0\}$ can be replaced by the entire $\Omega$. This is clear, in our assumptions on $a$, for the left-hand side. This is also the case of the right-hand side, since 
	\[
		\langle \nabla \Delta_a u, \nabla u \rangle = - f'(u)|\nabla u|^2 \in C(\Omega)
	\]
	and by Remark \ref{rem_stampa}. The case of \eqref{bochner_betterfornablau} is analogous.
\end{proof}

Next, we use the above Bochner's formula to get  a geometric Poincar\'e inequality. A related inequality was obtained in \cite{DPV}. The first step is the following proposition.

\begin{proposition} \label{prop_Poincare_loc}
	Let $\Omega \subseteq M$ be a $C^2$ domain with interior normal $\eta$, and let $u \in C^3(\Omega) \cap C^2(\overline\Omega)$ be a solution to 
	\[
		\Delta_a u + f(u) = 0,
	\]
	where $a$ satisfies \eqref{assu_A_superstrong} and $f \in C^1(\R)$. Let $\Omega_0\subseteq\Omega$ be a domain with locally finite perimeter in $\Omega$ such that 
	\begin{equation}\label{eq_patialomega0}
	\partial\Omega_0 \subseteq \partial\Omega \cup \critu. 
	\end{equation}
	If $w \in C^2(\Omega) \cap C^1(\overline{\Omega})$ solves
	\begin{equation} \label{poinloc0}
		w \ge 0 \ \ \text{in } \, \Omega_0, \qquad \diver \big( A(\nabla u) \nabla w\big) + f'(u)w \le 0 \qquad \text{weakly in } \, \Omega_0,
	\end{equation}
	then, setting $\lambda_j(t)$ as in \eqref{eigenvalues}, for each $\varphi \in \lip_c(M)$ and $\eps>0$ it holds
	\begin{equation} \label{poinloc}
		\begin{split}
			\int_{\Omega_0} |\nabla u|^2 \langle A(\nabla u)\nabla\varphi,\nabla\varphi\rangle & \geq \int_{\Omega_0} \varphi^2 \left\{ \lambda_1|\nabla^\top|\nabla u||^2 + \lambda_2 [|\nabla u|^2|\II|^2 + \Ric(\nabla u,\nabla u)] \right\} \\
			& \phantom{=\;} + \int_{\Omega_0} (w+\eps)^2 \langle A(\nabla u)\nabla\left(\frac{\varphi|\nabla u|}{w+\eps}\right),\nabla\left(\frac{\varphi|\nabla u|}{w+\eps}\right)\rangle \\
			& \phantom{=\;} - \int_{\Omega_0} \frac{\eps \varphi^2}{w+\eps} f'(u)|\nabla u|^2 \\
			& \phantom{=\;} + \int_{\overline\Omega_0\cap\partial\Omega} \varphi^2 \langle A(\nabla u)\nabla\frac{|\nabla u|^2}{2} - \frac{|\nabla u|^2}{w+\eps} A(\nabla u)\nabla w,\eta\rangle \,
		\end{split}
	\end{equation}
	where $\lambda_j = \lambda_j(|\nabla u|)$.
\end{proposition}

\begin{proof}
	Set $w_\eps = w+\eps$, $z = |\nabla u|/w_\eps$ and consider the vector field
	\[
		V = w_\eps |\nabla u| A(\nabla u) \nabla z \equiv \frac{1}{2} A(\nabla u)\nabla|\nabla u|^2 - \frac{|\nabla u|^2}{w_\eps} A(\nabla u)\nabla w
	\]
	which is continuous in $\overline{\Omega}$ and locally Lipschitz continuous in $\Omega$. From \eqref{bochner_betterfornablau} and \eqref{poinloc0} we have
	\[
		\diver\, V = \lambda_1|\nabla^\top|\nabla u||^2 + \lambda_2 [|\nabla u|^2|\II|^2 + \Ric(\nabla u,\nabla u)] - \frac{\eps}{w_\eps} f'(u) |\nabla u|^2 + w_\eps^2 \langle A(\nabla u)\nabla z,\nabla z\rangle \, 
	\]
	weakly in $\Omega$. Let $\varphi\in\lip_c(M)$ be given. We have
	\[
		\diver(\varphi^2 V) = \varphi^2 \diver\, V + 2 w_\eps \varphi |\nabla u| \langle A(\nabla u)\nabla z, \nabla \varphi\rangle
	\]
	and from the identity
	\[
	\begin{array}{lcl}	
	\disp \varphi^2 w_\eps^2 \langle A(\nabla u)\nabla z,\nabla z\rangle + 2 \varphi w_\eps |\nabla u| \langle A(\nabla u)\nabla z,\nabla \varphi\rangle & = & \disp w_\eps^2 \langle A(\nabla u)\nabla (\varphi z),\nabla(\varphi z)\rangle \\[0.3cm]
	& & \disp - |\nabla u|^2 \langle A(\nabla u)\nabla\varphi,\nabla\varphi\rangle
	\end{array}
	\]
	we get
	\begin{equation}
		\begin{split}
			\diver\,(\varphi^2 V) & \ge \varphi^2 \left\{\lambda_1|\nabla^\top|\nabla u||^2 + \lambda_2 [|\nabla u|^2|\II|^2 + \Ric(\nabla u,\nabla u)]\right\} - \frac{\eps \varphi^2}{w_\eps} f'(u)|\nabla u|^2 \\
			& \phantom{=\;} + w_\eps^2 \langle A(\nabla u)\nabla (\varphi z),\nabla(\varphi z)\rangle - |\nabla u|^2 \langle A(\nabla u)\nabla\varphi,\nabla\varphi\rangle \, .
		\end{split}
	\end{equation}
	The vector field $W = \varphi^2 V$ vanishes on $\partial\Omega_0 \cap \Omega$ and is compactly supported in $\overline{\Omega}$, so by the divergence theorem (see Lemma \ref{lem_divergence} in the appendix) we have
	\begin{align*}
		0 & \geq \int_{\Omega_0} \varphi^2 \left\{\lambda_1|\nabla^\top|\nabla u||^2 + \lambda_2 [|\nabla u|^2|\II|^2 + \Ric(\nabla u,\nabla u)]\right\} - \int_{\Omega_0} \frac{\eps \varphi^2}{w_\eps} f'(u)|\nabla u|^2 \\
		& \phantom{=\;} + \int_{\Omega_0} w_\eps^2 \langle A(\nabla u)\nabla (\varphi z),\nabla(\varphi z)\rangle - \int_{\Omega_0} |\nabla u|^2 \langle A(\nabla u)\nabla\varphi,\nabla\varphi\rangle \\
		& \phantom{=\;} + \int_{\overline{\Omega}_0\cap\partial\Omega} \varphi^2 \langle V,\eta\rangle
	\end{align*}
	which is \eqref{poinloc}.
	
%	Let $K\subseteq \{x\in M : \dist(x, {\rm supp}\, \varphi) < 1\}$. Let
%	\[
%		\Psi \ \ : \ \ \R \times \partial\Omega \to M, \qquad \Psi(t,x) = \exp_x(t\eta_x)
%	\]
%	Since $\partial\Omega$ is $C^2$ regular and $K$ is relatively compact and open in $M$, there exists $\tau_0\in(0,1]$ such that the restriction of $\Psi$ to the open set
%	\[
%		E := (-\tau_0,\tau_0) \times (K\cap\partial\Omega)
%	\]
%	is injective and realizes a $C^2$ diffeomorphism onto $\Psi(E)$. %Since $\partial\Omega$ is $C^2$ regular, there exists $\tau_0>0$ such that $r$ is $C^2$ in the tubular neighbourhood $\{x\in M : r(x) < \tau_0\}$ of $\partial\Omega$.
%	From our choice of $\tau_0$, we have that the distance $r$ from $\partial \Omega$ is $C^2$ regular in the set
%	\[
%		U := K \cap \{x\in\Omega : 0 < r(x) < \tau_0\}
%	\]
%	and
%	\begin{equation} \label{rPsi}
%		U\subseteq\Psi(E) \qquad \text{and} \qquad r = \pi_\R \circ (\Psi_{|E})^{-1} \quad \text{in } \, U
%	\end{equation}
%	where $\pi_\R : \R \times \partial\Omega \to \R$ denotes the canonical projection onto the first factor. 
%%		Indeed, by definition of $r$ and by completeness of $M$ and closedness of $\partial\Omega$, for each $x\in U$ there exists a unit speed geodesic arc $\gamma_x : [0,r(x)] \to M$ such that
%%		\[
%%		\gamma_x(0) = x, \qquad \gamma_x([0,r(x))) \subseteq \Omega, \qquad \gamma_x(r(x)) \in \partial\Omega \qquad \text{and} \qquad \dot\gamma_x(r(x)) = \eta_{\gamma_x(r(x))}.
%%		\]
%%		For ease of notation, let us set $y = \gamma_x(r(x))$. Since $\dist(x,y) \leq \mathrm{length}(\gamma_x) = r(x) < \tau_0 \leq 1$, we have $y \in K$ by definition of $K$. But then $(r(x),y) \in E$ and so $\gamma_x(s) = \Psi(r(x)-s,y)$ for all $s\in[0,r(x)]$. In particular, $x = \Psi(r(x),y) \in \Psi(E)$ and this proves \eqref{rPsi}.
%	For $\tau \in (0, \tau_0/2)$, consider the cut-off function $\psi\in\lip(M)$ given by 
%	\begin{equation} \label{def_psi}
%		\psi(x) = \left\{ \begin{array}{ll}
%			0 & \quad \text{if } \, r(x) < \tau \\[0.1cm]
%			\dfrac{r(x)-\tau}{\tau} & \quad \text{if } \, r(x) \in [\tau, 2\tau] \\[0.3cm]
%			1 & \quad \text{if } \, r(x) > 2\tau
%		\end{array}\right.
%	\end{equation}
%	and the vector field
%	\begin{equation} \label{def_V}
%		V = \frac{\varphi^2|\nabla u|^2}{w+\eps} A(\nabla u)\nabla w \, ,
%	\end{equation}
%	which is locally Lipschitz in $\Omega$ and continuous in $\overline\Omega$, with compact support in $\overline{\Omega}$. Applying the divergence theorem to the vector field $\psi V$, which is globally Lipschitz and compactly supported in $\Omega$, and observing that $\psi V \equiv 0$ on the entire $\partial \Omega_0$ because of \eqref{eq_patialomega0}, by the divergence theorem (see \tcr{CITAZIONE}) we get
%	\begin{equation} \label{div_V}
%		0 = \int_{\Omega_0} \diver (\psi V) = \int_{\Omega_0} \langle V, \nabla \psi \rangle  + \int_{\Omega_0} \psi \, \diver V \, .
%	\end{equation}
%	From \eqref{poinloc0} we obtain
%	\[
%	\diver V \le - \frac{\varphi^2|\nabla u|^2}{w+\eps} f'(u)w + \langle A(\nabla u)\nabla w,\nabla\left(\frac{\varphi^2|\nabla u|^2}{w+\eps}\right)\rangle
%	\]
%	so \eqref{div_V} yields
%	\begin{equation} \label{poinloc1}
%		\int_{\Omega_0} \psi \frac{\varphi^2 w}{w+\eps} f'(u)|\nabla u|^2 \le \int_{\Omega_0} \psi \langle A(\nabla u)\nabla w,\nabla\left(\frac{\varphi^2|\nabla u|^2}{w+\eps}\right)\rangle + \int_{\Omega_0} \langle V, \nabla \psi \rangle.
%	\end{equation}
%%	We rewrite the RHS of \eqref{poinloc1} as
%%	\begin{equation} \label{poinloc2}
%%		\int_{\Omega_0} \psi \langle A(\nabla u)\nabla w,\nabla\left(\frac{\varphi^2|\nabla u|^2}{w+\eps}\right)\rangle + \int_{\Omega_0} \frac{\varphi^2|\nabla u|^2}{w+\eps} \langle A(\nabla u)\nabla w,\nabla\psi\rangle,
%%	\end{equation}
%	Using the definition of $\psi$ and the coarea formula,
%	\begin{equation} \label{poinloc3}
%		\int_{\Omega_0} \langle V, \nabla\psi\rangle = \frac{1}{\tau} \int_\tau^{2\tau} \int_{\Omega_0\cap\{r=t\}}  \langle V, \nabla r\rangle \, \di t \, .
%	\end{equation}
%	We claim that passing to limits as $\tau\to0$ it holds
%		\begin{align}
%			\label{poinloc3a}
%			\int_{\Omega_0} \psi \frac{\varphi^2 w}{w+\eps} f'(u)|\nabla u|^2 & \to \int_{\Omega_0} \frac{\varphi^2 w}{w+\eps} f'(u)|\nabla u|^2 \\
%			\label{poinloc3b}
%			\int_{\Omega_0} \psi \langle A(\nabla u)\nabla w,\nabla\left(\frac{\varphi^2|\nabla u|^2}{w+\eps}\right)\rangle & \to \int_{\Omega_0} \langle A(\nabla u)\nabla w,\nabla\left(\frac{\varphi^2|\nabla u|^2}{w+\eps}\right)\rangle \\
%			\label{poinloc3c}
%			\frac{1}{\tau} \int_\tau^{2\tau} \int_{\Omega_0\cap\{r=t\}} 
%			\langle V, \nabla r\rangle \, \di t & \to \int_{\overline{\Omega}_0\cap\partial\Omega} \langle V,\eta\rangle	
%		\end{align}
%		so that from \eqref{poinloc1}-\eqref{poinloc3} we obtain
%		\begin{equation} \label{poinloc4}
%			\int_{\Omega_0} \frac{\varphi^2 w}{w+\eps} f'(u)|\nabla u|^2 \le \int_{\Omega_0} \langle A(\nabla u)\nabla w,\nabla\left(\frac{\varphi^2|\nabla u|^2}{w+\eps}\right)\rangle + \int_{\overline{\Omega}_0\cap\partial\Omega} \langle V,\eta\rangle \, .
%		\end{equation}
%		Claims \eqref{poinloc3a} and \eqref{poinloc3b} follow by the Lebesgue convergence theorem, so we focus on justifying \eqref{poinloc3c}, which requires a bit of care. Fix $b> 0$ and consider 
%		\[
%		F_b = K \cap \partial \Omega \cap \{|\nabla u| \ge b\}, \qquad F^b = K \cap \partial \Omega \cap \{|\nabla u| < b\}
%		\] 
%		Define also $F_b(t) = \Psi(t, F_b)$ and $F^b(t) = \Psi(t, F^b)$, which are disjoint sets being $\Psi(t,\cdot)$ a diffeomorphism for $t \in (0, \tau_0)$. Then, 
%		\[
%		\begin{array}{lcl}
%		\disp \int_{\Omega_0\cap\{r=t\}} \langle V, \nabla r\rangle =  \int_{\Omega_0 \cap F_b(t)} \langle V, \nabla r\rangle + \int_{\Omega_0\cap F^b(t)} \langle V, \nabla r\rangle .
%		\end{array}
%		\]
%		The second integral in the right-hand side can be estimated as follows: using that $|\nabla u| < b$ on $F^b$, $|V| \le C_\eps|\nabla u|^2$ and $\Psi$ is smooth, it holds
%		\[
%		\left| \int_{\Omega_0\cap F^b(t)} \langle V,  \nabla r\rangle \right| \le C_\eps (b+C_1t)^2 \haus^{n-1}(F^b(t)).
%		\]
%		Being $\haus^{n-1}(F^b(t))$ uniformly bounded as $t \in (0, \tau_0)$, we deduce for some constant $C_\eps'$ that 
%		\begin{equation}\label{eq_inte1}
%		\limsup_{\tau\to0} \frac{1}{\tau} \int_\tau^{2\tau} \left| \int_{\Omega_0\cap F^b(t)} \langle V, \nabla r \rangle \right|\di t \le C'_\eps b^2.
%		\end{equation}
%		Regarding the first integral, the inequality $|\nabla u| \ge b$ on $F_b$ guarantees that, for $t$ small enough, $\nabla u\neq 0$ everywhere on $F_b(t)$ and 
%		\[
%		\Omega_0 \cap F_b(t) \equiv \Psi(t, \overline\Omega_0 \cap F_b). 
%		\]
%		(in other words, moving $\overline\Omega_0 \cap F_b$ inwards via the flow $\Psi_t$ one does not escape $\Omega_0$)
%		The family of hypersurfaces $t \mapsto \Omega_0 \cap F_b(t)$ is thus continuous in the $C^1$ topology and converges to $\overline\Omega_0 \cap F_b$ as $t\to0$, and $V$ is continuous on $\overline\Omega$. Therefore, 
%		\[
%		g(t) = \int_{\Omega \cap F_b(t)} \langle V, \nabla r \rangle
%		\]
%		is continuous in a neighbourhood of $t=0$ and converges to
%		\[
%		\int_{\overline{\Omega}_0\cap F_b} \langle V, \eta\rangle
%		\]
%		as $t\to0$. Whence, 
%		\begin{equation}\label{eq_inte2}
%		\lim_{\tau\to0} \frac{1}{\tau} \int_\tau^{2\tau} \int_{\Omega_0\cap F_b(t)} \langle V, \nabla r \rangle \di t = \int_{\overline\Omega_0\cap F_b} \langle V, \eta\rangle.
%		\end{equation}
%		by the mean value theorem. Combining \eqref{eq_inte1} and \eqref{eq_inte2} and letting $b \to 0$ we get \eqref{poinloc3c}.	
%%
%%
%%
%%
%%	The set $K$ defined at the beginning of the proof is relatively compact, so it intersects only finitely many components of $\partial\Omega$, call them $\partial_1\Omega,\dots,\partial_n\Omega$. We set
%%	\[
%%	F_i = K \cap \partial_i\Omega \qquad \text{for } \, i = 1,\dots,n
%%	\]
%%	and for each $t\in(0,\tau_0)$ we also define
%%	\[
%%	F^t_i := \Psi(t,F_i) \qquad \text{for } \, i = 1,\dots,n .
%%	\]
%%	The sets $F_1,\dots,F_n$ are pairwise disjoint and their union equals $K\cap\partial\Omega$, so for each $t\in(0,\tau_0)$ the sets $F^t_i$ are also pairwise disjoint and their union covers $K\cap\Omega\cap\{r=t\}$ by \eqref{rPsi}. Then, for each $t\in(0,\tau_0)$ we have
%%	\[
%%	\int_{\Omega_0\cap\{r=t\}} \frac{\varphi^2|\nabla u|^2}{w+\eps} \langle A(\nabla u)\nabla w, \nabla r\rangle = \sum_{i=1}^n \int_{\Omega_0 \cap F^t_i} \frac{\varphi^2|\nabla u|^2}{w+\eps} \langle A(\nabla u)\nabla w, \nabla r\rangle \, .
%%	\]
%%	Claim \eqref{poinloc3c} will then follow by showing that
%%	\begin{equation} \label{poinloc3i}
%%		\lim_{\tau\to0} \lim_{\delta\to0} \frac{1}{\tau} \int_\delta^{\delta+\tau} \left( \int_{\Omega_0\cap F^t_i} \frac{\varphi^2|\nabla u|^2}{w+\eps} \langle A(\nabla u)\nabla w, \nabla r\rangle \right) \di t = \int_{\overline{\Omega}_0\cap F_i} \frac{\varphi^2|\nabla u|^2}{w+\eps} \langle A(\nabla u)\nabla w, \nabla r\rangle
%%	\end{equation}
%%	for $i=1,\dots,n$. Let $i$ be given. If $\partial_\eta u \equiv 0$ on $\partial_i\Omega$ then for some constant $C>0$ we have $|\nabla u| \leq Ct$ on $F^t_i$ for each $t\in(0,\tau_0)$, since the function $|\nabla u|$ is locally Lipschitz in $\overline{\Omega}$ and vanishes on the bounded set $F_i$, so
%%	\begin{align*}
%%		\left|\int_{F^t_i} \frac{\varphi^2|\nabla u|^2}{w+\eps} \langle A(\nabla u)\nabla w, \nabla r\rangle\right| \leq C' t^2 \mathscr{H}^{n-1}(F^t_i) \leq C'' t^2 \to 0 \qquad \text{as } \, t \to 0
%%	\end{align*}
%%	as $\mathscr{H}^{n-1}(F^t_i)$ remains bounded when $t$ approaches $0$ and then \eqref{poinloc3i} follows since both sides vanish. On the other hand, if $\partial_\eta u\not \equiv  0$ on $\partial_i\Omega$ then for all sufficiently small $t>0$ we have $\nabla u\neq 0$ everywhere on $F^t_i$, hence $\Omega_0 \cap F^t_i = \Omega \cap F^t_i$. The family of hypersurfaces $t \mapsto F^t_i$ is continuous in the $C^1$ topology and converges to $F_i$ as $t\to0$ and the vector field
%%	\[
%%	\frac{\varphi^2|\nabla u|^2}{w+\eps} A(\nabla u)\nabla w
%%	\]
%%	is continuous in $\overline{\Omega}$. Hence, the function
%%	\[
%%	g(t) = \int_{\Omega \cap F^t_i} \frac{\varphi^2|\nabla u|^2}{w+\eps} \langle A(\nabla u)\nabla w, \nabla r\rangle
%%	\]
%%	is continuous on $(0,\tau_0)$ and converges to
%%	\[
%%	\int_{F_i} \frac{\varphi^2|\nabla u|^2}{w+\eps} \langle A(\nabla u)\nabla w, \nabla \eta\rangle \equiv \int_{\overline{\Omega}_0\cap F_i} \frac{\varphi^2|\nabla u|^2}{w+\eps} \langle A(\nabla u)\nabla w, \nabla \eta\rangle
%%	\]
%%	as $t\to0$, so in this case \eqref{poinloc3i} follows by the integral mean value theorem.
%%	
%%	
%%	
%%	\tcb{[Versione precedente] When we let $\delta\to0$ and then $\tau\to0$, the RHS of \eqref{poinloc3} converges to
%%		\[
%%		\int_{\overline{\Omega}_0\cap\partial\Omega} \frac{\varphi^2|\nabla u|^2}{w+\eps} \langle A(\nabla u)\nabla w,\eta\rangle \, .
%%		\]
%%		while the LHS of \eqref{poinloc1} and the first integral in \eqref{poinloc2} converge to
%%		\[
%%		\int_{\Omega_0} \varphi^2 \frac{w}{w+\eps} f'(u)|\nabla u|^2 \qquad \text{and} \qquad \int_{\Omega_0} \langle A(\nabla u)\nabla w,\nabla\left(\frac{\varphi^2|\nabla u|^2}{w+\eps}\right)\rangle
%%		\]
%%		respectively, by the dominated convergence theorem. Hence, we obtain
%%		\begin{equation} \label{poinloc4}
%%			\int_{\Omega_0} \frac{\varphi^2 w}{w+\eps} f'(u)|\nabla u|^2 \le \int_{\Omega_0} \langle A(\nabla u)\nabla w,\nabla\left(\frac{\varphi^2|\nabla u|^2}{w+\eps}\right)\rangle + \int_{\overline{\Omega}_0\cap\partial\Omega} \frac{\varphi^2|\nabla u|^2}{w+\eps} \langle A(\nabla u)\nabla w,\eta\rangle \, .
%%	\end{equation}}
%	
%	We next consider the vector field
%	\begin{equation} \label{def_W}
%		W = \frac{1}{2} \varphi^2 A(\nabla u)\nabla|\nabla u|^2
%	\end{equation}
%	which, as $V$, is locally Lipschitz in $\Omega$ and continuous on $\overline\Omega$. Applying the divergence theorem to $\psi W$, which is Lipschitz and compactly supported in $\Omega$ and vanishes on $\partial \Omega_0$, we deduce
%	\begin{equation} \label{div_W}
%		0 = \int_{\Omega_0} \diver (\psi W) = \int_{\Omega_0} \psi \, \diver W + \int_{\Omega_0} \langle W, \nabla \psi \rangle 
%%		 -\int_{\overline{\Omega}_0\cap\partial\Omega} \langle W,\eta\rangle
%	\end{equation}
%	The term with $\langle W, \nabla \psi \rangle$ can be treated exactly as we did before for $\langle V, \nabla \psi \rangle$ (note that $|W| \le C|\nabla u|$, which is enough to follow the argument verbatim), yielding to 
%	\[
%		\lim_{\tau \to 0} \int_{\Omega_0} \langle W, \nabla \psi \rangle =  \int_{\overline{\Omega}_0\cap\partial\Omega} \langle W,\eta\rangle.
%	\] 
%	By the Bochner formula \eqref{bochner_betterfornablau} and using that $\Delta_a u = -f(u)$ we have
%	\begin{align*}
%		\diver\, W & = \varphi^2 |\nabla u|\diver(A(\nabla u)\nabla|\nabla u|) + \langle A(\nabla u)\nabla|\nabla u|, \nabla\left(\varphi^2|\nabla u|\right)\rangle \\
%		& \geq \varphi^2 \left\{ \lambda_1|\nabla^\top|\nabla u||^2 + \lambda_2 [|\nabla u|^2|\II|^2 + \Ric(\nabla u,\nabla u)] \right\} - \varphi^2 f'(u)|\nabla u|^2 \\
%		& \phantom{=\;} + \langle A(\nabla u)\nabla|\nabla u|, \nabla\left(\varphi^2|\nabla u|\right)\rangle
%	\end{align*}
%	in $\Omega$, so from \eqref{div_W} we get
%	\begin{equation} \label{poinloc4b}
%		\begin{split}
%			\int_{\Omega_0} \varphi^2 f'(u)|\nabla u|^2 & \geq \int_{\Omega_0} \varphi^2 \left\{ \lambda_1|\nabla^\top|\nabla u||^2 + \lambda_2 [|\nabla u|^2|\II|^2 + \Ric(\nabla u,\nabla u)] \right\} \\
%			& \phantom{=\;} + \int_{\Omega_0} \langle A(\nabla u)\nabla|\nabla u|, \nabla\left(\varphi^2|\nabla u|\right)\rangle + \int_{\overline\Omega_0\cap\partial\Omega} \langle W,\eta\rangle \, .
%		\end{split}
%	\end{equation}
%	Subtracting \eqref{poinloc4} from \eqref{poinloc4b} we obtain
%	\begin{equation} \label{poinloc5}
%		\begin{split}
%			\int_{\Omega_0} \frac{\eps \varphi^2}{w+\eps} f'(u)|\nabla u|^2 & \geq \int_{\Omega_0} \varphi^2 \left\{ \lambda_1|\nabla^\top|\nabla u||^2 + \lambda_2 [|\nabla u|^2|\II|^2 + \Ric(\nabla u,\nabla u)] \right\} \\
%			& \phantom{=\;} + \int_{\Omega_0} \langle A(\nabla u)\nabla|\nabla u|, \nabla(\varphi^2|\nabla u|)\rangle - \int_{\Omega_0} \langle A(\nabla u)\nabla w,\nabla\left(\frac{\varphi^2|\nabla u|^2}{w+\eps}\right)\rangle \\
%			& \phantom{=\;} + \int_{\overline\Omega_0\cap\partial\Omega} \langle W-V,\eta\rangle \, .
%		\end{split}
%	\end{equation}
%	Direct computations yield
%	\begin{align*}
%		\langle A(\nabla u)\nabla|\nabla u|, \nabla(\varphi^2|\nabla u|)\rangle & = \langle A(\nabla u)\nabla(\varphi|\nabla u|),\nabla(\varphi|\nabla u|)\rangle \\
%		& \phantom{=\;} - |\nabla u|^2 \langle A(\nabla u)\nabla\varphi,\nabla\varphi\rangle
%	\end{align*}
%	and
%	\begin{align*}
%		\langle A(\nabla u)\nabla w,\nabla\left(\frac{\varphi^2|\nabla u|^2}{w+\eps}\right)\rangle & = \langle A(\nabla u)\nabla(\varphi|\nabla u|),\nabla(\varphi|\nabla u|)\rangle \\
%		& \phantom{=\;} - (w+\eps)^2 \langle A(\nabla u)\nabla\left(\frac{\varphi|\nabla u|}{w+\eps}\right),\nabla\left(\frac{\varphi|\nabla u|}{w+\eps}\right)\rangle \, .
%	\end{align*}
%	%	\begin{align*}
%		%		\langle A(\nabla u)\nabla w,\nabla\left(\frac{\varphi^2|\nabla u|^2}{w+\eps}\right)\rangle & = \varphi|\nabla u|\langle A(\nabla u)\nabla w,\nabla\left(\frac{\varphi|\nabla u|}{w+\eps}\right)\rangle \\
%		%		& \phantom{=\;} + \frac{\varphi|\nabla u|}{w+\eps} \langle A(\nabla u)\nabla w,\nabla(\varphi|\nabla u|)\rangle \\
%		%		& = -(w+\eps)^2 \langle A(\nabla u)\nabla\left(\frac{\varphi|\nabla u|}{w+\eps}\right),\nabla\left(\frac{\varphi|\nabla u|}{w+\eps}\right)\rangle \\
%		%		& \phantom{=\;} + \langle A(\nabla u)\nabla(\varphi|\nabla u|),\nabla(\varphi|\nabla u|)\rangle \, .
%		%	\end{align*}
%	Substituting these expressions in \eqref{poinloc5} and then rearranging terms we finally obtain \eqref{poinloc}.
%	%	\begin{align*}
%		%		0 & \geq \int_{\Omega_0} \frac{\eps \varphi^2}{w+\eps} f'(u)|\nabla u|^2 + \int_{\Omega_0} \varphi^2 \left\{ \lambda_1|\nabla^\top|\nabla u||^2 + \lambda_2 [|\nabla u|^2|\II|^2 + \Ric(\nabla u,\nabla u)] \right\} \\
%		%		& \phantom{=\;} + \int_{\Omega_0} (w+\eps)^2 \langle A(\nabla u)\nabla\left(\frac{\varphi|\nabla u|}{w+\eps}\right),\nabla\left(\frac{\varphi|\nabla u|}{w+\eps}\right)\rangle - \int_{\Omega_0} |\nabla u|^2 \langle A(\nabla u)\nabla\varphi,\nabla\varphi\rangle \\
%		%		& \phantom{=\;} + \int_{\overline\Omega_0\cap\partial\Omega} \varphi^2 \langle A(\nabla u)\nabla\frac{|\nabla u|^2}{2},\eta\rangle - \int_{\overline\Omega_0\cap\partial\Omega} \frac{\varphi^2|\nabla u|^2}{w+\eps} \langle A(\nabla u)\nabla w,\eta\rangle
%		%	\end{align*}
%	%	that is, by rearranging terms, \eqref{poinloc}.
\end{proof}

%\begin{proposition}\label{prop_poincare}
%	Let $\Omega \subseteq M$ be a $C^2$ domain with interior normal $\eta$, and let $u \in C^3(\Omega) \cap C^2(\overline\Omega)$ be a stable solution to 
%	\[
%		\Delta_a u + f(u) = 0,
%	\]
%	where $a$ satisfies \eqref{assu_A_superstrong} and $f \in C^1(\R)$. If $w \in C^2(\Omega) \cap C^1(\overline{\Omega})$ solves
%	\[
%		w \ge 0 \ \ \text{in } \, \Omega, \qquad \diver \big( A(\nabla u) \nabla w\big) + f'(u)w \le 0 \quad \text{weakly in } \, \Omega
%	\]
%	then, setting $\lambda_j(t)$ as in \eqref{eigenvalues}, for each $\varphi \in \lip_c(M)$ and $\eps>0$ it holds
%	\begin{equation}\label{piufinalissima_teo}
%		\begin{array}{l}
%		\disp \int_{\Omega} \frac{\varphi^2 w}{w+\eps}\Big\{\lambda_1 \big|\nabla^\top   u_\nu\big|^2 +\lambda_2\Big[u_\nu^2|\II|^2 +\Ric(\nabla u, \nabla u)\Big]\Big\} \\[0,5cm]
%		\le \disp \int_{\partial \Omega}\frac{\varphi^2}{w+\eps} \left\{  |\nabla u|^2  \langle A(\nabla u)\nabla w, \eta\rangle - w\langle A(\nabla u) \nabla \frac{|\nabla u|^2}{2} , \eta \rangle \right\} \\[0.5cm]
%		\quad \disp - \frac{\eps}{2} \int_{\Omega} \frac{\varphi^2}{(w+\eps)^2} \langle A(\nabla u)\nabla w, \nabla |\nabla u|^2 \rangle \\[0.5cm]
%		\disp + \int_{\Omega} \frac{w}{w+\eps}|\nabla u|^2 \langle A(\nabla u) \nabla \varphi, \nabla \varphi \rangle \\[0.5cm]
%		\disp - \int_{\Omega} (w+\eps)^2 \langle A(\nabla u) \nabla\left(\frac{\varphi|\nabla u|}{w+\eps}\right),\nabla\left(\frac{\varphi|\nabla u|}{w+\eps}\right) \rangle \\[0.5cm]
%		\disp + \int_{\Omega} \frac{\eps}{w+\eps} \disp \langle A(\nabla u) \nabla \big(\varphi |\nabla u|\big), \nabla \big(\varphi |\nabla u|\big)\rangle
%	\end{array}
%\end{equation}
%%$$
%%\begin{array}{l}
%%\disp \int_\Omega \lambda_2(|\nabla u|) \big[ |\nabla u|^2|II|^2 + \Ric(\nabla u, \nabla u)\big]\frac{
%%\varphi^2w}{w+\eps} + \int_\Omega \lambda_1(|\nabla u|) \big|\nabla^\top  |\nabla u|\big|^2\varphi^2 \\[0.5cm]
%%\quad \le \disp \int_{\partial\Omega}\frac{\varphi^2}{w+\eps} \left[w \langle A(\nabla u) \nabla \left(\frac{|\nabla u|^2}{2}\right), \eta \rangle - |\nabla u|^2 \langle A(\nabla u)\nabla w, \eta \rangle \right] \\[0.5cm] 
%%\disp + \eps \int_{\Omega} \frac{\varphi}{w+\eps} \langle A(\nabla u)\nabla \varphi, \nabla |\nabla u|^2 \rangle - \frac 12 \int_\Omega \varphi^2 \langle A(\nabla u) \nabla |\nabla u|^2, \nabla \left(\frac{w}{w+\eps}\right)\rangle \\[0.5cm]
%%\disp \disp + \int_\Omega |\nabla u|^2 \langle A(\nabla u)\nabla \varphi, \nabla \varphi \rangle - \int_\Omega (w+\eps)^2 \langle A(\nabla u) \left(\frac{\varphi|\nabla u|}{w+\eps}\right),\left(\frac{\varphi|\nabla u|}{w+\eps}\right) \rangle.
%%\end{array}
%%\end{equation}
%where $\lambda_j = \lambda_j(|\nabla u|)$.
%\end{proposition}
%%
%%
%\begin{proof}
%%We apply the Gauss-Green theorem on $\Omega_0$ to the Lipschitz vector field
%	%	\[
%	%		V = \frac{\varphi^2|\nabla u|^2}{w+\eps} A(\nabla u)\nabla w - \frac{1}{2} \frac{\varphi^2 w}{w+\eps} A(\nabla u)\nabla|\nabla u|^2
%	%	\]
%	%	Let $r(x) = \mathrm{dist}(x, \partial \Omega)$, and let $\tau_0$ be such that $r$ is $C^1$ in $\{r\in [0, \tau_0)\} \cap \overline\Omega$. For $\tau, \delta \in (0, \tau_0/2)$, consider the cut-off function $\psi\in\lip(M)$ given by 
%	%	\[
%	%	\psi(x) = \left\{ \begin{array}{ll}
%		%		0 & \quad \text{if } \, r(x) \in (0,\delta) \\[0.3cm]
%		%		\frac{r(x)-\delta}{\tau} & \quad \text{if } \, r(x) \in [\delta, \delta+\tau] \\[0.3cm]
%		%		1 & \quad \text{if } \, r(x) > \delta + \tau.
%		%	\end{array}\right.
%	%	\]
%	%	We integrate $\diver\big( A(\nabla u) \nabla w\big) + f'(u) w \le 0$ against the test function $\psi\phi^2/(w+\eps) \in \lip_c(\Omega)$ to deduce
%	%	\begin{equation}\label{step1}
%		%		\begin{array}{lcl}
%			%			0 & \le  & \disp - \int f'(u) \frac{\psi\phi^2w}{w+\eps} + \int \langle A(\nabla u)\nabla w, \nabla \left( \frac{\psi\phi^2}{w+\eps}\right)\rangle  \\[0.5cm]
%			%			& = & \disp - \int f'(u) \frac{\psi\phi^2w}{w+\eps} + 2 \int \frac{\psi \phi}{w+\eps}\psi \langle A(\nabla u)\nabla w, \nabla \phi\rangle + \int \frac{\phi^2}{w+\eps} \langle A(\nabla u)\nabla w, \nabla \psi \rangle \\[0.5cm]
%			%			& & \disp - \int \frac{\psi \phi^2}{(w+\eps)^2}\psi \langle A(\nabla u) \nabla w, \nabla w \rangle.
%			%		\end{array}
%		%	\end{equation}
%	%	Inserting the identity
%	%	\begin{equation}\label{picone}
%		%		\begin{array}{l}
%			%			\disp (w+\eps)^2 \langle A(\nabla u)\nabla \left(\frac{\phi}{w+\eps}\right),\nabla\left(\frac{\phi}{w+\eps}\right) \rangle \\[0.5cm]
%			%			\disp = \langle A(\nabla u) \nabla \phi, \nabla \phi \rangle + \frac{\phi^2}{(w+\eps)^2}\langle A(\nabla u) \nabla w, \nabla w \rangle - 2 \frac{\phi}{w+\eps} \langle A(\nabla u)\nabla w,\nabla \phi\rangle 
%			%		\end{array}
%		%	\end{equation}
%	%	into \eqref{step1} and using the definition of $\psi$ and the Coarea formula we obtain
%	%	\begin{equation}\label{step2}
%		%		\begin{array}{lcl}
%			%			0 & \le  & \disp  - \int f'(u) \frac{\psi\phi^2w}{w+\eps} + \frac{1}{\tau} \int_{\delta}^{\delta+ \tau} \frac{\phi^2}{w+\eps} \langle A(\nabla u)\nabla w, \nabla r\rangle \\[0.5cm]
%			%			& & \disp + \int \psi \langle A(\nabla u) \nabla \phi, \nabla \phi \rangle - \int \psi(w+\eps)^2 \langle A(\nabla u) \nabla\left(\frac{\phi}{w+\eps}\right),\nabla\left(\frac{\phi}{w+\eps}\right) \rangle. 
%			%		\end{array}
%		%	\end{equation}
%	%	The desired result follows by letting $\delta \ra 0$ and then $\tau \ra 0$, taking into account that the vector field 
%	%	$$
%	%	\frac{\phi^2}{w+\eps} A(\nabla u)\nabla w
%	%	$$
%	%	is continuous in $\overline{\Omega}$ and that $\nabla r = \eta$ on $\partial \Omega$.
%	Let $r\in\lip(M)$ be the distance function from $\partial\Omega$, that is, $r(x) = \mathrm{dist}(x, \partial \Omega)$ for each $x\in M$. Since $\partial\Omega$ is $C^3$ regular, there exists $\tau_0>0$ such that $r$ is $C^3$ in the tubular neighbourhood $\{x\in M : r(x) < \tau_0\}$ of $\partial\Omega$. For $\delta,\tau \in (0, \tau_0/2)$, consider the cut-off function $\psi\in\lip(M)$ given by 
%	\begin{equation} \label{defpsi}
%	\psi(x) = \left\{ \begin{array}{ll}
%		0 & \quad \text{if } \, r(x) < \delta \\[0.3cm]
%		\dfrac{r(x)-\delta}{\tau} & \quad \text{if } \, r(x) \in [\delta, \delta+\tau] \\[0.3cm]
%		1 & \quad \text{if } \, r(x) > \delta + \tau
%	\end{array}\right.
%	\end{equation}
%	%	For each sufficiently small $\tau,\delta>0$ let $\psi\in\lip_c(\Omega)$ be the cut-off function as in \tcr{the proof of Proposition 27}. By Lemma \ref{Om_finper} the set $\Omega_0$ has locally finite perimeter, so we apply the Gauss-Green theorem to the Lipschitz vector field
%	and the compactly supported Lipschitz vector field $V$ defined in $\overline{\Omega}$ by
%	\begin{equation} \label{defV}
%		V = \frac{\psi\varphi^2|\nabla u|^2}{w+\eps} A(\nabla u)\nabla w \, .
%	\end{equation}
%	Since $V$ is compactly supported, by the divergence theorem we have
%	\begin{equation} \label{divV}
%		\int_{\Omega} \diver\, V = 0 \, .
%	\end{equation}
%	Since $w$ satisfies
%	\begin{equation} \label{poinalt0}
%		\diver(A(\nabla u)\nabla w) + f'(u)w \le 0 \qquad \text{weakly in } \, \Omega
%	\end{equation}
%	we have
%	\[
%		\diver\,V \le - \frac{\psi\varphi^2|\nabla u|^2}{w+\eps} f'(u)w + \langle A(\nabla u)\nabla w,\nabla\left(\frac{\psi\varphi^2|\nabla u|^2}{w+\eps}\right)\rangle \, ,
%	\]
%	so \eqref{divV} yields
%	\begin{equation} \label{poinalt1}
%		\int_\Omega \psi \frac{\varphi^2 w}{w+\eps} f'(u)|\nabla u|^2 \le  \int_\Omega \langle A(\nabla u)\nabla w,\nabla\left(\frac{\psi\varphi^2|\nabla u|^2}{w+\eps}\right)\rangle \, .
%	\end{equation}
%	We rewrite the RHS of \eqref{poinalt1} as
%	\begin{equation} \label{poinalt2}
%		\int_\Omega \psi \langle A(\nabla u)\nabla w,\nabla\left(\frac{\varphi^2|\nabla u|^2}{w+\eps}\right)\rangle + \int_\Omega \frac{\varphi^2|\nabla u|^2}{w+\eps} \langle A(\nabla u)\nabla w,\nabla\psi\rangle \, .
%	\end{equation}
%	Using the definition of $\psi$ and the coarea formula we have
%	\begin{equation} \label{poinalt3}
%		\int_\Omega \frac{\varphi^2|\nabla u|^2}{w+\eps} \langle A(\nabla u)\nabla w,\nabla\psi\rangle = \frac{1}{\tau} \int_\delta^{\delta+\tau} \left( \int_{\Omega\cap\{r=t\}} \frac{\varphi^2|\nabla u|^2}{w+\eps} \langle A(\nabla u)\nabla w, \nabla r\rangle \right) \di t \, .
%	\end{equation}
%	When letting $\delta\to0$ and then $\tau\to0$, the RHS of \eqref{poinalt3} converges to
%	\[
%		\int_{\partial\Omega} \frac{\varphi^2|\nabla u|^2}{w+\eps} \langle A(\nabla u)\nabla w,\eta\rangle
%	\]
%	while the LHS of \eqref{poinalt1} and the first integral in \eqref{poinalt2} converge to
%	\[
%		\int_\Omega \varphi^2 \frac{w}{w+\eps} f'(u)|\nabla u|^2 \qquad \text{and} \qquad \int_\Omega \langle A(\nabla u)\nabla w,\nabla\left(\frac{\varphi^2|\nabla u|^2}{w+\eps}\right)\rangle
%	\]
%	respectively, by the dominated convergence theorem. Hence, we obtain
%	\begin{equation} \label{poinalt4}
%		\int_\Omega \frac{\varphi^2 w}{w+\eps} f'(u)|\nabla u|^2 \le \int_\Omega \langle A(\nabla u)\nabla w,\nabla\left(\frac{\varphi^2|\nabla u|^2}{w+\eps}\right)\rangle + \int_{\partial\Omega} \frac{\varphi^2|\nabla u|^2}{w+\eps} \langle A(\nabla u)\nabla w,\eta\rangle \, .
%	\end{equation}
%	%	that we rewrite as
%	%	\begin{align*}
%		%		\int_{\Omega_0} \langle A(\nabla u)\nabla(\varphi|\nabla u|),\nabla(\varphi|\nabla u|)\rangle & = \int_{\Omega_0} \varphi^2 \frac{w}{w+\eps} f'(u)|\nabla u|^2 - \int_{\Gamma_0} \frac{\varphi^2|\nabla u|^2}{w+\eps} \langle A(\nabla u)\nabla w,\eta\rangle \\
%		%		& \phantom{=\;} + \int_{\Omega_0} (w+\eps)^2 \langle A(\nabla u)\nabla\left(\frac{\varphi|\nabla u|}{w+\eps}\right),\nabla\left(\frac{\varphi|\nabla u|}{w+\eps}\right)\rangle \, .
%		%	\end{align*}
%	
%	We next apply the divergence theorem to the Lipschitz vector field
%	\begin{equation} \label{defW}
%		W = \frac{1}{2} \frac{\varphi^2 w}{w+\eps} A(\nabla u)\nabla|\nabla u|^2
%	\end{equation}
%	to obtain
%	\begin{equation} \label{divW}
%		\int_\Omega \diver\, W = -\int_{\partial\Omega} \langle W,\eta\rangle
%	\end{equation}
%	where we used the fact that $W=0$ on the portion $\partial\Omega_0\backslash\partial\Omega \subseteq \critu$ of $\partial\Omega_0$. By the Bochner formula \eqref{bochner_betterfornablau} and using that $\Delta_a u = -f(u)$ we have
%	\begin{align*}
%		\diver\, W & = \frac{\varphi^2 w}{w+\eps} |\nabla u|\diver(A(\nabla u)\nabla|\nabla u|) + \langle A(\nabla u)\nabla|\nabla u|, \nabla\left(\frac{\varphi^2|\nabla u|w}{w+\eps}\right)\rangle \\
%		& = \frac{\varphi^2 w}{w+\eps} \left\{ \lambda_1|\nabla^\top|\nabla u||^2 + \lambda_2 [|\nabla u|^2|\II|^2 + \Ric(\nabla u,\nabla u)] \right\} - \varphi^2 \frac{w}{w+\eps} f'(u)|\nabla u|^2 \\
%		& \phantom{=\;} + \langle A(\nabla u)\nabla|\nabla u|, \nabla\left(\frac{\varphi^2|\nabla u|w}{w+\eps}\right)\rangle
%	\end{align*}
%	in $\Omega$, so from \eqref{divW} we get
%	\begin{align*}
%		\int_\Omega \frac{\varphi^2 w}{w+\eps} f'(u)|\nabla u|^2 & = \int_\Omega \frac{\varphi^2 w}{w+\eps} \left\{ \lambda_1|\nabla^\top|\nabla u||^2 + \lambda_2 [|\nabla u|^2|\II|^2 + \Ric(\nabla u,\nabla u)] \right\} \\
%		& \phantom{=\;} + \int_\Omega \langle A(\nabla u)\nabla|\nabla u|, \nabla\left(\frac{\varphi^2|\nabla u|w}{w+\eps}\right)\rangle \\
%		& \phantom{=\;} + \frac{1}{2} \int_{\partial\Omega} \varphi^2 \frac{w}{w+\eps} \langle A(\nabla u)\nabla|\nabla u|^2,\eta\rangle \, .
%	\end{align*}
%	Combining this with \eqref{poinalt4} we obtain
%	\begin{equation} \label{poinalt5}
%		\begin{split}
%			0 & \ge \int_\Omega \frac{\varphi^2 w}{w+\eps} \left\{ \lambda_1|\nabla^\top|\nabla u||^2 + \lambda_2 [|\nabla u|^2|\II|^2 + \Ric(\nabla u,\nabla u)] \right\} \\
%			& \phantom{=\;} + \int_\Omega \langle A(\nabla u)\nabla|\nabla u|, \nabla\left(\frac{\varphi^2|\nabla u|w}{w+\eps}\right)\rangle - \langle A(\nabla u)\nabla w,\nabla\left(\frac{\varphi^2|\nabla u|^2}{w+\eps}\right)\rangle \\
%			& \phantom{=\;} + \int_{\partial\Omega} \frac{\varphi^2}{w+\eps} \langle w A(\nabla u)\nabla\frac{|\nabla u|^2}{2} - |\nabla u|^2 A(\nabla u)\nabla w,\eta\rangle \, .
%		\end{split}
%	\end{equation}
%	Direct computations yield
%	\begin{align*}
%		\langle A(\nabla u)\nabla|\nabla u|, \nabla\left(\frac{\varphi^2|\nabla u|w}{w+\eps}\right)\rangle & = \varphi^2 |\nabla u| \langle A(\nabla u)\nabla|\nabla u|,\nabla\frac{w}{w+\eps}\rangle \\
%		& \phantom{=\;} + \frac{w}{w+\eps} \langle A(\nabla u)\nabla|\nabla u|,\nabla(\varphi^2|\nabla u|)\rangle \\
%		& = - \frac{\eps}{2} \frac{\varphi^2}{(w+\eps)^2} \langle A(\nabla u)\nabla|\nabla u|^2,\nabla w\rangle \\
%		& \phantom{=\;} + \frac{w}{w+\eps} \langle A(\nabla u)\nabla(\varphi|\nabla u|),\nabla(\varphi|\nabla u|)\rangle \\
%		& \phantom{=\;} - \frac{w}{w+\eps} |\nabla u|^2 \langle A(\nabla u)\nabla\varphi,\nabla\varphi\rangle
%	\end{align*}
%	and
%	\begin{align*}
%		\langle A(\nabla u)\nabla w,\nabla\left(\frac{\varphi^2|\nabla u|^2}{w+\eps}\right)\rangle & = \varphi|\nabla u|\langle A(\nabla u)\nabla w,\nabla\left(\frac{\varphi|\nabla u|}{w+\eps}\right)\rangle \\
%		& \phantom{=\;} + \frac{\varphi|\nabla u|}{w+\eps} \langle A(\nabla u)\nabla w,\nabla(\varphi|\nabla u|)\rangle \\
%		& = -(w+\eps)^2 \langle A(\nabla u)\nabla\left(\frac{\varphi|\nabla u|}{w+\eps}\right),\nabla\left(\frac{\varphi|\nabla u|}{w+\eps}\right)\rangle \\
%		& \phantom{=\;} + \langle A(\nabla u)\nabla(\varphi|\nabla u|),\nabla(\varphi|\nabla u|)\rangle
%	\end{align*}
%Plugging into \eqref{poinalt5} we finally get \eqref{piufinalissima_teo}. 
%\end{proof}





%
%
%
%%%%%%%%%%%%%%%%     OLD PROOF
%
%
%
%\begin{proof}
%We begin by testing \eqref{eq_picone} with $\varphi|\nabla u|$:
%\begin{equation}\label{eq_picone_2}
%\begin{array}{l}
%\disp \int_\Omega \Big[ \langle A(\nabla u) \nabla \big(\varphi |\nabla u|\big), \nabla \big(\varphi |\nabla u|\big)\rangle - f'(u) \frac{\varphi^2|\nabla u|^2 w}{w+\eps}\Big] \\[0.5cm]
%\quad \ge - \disp \int_{\partial \Omega} \frac{\varphi^2|\nabla u|^2}{w+\eps} \langle A(\nabla u)\nabla w, \eta\rangle + \int_\Omega (w+\eps)^2 \langle A(\nabla u) \nabla\left(\frac{\varphi|\nabla u|}{w+\eps}\right),\nabla\left(\frac{\varphi|\nabla u|}{w+\eps}\right) \rangle. 
%\end{array}
%\end{equation}
%We integrate the Bochner formula \eqref{bochner_better} against the Lipschitz function 
%$$
%\bar \varphi^2 = \varphi^2 \frac{w}{w+\eps} 
%$$
%to deduce 
%$$
%\begin{array}{l}
%\disp \int_{\Omega} \bar \varphi^2\Big\{\lambda_1 u_{\nu\nu}^2 + (\lambda_1+\lambda_2) \big|\nabla^\top   u_\nu\big|^2 +\lambda_2\Big[u_\nu^2|\II|^2 +\Ric(\nabla u, \nabla u)\Big]\Big\} \\[0,5cm]
%= \disp \frac{1}{2} \int_\Omega \bar \varphi^2 \diver\big( A(\nabla u) \nabla |\nabla u|^2 \big) + \int_\Omega \bar \varphi^2 f'(u)|\nabla u|^2 \\[0.5cm]
%= \disp - \int_{\partial \Omega} \bar{\varphi}^2 \langle A(\nabla u) \nabla \frac{|\nabla u|^2}{2} , \eta \rangle - \frac{1}{2} \int_\Omega \langle A(\nabla u) \nabla(\bar \varphi^2), \nabla |\nabla u|^2 \rangle + \disp \int_\Omega \bar \varphi^2 f'(u)|\nabla u|^2 \\[0.5cm]
%%\end{array}
%%$$
%%And using the expression of $\bar \varphi$, 
%%$$
%%\begin{array}{l}
%%\disp \int_{\Omega} \bar \varphi^2\Big\{\lambda_1 u_{\nu\nu}^2 + (\lambda_1+\lambda_2) \big|\nabla^\top   u_\nu\big|^2 +\lambda_2\Big[u_\nu^2|\II|^2 +\Ric(\nabla u, \nabla u)\Big]\Big\} \\[0,5cm]
%%= \disp \int_{\partial \Omega} \bar{\varphi}^2 \langle A(\nabla u) \nabla \frac{|\nabla u|^2}{2} , \eta \rangle - \frac{1}{2} \int_\Omega \frac{w}{w+\eps}\langle A(\nabla u) \nabla(\varphi^2), \nabla |\nabla u|^2 \rangle \\[0.5cm]
%%\quad \disp - \int_\Omega \varphi^2 \langle A(\nabla u) \nabla\left( \frac{w}{w+\eps}\right), \nabla |\nabla u|^2 \rangle + \int_\Omega \bar \varphi^2 f'(u)|\nabla u|^2 \\[0.5cm]
%= \disp - \int_{\partial \Omega} \bar{\varphi}^2 \langle A(\nabla u) \nabla \frac{|\nabla u|^2}{2} , \eta \rangle - \frac{1}{2} \int_\Omega \frac{w}{w+\eps}\langle A(\nabla u) \nabla(\varphi^2), \nabla |\nabla u|^2 \rangle \\[0.5cm]
%\quad \disp - \eps \int_\Omega \frac{\varphi^2}{(w+\eps)^2} \langle A(\nabla u)\nabla w, \nabla |\nabla u|^2 \rangle + \int_\Omega \bar \varphi^2 f'(u)|\nabla u|^2\\[0.5cm]
%\end{array}
%$$
%From $3)$ in Lemma \ref{lem_basicomp}, the following identity holds:
%\begin{equation}\label{fir}
%\begin{array}{l}
%\disp \langle A(\nabla u) \nabla \big(\varphi |\nabla u|\big), \nabla \big(\varphi |\nabla u|\big)\rangle \\[0.5cm]
%\disp = \varphi^2\langle A(\nabla u)\nabla |\nabla u|, \nabla |\nabla u| \rangle + |\nabla u|^2 \langle A(\nabla u) \nabla \varphi, \nabla \varphi\rangle \\[0.5cm]
%\disp \quad + \frac{1}{2} \langle A(\nabla u) \nabla (\varphi^2), \nabla |\nabla u|^2 \rangle \\[0.5cm]
%\disp = \varphi^2\Big[ \lambda_2 |\nabla^\top   u_\nu|^2 + \lambda_1 u_{\nu\nu}^2\Big] + |\nabla u|^2 \langle A(\nabla u) \nabla \varphi, \nabla \varphi \rangle \\[0.5cm]
%\disp \quad + \frac{1}{2} \langle A(\nabla u) \nabla (\varphi^2), \nabla |\nabla u|^2 \rangle 
%\end{array}
%\end{equation}
%Inserting into the above,
%$$
%\begin{array}{l}
%\disp \int_{\Omega} \frac{\varphi^2 w}{w+\eps}\Big\{\lambda_1 u_{\nu\nu}^2 + (\lambda_1+\lambda_2) \big|\nabla^\top   u_\nu\big|^2 +\lambda_2\Big[u_\nu^2|\II|^2 +\Ric(\nabla u, \nabla u)\Big]\Big\} \\[0,5cm]
%= \disp - \int_{\partial \Omega} \frac{\varphi^2w}{w+\eps} \langle A(\nabla u) \nabla \frac{|\nabla u|^2}{2} , \eta \rangle  \\[0.5cm]
%\quad \disp - \eps \int_\Omega \frac{\varphi^2}{(w+\eps)^2} \langle A(\nabla u)\nabla w, \nabla |\nabla u|^2 \rangle \\[0.5cm]
%\disp + \int_\Omega \frac{\varphi^2w}{w+\eps}\Big[ \lambda_2 |\nabla^\top   u_\nu|^2 + \lambda_1 u_{\nu\nu}^2\Big] + \int_\Omega \frac{|\nabla u|^2 w}{w+\eps} \langle A(\nabla u) \nabla \varphi, \nabla \varphi \rangle \\[0.5cm]
%\disp - \int_\Omega \frac{w}{w+\eps} \disp \langle A(\nabla u) \nabla \big(\varphi |\nabla u|\big), \nabla \big(\varphi |\nabla u|\big)\rangle + \int_\Omega \frac{\varphi^2w}{w+\eps} f'(u)|\nabla u|^2
%\end{array}
%$$
%And simplifying, 
%$$
%\begin{array}{l}
%\disp \int_{\Omega} \frac{\varphi^2 w}{w+\eps}\Big\{\lambda_1 \big|\nabla^\top   u_\nu\big|^2 +\lambda_2\Big[u_\nu^2|\II|^2 +\Ric(\nabla u, \nabla u)\Big]\Big\} \\[0,5cm]
%= \disp -\int_{\partial \Omega} \frac{\varphi^2w}{w+\eps} \langle A(\nabla u) \nabla \frac{|\nabla u|^2}{2} , \eta \rangle  \\[0.5cm]
%\quad \disp - \eps \int_\Omega \frac{\varphi^2}{(w+\eps)^2} \langle A(\nabla u)\nabla w, \nabla |\nabla u|^2 \rangle \\[0.5cm]
%\disp + \int_\Omega \frac{|\nabla u|^2 w}{w+\eps} \langle A(\nabla u) \nabla \varphi, \nabla \varphi \rangle \\[0.5cm]
%\disp - \int_\Omega \frac{w}{w+\eps} \disp \langle A(\nabla u) \nabla \big(\varphi |\nabla u|\big), \nabla \big(\varphi |\nabla u|\big)\rangle + \int_\Omega \frac{\varphi^2w}{w+\eps} f'(u)|\nabla u|^2 \\[0.5cm]
%= \disp - \int_{\partial \Omega} \frac{\varphi^2w}{w+\eps} \langle A(\nabla u) \nabla \frac{|\nabla u|^2}{2} , \eta \rangle  \\[0.5cm]
%\quad \disp - \eps \int_\Omega \frac{\varphi^2}{(w+\eps)^2} \langle A(\nabla u)\nabla w, \nabla |\nabla u|^2 \rangle \\[0.5cm]
%\disp + \int_\Omega \frac{w}{w+\eps}|\nabla u|^2 \langle A(\nabla u) \nabla \varphi, \nabla \varphi \rangle \\[0.5cm]
%\disp - \int_\Omega \left\{\disp \langle A(\nabla u) \nabla \big(\varphi |\nabla u|\big), \nabla \big(\varphi |\nabla u|\big)\rangle -\frac{\varphi^2w}{w+\eps} f'(u)|\nabla u|^2 \right\} \\[0.5cm]
%\disp + \int_\Omega \frac{\eps}{w+\eps} \disp \langle A(\nabla u) \nabla \big(\varphi |\nabla u|\big), \nabla \big(\varphi |\nabla u|\big)\rangle.
%\end{array}
%$$
%Inserting \eqref{eq_picone_2}, 
%$$
%\begin{array}{l}
%\disp \int_{\Omega} \frac{\varphi^2 w}{w+\eps}\Big\{\lambda_1 \big|\nabla^\top   u_\nu\big|^2 +\lambda_2\Big[u_\nu^2|\II|^2 +\Ric(\nabla u, \nabla u)\Big]\Big\} \\[0,5cm]
%\le \disp - \int_{\partial \Omega} \frac{\varphi^2w}{w+\eps} \langle A(\nabla u) \nabla \frac{|\nabla u|^2}{2} , \eta \rangle  \\[0.5cm]
%\quad \disp - \eps \int_\Omega \frac{\varphi^2}{(w+\eps)^2} \langle A(\nabla u)\nabla w, \nabla |\nabla u|^2 \rangle \\[0.5cm]
%\disp + \int_\Omega \frac{w}{w+\eps}|\nabla u|^2 \langle A(\nabla u) \nabla \varphi, \nabla \varphi \rangle \\[0.5cm]
%\disp + \int_{\partial \Omega} \frac{\varphi^2|\nabla u|^2}{w+\eps} \langle A(\nabla u)\nabla w, \eta \rangle - \int_\Omega (w+\eps)^2 \langle A(\nabla u)\nabla \left(\frac{\varphi|\nabla u|}{w+\eps}\right),\nabla\left(\frac{\varphi|\nabla u|}{w+\eps}\right) \rangle \\[0.5cm]
%\disp + \int_\Omega \frac{\eps}{w+\eps} \disp \langle A(\nabla u) \nabla \big(\varphi |\nabla u|\big), \nabla \big(\varphi |\nabla u|\big)\rangle
%\end{array}
%$$
%That is, 
%$$
%\begin{array}{l}
%\disp \int_{\Omega} \frac{\varphi^2 w}{w+\eps}\Big\{\lambda_1 \big|\nabla^\top   u_\nu\big|^2 +\lambda_2\Big[u_\nu^2|\II|^2 +\Ric(\nabla u, \nabla u)\Big]\Big\} \\[0,5cm]
%\le \disp - \int_{\partial \Omega}\frac{\varphi^2}{w+\eps} \left\{ w\langle A(\nabla u) \nabla \frac{|\nabla u|^2}{2} , \eta \rangle - |\nabla u|^2  \langle A(\nabla u)\nabla w, \eta\rangle \right\} \\[0.5cm]
%\quad \disp - \eps \int_\Omega \frac{\varphi^2}{(w+\eps)^2} \langle A(\nabla u)\nabla w, \nabla |\nabla u|^2 \rangle \\[0.5cm]
%\disp + \int_\Omega \frac{w}{w+\eps}|\nabla u|^2 \langle A(\nabla u) \nabla \varphi, \nabla \varphi \rangle \\[0.5cm]
%\disp - \int_\Omega (w+\eps)^2 \langle A(\nabla u)\nabla \left(\frac{\varphi|\nabla u|}{w+\eps}\right),\nabla\left(\frac{\varphi|\nabla u|}{w+\eps}\right) \rangle \\[0.5cm]
%\disp + \int_\Omega \frac{\eps}{w+\eps} \disp \langle A(\nabla u) \nabla \big(\varphi |\nabla u|\big), \nabla \big(\varphi |\nabla u|\big)\rangle,
%\end{array}
%$$
%that was to be proved.
%\end{proof}
%
%
%
%%%%%%%%%%%%%%%    END OLD PROOF
%
%
%
In the next Lemma, we deal with the boundary term and extend to a quasilinear setting a ``magical identity"  first discovered in \cite{Arma,Arma2}, see also \cite{FarMarVal,cmr}.


\begin{lemma}\label{lem_bordo}
	Let $\Omega\subseteq M$ be a $C^2$ domain and let $u \in C^2(\overline\Omega)$ be such that both $u$ and $c_j \doteq \partial_\eta u$ are constant on a connected component $\partial_j \Omega$ of $\partial\Omega$, and let $a$ satisfy \eqref{assu_A}. Then, for each vector field $X$ in a neighbourhood of $\partial_j \Omega$, the function $w= \langle \nabla u, X\rangle$ satisfies
	\begin{equation}\label{identity_quasilinear}
		\langle w A(\nabla u) \nabla \frac{|\nabla u|^2}{2} - |\nabla u|^2  A(\nabla u)\nabla w, \eta\rangle = - \lambda_1(|\nabla u|)|\nabla u|^2 c_j \langle \eta, \nabla_{\eta} X \rangle
	\end{equation}
	on $\partial_j \Omega$. In particular, if $X$ is Killing then the right hand side of \eqref{identity_quasilinear} vanishes.
\end{lemma}

\begin{proof}
	The identity is obvious if $\partial_\eta u = 0$, since in this case $\nabla u \equiv 0$ on $\partial_j \Omega$. Thus, assume $c_j \neq 0$, so that $c_j \eta = \nabla u$. Since $A(\nabla u)$ is symmetric and $A(\nabla u)\eta = \lambda_1(|\nabla u|) \eta$, we deduce
	\[
		\begin{array}{l}
			\disp \langle w A(\nabla u) \nabla \frac{|\nabla u|^2}{2} - |\nabla u|^2  A(\nabla u)\nabla w, \eta\rangle \\[0.3cm]
			= \ \disp \langle w \nabla \frac{|\nabla u|^2}{2} - |\nabla u|^2 \nabla w, \eta\rangle \lambda_1(|\nabla u|) \qquad \text{on } \, \partial_j \Omega.
	\end{array}
	\]
	The conclusion therefore follows from the identity 
	\begin{equation}\label{eq_FMV}
		\langle w \nabla \frac{|\nabla u|^2}{2} - |\nabla u|^2 \nabla w, \eta \rangle =  - |\nabla u|^2 c_j \langle \eta, \nabla_\eta X \rangle \qquad \text{on } \partial_j \Omega,
	\end{equation} 
	which was proved\footnote{There is an uninfluent sign mistake there, since the right-hand side of \eqref{eq_FMV} is written as $- |\nabla u|^3 \langle \eta, \nabla_\eta X \rangle$.} in \cite[Lemma 32]{FarMarVal}.
\end{proof}

We are ready to pass to limits as $\eps \to 0$ in Proposition \ref{prop_Poincare_loc}.

\begin{theorem} \label{teo_Poincare_loc}
	Let $\Omega \subseteq M$ be a $C^2$ domain and let $u \in C^3(\Omega) \cap C^2(\overline\Omega)$ be a stable solution to 
	\[
	\left\{ \begin{array}{ll}
		\Delta_a u + f(u)=0 & \quad \text{in } \, \Omega, \\[0.2cm]
		u, \partial_\eta u \, \text{ locally constant} & \quad \text{on } \, \partial \Omega, 
	\end{array}
	\right.
	\]
	where $f \in C^1(\R)$ and $a$ satisfies \eqref{assu_A_superstrong}. Let $\Omega_0\subseteq\Omega$ be a domain with locally finite perimeter in $\Omega$ such that $\partial\Omega_0 \subseteq \partial\Omega \cup \critu$, let $X$ be a Killing vector field in $\overline\Omega$ and suppose that $w = \langle \nabla u, X \rangle$ satisfies $w \ge 0$, $w \not\equiv 0$ in $\Omega_0$. Then for each $\varphi \in \lip_c(M)$ it holds
	\begin{equation} \label{piufinalissima_teo3}
		\begin{split}
			\int_{\Omega_0} |\nabla u|^2 \langle A(\nabla u)\nabla\varphi,\nabla\varphi\rangle & \geq \int_{\Omega_0} \varphi^2 \left\{ \lambda_1|\nabla^\top|\nabla u||^2 + \lambda_2 [|\nabla u|^2|\II|^2 + \Ric(\nabla u,\nabla u)] \right\} \\
			& \phantom{=\;} + \int_{\Omega_0} w^2 \langle A(\nabla u)\nabla\left(\frac{\varphi|\nabla u|}{w}\right),\nabla\left(\frac{\varphi|\nabla u|}{w}\right)\rangle \, .
		\end{split}
	\end{equation}
\end{theorem}

\begin{proof}
	Since $w$ satisfies $\diver(A(\nabla u)\nabla w) + f'(u) w = 0$ in $\Omega$ and $w\geq 0$, $w\not\equiv 0$ in $\Omega_0$, by the Harnack inequality (see Remark \ref{rem_harnackHopf}) we have $w>0$ in $\Omega_0$. From Lemma \ref{lem_bordo} (applied to any $C^1$ extension of $X$ in a neighbourhood of $\partial\Omega$) we know that
	\begin{equation} \label{w_bordo}
		\langle w A(\nabla u) \nabla \frac{|\nabla u|^2}{2},\eta\rangle = \langle |\nabla u|^2  A(\nabla u)\nabla w, \eta\rangle \qquad \text{on } \, \partial\Omega \, .
	\end{equation}
	We claim that
	\begin{equation} \label{w_bordo2}
		\{ x \in \overline{\Omega}_0 \cap \partial\Omega : w(x) = 0 \} \subseteq \{ x \in \overline{\Omega} : \nabla u(x) = 0 \} \, .
	\end{equation}
	Indeed, suppose by contradiction that there exists $x\in\overline{\Omega}_0 \cap \partial\Omega$ such that $w(x) = 0$ and $\nabla u(x)\neq 0$. Since $\partial\Omega_0 \subseteq \critu \cup \partial\Omega$ and by assumption we have $\nabla u\neq 0$ at $x$, there exists a neighbourhood $U$ of $x$ in $M$ such that $U\cap\critu = \emptyset$ and therefore $U\cap\partial\Omega_0 = U\cap\partial\Omega$. Since $\partial\Omega$ is $C^2$ regular, $\Omega_0$ satisfies an interior ball condition at $x$. Having $w>0$ in $\Omega_0$ and $w(x)=0$, by Hopf's lemma (see again Remark \ref{rem_harnackHopf}) we deduce $\partial_\eta w(x) > 0$. Then, we have
	\[
	\langle|\nabla u|^2 A(\nabla u)\nabla w,\eta\rangle = \lambda_1(|\nabla u|) |\nabla u|^2 \partial_\eta w \neq 0 \qquad \text{at } \, x
	\]
	where we used that $A(\nabla u)\eta = \lambda_1(|\nabla u|)\eta$ since $\eta$ is proportional to $\nabla u$ at $x$. But then the left-hand side of \eqref{w_bordo} must also be nonzero at $x$, contradicting the assumption $w(x) = 0$.
	
	Let $\varphi\in\lip_c(M)$ be given. From Proposition \ref{prop_Poincare_loc} and \eqref{w_bordo} we have
	\begin{equation} \label{poinalt}
		\begin{split}
			\int_{\Omega_0} |\nabla u|^2 \langle A(\nabla u)\nabla\varphi,\nabla\varphi\rangle & = \int_{\Omega_0} \varphi^2 \left\{ \lambda_1|\nabla^\top|\nabla u||^2 + \lambda_2 [|\nabla u|^2|\II|^2 + \Ric(\nabla u,\nabla u)] \right\} \\
			& \phantom{=\;} + \int_{\Omega_0} (w+\eps)^2 \langle A(\nabla u)\nabla\left(\frac{\varphi|\nabla u|}{w+\eps}\right),\nabla\left(\frac{\varphi|\nabla u|}{w+\eps}\right)\rangle \\
			& \phantom{=\;} - \int_{\Omega_0} \frac{\eps \varphi^2}{w+\eps} f'(u)|\nabla u|^2 \\
			& \phantom{=\;} + \int_{\overline\Omega_0\cap\partial\Omega} \frac{\eps\varphi^2}{w+\eps} \langle A(\nabla u)\frac{\nabla|\nabla u|^2}{2},\eta\rangle
		\end{split}
	\end{equation}
	for any $\eps>0$. We now let $\eps\to0$ to obtain \eqref{piufinalissima_teo3} from \eqref{poinalt}. Since $|\eps/(w+\eps)|<1$ in $\Omega_0$, by the Lebesgue convergence theorem we have
	\[
	\int_{\Omega_0} \frac{\eps \varphi^2}{w+\eps} f'(u)|\nabla u|^2 \to 0 \qquad \text{as } \, \eps \to 0 \, .
	\]
	Setting $E_0 = \{w=0\} \cap \overline\Omega_0\cap\partial\Omega$ and $E_1 = \{w>0\} \cap \overline\Omega_0\cap\partial\Omega$, we have
	\begin{align*}
		\int_{\overline\Omega_0\cap\partial\Omega} \frac{\eps\varphi^2}{w+\eps} \langle A(\nabla u)\frac{\nabla|\nabla u|^2}{2},\eta\rangle & = \int_{E_0} \varphi^2 \langle A(\nabla u)\frac{\nabla|\nabla u|^2}{2},\eta\rangle \\
		& \phantom{=\;} + \int_{E_1} \frac{\eps\varphi^2}{w+\eps} \langle A(\nabla u)\frac{\nabla|\nabla u|^2}{2},\eta\rangle \, .
	\end{align*}
	When letting $\eps\to0$, the integral over $E_1$ converges to $0$ by the Lebesgue dominated convergence theorem. On the other hand, by \eqref{w_bordo2} we have $\nabla u = 0$ and therefore $\nabla|\nabla u|^2 = 0$ on $E_0$, so the integral over $E_0$ also vanishes and we get
	\[
	\int_{\overline\Omega_0\cap\partial\Omega} \frac{\eps\varphi^2}{w+\eps} \langle A(\nabla u)\frac{\nabla|\nabla u|^2}{2},\eta\rangle \to 0 \qquad \text{as } \, \eps \to 0 \, .
	\]
	Finally, by explicit computation of $\nabla(\varphi|\nabla u|/(w+\eps))$ and Fatou's lemma we have
	\[
	\liminf_{\eps\to0} \int_{\Omega_0} (w+\eps)^2 \langle A(\nabla u)\nabla\left(\frac{\varphi|\nabla u|}{w+\eps}\right),\nabla\left(\frac{\varphi|\nabla u|}{w+\eps}\right)\rangle \geq \int_{\Omega_0} w^2 \langle A(\nabla u)\nabla\left(\frac{\varphi|\nabla u|}{w}\right),\nabla\left(\frac{\varphi|\nabla u|}{w}\right)\rangle
	\]
	and this concludes the proof.
\end{proof}

%\tcb{We are ready to pass to limits as $\eps \to 0$ in Proposition \ref{prop_poincare}. The issue is subtle if $w=0$ somewhere on $\partial \Omega$.} 
%
%
%
%\begin{theorem}\label{teo_Poincare}
%Let $\Omega \subseteq M$ be a $C^3$ domain, and let $u \in C^3(\Omega) \cap C^{2,\alpha}(\overline\Omega)$ be a stable solution to 
%$$
%\left\{ \begin{array}{ll}
%	\Delta_a u + f(u)=0 & \quad \text{in } \, \Omega, \\[0.2cm]
%	u, \ \partial_\eta u & \quad \text{locally constant on } \, \partial \Omega, 
%\end{array}\right. 
%$$
%where $f \in C^1(\R)$ and $a$ satisfies \eqref{assu_A_superstrong}. Let $X$ be a Killing vector field in $\overline\Omega$ and suppose that $w = \langle \nabla u, X \rangle$ satisfies $w \ge 0$, $w \not\equiv 0$ in $\Omega$. Then, for each $\varphi \in \lip_c(M)$ it holds
%\begin{equation}\label{piufinalissima_teo_2}
%\begin{array}{l}
%\disp \int_{\Omega} \varphi^2 \Big\{\lambda_1 \big|\nabla^\top   u_\nu\big|^2 +\lambda_2\Big[u_\nu^2|\II|^2 +\Ric(\nabla u, \nabla u)\Big]\Big\} \\[0,5cm]
%\disp + \int_{\Omega} w^2 \langle A(\nabla u) \nabla \left(\frac{\varphi|\nabla u|}{w}\right), \nabla \left(\frac{\varphi|\nabla u|}{w}\right) \rangle \\[0.5cm]
%\le \disp \int_{\Omega} |\nabla u|^2 \langle A(\nabla u) \nabla \varphi, \nabla \varphi \rangle.
%\end{array}
%\end{equation}
%\end{theorem}
%
%%Before proving the result, we observe that if $w\ge 0$ and $w \not \equiv 0$ on the entire $\Omega$, then by the Harnack inequality (see Remark \ref{rem_harnackHopf}) $w>0$ in $\Omega$. Therefore, $\Omega_0 \equiv \Omega$ and \eqref{piufinalissima_teo_2} holds on the entire $\Omega$, leading to the following: 
%%
%%\begin{corollary}\label{cor_Poincare}
%%	Let $\Omega \subseteq M$ be a $C^3$ domain, and let $u \in C^3(\Omega) \cap C^{2,\alpha}(\overline\Omega)$ be a stable solution to 
%%	$$
%%	\left\{ \begin{array}{ll}
%%		\Delta_a u + f(u)=0 & \quad \text{in } \, \Omega, \\[0.2cm]
%%		u, \ \partial_\eta u & \quad \text{locally constant on } \, \partial \Omega, 
%%	\end{array}\right. 
%%	$$
%%	where $f \in C^1(\R)$ and $a$ satisfies \eqref{assu_A_superstrong}.
%%	Let $X$ be a Killing vector field in $\overline\Omega$, and suppose that $w = \langle \nabla u, X \rangle$ satisfies $w \ge 0$, $w \not\equiv 0$ in $\Omega$. Then, for each $\varphi \in \lip_c(M)$ it holds
%%	\begin{equation}\label{piufinalissima_teo_2}
%%		\begin{array}{l}
%%			\disp \int_{\Omega} \varphi^2 \Big\{\lambda_1 \big|\nabla^\top   u_\nu\big|^2 +\lambda_2\Big[u_\nu^2|\II|^2 +\Ric(\nabla u, \nabla u)\Big]\Big\} \\[0,5cm]
%%			\disp + \int_{\Omega} w^2 \langle A(\nabla u) \nabla \left(\frac{\varphi|\nabla u|}{w}\right), \nabla \left(\frac{\varphi|\nabla u|}{w}\right) \rangle \\[0.5cm]
%%			\le \disp \int_{\Omega} |\nabla u|^2 \langle A(\nabla u) \nabla \varphi, \nabla \varphi \rangle.
%%		\end{array}
%%	\end{equation}
%%\end{corollary}
%
%
%
%
%\begin{proof}%[Proof of Theorem \ref{piufinalissima_teo_2}]
%By Lemma \ref{lem_linearKilling}, $w$ solves
%$$
%\diver \big( A(\nabla u) \nabla w \big) + f'(u)w = 0 \qquad \text{weakly in } \, \Omega,
%$$
%so by the Harnack inequality (see Remark \ref{rem_harnackHopf}) $w>0$ in $\Omega$.
%%Since $w \in C^2(\Omega)$ and $a$ satisfies \eqref{assu_A_superstrong}, the strong maximum principle in \cite[Theorem 2.8.1]{pucci_serrin} implies that $w>0$ on $\Omega$. 
%For each $\eps>0$, we therefore deduce the validity of the geometric Poincar\'e inequality \eqref{piufinalissima_teo}. Next, by Lemma \ref{lem_bordo} and since $X$ is Killing, the boundary term in \eqref{piufinalissima_teo} vanishes identically, and thus 
%\begin{equation}\label{poinc_part}
%\begin{array}{l}
%\disp \int_{\Omega} \frac{\varphi^2 w}{w+\eps}\Big\{\lambda_1 \big|\nabla^\top   u_\nu\big|^2 +\lambda_2\Big[u_\nu^2|\II|^2 +\Ric(\nabla u, \nabla u)\Big]\Big\} \\[0,5cm]
%\le \disp - \frac{\eps}{2} \int_{\Omega} \frac{\varphi^2}{(w+\eps)^2} \langle A(\nabla u)\nabla w, \nabla |\nabla u|^2 \rangle \\[0.5cm]
%\disp + \int_{\Omega} \frac{w}{w+\eps}|\nabla u|^2 \langle A(\nabla u) \nabla \varphi, \nabla \varphi \rangle \\[0.5cm]
%\disp - \int_{\Omega} (w+\eps)^2 \langle A(\nabla u) \nabla \left(\frac{\varphi|\nabla u|}{w+\eps}\right), \nabla \left(\frac{\varphi|\nabla u|}{w+\eps}\right) \rangle \\[0.5cm]
%\disp + \int_{\Omega} \frac{\eps}{w+\eps} \disp \langle A(\nabla u) \nabla \big(\varphi |\nabla u|\big), \nabla \big(\varphi |\nabla u|\big)\rangle
%\end{array}
%\end{equation}
%Since $|\eps /w+\eps| \le 1$ and converges to zero as $\eps \ra 0$, a direct application of Lebesgue convergence theorem shows
%\begin{equation}\label{uno_eps}
%\disp \int_{\Omega} \frac{\eps}{w+\eps} \disp \langle A(\nabla u) \nabla \big(\varphi |\nabla u|\big), \nabla \big(\varphi |\nabla u|\big)\rangle \ra 0 \qquad \text{as } \, \eps \ra 0.
%\end{equation}
%\tcb{Next, we investigate the term
%$$
%(\star) \doteq \int_{\Omega} \frac{\eps \varphi^2}{(w+\eps)^2} \langle A(\nabla u)\nabla w, \nabla |\nabla u|^2 \rangle.
%$$
%Because of the positivity of $w$ in $\Omega$, the integrand converges pointwise to zero a.e. as $\eps \ra 0$. The conlusion $(\star) \to 0$ as $\eps \to 0$ is then easy if $w>0$ in $\overline\Omega$, by Lebesgue convergence theorem, but if $w$ vanishes somewhere on $\cap \partial \Omega$ it needs some care. 
%We first prove that for each $K \Subset \overline\Omega$ compact it holds 
%\begin{equation}\label{claim_partial}
%\big| \langle A(\nabla u) \nabla |\nabla u|^2, \nabla w \rangle\big| \le C'_K |w| \qquad \text{on } \, K \cap \partial \Omega.
%\end{equation}
%The inequality holds on components $\partial_j\Omega$ for which $\partial_\eta u = 0$. Indeed, in this case $\nabla u = 0$ on $\partial_j \Omega$ and thus both $w$ and $\nabla|\nabla u|^2$ vanish there. Assume that $\partial_\eta u \neq 0$, so $\eta = \pm \nabla u/|\nabla u|$. Since $|\nabla u|$ is locally constant on $\partial \Omega$ and $A(\nabla u)\eta = \lambda_1(|\nabla u|) \eta$, we deduce on $K \cap \partial \Omega$
%\begin{equation}\label{primo_esti}
%\begin{array}{lcl}
%\disp \big|\langle A(\nabla u) \nabla |\nabla u|^2, \nabla w \rangle \big| & = & \disp \big|\partial_\eta(|\nabla u|)^2\big| \big| \langle A(\nabla u) \eta, \nabla w \rangle \big| \\[0.3cm]
%& = & \disp \disp \lambda_1(|\nabla u|) \big|\partial_\eta(|\nabla u|)^2\big| \big| \langle \eta, \nabla w \rangle\big| \le 2\lambda_1(|\nabla u|) |\nabla u| |\nabla \di u| \big| \langle \eta, \nabla w \rangle\big| \\[0.3cm]
%& \le & C''_K \big| \langle \eta, \nabla w \rangle\big|,
%\end{array}
%\end{equation}
%where 
%$$
%C''_K = \left[\sup_{K \cap \partial \Omega} |\nabla \di u|\right] \cdot \left[\sup \left\{ t \lambda_1(t) \ : \ t \in \Big[0, \sup_{K \cap \partial \Omega} |\nabla u|\Big]\right\}\right].  
%$$
%Next, we estimate
%$$
%\langle \eta, \nabla w \rangle = \eta(w) = \eta \big( \langle \nabla u, X \rangle \big) = \langle \nabla_\eta \nabla u, X \rangle + \langle \nabla u, \nabla_\eta X \rangle.
%$$
%Since  $X$ is Killing, the second term vanishes, and since $|\nabla u|$ is constant on $\partial_j \Omega$ we infer 
%$$
%\langle \eta, \nabla w \rangle  =  \disp \langle \nabla_\eta \nabla u, X \rangle = \nabla \di u(\eta, X) =\pm  \langle \nabla |\nabla u|, X \rangle =\pm \partial_\eta\left( |\nabla u|\right) \langle \eta,X \rangle =  \frac{\partial_\eta\big( |\nabla u|\big)}{|\nabla u|} w.
%$$
%Inserting into \eqref{primo_esti} we deduce 
%$$
%\disp \big|\langle A(\nabla u) \nabla |\nabla u|^2, \nabla w \rangle \big| \le C_K''|\nabla |\nabla u|| |w| \le C_K' |w|,
%$$
%and \eqref{claim_partial} holds.\\
%%
%Next, let $K$ be the closure of a neighbourhood where Fermi coordinates
%$$
%\Psi : V \times [0,\delta] \ \ \longrightarrow \ \  K, \qquad \Psi(y,t) = \exp^\perp\big(t \eta(y)\big)
%$$
%are defined for some $V \Subset \partial \Omega$. We claim that there exists a constant $C_K > 0$ such that 
%\begin{equation}\label{claim}
%	\big| \langle A(\nabla u) \nabla |\nabla u|^2, \nabla w \rangle\big| \le C_K |w|t^{\alpha-1} \qquad \text{on } \, K,
%\end{equation}
%where $\alpha$ is the H\"older exponent of the second derivative of $u$ in $(ii)$ on $K$. 
%%To prove \eqref{claim}, observe that 
%%$$
%%|A(\nabla u) \nabla |\nabla u|^2| \le 2 |\nabla \di u| |\nabla u| \lambda_{\max}(|\nabla u|) \in L^\infty(K).
%%$$
%%Thus, by the positivity of $w$ on $\Omega$ and a compactness argument, it is sufficient to prove \eqref{claim} in a small 
%Since by \eqref{assu_A_superstrong} the matrix $A(\nabla u)$ has locally $C^{1,1}$ coefficients and it is uniformly elliptic, and since $w = \langle \nabla u, X \rangle \in C^2(\Omega) \cap C^1(\overline\Omega)$, we can apply the Hopf Lemma \cite[Theorem 2.8.3]{pucci_serrin} to obtain 
%\[	
%w(y,t) \ge Ct \qquad \text{on } \, K.
%\]
%for some constant $C>0$. On the other hand, $w \in C^1(K)$ and the function 
%$$
%g(y,t) = \langle A(\nabla u) \nabla |\nabla u|^2, \nabla w \rangle (\Psi(y,t))
%$$
%is in $C^{0,\alpha}(K)$ \tcr{CHECK: it is the product of two $C^{0,\alpha}$ functions and one $C^{0,1}$ function}. Let $L$ be a common upper bound for the $\alpha$-H\"older seminorm of $g$ and $w$. From \eqref{claim_partial}, we obtain
%$$
%\begin{array}{lcl}
%|g(y,t)| & \le & \disp |g(y,0)| + Lt^\alpha \le C'_K|w(y,0)| + Lt^\alpha \le C'_K(|w(y,t)|+ Lt^\alpha) + Lt^\alpha \\[0.3cm]
%& \le & \disp C_K|w(y,t)|t^{\alpha-1}
%\end{array}
%$$
%proving \eqref{claim} in $K$. We therefore  deduce, on compact subsets $K$ which are domains of a Fermi chart,
%$$
%\left|\frac{\eps \varphi^2}{(w+\eps)^2} \langle A(\nabla u)\nabla w, \nabla |\nabla u|^2 \rangle\right| \le \varphi^2 C_K \frac{\eps wt^{\alpha-1}}{(w+\eps)^2} \le C_K\varphi^2t^{\alpha-1} \in L^1(K)
%$$
%On the other hand, if $K \cap \partial \Omega = \emptyset$ the stronger 
%\begin{equation}\label{claim_2}
%	\big| \langle A(\nabla u) \nabla |\nabla u|^2, \nabla w \rangle\big| \le C_K |w| \qquad \text{holds in } \, K
%\end{equation}
%for some constant $C_K$. Summarizing, we can apply Lebesgue convergence to conclude that $(\star) \to 0$ as $\eps \to 0$, as required.  
%%$$
%%\eps \int_\Omega \frac{\varphi^2}{(w+\eps)^2} \langle A(\nabla u)\nabla w, \nabla |\nabla u|^2 \rangle \ra 0 \qquad \text{as } \, \eps \ra 0.
%%$$
%}
%Letting $\eps \ra 0$ in \eqref{poinc_part} we get
%\begin{equation}\label{poinc_part_2}
%\begin{array}{l}
%\disp \int_{\Omega} \varphi^2\Big\{\lambda_1 \big|\nabla^\top   u_\nu\big|^2 +\lambda_2\Big[u_\nu^2|\II|^2 +\Ric(\nabla u, \nabla u)\Big]\Big\} \\[0,5cm]
%\disp + \liminf_{\eps \ra 0} \int_\Omega (w+\eps)^2 \langle A(\nabla u)  \nabla \left(\frac{\varphi|\nabla u|}{w+\eps}\right), \nabla \left(\frac{\varphi|\nabla u|}{w+\eps}\right) \rangle \\[0.5cm]
%\le \disp \int_\Omega |\nabla u|^2 \langle A(\nabla u) \nabla \varphi, \nabla \varphi \rangle.
%\end{array}
%\end{equation}
%To conclude, observe that for $\delta >0$ and setting $\Omega_\delta = \Omega \cap \{w> \delta\}$, 
%$$
%\begin{array}{l}
%\disp \liminf_{\eps \ra 0} \int_{\Omega} (w+\eps)^2 \langle A(\nabla u) \nabla \left(\frac{\varphi|\nabla u|}{w+\eps}\right), \nabla \left(\frac{\varphi|\nabla u|}{w+\eps}\right) \rangle \\[0.5cm]
%\disp \ge \disp \liminf_{\eps \ra 0} \int_{\Omega_\delta} (w+\eps)^2 \langle A(\nabla u) \nabla \left(\frac{\varphi|\nabla u|}{w+\eps}\right), \nabla \left(\frac{\varphi|\nabla u|}{w+\eps}\right) \rangle \\[0.5cm]
%\disp = \int_{\Omega_\delta} w^2 \langle A(\nabla u) \nabla \left(\frac{\varphi|\nabla u|}{w}\right), \nabla \left(\frac{\varphi|\nabla u|}{w}\right) \rangle,
%\end{array}
%$$
%and letting $\delta \ra 0$ we obtain 
%$$
%\begin{array}{l}
%\disp \liminf_{\eps \ra 0} \int_{\Omega} (w+\eps)^2 \langle A(\nabla u)  \nabla \left(\frac{\varphi|\nabla u|}{w+\eps}\right), \nabla \left(\frac{\varphi|\nabla u|}{w+\eps}\right) \rangle \ge \disp \int_{\Omega} w^2 \langle A(\nabla u) \nabla \left(\frac{\varphi|\nabla u|}{w}\right), \nabla\left(\frac{\varphi|\nabla u|}{w}\right) \rangle 
%\end{array}
%$$
%inserting into \eqref{poinc_part_2} we deduce the desired inequality.
%\end{proof}







%For later purposes, we shall localize inequality \eqref{piufinalissima_teo_2} to \tcr{connected components of $\Omega\backslash\critu$, where $\critu = \{x\in\Omega : \}$} subsets $\Omega_0 \subseteq \Omega$ whose boundary consists of two disjoint closed pieces, one contained in $\partial \Omega$ and one in the set of interior critical points of $u$. 
%
%\begin{proposition}\label{prop_Poincare_loc}
%	Let $\Omega \subseteq M$ be a $C^3$ domain, and let $u \in C^3(\Omega) \cap C^{2,\alpha}(\overline\Omega)$ be a stable solution to 
%	\[
%		\left\{ \begin{array}{ll}
%			\Delta_a u + f(u)=0 & \quad \text{in } \, \Omega, \\[0.2cm]
%			u, \, \partial_\eta u \text{ locally constant} & \quad \text{on } \, \partial \Omega, 
%		\end{array}\right. 
%	\]
%	where $f \in C^1(\R)$ and $a$ satisfies \eqref{assu_A_superstrong}.
%	%Define $\Gamma = \{x \in \Omega : \nabla u(x) = 0\}$, and let $\Omega_0$ be a connected component of $\Omega\backslash\Gamma$. 
%	Let $\Omega_0\subseteq\Omega$ be an open set of locally finite perimeter in $M$ such that $\partial\Omega_0 = \Gamma_0 \cup \Gamma_1$ for some closed sets $\Gamma_0$, $\Gamma_1$ satisfying
%	\[
%		\Gamma_0 \subseteq \{ x \in \Omega : \nabla u(x) = 0 \} \, , \qquad \Gamma_1 \subseteq \partial\Omega \, .
%	\]
%	Let $X$ be a Killing vector field on $\overline\Omega_0$, and suppose that $w = \langle \nabla u, X \rangle$ satisfies $w \ge 0$, $w \not\equiv 0$ in $\Omega_0$. Then, for each $\varphi\in\lip_c(M)$ it holds
%	\begin{equation}\label{piufinalissima_teo_3}
%		\begin{split}
%			\int_{\Omega_0} |\nabla u|^2 \langle A(\nabla u) \nabla \varphi, \nabla \varphi \rangle & \geq \int_{\Omega_0} \varphi^2 \left\{\lambda_1 \left|\nabla^\top|\nabla u|\right|^2 +\lambda_2\left[|\nabla u|^2|\II|^2 +\Ric(\nabla u, \nabla u)\right]\right\} \\
%			& \phantom{=\;} + \int_{\Omega_0} w^2 \langle A(\nabla u) \nabla \left(\frac{\varphi|\nabla u|}{w}\right), \nabla \left(\frac{\varphi|\nabla u|}{w}\right) \rangle \, .
%		\end{split}
%	\end{equation}
%\end{proposition}
%
%\begin{proof}
%	The argument is similar to that developed in the proofs of Proposition \ref{prop_poincare} and Theorem \ref{teo_Poincare}, so we only highlight the key differences. Let $\varphi\in\lip_c(M)$ be given and let $\eps>0$. Denoting by $r$ the distance function from $\partial\Omega$ in $M$, for sufficiently small $\delta,\tau>0$ let $\psi$ be defined as in \eqref{defpsi} and consider the compactly supported Lipschitz vector field $V$ defined as in \eqref{defV}. Since $\Omega_0$ has locally finite perimeter in $M$ we apply the divergence theorem to obtain
%	\begin{equation} \label{poinloc1}
%		\int_{\Omega_0} \diver\,V = 0 \, ,
%	\end{equation}
%	where we used that $V=0$ everywhere on $\partial\Omega_0$, as $|\nabla u|=0$ on $\Gamma_0$ and $\psi=0$ on $\Gamma_1$. Since in the present setting $w$ satisfies \eqref{poinalt0} with equality sign, from \eqref{poinloc1} we obtain integral identities analogous to inequalities \eqref{poinalt1}, \eqref{poinalt2} and \eqref{poinalt3}, in which $\Omega$ is replaced everywhere by $\Omega_0$ as domain of integration. Then, letting $\delta\to0$ and then $\tau\to0$, and noting that $\Omega_0\cap\{r=t\}$ converges to $\overline{\Omega}_0\cap\partial\Omega = \Gamma_1$ as $t\to0$, we obtain
%	\begin{equation} \label{poinloc2}
%		\int_{\Omega_0} \frac{\varphi^2 w}{w+\eps} f'(u)|\nabla u|^2 = \int_{\Omega_0} \langle A(\nabla u)\nabla w,\nabla\left(\frac{\varphi^2|\nabla u|^2}{w+\eps}\right)\rangle + \int_{\Gamma_1} \frac{\varphi^2|\nabla u|^2}{w+\eps} \langle A(\nabla u)\nabla w,\eta\rangle \, .
%	\end{equation}
%	Consider now the compactly supported Lipschitz vector field $W$ defined as in \eqref{defW}. By the divergence theorem we have
%	\[
%		\int_{\Omega_0} \diver\, W = - \int_{\Gamma_1} \langle W,\eta\rangle
%	\]
%	where we used that $W$ vanishes on the portion $\Gamma_0$ of $\partial\Omega_0$. From this point on, by repeating the same straightforward computations performed in the remainder of the proof of Proposition \ref{prop_poincare} we arrive at
%	\begin{align*}
%		\int_{\Omega_0} \frac{w}{w+\eps}|\nabla u|^2 \langle A(\nabla u) \nabla \varphi, \nabla \varphi \rangle & = \int_{\Omega_0} \frac{\varphi^2 w}{w+\eps}\{\lambda_1 |\nabla^\top|\nabla u||^2 +\lambda_2[|\nabla u|^2|\II|^2 +\Ric(\nabla u, \nabla u)]\} \\
%		& \phantom{=\;} + \int_{\Omega_0} (w+\eps)^2 \langle A(\nabla u) \nabla\left(\frac{\varphi|\nabla u|}{w+\eps}\right),\nabla\left(\frac{\varphi|\nabla u|}{w+\eps}\right) \rangle \\
%		& \phantom{=\;} - \int_{\Omega_0} \frac{\eps}{w+\eps} \disp \langle A(\nabla u) \nabla \big(\varphi |\nabla u|\big), \nabla \big(\varphi |\nabla u|\big)\rangle \\
%		& \phantom{=\;} + \frac{\eps}{2} \int_{\Omega_0} \frac{\varphi^2}{(w+\eps)^2} \langle A(\nabla u)\nabla w, \nabla |\nabla u|^2 \rangle \\
%		& \phantom{=\;} - \int_{\Gamma_1}\frac{\varphi^2}{w+\eps} \langle |\nabla u|^2 A(\nabla u)\nabla w - w A(\nabla u) \nabla \frac{|\nabla u|^2}{2} , \eta \rangle
%	\end{align*}
%	The integral over $\Gamma_1\subseteq\partial\Omega$ vanishes thanks to Lemma \ref{lem_bordo}. Then, letting $\eps\to0$ and arguing as in the proof of Theorem \ref{teo_Poincare} we finally obtain \eqref{piufinalissima_teo_3}.
%\end{proof}









%\tcr{CORREGGERE}
%
%We are ready to prove Theorem \ref{teo_main}, which follows from the next more general result:
%
%\begin{proposition}\label{prop_Poincare}
%	Let $(M^, \metric)$ be a complete Riemannian manifold, let $\Omega$ be an open subset with $C^3$-boundary, and let $u \in C^3(\overline\Omega)$ be a stable solution to 
%	$$
%	\diver\big(a(|\nabla u|) \nabla u\big) + f(u)=0. 
%	$$
%	where $a$ satisfy \eqref{assu_A}, 
%	$$
%	t \max\left \{ \lambda_1(t), \lambda_2(t)\right\}   \in L^\infty_\loc(\R^+_0)
%	$$
%	and $f \in C^1(\R)$. Let $X$ be a Killing vector field, and suppose that $w = \langle \nabla u, X \rangle$ satisfies $w \ge 0$, $w \not\equiv 0$ on $\Omega$. Then, for each $\varphi \in \lip_c(M)$ it holds
%	\begin{equation}\label{piufinalissima_teo}
%		\begin{array}{l}
%			\disp \int_{\Omega} \varphi^2 \Big\{\lambda_1 \big|\nabla^\top   u_\nu\big|^2 +\lambda_2\Big[u_\nu^2|\II|^2 +\Ric(\nabla u, \nabla u)\Big]\Big\} \\[0,5cm]
%			\disp + \int_\Omega w^2 \langle A(\nabla u) \nabla \left(\frac{\varphi|\nabla u|}{w}\right), \nabla \left(\frac{\varphi|\nabla u|}{w}\right) \rangle \\[0.5cm]
%			\le \disp \int_\Omega |\nabla u|^2 \langle A(\nabla u) \nabla \varphi, \nabla \varphi \rangle.
%		\end{array}
%	\end{equation}
%\end{proposition}



%\begin{proposition} \label{prop_splitting}
%	Let $(M^n,\metric)$ be a complete Riemannian manifold, $\Omega$ an open subset with \tcr{$C^1$} %smooth
%	boundary and let $u\in C^2(\Omega) \cap C^1(\overline{\Omega})$ be a solution to
%	\[
%		\Delta_a u + f(u) = 0
%	\]
%	satisfying $\nabla u\neq 0$ everywhere in $\Omega$,
%	\[
%		\big|\nabla^\top  |\nabla u|\big|^2 \equiv 0 \, , \qquad |\II|^2 \equiv 0 \qquad \text{in } \, \Omega
%	\]
%	and such that \tcr{$u$ is locally constant on $\partial\Omega$.} %and $\partial_\eta u$ are locally constant on $\partial\Omega$. %\tcr{and assume that $\overline{u(\partial\Omega)}$ is not an interval [serve davvero?]}.
%	Then,
%	\begin{itemize}
%		\item [i)] there exists an isometry $\Psi : I \times N \to \Omega$, where $I\subseteq\R$ is an open interval and $N$ is a complete, totally geodesic hypersurface of $M$.
%		\item [ii)] $u(\Psi(t,x)) = u_0(t)$ for all $(t,x)\in I\times N$, for a suitable function $u_0 : I \to \R$ such that $u_0'>0$.
%	\end{itemize}
%\end{proposition}
%
%\begin{proof}
%	By the identities 1) and 2) in \eqref{identities}, the only nontrivial component of $\nabla \di u$ is the one in the direction of $\nu = \nabla u/|\nabla u|$:
%	\begin{equation}\label{eq_Ddu}
%		\nabla \di u = \nabla \di u\left(\nu,\nu\right) \frac{\di u}{|\nabla u|} \otimes \frac{\di u}{|\nabla u|} \qquad \text{on } \, \Omega.
%	\end{equation}
%	A straightforward computation allows us to deduce from \eqref{eq_Ddu} that $|\nabla u|$ is locally constant on level sets of $u$, that integral curves of $\nu$ are geodesics in $M$ and that level sets of $u$ are totally geodesic in $\Omega$. Fix $b\in u(\Omega)\backslash u(\partial\Omega)$, which exists since $u(\partial\Omega)$ is at most countable and $u$ is non-constant. %Fix  $b\in u(\Omega)\backslash \overline{u(\partial\Omega)}$. Note that such a $b$ exists: indeed, we have
%	%\[
%	%	\overline{u(\partial\Omega)} \subseteq \overline{u(\overline\Omega)} = \overline{u(\Omega)} 
%	%\] 
%	%where $\overline{u(\Omega)}$ is a closed interval while, by assumption, $\overline{u(\partial\Omega)}$ is not an interval. Hence, there exists an interior point $b$ of $\overline{u(\Omega)}$ (that is, a point of $u(\Omega)$) not belonging to $\overline{u(\partial\Omega)}$, as desired.
%	Let $N \subseteq \Omega$ be a connected component of $\{u = b\}$. Observe that $N$ is closed in $\overline{\Omega}$ (by continuity of $u$ in $\overline{\Omega}$), hence it is closed in $M$. %By the implicit function theorem, \tcr{$N$ is properly embedded in $\Omega$ [ci serve?]}. 
%	%Moreover, by considering an interval $I_0 = (b-\eps,b+\eps)$ disjoint from $\overline{u(\partial\Omega)}$ and 
%	%$\psi \in C^\infty_c(I_0)$ with $\psi(t) \equiv t$ in a neighbourhood of $b$, the function $\psi(u)$ can be extended to the entire $M$. Again by the implicit function theorem $\{\psi(u)=b\}$, thus $N$, is properly embedded in the entire $M$, \tcr{and is therefore a complete manifold without boundary [è necessario che $N$ sia embedded per dedurre ciò?]}.
%	Furthermore, $N$ is also embedded and totally geodesic in $\Omega$, therefore in $M$. In particular, $N$ is also complete. Let
%	\[
%		\Psi : \R \times N \to M : (t,x) \mapsto \exp_x(t\nu_x) \, .
%	\]
%	Note that $\Psi$ is well defined and surjective by completeness of $M$ and closedness of $N$, %and the proper embeddedness of $N$,
%	and that $\Psi^{-1}(\Omega)$ is open in $\R\times N$. For each $x\in N$ we denote by $I_x \subseteq \R$ the maximal open interval containing $0$ such that $\Psi(t,x) \in \Omega$ for all $t\in I_x$ and we denote by $\gamma_x : I_x \to \Omega$ the geodesic curve given by
%	\[
%		\gamma_x(t) = \Psi(t,x) \qquad \text{for all } \, t \in I_x \, .
%	\]
%	For each $x\in N$ we also set
%	\[
%		t_\ast(x) = \inf I_x \, , \qquad t^\ast(x) = \sup I_x \, .
%	\]
%	%Note that if $t^\ast(x) = \infty$ then $\Psi(t,x) \in \Omega$ for all times $t\in\R^+_0$, while if $t^\ast(x) < \infty$ then $\Psi(t^\ast(x),x)\in\partial\Omega$, and in fact in this second case the value $t = t^\ast(x)$ is the first (positive) time at which the geodesic curve $\R^+_0 \ni t \mapsto \Psi(t,x)$ hits $\partial\Omega$.
%	Note that if $t_\ast(x) > -\infty$ then $\Psi(t_\ast(x),x) \in \partial\Omega$. Similarly, $\Psi(t^\ast(x),x) \in \partial\Omega$ if $t^\ast(x) < \infty$.
%	
%	We define
%	\[
%		\mathscr D = \{ (t,x) : x \in N, \, t \in I_x \} \subseteq \Psi^{-1}(\Omega) \, .
%	\]
%	Our ultimate goal is to show that $\mathscr D = I \times N$ for a fixed open interval $I\subseteq\R$, that $\Psi : \mathscr D \to \Omega$ is an isometry when $\mathscr D$ is equipped with the product metric $\di t^2 + \metric_{|N}$ and that $u$ is of the form $u(\Psi(t,x)) = u_0(t)$ for some function $u_0 : I \to \R$.
%	
%	First, we show that $\mathscr D$ is open in $\Psi^{-1}(\Omega)$, and therefore also in $\R\times N$. Suppose, by contradiction, that this is false, that is, suppose that $\Psi^{-1}(\Omega)\backslash\mathscr D$ is not closed in $\Psi^{-1}(\Omega)$. Then there exist $(t,x) \in \mathscr D$ and a sequence $\{(t_n,x_n)\} \in \Psi^{-1}(\Omega)\backslash\mathscr D$ such that $(t_n,x_n) \to (t,x)$. Suppose, without loss of generality, that $t>0$. Then we also have $t_n > 0$ for all sufficiently large $n\in\N$. Since $(t_n,x_n)\not\in\mathscr D$, that is, $t_n \not\in I_{x_n}$, this implies that $t^\ast(x_n) \leq t_n$ for all sufficiently large $n$. Since $0 \leq t^\ast(x_n) \leq t_n$ and $t_n \to t < \infty$, the sequence $t^\ast(x_n)$ is bounded, so there exists a subsequence $x_{n_k}$ such that $t^\ast(x_{n_k})$ converges to some limit $\bar t \in [0,t]$. By continuity, we have $\Psi(t^\ast(x_{n_k}),x_{n_k}) \to \Psi(\bar t,x)$ as $k\to\infty$. Since $\bar t \in [0,t] \subseteq I_x$, we have $\Psi(\bar t,x) \in \Omega$. But then $\{\Psi(t^\ast(x_{n_k}),x_{n_k})\}$ is a sequence of points of $\partial\Omega$ converging to a point of $\Omega$, absurd.
%	
%	From \eqref{eq_Ddu} we deduce that the vector field $\nu$ is parallel in $\Omega$. We let
%	\[
%		\Phi : E \subseteq \R \times \Omega \to \Omega
%	\]
%	be the maximally defined flow of $\nu$. By the very definition of $\Psi$ and $\mathscr D$ and by parallelism of $\nu$, we have $\mathscr D \subseteq E$ and
%	\[
%		\Psi(t,x) = \Phi(t,\Psi(0,x)) \equiv \Phi(t,x) \qquad \text{for all } \, (t,x) \in \mathscr D \, .
%	\]
%	%Since $\Psi$ maps some neighbourhood of $\{0\}\times N$ in $\mathscr D$ diffeomorphically onto a neighbourhood of $N$ in $\Omega$, we infer that the restriction $\Psi_{|\mathscr D} : \mathscr D \to \Omega$ is an open map and a local diffeomorphism onto its image (in fact, we will show that it is also injective, hence a global diffeomorphism). In particular, $\Psi(\mathscr D)$ is open in $\Omega$.
%	Moreover, since $\nu$ is a parallel vector field and since $\Psi \equiv \Phi_\nu$ on $\mathscr D$, the pullback $\Psi_{|\mathscr D}^\ast \metric$ of the metric $\metric$ of $\Omega$ induced by $\Psi_{|\mathscr D}$ on $\mathscr D$ is the product metric
%	\[
%		\di t^2 + \metric_{|N} \, .
%	\]
%	Indeed, for each $(t,x)\in\mathscr D$ we have
%	\begin{align*}
%		(\Psi_{|\mathscr D}^\ast \metric)(\partial_t,\partial_t) & = \langle\Psi_\ast\partial_t,\Psi_\ast\partial_t\rangle \underset{(i)}{=} \langle\nu,\nu\rangle = 1 \\
%		(\Psi_{|\mathscr D}^\ast \metric)(\partial_t,X) & = \langle\Psi_\ast\partial_t,\Psi_\ast X\rangle = \langle(\Phi_t)_\ast\nu,(\Phi_t)_\ast X\rangle \underset{(ii)}{=} \langle\nu,X\rangle = 0 \\
%		(\Psi_{|\mathscr D}^\ast \metric)(X,Y) & = \langle\Psi_\ast X,\Psi_\ast Y\rangle = \langle(\Phi_t)_\ast X,(\Phi_t)_\ast Y\rangle \underset{(iii)}{=} \langle X,Y\rangle = \langle X,Y\rangle_N
%	\end{align*}
%	for any $X,Y\in TN$, where $(i)$ holds because $\Psi$ is a restriction of the exponential map and $(ii)$-$(iii)$ hold because %For each $t\in\R$, denoting by $\Omega_t = \{ y \in \Omega : (t,y) \in E \}$ the (open) subset of $\Omega$ of those points for which the flow of $\nu$ starting from $x$ is defined up to time $t$, we have that
%	the map $y \mapsto \Phi_t(y) \doteq \Phi(t,y)$ is a local isometry by parallelism of $\nu$. In particular, $\Psi_{|\mathscr D}$ is an open map and a local diffeomorphism onto its image (in fact, we will show that it is also injective, hence a global diffeomorphism) and $\Psi(\mathscr D)$ is open in $\Omega$.
%		
%	Let $c>0$ be the constant value of $|\nabla u|$ on $N$ and let $y_0 : J \to \R$ be the maximal solution to the Cauchy problem
%	\[
%		\begin{cases}
%			[a(y')+y'a'(y')] y'' + f(y) = 0 \\
%			y(0) = b \\
%			y'(0) = c
%		\end{cases}
%	\]
%	subject to the constraint
%	\[
%		y_0'(t) \neq 0 \qquad \text{for all } \, t \in J \, .
%	\]
%	We claim that $I_x \subseteq J$ for all $x\in N$ and that $u(\Psi(t,x)) = y_0(t)$ for all $x\in N$, $t\in I_x$. For a given $x\in N$, consider the function $u_x = u\circ\gamma_x : I_x \to \R$. We have
%	\[
%		u_x' = |\nabla u|\circ\gamma_x > 0 \, , \qquad u_x'' = \nabla\di u(\nu,\nu) \circ \gamma_x
%	\]
%	so $u_x$ satisfies
%	\[
%		[a(u_x') + u_x' a'(u_x')] u_x'' + f(u_x) = 0 \qquad \text{on } \, I_x
%	\]
%	and $u_x(0) = b$, $u_x'(0) = c$. Hence, the claim follows by uniqueness of the local solution to the Cauchy problem for the ODE
%	\[
%		[a(y')+y'a'(y')]y'' + f(y) = 0
%	\]
%	for any initial datum $(y(t_0),y'(t_0)) \in \R\times\R^+$.
%	
%	We now show that the intervals $I_x$, $x\in N$, are all equal to some given interval $I\subseteq\R$ independent from $x$. This is equivalent to saying that for each $t\in\R$ the set
%	\[
%		N_t = \{ x \in N : t \in I_x \}
%	\]
%	is either empty or equal to $N$. We prove this latter claim. Let $t\in\R$ be such that $N_t\neq\emptyset$, then we prove that $N_t = N$ by showing that $N_t$ is both open and closed in the connected $N$. The first property is an immediate consequence of openness of $\mathscr D$: since $\mathscr D$ is open in $\R\times N$, the intersection $\mathscr D \cap (\{t\}\times N)$ is also open in $\{t\}\times N$; since $\mathscr D \cap (\{t\}\times N) \simeq N_t$ and $\{t\}\times N \simeq N$ via the trivial homeomorphism $\psi : \{t\}\times N \to N : (t,x) \mapsto x$, it follows that $N_t$ is open in $N$. %set $N_t = \{ x \in N : (t,x) \in \}$ Without loss of generality, we assume $t>0$ (otherwise, all arguments in the following can be equivalently rephrased by replacing $t^\ast$ with $t_1$ at all occurrences). Suppose, by contradiction, that $N_t$ is not open in $N$. Then there exist $x\in N_t$ and a sequence $\{x_n\} \subseteq N\backslash N_t$ such that $x_n \to x$. Since $x_n \not \in N_t$, we have $t^\ast(x_n) \leq t < \infty$ for each $n\in\N$. The sequence $\{p_n\} = \{\Psi(t^\ast(x_n),x_n)\}$ is contained in $\partial\Omega$. Since $0 < t^\ast(x_n) \leq t$ for each $n$, there exists a subsequence $\{x_{n_k}\}$ such that $t^\ast(x_{n_k})$ converges to some $\bar t \in [0,t]$. Then by continuity of $\Psi$ we have $\partial\Omega \ni p_{n_k} \to \Psi(\bar t,x) \in \Omega$, contradiction. This shows that $N_t$ is open.
%	We now prove that $N_t$ is closed in $N$. Without loss of generality, assume that $t>0$. Let $\{x_n\} \subseteq N_t$ be a sequence converging to some point $x\in N$. Suppose, by contradiction, that $x\not\in N_t$. Then we have $t^\ast(x) \leq t < t^\ast(x_n)$. By continuity,
%	\[
%		|\nabla u|(\Psi(t^\ast(x),x)) = \lim_{n\to\infty} |\nabla u|(\Psi(t^\ast(x),x_n)) = \lim_{n\to\infty} u_{x_n}'(t^\ast(x)) = y_0'(t^\ast(x)) > 0.
%	\]
%	Therefore, $\nabla u\neq 0$ at $\Psi(t^\ast(x),x)$. Since $\Psi(t^\ast(x),x) \in \partial\Omega$ and $u$ is locally constant on $\partial\Omega$, the gradient $\nabla u = |\nabla u|\dot\gamma_x$ must be orthogonal to $\partial\Omega$ at $\Psi(t^\ast(x),x)$, that is, the curve $s \mapsto \Psi(s,x) \equiv \gamma_x(s)$ meets $\partial\Omega$ orthogonally at time $s = t^\ast(x)$. Since $\partial\Omega$ is an embedded smooth hypersurface of $M$, for each large enough $n\in\N$ the curve $s \mapsto \Psi(s,x_n)$ must intersect $\partial\Omega$ for some value $s_n\in\R$ of $s$ such that $s_n \to t^\ast(x)$ as $n\to\infty$. Since $t^\ast(x) \leq t < t^\ast(x_n) \leq s_n$, by comparison it must be $t^\ast(x_n) \to t^\ast(x)$ as $n\to\infty$, hence $\Psi(t^\ast(x_n),x_n) \to \Psi(t^\ast(x),x)$. Since $\partial\Omega$ is smooth and embedded, for all sufficiently large $n\in\N$ the point $\Psi(t^\ast(x_n),x_n)$ belongs to the same connected component of $\partial\Omega$ containing $\Psi(t^\ast(x),x)$. Since $u$ is constant along that component of $\partial\Omega$, it must be
%	\[
%		u(\Psi(t^\ast(x),x)) = u(\Psi(t^\ast(x_n),x_n))
%	\]
%	for all sufficiently large $n\in\N$. By continuity, we have
%	\[
%		u(\Psi(t^\ast(x),x)) = \lim_{s\to t^\ast(x)} u(\Psi(s,x)) = \lim_{s\to t^\ast(x)} y_0(s) = y_0(t^\ast(x)) \, .
%	\]
%	On the other hand, since $y_0$ is strictly increasing and $t^\ast(x) \leq t < t^\ast(x_n)$, we also have
%	\[
%		y_0(t^\ast(x)) < \lim_{s\to t^\ast(x_n)} y_0(s) = \lim_{s\to t^\ast(x_n)} u(\Psi(s,x_n)) = u(\Psi(t^\ast(x_n),x_n))
%	\]
%	for each $n\in\N$, contradiction. This shows that $N_t$ is also closed. By connectedness of $N$ and since $N_t\neq\emptyset$ we deduce $N_t=N$.
%
%	So far, we have proved that
%	\[
%		\mathscr D = I \times N
%	\]
%	for some open interval $I\subseteq J$, that $\Psi_{|\mathscr D} : \mathscr D \to \Omega$ is a local isometry (when $\mathscr D$ is equipped with the product metric), that $\Psi(\mathscr D)$ is open in $\Omega$ and that $u$ satisfies $u(\Psi(t,x)) = y_0(t)$ for all $(t,x) \in \mathscr D$ for a fixed function $y_0 : J \to \R$. Since $y_0$ is strictly increasing on $J$ and integral curves of $\nu$ are disjoint, we further deduce that $\Psi_{|\mathscr D}$ is injective, hence a global isometry onto its image. We now prove that $\Psi(\mathscr D)$ is also closed in $\Omega$, thus concluding $\Psi(\mathscr D) = \Omega$ by connectedness of $\Omega$. Let $\{p_n\}\subseteq \Psi(\mathscr{D})$ be a given sequence converging to some point $p \in \Omega$. We have to show that $p \in \Psi(\mathscr{D})$. We consider the (unique) sequence $\{(t_n,x_n)\} \subseteq \mathscr D$ such that $p_n = \Psi(t_n,x_n)$ for each $n\in\N$. Let $U\subseteq\Omega$ be a geodesically convex open neighbourhood of $p$, which exists by openness of $\Omega$. For all $n,m\in\N$ large enough we have $p_n,p_m\in U$ and there exists a minimizing geodesic $\gamma = \gamma_{n,m} : [0,1] \to U$ such that $\gamma(0) = p_n$ and $\gamma(1) = p_m$. Let
%	\[
%		v = \dot\gamma(0) - \langle\dot\gamma(0),\nu\rangle \nu
%	\]
%	be the component of $\dot\gamma(0)$ orthogonal to $\nu$ and let $\sigma : [0,1] \to N$ be the geodesic with initial point $x_n$ and initial velocity $(\Psi(t_n,\,\cdot\,)^{-1})_\ast v \in T_{x_n} N$, which exists by completeness of $N$. Consider the curve $c : [0,1] \to \Omega$ defined by
%	\[
%		c(s) = \Psi(t_n+s\langle\dot\gamma(0),\nu\rangle,\sigma(s)) \, .
%	\]
%	We claim that $c(s) = \gamma(s)$ for all $s\in[0,1]$. Indeed, by construction we have $c(0) = \gamma(0)$ and $\dot c(0) = \dot\gamma(0)$. Let
%	\[
%		\bar s = \sup\{ s_0 \in [0,1] : t_n + s\langle\dot\gamma(0),\nu\rangle \in I \text{ for all } s \in [0,s_0] \} \, .
%	\]
%	The restriction $c_{|[0,\bar s]} : [0,\bar s] \to \Omega$ is a geodesic, and by uniqueness of geodesics we have $c(s) = \gamma(s)$ for all $s\in[0,\bar s]$. Clearly $\bar s > 0$, since $t_n \in I$ and $I$ is open. In fact, we have $\bar s = 1$, since otherwise $t_n + \bar s\langle\dot\gamma(0),\nu\rangle$ would be a boundary point for $I$, yielding $\gamma(\bar s) = c(\bar s) \in \partial\Omega$, contradiction. So, we have $c(s) = \gamma(s)$ for all $s\in[0,1]$, as claimed. In particular,
%	\[
%		\Psi(t_n + \langle\dot\gamma(0),\nu\rangle,\sigma(1)) = c(1) = \gamma(1) = \Psi(t_m,x_m)
%	\]
%	and by injectivity of $\Psi_{|\mathscr D}$ this implies
%	\[
%		t_n + \langle\dot\gamma(0),\nu\rangle = t_m \qquad \text{and} \qquad \sigma(1) = x_m \, .
%	\]
%	Since $\gamma([0,1])$ lies entirely in $\Psi(\mathscr D)$, we have
%	\[
%		\dist_\Omega(p_n,p_m)^2 = |\dot\gamma(s)|^2 = \langle\dot\gamma(0),\nu\rangle^2 + |\dot\sigma(s)|^2 \geq |t_n-t_m|^2 + \dist_N(x_n,x_m)^2
%	\]
%%	The level sets of $u$ provide a foliation of $U$ by totally geodesic hypersurfaces with parallel normal vector field $\nu$. This, as above, induces a splitting of the metric $\metric$ on $U$ as a product metric. Up to restricting $U$, we can assume that it is isometric to a product
%%	\[
%%		((-\eps,\eps) \times \Sigma, \di y^2 + \metric_\Sigma) \, .
%%	\]
%%	Let us denote by $f : U \to (-\eps,\eps)\times\Sigma$ this isometry, so that $f_\ast \nu = \partial_y$. Then
%%	\[
%%		\gamma_{n,m}(s) = f^{-1}(c_1+c_2s,c(s))
%%	\]
%%	for some $c_1,c_2\in\R$ and some geodesic $c : [0,1] \to \Sigma$. The curve $\gamma : [0,1] \to U$ given by
%%	\[
%%		\gamma(s) = f^{-1}(c_1,c(s))
%%	\]
%%	is also a geodesic, with initial point $\gamma(0) = \gamma_{n,m}(0) = p_n = \Psi(t_n,x_n)$ and initial velocity $\dot\gamma(0) \in \nu^\bot$. Since $\Psi(t_n,N)$ is totally geodesic in $\Omega$, we also have
%%	\[
%%		\gamma(s) = \Psi(t_n,\sigma(s))
%%	\]
%%	where $\sigma : [0,1] \to N$ is the geodesic of $N$ with initial point $x_n$ and initial velocity $(\Psi(t_n,\,\cdot\,)^{-1})_\ast\dot\gamma(0) \in T_{x_n} N$, which exists by completeness of $N$. Then \tcr{[qualche dettaglio in più]} the geodesic $\gamma_{n,m}$ is of the form
%%	\[
%%		\gamma_{n,m}(s) = \Psi((1-s)t_n+s t_m,\sigma_{n,m}(s))
%%	\]
%%	and we have
%%	\[
%%		\dist_\Omega(p_n,p_m)^2 = |t_n-t_m|^2 + \dist_N(x_n,x_m)^2 \, .
%%	\]
%	Hence, $\{t_n\}$ and $\{x_n\}$ are Cauchy sequences in $\R$ and $N$, respectively, so we have the existence of $(t,x) \in \R\times N$ such that $(t_n,x_n) \to (t,x)$. By continuity of $\Psi$ we also have $p_n = \Psi(t_n,x_n) \to \Psi(t,x)$, so $p = \Psi(t,x)$. If $t\not\in I$ then we would have either $t = \sup I$ or $t = \inf I$; in both cases, $\Psi(t,x_n)$ would be a sequence of points of $\partial\Omega$ converging to the point $p\in\Omega$, absurd. Hence, it must be $t\in I$ and so $p\in\Psi(\mathscr D)$, as desired.
%\end{proof}
%
%\textbf{Alternativa: splitting di ogni componente connessa di $\Omega\backslash\{\nabla u = 0\}$}

\begin{remark}
	In the above proof, since $w>0$ in $\Omega_0$ then $\critu\cap\Omega_0 = \emptyset$, whence $\Omega_0$ is a connected component of $\Omega\backslash\critu$.
\end{remark}

Theorem \ref{teo_Poincare_loc} will be used later to infer that solutions with moderate energy growth satisfy $\nabla^\top|\nabla u| \equiv 0$ and $\II \equiv 0$ in $\Omega_0$. Once this is achieved, the next proposition provides a general splitting result, of independent interest.



\begin{proposition}[\textbf{The splitting lemma}] \label{prop_splitting_alt}
	Let $(M,\metric)$ be a complete Riemannian manifold and $\Omega\subseteq M$ a $C^1$ domain. Let $a$ satisfy \eqref{assu_A}, $f \in C^1(\R)$ and let $u\in C^1(\overline{\Omega})$ be a non-constant weak solution to
	 
%	$a\in C(\R^+_0) \cap C^1(\R^+)$ satisfy
%	\[
%	ta'(t) + a(t) > 0 \qquad \text{for all } \, t \in \R^+ \, ,
%	\]
	\[
		\left\{
			\begin{array}{ll}
				\Delta_a u + f(u) = 0 & \qquad \text{in } \, \Omega \\[0.2cm]
				u \text{ locally constant} & \qquad \text{on } \, \partial\Omega \, .
			\end{array}
		\right.
	\]
	If $\Omega_0$ is a connected component of $\Omega\backslash\critu$ such that
	\begin{equation} \label{Ddu_trivial}
		u \in C^2(\Omega_0), \qquad \big|\nabla^\top  |\nabla u|\big|^2 \equiv 0 \qquad \text{and} \qquad |\II|^2 \equiv 0 \qquad \text{in } \, \Omega_0
	\end{equation}
	then:
	\begin{itemize}
		\item [i)] there exists an isometry $\Psi : I \times N \to \Omega_0$, where $I\subseteq\R$ is an open interval and $N$ is a complete, totally geodesic hypersurface of $M$;
		\item [ii)] $u(\Psi(t,x)) = u_0(t)$ for all $(t,x)\in I\times N$, for a suitable function $u_0 : I \to \R$ such that $u_0'>0$.
	\end{itemize}
\end{proposition}

\begin{remark}
	Observe that we assume no regularity at all on the the portion $\partial \Omega_0 \cap \critu$. The fact that the latter is totally geodesic, in particular smooth, is a consequence of the proof.	
\end{remark}

\begin{proof}
	By the identities 1) and 2) in \eqref{identities} and by \eqref{Ddu_trivial}, the only nontrivial component of $\nabla \di u$ in $\Omega_0$ is the one in the direction of $\nu = \nabla u/|\nabla u|$:
	\begin{equation}\label{eq_Ddu}
		\nabla \di u = \nabla \di u\left(\nu,\nu\right) \frac{\di u}{|\nabla u|} \otimes \frac{\di u}{|\nabla u|} \qquad \text{in } \, \Omega_0 .
	\end{equation}
	A straightforward computation allows us to deduce from \eqref{eq_Ddu} that $|\nabla u|$ is locally constant on level sets of $u$ intersected with $\Omega_0$ and that $\nu$ is a parallel vector field. Indeed, from the standard identity $\di|\nabla u|^2 = 2\nabla\di u(\nabla u,\,\cdot\,)$ we have $\di|\nabla u| = \nabla\di u(\nu,\,\cdot)$ in $\{\nabla u\neq 0\}$, so by \eqref{eq_Ddu} we get
	\[
		\nabla|\nabla u| = \nabla\di u(\nu,\nu) \nu \qquad \text{in } \, \Omega_0
	\]
	from which we have
	\begin{align*}
		\langle\nabla_X \nu,Y\rangle & = \frac{\langle\nabla_X \nabla u,Y\rangle}{|\nabla u|} - \frac{\langle X,\nabla|\nabla u|\rangle\langle\nabla u,Y\rangle}{|\nabla u|^2} \\
		& = \frac{\nabla\di u(X,Y)}{|\nabla u|} - \frac{\nabla\di u(\nu,\nu)\langle\nu, X\rangle\langle\nabla u,Y\rangle}{|\nabla u|^2} \\
		& = \frac{1}{|\nabla u|} \left(\nabla\di u(X,Y) - \nabla\di u(\nu,\nu)\langle\nu,X\rangle\langle\nu,Y\rangle\right) = 0
	\end{align*}
	for any pair of tangent vectors $X,Y\in T_x M$, $x\in\Omega_0$, that is, $\nabla\nu = 0$ in $\Omega_0$. In particular, integral curves of $\nu$ are geodesics in $M$. By the very definition of $\II$ and by \eqref{Ddu_trivial} we also have that level sets of $u$ intersected with $\Omega_0$ are totally geodesic hypersurfaces in $M$. 
	%A straightforward computation allows us to deduce from \eqref{eq_Ddu} that $|\nabla u|$ is locally constant on level sets of $u$ \tcr{intersected with $\Omega_0$}, that integral curves of $\nu$ are geodesics in $M$ that level sets of $u$ \tcr{intersected with $\Omega_0$} are totally geodesic in $M$.
	
	Fix $b\in u(\Omega_0)\backslash u(\partial\Omega)$, which exists since $u(\partial\Omega)$ is at most countable and $u$ is non-constant. Let $N \subseteq \Omega_0$ be a connected component of $\Sigma_b \doteq \{x\in\Omega_0 : u(x) = b\}$ and let $c>0$ be the constant value of $|\nabla u|$ on $N$. Since $\Sigma_b$ is closed in $\Omega_0$ and $N$ is a connected component of $\Sigma_b$, we have that $N$ is closed in $\Sigma_b$, hence in $\Omega_0$. We claim that $N$ is also closed in $M$. This is equivalent to saying that the closure $\overline{N}$ of $N$ in $M$ is contained in $\Omega_0$, that is, that it does not intersect $\partial\Omega_0\equiv\overline{\Omega}_0\backslash\Omega_0$. Indeed, by continuity of $u$ and $|\nabla u|$ in $\overline{\Omega}$ we have $u\equiv b$ and $|\nabla u|\equiv c$ on $\overline{N}$; since $\partial\Omega_0$ is contained in $\critu \cup \partial\Omega$, $u\neq b$ on $\partial\Omega$ and $|\nabla u| = 0 \neq c$ on $\critu$, we have $\overline{N}\cap\partial\Omega_0 = \emptyset$ as claimed. Furthermore, $N$ is also embedded and totally geodesic in $\Omega$, therefore in $M$. In particular, $N$ is also complete. Let
	\[
		\Psi : \R \times N \to M : (t,x) \mapsto \exp_x(t\nu_x) \, .
	\]
	Note that $\Psi$ is well defined and surjective by completeness of $M$ and closedness of $N$, and that $\Psi^{-1}(\Omega_0)$ is open in $\R\times N$. For each $x\in N$ we denote by $I_x \subseteq \R$ the maximal open interval containing $0$ such that $\Psi(t,x) \in \Omega_0$ for all $t\in I_x$ and we denote by $\gamma_x : I_x \to \Omega_0$ the geodesic curve given by
	\[
		\gamma_x(t) = \Psi(t,x) \qquad \text{for all } \, t \in I_x \, .
	\]
	For each $x\in N$ we also set
	\[
		t_\ast(x) = \inf I_x \, , \qquad t^\ast(x) = \sup I_x \, .
	\]
	Note that if $t_\ast(x) > -\infty$ then $\Psi(t_\ast(x),x) \in \partial\Omega_0$. Similarly, $\Psi(t^\ast(x),x) \in \partial\Omega_0$ if $t^\ast(x) < \infty$.
	
	We define
	\[
		\mathscr D = \{ (t,x) : x \in N, \, t \in I_x \} \subseteq \Psi^{-1}(\Omega_0) \, .
	\]
	Our ultimate goal is to show that $\mathscr D = I \times N$ for a fixed open interval $I\subseteq\R$, that $\Psi : \mathscr D \to \Omega_0$ is an isometry when $\mathscr D$ is equipped with the product metric $\di t^2 + \metric_{|N}$ and that $u$ is of the form $u(\Psi(t,x)) = u_0(t)$ for some function $u_0 : I \to \R$.
	
	First, we show that $\mathscr D$ is open in $\Psi^{-1}(\Omega_0)$, and therefore also in $\R\times N$. Suppose, by contradiction, that this is false, that is, suppose that $\Psi^{-1}(\Omega_0)\backslash\mathscr D$ is not closed in $\Psi^{-1}(\Omega_0)$. Then there exist $(t,x) \in \mathscr D$ and a sequence $\{(t_n,x_n)\} \in \Psi^{-1}(\Omega_0)\backslash\mathscr D$ such that $(t_n,x_n) \to (t,x)$. Suppose, without loss of generality, that $t>0$. Then we also have $t_n > 0$ for all sufficiently large $n\in\N$. Since $(t_n,x_n)\not\in\mathscr D$, that is, $t_n \not\in I_{x_n}$, this implies that $t^\ast(x_n) \leq t_n$ for all sufficiently large $n$. Since $0 \leq t^\ast(x_n) \leq t_n$ and $t_n \to t < \infty$, the sequence $t^\ast(x_n)$ is bounded, so there exists a subsequence $x_{n_k}$ such that $t^\ast(x_{n_k})$ converges to some limit $\bar t \in [0,t]$. By continuity, we have $\Psi(t^\ast(x_{n_k}),x_{n_k}) \to \Psi(\bar t,x)$ as $k\to\infty$. Since $\bar t \in [0,t] \subseteq I_x$, we have $\Psi(\bar t,x) \in \Omega$. But then $\{\Psi(t^\ast(x_{n_k}),x_{n_k})\}$ is a sequence of points of $\partial\Omega$ converging to a point of $\Omega$, absurd.
	
	From \eqref{eq_Ddu} we deduce that the vector field $\nu$ is parallel in $\Omega_0$. We let
	\[
		\Phi : E \subseteq \R \times \Omega_0 \to \Omega_0
	\]
	be the maximally defined flow of $\nu$. By the very definition of $\Psi$ and $\mathscr D$ and by parallelism of $\nu$, we have $\mathscr D \subseteq E$ and
	\[
		\Psi(t,x) = \Phi(t,\Psi(0,x)) \equiv \Phi(t,x) \qquad \text{for all } \, (t,x) \in \mathscr D \, .
	\]
	Moreover, since $\nu$ is a parallel vector field and since $\Psi \equiv \Phi_\nu$ on $\mathscr D$, the pullback $\Psi_{|\mathscr D}^\ast \metric$ of the metric $\metric$ of $\Omega_0$ induced by $\Psi_{|\mathscr D}$ on $\mathscr D$ is the product metric
	\[
		\di t^2 + \metric_{|N} \, .
	\]
	Indeed, for each $(t,x)\in\mathscr D$ we have
	\begin{align*}
		(\Psi_{|\mathscr D}^\ast \metric)(\partial_t,\partial_t) & = \langle\Psi_\ast\partial_t,\Psi_\ast\partial_t\rangle \underset{(i)}{=} \langle\nu,\nu\rangle = 1 \\
		(\Psi_{|\mathscr D}^\ast \metric)(\partial_t,X) & = \langle\Psi_\ast\partial_t,\Psi_\ast X\rangle = \langle(\Phi_t)_\ast\nu,(\Phi_t)_\ast X\rangle \underset{(ii)}{=} \langle\nu,X\rangle = 0 \\
		(\Psi_{|\mathscr D}^\ast \metric)(X,Y) & = \langle\Psi_\ast X,\Psi_\ast Y\rangle = \langle(\Phi_t)_\ast X,(\Phi_t)_\ast Y\rangle \underset{(iii)}{=} \langle X,Y\rangle = \langle X,Y\rangle_N
	\end{align*}
	for any $X,Y\in TN$, where $(i)$ holds because $\Psi$ is a restriction of the exponential map and $(ii)$-$(iii)$ hold because the map $y \mapsto \Phi_t(y) \doteq \Phi(t,y)$ is a local isometry by parallelism of $\nu$. In particular, $\Psi_{|\mathscr D}$ is an open map and a local diffeomorphism onto its image (in fact, we will show that it is also injective, hence a global diffeomorphism) and $\Psi(\mathscr D)$ is open in $\Omega_0$.
	
	Let $c>0$ be the constant value of $|\nabla u|$ on $N$ and let $y_0 : J \to \R$ be the maximal solution to the Cauchy problem
	\[
	\begin{cases}
		[a(y')+y'a'(y')] y'' + f(y) = 0 \\
		y(0) = b \\
		y'(0) = c
	\end{cases}
	\]
	subject to the constraint
	\[
		y_0'(t) \neq 0 \qquad \text{for all } \, t \in J \, .
	\]
	We claim that $I_x \subseteq J$ for all $x\in N$ and that $u(\Psi(t,x)) = y_0(t)$ for all $x\in N$, $t\in I_x$. For a given $x\in N$, consider the function $u_x = u\circ\gamma_x : I_x \to \R$. We have
	\[
		u_x' = |\nabla u|\circ\gamma_x > 0 \, , \qquad u_x'' = \nabla\di u(\nu,\nu) \circ \gamma_x
	\]
	so $u_x$ satisfies
	\[
		[a(u_x') + u_x' a'(u_x')] u_x'' + f(u_x) = 0 \qquad \text{on } \, I_x
	\]
	and $u_x(0) = b$, $u_x'(0) = c$. Hence, the claim follows by uniqueness of the local solution to the Cauchy problem for the ODE
	\[
		[a(y')+y'a'(y')]y'' + f(y) = 0
	\]
	for any initial datum $(y(t_0),y'(t_0)) \in \R\times\R^+$.
	
	We now show that the intervals $I_x$, $x\in N$, are all equal to some given interval $I\subseteq\R$ independent from $x$. This is equivalent to saying that for each $t\in\R$ the set
	\[
		N_t = \{ x \in N : t \in I_x \}
	\]
	is either empty or equal to $N$. We prove this latter claim. Let $t\in\R$ be such that $N_t\neq\emptyset$, then we prove that $N_t = N$ by showing that $N_t$ is both open and closed in the connected $N$. The first property is an immediate consequence of openness of $\mathscr D$: since $\mathscr D$ is open in $\R\times N$, the intersection $\mathscr D \cap (\{t\}\times N)$ is also open in $\{t\}\times N$; since $\mathscr D \cap (\{t\}\times N) \simeq N_t$ and $\{t\}\times N \simeq N$ via the trivial homeomorphism $\psi : \{t\}\times N \to N : (t,x) \mapsto x$, it follows that $N_t$ is open in $N$. We now prove that $N_t$ is closed in $N$. Without loss of generality, assume that $t>0$. Let $\{x_n\} \subseteq N_t$ be a sequence converging to some point $x\in N$. Suppose, by contradiction, that $x\not\in N_t$. Then we have $t^\ast(x) \leq t < t^\ast(x_n)$. In particular, $t^\ast(x) < \infty$ and so $\Psi(t^\ast(x),x) \in \partial\Omega_0 \subseteq \critu\cup\partial\Omega$. By continuity,
	\[
		|\nabla u|(\Psi(t^\ast(x),x)) = \lim_{n\to\infty} |\nabla u|(\Psi(t^\ast(x),x_n)) = \lim_{n\to\infty} u_{x_n}'(t^\ast(x)) = y_0'(t^\ast(x)) > 0.
	\]
	Therefore, $\nabla u\neq 0$ at $\Psi(t^\ast(x),x)$ and we have $\Psi(t^\ast(x),x) \in \partial\Omega$. Since $u$ is locally constant on $\partial\Omega$, the gradient $\nabla u = |\nabla u|\dot\gamma_x$ must be orthogonal to $\partial\Omega$ at $\Psi(t^\ast(x),x)$, that is, the curve $s \mapsto \Psi(s,x) \equiv \gamma_x(s)$ meets $\partial\Omega$ orthogonally at time $s = t^\ast(x)$. Since $\partial\Omega$ is an embedded $C^1$ hypersurface of $M$, for each large enough $n\in\N$ the curve $s \mapsto \Psi(s,x_n)$ must intersect $\partial\Omega$ for some value $s_n\in\R$ of $s$ such that $s_n \to t^\ast(x)$ as $n\to\infty$. Since $t^\ast(x) \leq t < t^\ast(x_n) \leq s_n$, by comparison it must be $t^\ast(x_n) \to t^\ast(x)$ as $n\to\infty$, hence $\Psi(t^\ast(x_n),x_n) \to \Psi(t^\ast(x),x)$. Since $\partial\Omega$ is $C^1$ and embedded, for all sufficiently large $n\in\N$ the point $\Psi(t^\ast(x_n),x_n)$ belongs to the same connected component of $\partial\Omega$ containing $\Psi(t^\ast(x),x)$. Since $u$ is constant along that component of $\partial\Omega$, it must be
	\[
		u(\Psi(t^\ast(x),x)) = u(\Psi(t^\ast(x_n),x_n))
	\]
	for all sufficiently large $n\in\N$. By continuity, we have
	\[
		u(\Psi(t^\ast(x),x)) = \lim_{s\to t^\ast(x)} u(\Psi(s,x)) = \lim_{s\to t^\ast(x)} y_0(s) = y_0(t^\ast(x)) \, .
	\]
	On the other hand, since $y_0$ is strictly increasing and $t^\ast(x) \leq t < t^\ast(x_n)$, we also have
	\[
		y_0(t^\ast(x)) < \lim_{s\to t^\ast(x_n)} y_0(s) = \lim_{s\to t^\ast(x_n)} u(\Psi(s,x_n)) = u(\Psi(t^\ast(x_n),x_n))
	\]
	for each $n\in\N$, contradiction. This shows that $N_t$ is also closed. By connectedness of $N$ and since $N_t\neq\emptyset$ we deduce $N_t=N$.
	
	So far, we have proved that
	\[
		\mathscr D = I \times N
	\]
	for some open interval $I\subseteq J$, that $\Psi_{|\mathscr D} : \mathscr D \to \Omega_0$ is a local isometry (when $\mathscr D$ is equipped with the product metric), that $\Psi(\mathscr D)$ is open in $\Omega_0$ and that $u$ satisfies $u(\Psi(t,x)) = y_0(t)$ for all $(t,x) \in \mathscr D$ for a fixed function $y_0 : J \to \R$. Since $y_0$ is strictly increasing on $J$ and integral curves of $\nu$ are disjoint, we further deduce that $\Psi_{|\mathscr D}$ is injective, hence a global isometry onto its image. We now prove that $\Psi(\mathscr D)$ is also closed in $\Omega_0$, thus concluding $\Psi(\mathscr D) = \Omega_0$ by connectedness of $\Omega_0$. Let $\{p_n\}\subseteq \Psi(\mathscr{D})$ be a given sequence converging to some point $p \in \Omega_0$. We have to show that $p \in \Psi(\mathscr{D})$. We consider the (unique) sequence $\{(t_n,x_n)\} \subseteq \mathscr D$ such that $p_n = \Psi(t_n,x_n)$ for each $n\in\N$. Let $U\subseteq\Omega_0$ be a geodesically convex open neighbourhood of $p$, which exists by openness of $\Omega_0$. For all $n,m\in\N$ large enough we have $p_n,p_m\in U$ and there exists a minimizing geodesic $\gamma = \gamma_{n,m} : [0,1] \to U$ such that $\gamma(0) = p_n$ and $\gamma(1) = p_m$. Let
	\[
		v = \dot\gamma(0) - \langle\dot\gamma(0),\nu\rangle \nu
	\]
	be the component of $\dot\gamma(0)$ orthogonal to $\nu$ and let $\sigma : [0,1] \to N$ be the geodesic with initial point $x_n$ and initial velocity $(\Psi(t_n,\,\cdot\,)^{-1})_\ast v \in T_{x_n} N$, which exists by completeness of $N$. Consider the curve $c : [0,1] \to M$ defined by
	\[
		c(s) = \Psi(t_n+s\langle\dot\gamma(0),\nu\rangle,\sigma(s)) \, .
	\]
	We claim that $c(s) = \gamma(s)$ for all $s\in[0,1]$. Indeed, by construction we have $c(0) = \gamma(0)$ and $\dot c(0) = \dot\gamma(0)$. Let
	\[
		\bar s = \sup\{ s_0 \in [0,1] : t_n + s\langle\dot\gamma(0),\nu\rangle \in I \text{ for all } s \in [0,s_0] \} \, .
	\]
	The restriction $c_{|[0,\bar s]} : [0,\bar s] \to \Omega_0$ is a geodesic, and by uniqueness of geodesics we have $c(s) = \gamma(s)$ for all $s\in[0,\bar s]$. Clearly $\bar s > 0$, since $t_n \in I$ and $I$ is open. In fact, we have $\bar s = 1$, since otherwise $t_n + \bar s\langle\dot\gamma(0),\nu\rangle$ would be a boundary point for $I$, yielding $\gamma(\bar s) = c(\bar s) \in \partial\Omega_0$, contradiction. So, we have $c(s) = \gamma(s)$ for all $s\in[0,1]$, as claimed. In particular,
	\[
		\Psi(t_n + \langle\dot\gamma(0),\nu\rangle,\sigma(1)) = c(1) = \gamma(1) = \Psi(t_m,x_m)
	\]
	and by injectivity of $\Psi_{|\mathscr D}$ this implies
	\[
		t_n + \langle\dot\gamma(0),\nu\rangle = t_m \qquad \text{and} \qquad \sigma(1) = x_m \, .
	\]
	Since $\gamma([0,1])$ lies entirely in $\Psi(\mathscr D)$, we have
	\[
		\dist_\Omega(p_n,p_m)^2 = |\dot\gamma(s)|^2 = \langle\dot\gamma(0),\nu\rangle^2 + |\dot\sigma(s)|^2 \geq |t_n-t_m|^2 + \dist_N(x_n,x_m)^2 .
	\]
	Hence, $\{t_n\}$ and $\{x_n\}$ are Cauchy sequences in $\R$ and $N$, respectively, so we have the existence of $(t,x) \in \R\times N$ such that $(t_n,x_n) \to (t,x)$. By continuity of $\Psi$ we also have $p_n = \Psi(t_n,x_n) \to \Psi(t,x)$, so $p = \Psi(t,x)$. If $t\not\in I$ then we would have either $t = \sup I$ or $t = \inf I$; in both cases, $\Psi(t,x_n)$ would be a sequence of points of $\partial\Omega_0$ converging to the point $p\in\Omega_0$, absurd. Hence, it must be $t\in I$ and so $p\in\Psi(\mathscr D)$, as desired.
\end{proof}

We are ready to prove Theorem \ref{teo_splitting_intro} in the Introduction, in the following more general form: 


\begin{theorem} \label{teo_splitting}
	Let $(M^n, \metric)$ be a complete Riemannian manifold. Let $\Omega$ be a $C^2$ domain such that $\Ric\geq 0$ in $\Omega$. Let $a$ satisfy \eqref{assu_A_superstrong}, $f \in C^1(\R)$ and let $u \in C^3(\Omega) \cap C^2(\overline\Omega)$ be a non-constant, stable solution to
	\[
	\left\{
	\begin{array}{l@{\qquad}l}
		\Delta_a u + f(u) = 0 & \text{in } \Omega \\[0.3cm]
		u, \partial_\eta u \text{ locally constant} & \text{on } \partial\Omega
	\end{array}
	\right.
	\]
	%\tcr{and such that $\overline{u(\partial\Omega)}$ is not an interval}.
	with moderate energy growth. Let $\Omega_0$ be a connected component of $\Omega\backslash\critu$ with locally finite perimeter in $\Omega$. Assume that $X$ is a bounded Killing vector field in  $\overline\Omega$ such that
	\begin{equation}\label{eq_monotonoc}
		\langle \nabla u, X \rangle > 0 \qquad \text{in } \, \Omega_0.
	\end{equation}
	%	and that $u$ has moderate energy growth. 
	%	\[
	%		\int_{\Omega \cap B_R} |\nabla u|^2 \lambda_{\max}(|\nabla u|) = O(R^2 \log R) \qquad \text{as } \, R \to \infty \, .
	%	\]
	Then,
	\begin{itemize}
		\item [i)] there exists an isometry $\Psi : I \times N \to \Omega_0$, where $I\subseteq\R$ is an open interval and $N$ is a complete, totally geodesic hypersurface of $M$ with $\Ric_N \geq 0$;
		\item [ii)] $u(\Psi(t,x)) = u_0(t)$ for all $(t,x)\in I\times N$, for some $u_0 : I \to \R$ with $u_0'>0$;
		\item [iii)] $\langle \nu,X \rangle$ is a positive constant in $\Omega_0$, where $\nu = \frac{\nabla u}{|\nabla u|}$.
	\end{itemize}
\end{theorem}

\begin{proof}[Proof of Theorem \ref{teo_splitting}]
Setting $w \doteq \langle \nabla u, X \rangle$, by Theorem \ref{teo_Poincare_loc} we have
	\begin{align*}
		\int_{\Omega_0} |\nabla u|^2 \langle A(\nabla u) \nabla \varphi, \nabla \varphi \rangle & \geq \int_{\Omega_0} \varphi^2 \Big\{\lambda_1 \big|\nabla^\top  |\nabla u|\big|^2 +\lambda_2\Big[|\nabla u|^2|\II|^2 +\Ric(\nabla u, \nabla u)\Big]\Big\} \\[0.2cm]
		& \phantom{=\;} + \int_{\Omega_0} w^2 \langle A(\nabla u) \nabla \left(\frac{\varphi|\nabla u|}{w}\right), \nabla \left(\frac{\varphi|\nabla u|}{w}\right) \rangle
	\end{align*}
	for all $\varphi\in\lip_c(M)$. 
%	By Lemma \ref{lem_phij} there exists a sequence $\{\varphi_j\} \subseteq \lip_c(M)$ such that $\varphi_j \to 1$ in $W^{1,\infty}_\loc(M)$ and
%	\[
%		\int_\Omega |\nabla u|^2 \lambda_{\max}(|\nabla u|) |\nabla\varphi_j|^2 \to 0
%	\]
%	as $j\to\infty$. 
	Testing the inequality with $\varphi = \varphi_j$ in the definition of moderate energy growth, and letting $j \to \infty$, we obtain
	\[
		|\nabla^\top  |\nabla u|\big|^2 \equiv 0 \, , \qquad |\II|^2 \equiv 0 \, , \qquad \Ric(\nabla u, \nabla u) \equiv 0 \, , \qquad \nabla \frac{|\nabla u|}{w} \equiv 0 \qquad \text{in } \Omega_0 \, .
	\]
	Proposition \ref{prop_splitting_alt} ensures the validity of clauses $i)$ and $ii)$ of the statement, %(where the claim $\Ric_N \geq 0$ follows from the %fact that the product metric on $I\times %N\simeq\Omega$ has non negative Ricci tensor), 
	while $iii)$ follows by observing that
	\[
		\frac{|\nabla u|}{w} = \frac{|\nabla u|}{\langle\nabla u,X\rangle} = \frac{1}{\langle \nu,X\rangle} \, .
	\]
\end{proof}

\begin{proof}[Proof of Theorem \ref{teo_splitting_intro}]
	Condition \eqref{condi_Hcj_con_a} guarantees that 
	\[
	\langle \nabla u, X \rangle \ge 0, \qquad \langle \nabla u, X \rangle \not\equiv 0 \qquad \text{on } \, \partial \Omega,
	\]
	and by Corollary \ref{cor_lowenergy} $u$ has moderate energy growth. Hence, using Theorem \ref{teo_monotonicity} we deduce 
	\[
	\langle \nabla u, X \rangle > 0 \qquad \text{in } \, \Omega,
	\]
	and we conclude by Theorem \ref{teo_splitting} applied to $\Omega_0 = \Omega$.
\end{proof}

%\begin{remark}\label{rem_use_observ}
%	The choice of \eqref{eq_monotonoc} rather than the more appealing sufficient condition 
%	\[
%	\langle \nabla u, X \rangle \ge 0, \qquad \langle \nabla u, X \rangle \not\equiv 0 \qquad \text{on } \, \partial \Omega,
%	\]
%	is to apply Theorem \ref{teo_splitting} to settings where \eqref{eq_monotonoc} holds but $\partial \Omega \subseteq \{\nabla u = 0\}$. This will be useful later on. 
%\end{remark}




%\textbf{Dimostrazione del teorema di spezzamento da JFA}
%
%\begin{proof}
%By the parabolicity of $\overline \Omega$, every compact subset $K\subseteq\overline{\Omega}$ has zero capacity in the manifold with boundary $(\overline{\Omega},\sigma)$, that is,
%$$
%\inf\left\{ \int_\Omega |D\phi|^2 \di x : \phi\in \lip_c(\overline{\Omega}), \, \phi \geq 1 \, \text{ on } \, K \right\} = 0.
%$$
%In particular, by \cite[Thm. 1.5]{imperapigolasetti}, there exists a sequence $\{\varphi_j\} \subseteq \lip_c(\overline{\Omega})$ satisfying
%\[
%\varphi_j \ra 1 \ \ \text{ in } \, W^{1,\infty}_\loc(\overline{\Omega}), \qquad \int_\Omega |D\varphi_j|^2 \ra 0
%\] 
%as $j \ra \infty$. For each $j\geq 1$, we apply Lemma \ref{lem_poinc} to deduce
%\[
%\begin{array}{l}
%	\disp \int_\Sigma \left[ W^2\left( \|\nabla_\top    \|\nabla u\|\|^2 + \|\nabla u\|^2 \|A\|^2 \right) + \frac{\Ric(Du,Du)}{W^2} \right] \varphi_j^2 \di x_g + \\[0.5cm]
%	\disp \qquad \qquad + \int_\Sigma \frac{\bar v^2}{W^2} \left\| \nabla \left(\frac{\varphi_j\|\nabla u\|W}{\bar v}\right) \right\|^2 \di x_g \le \int_\Sigma \|\nabla u\|^2 \|\nabla \varphi_j \|^2 \di x_g.
%\end{array}
%\]
%The assumed boundedness of $|Du|$ guarantees that there exists a constant $C_0$ such that $W \le C_0$ on $\Omega$. Hence, 
%\[
%\int_\Sigma \|\nabla u\|^2 \|\nabla \varphi_j\|^2 \di x_g \le \int_\Omega |D\varphi_j|^2 W \di x \le C_0 \int_\Omega |D\varphi_j|^2 \di x \to 0,
%\]
%thus letting $j \ra \infty$ and using Fatou's Lemma we deduce 
%\begin{equation}\label{eq_persplitting}
%	\|\nabla_\top    \|\nabla u\|\|^2 + \|\nabla u\|^2 \|A\|^2 \equiv 0, \qquad \nabla\left( \frac{\|\nabla u\|W}{\bar v} \right) \equiv 0 \qquad \text{on } \, \Sigma.
%\end{equation}
%
%%To justify applicability of Lemma \ref{lem_poinc} we observe that, since $\psi$ extends to a diffeomorphism between manifolds with boundary $(\Omega,\partial\Omega)$ and $(U,\partial U)$, we can use it to push forward to $\overline{U}$ the functions $W$, $u$, $\bar v$ and the metric $g$. Since the geometry of $M$ outside of $\overline{\Omega}$ plays no role in the proof of Lemma \ref{lem_poinc}, the argument there can be repeated verbatim with $U\subseteq N$ in place of $\Omega\subseteq M$.
%
%From now on, the splitting proceeds similarly to \cite{fmv}, with a few extra topological arguments. First, from $\bar v > 0$ we deduce $\di u \neq 0$ at every point of $\Omega$. For $x \in \Sigma$, define $\nu = \nabla u/\|\nabla u\|$ and let $\{e_\alpha\}$ be an orthonormal frame tangent to $\{u = u(x)\}$ in $\Sigma$. Because of the identities 
%\[
%\langle \nabla \|\nabla u\|, e_j \rangle  = \nabla \di  u(\nu, e_j), \qquad A_{\alpha\beta} = \frac{\nabla \di _{\alpha\beta} u}{\|\nabla u\|},
%\]
%with $1 \le j \le m$ and $A$ the second fundamental form of the hypersurface $\{ u = u(x)\}$ in $\Sigma$, the first in \eqref{eq_persplitting} implies that the only nonzero component of $\nabla \di  u$ (hence, by \eqref{hessian_on_graph_u}, of $D^2u$), is the one in the direction of $Du$:
%\begin{equation}\label{eq_D2u}
%	D^2 u = D^2 u\left( \frac{Du}{|Du|},\frac{Du}{|Du|} \right) \frac{\di u}{|Du|} \otimes \frac{\di u}{|Du|} \qquad \text{on } \, \Omega.
%\end{equation}
%A straightforward computation allows us to deduce from \eqref{eq_D2u} that $|Du|$ is locally constant on level sets of $u$, that integral curves of $Du /|Du|$ are geodesics in $M$, and that level sets of $u$ are totally geodesic both in $\Sigma$ and in $\Omega$. In the limit, we infer that each component of $\partial \Omega$ is totally geodesic. Let $N \subseteq \Omega$ be a connected component of a level set of $u$, say of $\{u = b\}$ with $b \not\in \{b_1,b_2, \ldots, b_{j_0}\}$. By the implicit function theorem, note that $N$ is properly embedded both in $\Omega$ and in $M$, and is therefore a complete manifold without boundary. We denote with $\Phi(t,x)$ the flow of $D u / |Du|$ starting from $N$, defined on the connected set
%\[
%\mathscr{D} \subseteq \R \times N, \qquad \mathscr{D} = \big\{ (t,x) : x \in N, \ t \in (t_1(x), t_2(x))\big\},
%\]
%where, for every $x\in N$, $t_1(x)\in[-\infty,0)$ and $t_2(x)\in(0,+\infty]$ are the extrema of the largest open interval $I_x = (t_1(x),t_2(x))$ such that for every $t\in I_x$ the point $\Phi(t,x)$ is well defined and belongs to $\Omega$. If $t_1(x) > -\infty$ (respectively, if $t_2(x) < +\infty$) then the curve $t \mapsto \Phi(t,x)$ converges to a point of $\partial\Omega$ as $t \searrow t_1(x)$ (resp., $t\nearrow t_2(x)$) which we shall denote as $x_\ast = \Phi(t_1(x)^+,x)$ (resp., $x^\ast = \Phi(t_2(x)^-,x)$). The function $t_1$ is upper semi-continuous on $N$, that is, for every $x\in N$ we have
%$$
%\limsup_{n\to \infty} t_1(x_n) \leq t_1(x)
%$$
%for every sequence $\{x_n\} \subseteq N$ converging to $x$: otherwise, we could find $t\in(t_1(x),0]$ and a sequence $\{x_n\}$ converging to $x$ such that $t_1(x_n) \to t$, yielding $\partial\Omega \ni (x_n)_\ast \to \Phi(t,x) \in \Omega$, absurd. Similarly, the function $t_2$ is lower semi-continuous on $N$. Hence, $\mathscr{D}$ is open in $\R \times N$. From \eqref{eq_D2u} we deduce
%\[
%\left( D_V \frac{Du}{|Du|}, W \right) = 0 \qquad \forall \, V,W \in T\Omega,
%\]
%thus $Du /|Du|$ is a parallel vector field.	By standard theory, the induced metric on $\mathscr{D}$ by $\Phi$ is the product metric $\di t^2 + \sigma_{|N}$. Let $c_0>0$ be the constant value of $|Du|$ on $N$ and let $\beta$ be the maximal solution of the Cauchy problem
%\begin{equation}\label{eq_beta_1}
%	\begin{cases}
%		\beta' = H \dfrac{(1+\beta^2)^{3/2}}{\beta}, \\
%		\beta(b) = c_0.
%	\end{cases}
%\end{equation}
%Since $u$ is strictly increasing along the curves $t \mapsto \Phi(t,x)$ and $|Du|$ is locally constant on level sets of $u$, for every $x\in\Omega$ there exist a neighbourhood $U_x \subseteq \Omega$ and a smooth real function $\beta_x$ such that
%$$
%|Du| = \beta_x(u) \qquad \text{on } \, U_x.
%$$
%Since $Du / |Du|$ is parallel, on $U_x$ we have
%\begin{align*}
%	H & = \diver\left( \frac{Du}{\sqrt{1+|Du|^2}} \right) = \diver\left( \frac{|Du|}{\sqrt{1+|Du|^2}} \frac{Du}{|Du|} \right) = D_{Du/|Du|}\frac{|Du|}{\sqrt{1+|Du|^2}} \\
%	& = \beta_x(u) \left( \frac{\beta_x}{\sqrt{1+\beta_x^2}} \right)'(u) = \frac{\beta_x(u)\beta_x'(u)}{(1+\beta_x(u)^2)^{3/2}}
%\end{align*}
%that is, $\beta_x$ is a solution of the Cauchy problem
%$$
%\begin{cases}
%	y' = H \dfrac{(1+y^2)^{3/2}}{y}, \\
%	y(u(x)) = |Du(x)|.
%\end{cases}
%$$
%Without loss of generality, let us suppose that $\beta_x$ is the maximal solution of this problem. For $x\in N$, by uniqueness $\beta_x = \beta$, hence for every $x_1,x_2\in\Phi(\mathscr{D})$ belonging to the same curve $t \mapsto \Phi(t,x)$, $x\in N$, it must hold $\beta_{x_1} = \beta_{x_2}$. Therefore, $\beta_x = \beta$ for every $x\in\Phi(\mathscr{D})$, equivalently
%$$
%|Du| = \beta(u) \qquad \text{on } \, \Phi(\mathscr{D}).
%$$
%
%We claim that $\Phi(\mathscr{D}) = \Omega$. The map $\Phi$ is a diffeomorphism and $\mathscr{D}$ is open in $\R\times N$, so $\Phi(\mathscr{D})$ is open in $\Omega$. We check that $\Phi(\mathscr{D})$ is also closed in $\Omega$, thus deducing $\Phi(\mathscr{D}) = \Omega$ by connectedness of $\Omega$.
%
%%\tcr{The map $\Phi$ is a diffeomorphism, hence it is sufficient to show that $\mathscr{D}$ is open in $\R \times N$ to conclude that $\Phi(\mathscr{D})$ is open in $\Omega$. Suppose, by contradiction, that $(t_0,x_0) \in \mathscr{D}$ is not an interior point for $\mathscr{D}$. Then, there exists a sequence $\{x_n\}\subseteq N$ converging to $x_0$ and such that either}
%%$$
%%	\lim_{n\to+\infty} t_1(x_n) \geq t_0 \qquad \text{or} \qquad \lim_{n\to+\infty} t_2(x_n) \leq t_0.
%%$$
%%\tcr{We only deal with the first case, the second one being analogous. If the first case is verified, then $t_2(x_n) - 1/n > t_1(x_n) > -\infty$ for every $n\geq n_0$ for some sufficiently large $n_0$. We set}
%%$$
%%	\hat x_n = \lim_{t\to t_1(x_n)^+} \Phi(t,x_n) \in \partial\Omega \qquad \text{for every } \, n \geq n_0.
%%$$
%%\tcr{Since $(t_1(x_n)+1/n,x_n) \to (t_0,x_0)$ in $\mathscr D$ as $n\to\infty$, we have $\Phi(t_1(x_n)+1/n,x_n) \to \Phi(t_0,x_0)$ in $\Omega$ as $n\to\infty$. Since $d_\sigma(\Phi(t_1(x_n)+1/n,x_n),\hat x_n) \leq 1/n$, we also have $\hat x_n \to \Phi(t_0,x_0)$ in $M$. But $\hat x_n \in \partial\Omega$ for every $n\geq n_0$ and $\partial\Omega$ is closed, so we get $\Phi(t_0,x_0) \in \partial\Omega$, contradiction.}
%
%%We have already stated that $t_1$ is upper semi-continuous. We now claim that $t_1(N) \subseteq u(\partial\Omega)$ and that $t_1$ is continuous, hence constant as $N$ is connected and $u(\partial\Omega) = \{b_1,b_2,\dots,\}$ does not contain any intervals.
%
%First, we prove that $t_1$ and $t_2$ are constant on $N$. We show this for $t_1$, the proof for $t_2$ being analogous. Let us define an auxiliary function $\rho : \overline{u(\Phi(\mathscr D))} \to \R$ by setting
%$$
%\rho(s) = \int_b^s \frac{\di \sigma}{\beta(\sigma)} \qquad \text{for every } \, s \in \overline{u(\Phi(\mathscr D))}.
%$$
%Note that $\rho$ is continuous, strictly increasing and finite-valued since the above integral converges for every $s\in \overline{u(\Phi(\mathscr D))}$. Indeed, $\beta$ is continuous and strictly positive on $u(\Phi(\mathscr D))$,  and if $s\in\overline{u(\Phi(\mathscr D))}$ is such that $\beta(\sigma)\to 0$ as $\sigma\to s$, then by \eqref{eq_beta_1} necessarily $H\neq0$ and $\beta(\sigma) \sim \sqrt{2|H||\sigma-s|}$ as $\sigma\to s$, so $1/\beta$ is integrable in a neighbourhood of $s$. Also note that by integrating $\frac{\di}{\di t} u(\Phi) = |D u|(\Phi) = \beta(u(\Phi))$ we get
%\[
%t = \int_{b}^{u(\Phi(t,x))} \frac{\di \sigma}{\beta(\sigma)} = \rho(u(\Phi(t,x))) \qquad \text{for every } \, (t,x) \in \mathscr{D}.
%\]
%We show that $t_1$ is lower semi-continuous on $N$. Suppose, by contradiction, that for some $x\in N$ and for some sequence $\{x_n\}\subseteq N$ converging to $x$ we have
%$$
%\lim_{n\to \infty} t_1(x_n) < t_1(x).
%$$
%Fix $\bar t$ such that
%$$
%\lim_{n\to \infty} t_1(x_n) < \bar t < t_1(x), \qquad \rho^{-1}(\bar t) \not\in \{b_1, b_2,\dots \}.
%$$
%Then, $\{\Phi(\bar t,x_n)\} \subseteq \Omega$ converges to a point $\bar x$ of $\partial\Omega$. Along this sequence, $u$ has the constant value $\rho^{-1}(\bar t)$, so by continuity it must be $u(\bar x) = \rho^{-1}(\bar t)$. But $\rho^{-1}(\bar t) \not\in \{b_1,b_2,\dots\} = u(\partial\Omega)$ and we have reached a contradiction. Since we already showed that $t_1$ is upper semi-continuous, we conclude that $t_1$ is continuous on $N$. For every $x\in N$, we either have $t_1(x) = -\infty$ or $t_1(x) \in (-\infty,0)$. In the second case, the endpoint $x_\ast = \lim_{t\to t_1(x)+} \Phi(t,x)$ belongs to $\partial\Omega$ and by continuity $t_1(x) = \rho(u(x_\ast))$. So, $t_1(N) \subseteq \{\rho(b_1),\rho(b_2),\dots\} \cup \{-\infty\}$. Since this set contains no open intervals and $t_1$ is continuous on the connected set $N$, we conclude that $t_1$ is constant.
%
%Let $T_1\in[-\infty,0)$ and $T_2 \in (0,+\infty]$ be the constant values of $t_1$ and $t_2$ on $N$, so that
%$$
%\mathscr{D} = (T_1,T_2) \times N.
%$$
%For every $\bar t \in (T_1,T_2)$, the image $N_{\bar t} = \Phi(\{\bar t\}\times N) \subseteq \Omega$ is a connected open subset of the embedded submanifold $\{u = \rho^{-1}(\bar t)\} \subseteq \Omega$. The restriction $\Phi_{|\{\bar t\} \times N} : \{\bar t\} \times N \to N_{\bar t}$ is a local Riemannian isometry and $\{\bar t\} \times N$ is complete, so $\Phi_{|\{\bar t\} \times N}$ is a Riemannian covering map and therefore $N_{\bar t}$ is also complete with respect to its intrisinc geodesic distance, that we shall denote by $d_{\bar t}$ (see \cite[Lemma 5.6.4 and Proposition 5.6.3]{petersen}).
%
%% In particular, for every $p_1,p_2 \in \Phi(\{\bar t\}\times N)$ and for any $x_1 \in N$ such that $\Phi(\bar t,x_1) = p_1$, every geodesic curve $\gamma \subseteq \Phi(\{\bar t\}\times N)$ joining $p_1$ and $p_2$ can be lifted to a geodesic arc $\gamma_0 \subseteq N$ joining $x_1$ with some point $x_2\in N$ such that $\Phi(\bar t,x_2) = p_2$.}
%
%We prove that $\Phi(\mathscr{D})$ is closed in $\Omega$. % by showing that the limit of every sequence of points of $\Phi(\mathscr{D})$ converging in $\Omega$ also belongs to $\Phi(\mathscr{D})$.
%Let $\{p_n\}\subseteq \Phi(\mathscr{D})$ be a given sequence converging to some point $\bar p \in \Omega$. We have to show that $\bar p \in \Phi(\mathscr{D})$. Set $\bar t = \rho(u(\bar p))$, and observe that $u(\bar p) \in u(\overline{\Phi(\mathscr D)}) \subseteq \overline{u(\Phi(\mathscr D))}$, hence $\bar t$ is finite. For every $n$ we can find $(t_n,x_n) \in \mathscr{D}$ such that $p_n = \Phi(t_n,x_n)$. By continuity, $t_n = \rho(u(p_n)) \to \rho(u(\bar p)) = \bar t$, hence $T_1 \leq \bar t \leq T_2$. Both inequalities are strict, otherwise either $\{(x_n)_\ast\} = \{\Phi(T_1^+,x_n)\}$ or $\{(x_n)^\ast\} = \{\Phi(T_2^-,x_n))\}$ would be a sequence of points of $\partial\Omega$ converging to $\bar p\in\Omega$, absurd. Setting $q_n = \Phi(\bar t,x_n)$ for every $n$, we have that $\{q_n\}$ is a sequence of points of $N_{\bar t}$ converging to $\bar p$ in $M$, since $d_\sigma(p_n,q_n) \leq |\bar t - t_n| \to 0$. Hence, $\{q_n\}$ is a Cauchy sequence in $M$. By completeness of $M$, any two points $q_n$, $q_{n'}$ are joined by a minimizing geodesic arc in $M$. Since $(N_{\bar t},d_{\bar t})$ is complete and totally geodesic, every geodesic in $M$ joining two points of $N_{\bar t}$ must lie in $N_{\bar t}$. So, $\{q_n\}$ is a Cauchy sequence in $N_{\bar t}$ and therefore converges to some point $\bar q \in N_{\bar t}$. Since $N_{\bar t}$ is embedded in $M$, $\bar q = \bar p$ and we conclude $\bar p \in N_{\bar t} \subseteq \Phi(\mathscr{D})$, as desired. This shows that $\Phi(\mathscr{D})$ is closed in $\Omega$.
%
%%\tcr{Let $\{(t_n,x_n)\} \subseteq \mathscr{D}$ be given such that $\{\Phi(t_n,x_n)\}$ converges to some $\bar x \in \Omega$. Let $U\subseteq M$ be a geodesically convex neighbourhood of $\bar x$. Since $\Omega$ is open and $\bar x$ has a fundamental system of geodesically convex neighbourhoods in $M$, we can assume $U\subseteq\Omega$. Let $\bar t = \rho(u(\bar x))$. By continuity, $t_n = \rho(u(\Phi(t_n,x_n))) \to \bar t$, hence $T_1 \leq \bar t \leq T_2$. Both inequalities are strict, otherwise either $\{(x_n)_\ast\} \subseteq \partial\Omega$ or $\{(x_n)^\ast\}\subseteq\partial\Omega$ would converge to $\bar x\in\Omega$, absurd. Then,}
%%$$
%%	\{\bar t\} \times N \subseteq \mathscr{D}
%%$$
%%\tcr{Indeed, if it was either $\bar t = T_1$ or $\bar t = T_2$ then either $\{(x_n)_\ast\}$ or $\{(x_n)^\ast\}$ would be a sequence of points of $\partial\Omega$ converging to $\bar x\in\Omega$, absurd. }%, where}
%%$$
%%	\rho(s) = \int_b^s \frac{\di\sigma}{\beta(\sigma)}.
%%$$
%%\tcr{Then, we have $u(\Phi(\bar t,x)) = u(\bar x)$ for every $x\in N$ such that $(\bar t,x)\in\mathscr{D}$. Since $\rho(t_n) = u(\Phi(t_n,x_n)) \to u(\bar x)$ as $n\to\infty$, we have}
%%$$
%%	T_1 \leq \bar t \leq T_1
%%$$
%%\tcr{and both inequalities are strict, since otherwise we could extract from either $\{(x_n)_\ast\}$ or $\{(x_n)^\ast\}$ a sequence of points of $\partial\Omega$ converging to $\bar x$.} %, if $t_1(x_{n_k}) \to \bar t$ for some subsequence $\{x_{n_k}\} \subseteq \{x_n\}$ then we can show as above that the endpoints $\lim_{t\to t_1(x_{n_k})^+} \Phi(t,x_{n_k}) \in \partial\Omega$ converge to $\bar x$, contradicting the fact that $\bar x \in \Omega$.
%%\tcr{and we have $(\bar t,x_n) \in \mathscr{D}$ for every $n$. Since $d_\sigma(\Phi(\bar t,x_n),\Phi(t_n,x_n)) \leq |\bar t - t_n| \to 0$ as $n\to\infty$, we also have $\Phi(\bar t,x_n) \in U$ for every sufficiently large $n$. The set $U$ is geodesically convex and level sets of $u$ are totally geodesic, so for every sufficiently large $n,m$ the points $\Phi(\bar t,x_n)$ and $\Phi(\bar t,x_m)$ are joined by a minimizing geodesic arc $\gamma : [0,d_\sigma(\Phi(\bar t,x_n),\Phi(\bar t,x_m))] \to U \cap \{u=u(\bar x)\}$, which can be lifted to a geodesic arc $\gamma_0 : [0,d_\sigma(\Phi(\bar t,x_n),\Phi(\bar t,x_m))] \to N$ such that $\gamma(s) = \Phi(\bar t,\gamma_0(s))$. }
%%$$
%%	\{\bar t\}\times N \subseteq \mathscr{D}
%%$$
%%\tcr{Indeed, for some $\eps>0$ there exists a unique geodesic arc $\gamma_0 : [0,\eps] \to N$ such that $\gamma(s) = \Phi(\bar t,\gamma_0(s))$ for every $s\in[0,\eps]$, and by completeness of $N$ the maximal $\eps\leq d_\sigma(\Phi(\bar t,x_n),\Phi(\bar t,x_m))$ for which this is true must be $d_\sigma(\Phi(\bar t,x_n),\Phi(\bar t,x_m))$.
%%\tcr{Hence, since $\{\Phi(\bar t,x_n)\}$ is a Cauchy sequence in $M$, we have that $\{x_n\}\subseteq N$ contains a Cauchy subsequence. By completeness of $N$, there exists $x_0 \in N$ such that, up to a subsequence, $x_n \to x_0$ and then $\Phi(t_n,x_n) \to \Phi(\bar t,x_0)$. By uniqueness of limits, $\bar x = \Phi(\bar t,x_0) \in \Phi(\mathscr{D})$. This shows that $\Phi(\mathscr{D})$ is closed in $\Omega$.}
%%\tcr{Let $(t_0,x_0) \in \mathscr{D}$ be given. We have to show that there exists a neighbourhood $U \subseteq \Omega$ of $\Phi(t_0,x_0)$ such that $U\subseteq \Phi(\mathscr U)$ for some $\mathscr{U} \subseteq \mathscr{D}$. and let $u_0 = u(\Phi(t_0,x_0))$. Fix $\eps>0$ such that $t_1(x_0) < t_0 - \eps$, $t_0 + \eps < t_2(x_0)$ and set $u_\ast = u(\Phi(t_0-\eps,x_0))$, $u^\ast = u(\Phi(t_0+\eps,x_0))$. Since $t\mapsto u(\Phi(t,x_0))$ is strictly increasing, we have}
%%$$
%%	\lim_{t\to t_1(x_0)^+} u(t) < u_\ast < u_0 < u^\ast < \lim_{t\to t_2(x_0)^-}
%%$$
%%\tcr{and therefore}
%%\tcr{The function $\beta$ is well defined for every $\sigma \in [u(\Phi(t_0-\eps,x_0)),u(\Phi(t_0+\eps,x_0))]$ In particular, $t_1(x_0),t_2(x_0)$}
%
%As already stated, since $\Phi(\mathscr{D})$ is non-empty and both open and closed in the connected set $\Omega$, we have $\Phi(\mathscr{D}) = \Omega$. Thus, $\Phi$ realizes an isometry between $\Omega$ and the product manifold
%$$
%(T_1,T_2)\times N,
%$$
%and $N$ is parabolic because of Lemma \ref{lem_paraprod}. Furthermore, $u$ only depends on the variable $t$ because 
%$$
%u(\Phi(t,x)) = \rho^{-1}(t) \qquad \text{for every } \, (t,x) \in (T_1,T_2)\times N.
%$$
%In the chart $\Phi$, $u= u(t)$ is therefore a solution of 
%\[
%H = \diver \left( \frac{Du}{\sqrt{1+ |Du|^2}} \right) = \diver \left( \frac{u'\partial_t}{\sqrt{1+ (u')^2}} \right) = \left( \frac{u'}{\sqrt{1+ (u')^2}} \right)'.
%\]
%Integration of the ODE shows that the possibility $(-\infty, T_2) \times N$ is incompatible with the fact that $u$ is increasing and bounded from below, while the other models lead to the solutions $u_{c_1,H}$ in \eqref{eq_split}, up to reparametrizing $t \mapsto t - T_1$ and setting $T = T_2 - T_1$. To check $(ii)$, the second in \eqref{eq_persplitting} implies 
%\[
%\bar v = c W\|\nabla u\| \qquad \text{on } \, \Omega,
%\]	
%for some constant $c>0$. Using \eqref{nablau2}, the identity rewrites as $(X, Du) = c|Du|$, that is, $(X, \partial_t) = c$.
%\end{proof}

\section{Euclidean space and the critical set of $u$}\label{sec_critical}

Let $u \in C^3(\Omega)$ be a non-constant solution to $\Delta_a u + f(u) = 0$. In this section, we examine more closely the case of Euclidean space and provide some results that allow to ``glue" different solutions along boundaries on which $\nabla u$ vanishes. Towards this aim, we assume as usual $f \in C^1(\R)$ and \eqref{assu_A_superstrong}, that is, 
\[
a\in C^{1,1}(\R_0^+), \qquad a(t) > 0 \qquad \text{and} \qquad ta'(t) + a(t) > 0 \qquad \text{for all } \, t \in \R^+_0
\]
to guarantee a certain amount of regularity for the set of critical points $\critu$. In fact, each partial derivative $w_j = \partial_j u$ belongs to $C^2(\Omega)$ and satisfies the linearized equation
\begin{equation}\label{eq_partialder}
\diver(A(\nabla u)\nabla w_j) + f'(u) w_j = 0 \qquad \text{in } \, \Omega \, ,
\end{equation}
where in our assumptions $A(\nabla u)$ has $\lip_\loc$ coefficients and $f'(u) \in C(\Omega)$. Therefore, the unique continuation principle holds and so either $w_j \equiv 0$ or $w_j$ has finite order of vanishing at every point 
\[
x \in Z_j = \{ x \in \Omega : w_j(x) = 0 \} \supseteq \critu.
\]
In particular, if $u$ is constant in an open subset $\Omega_0$ of $\Omega$, then $u$ is constant in the entire $\Omega$. 



\begin{lemma} \label{Om_finper}
	Let $\Omega\subseteq\R^n$ be a $C^1$ domain, Let $a$ satisfy \eqref{assu_A_superstrong} and let $f \in C^1(\R)$. Consider a non-constant solution $u \in C^3(\Omega)$ to
	\[
	\Delta_a u + f(u) = 0 \qquad \text{in } \, \Omega \, .
	\]
	%	\[
	%	\left\{
	%	\begin{array}{l@{\qquad}l}
		%		\Delta_a u + f(u) = 0 & \text{in } \Omega \\[0.3cm]
		%		u, \partial_\eta u \text{ locally constant} & \text{on } \partial\Omega
		%	\end{array}
	%	\right.
	%	\]
	If $\Omega_0\subseteq\Omega$ is an open set such that $\partial\Omega_0 \subseteq \partial\Omega \cup \critu$
%	$\partial\Omega_0 = \Gamma_0 \cup \Gamma_1$ for some closed disjoint sets $\Gamma_0$, $\Gamma_1$ of $\R^n$ satisfying
%	\[
%		\Gamma_0 \subseteq \critu, \qquad \Gamma_1 \subseteq \partial\Omega
%	\]
	then $\mathscr{H}^{n-1}(K\cap\partial\Omega_0) < \infty$ for all compact sets $K\subseteq\Omega$. In particular, $\Omega_0$ is a set of locally finite perimeter in $\Omega$.
\end{lemma}

%\begin{remark}
%	We emphasize that in the assumptions of Lemma \ref{Om_finper} the sets $\Gamma_0$ and $\Gamma_1$ are closed and disjoint (since $\Gamma_0\subseteq\Omega$ and $\Gamma_1\subseteq\partial\Omega$), so in particular $\Gamma_0$ and $\partial\Omega$ are separated from each other. In other words, we are assuming that every connected component of $\partial\Omega_0$ is either a component of $\partial\Omega$ or is entirely contained in $\Omega$ (and thus separated from $\partial\Omega$), and in the latter case we are further assuming that $\nabla u = 0$ everywhere on that component.
%\end{remark}

\begin{proof}
%	Let $K\subseteq\R^n$ be a compact set. We have
%	\begin{align*}
%		\mathscr{H}^{n-1}(K\cap\partial\Omega_0) & \leq \mathscr{H}^{n-1}(K\cap\Gamma_0) + \mathscr{H}^{n-1}(K\cap\Gamma_1) \\
%		& \leq \mathscr{H}^{n-1}(K\cap\Gamma_0) + \mathscr{H}^{n-1}(K\cap\partial\Omega) \, .
%	\end{align*}
%	Since $\Omega$ has $C^1$ boundary, $\mathscr H^{n-1}(K\cap\partial\Omega) < \infty$. We shall prove that also $\mathscr H^{n-1}(K\cap\Gamma_0) < \infty$.
%	For $i=1,\dots,n$ let us set $w_i = \partial_i u$ and
%	\[
%	Z_i = \{ x \in \Omega : w_i(x) = 0 \} \, .
%	\]
%	Each function $w_i$ belongs to $C^2(\Omega)$ and satisfies the linear elliptic equation
%	\[
%	\diver(A(\nabla u)\nabla w_i) + f'(u) w_i = 0 \qquad \text{in } \, \Omega \, .
%	\]
	Let $K\subseteq\Omega$ be a compact set. Since $u$ is non-constant, there exists some $j\in\{1,\dots,n\}$ such that $w_j = \partial_j u \not\equiv 0$ in $\Omega$. As $\critu \subseteq Z_j = \{ w_j = 0 \}$, it suffices to show that $\mathscr H^{n-1}(K\cap Z_j) < \infty$.
	%The assumptions on $a$, $f$ and $u$ ensure that the coefficient matrix function $\Omega \ni x \mapsto A(\nabla u(x))$ is locally uniformly elliptic and Lipschitz continuous, and that $\Omega \ni x \mapsto f'(u(x))$ is continuous and locally bounded. Therefore, the unique continuation principle holds and so $w_j$ has finite order of vanishing at every point of $Z_j$.
	As observed at the beginning of this section, $w_j$ satisfies the linear elliptic equation \eqref{eq_partialder} and has finite order of vanishing at every point of $Z_j$. By \cite[Theorem 1.7]{HS89}, for each point $x\in Z_j$ there exist $\rho_x>0$ and $C=C(n,x)>0$ such that $B_{\rho_x}(x) \subseteq \Omega$ and
	\[
	\mathscr{H}^{n-1}(B_\rho(x) \cap Z_j) \leq C \rho^{n-1} \qquad \text{for all } \, 0 < \rho \leq \rho_x \, .
	\]
	By covering the compact set $K\cap Z_j$ with a finite number of balls $B_{\rho_{x_1}}(x_1),\dots,B_{\rho_{x_k}}(x_k)$ centered at points $x_1,\dots,x_k$ of $K\cap Z_j$ we deduce $\mathscr{H}^{n-1}(K \cap Z_j) < \infty$, as desired.
	
	By \cite[Theorem 4.5.11]{federer} (also \cite[Proposition 3.62]{ambrosiofuscopallara} together with the introductory remark), since $\mathscr H^{n-1}(K\cap\partial\Omega_0) < \infty$ for all compact $K\subseteq\Omega$ it follows that $\Omega_0$ has finite perimeter in every compact set $K\subseteq\Omega$, as claimed.
\end{proof}

%\begin{lemma} \label{lem_Hn-1}
%	Let $\Omega\subseteq\R^n$ be an open set, let $a\in C^{1,1}(\R_0^+)$ satisfy
%	\[
%	a(t) > 0 \qquad \text{and} \qquad ta'(t) + a(t) > 0 \qquad \text{for all } \, t \in \R^+_0
%	\]
%	and let $f \in C^1(\R)$. If $u \in C^3(\Omega)$ is a non-constant solution to
%	\[
%	\Delta_a u + f(u) = 0 \qquad \text{in } \Omega
%	\]
%	then
%	\[
%	\mathscr{H}^{n-1}(\{ x \in \Omega : \nabla u(x) = 0 \} \cap K) < \infty
%	\]
%	for all compact sets $K\subseteq\Omega$.
%	%	\[
%	%		Z_j = \{ x \in \Omega : \partial_j u(x) = 0 \} \qquad \text{for } \, j = 1,\dots,n \, .
%	%	\]
%	%	For each $j=1,\dots,n$ such that $Z_j \neq \Omega$, the set
%	%	\[
%	%		Z^0_j = \{ x \in Z_j : \partial^2_{ij} u(x) = 0 \text{ for } 1 \leq i \leq n \}
%	%	\]
%	%	is of Hausdorff dimension $\leq n-2$, the set $Z_j \backslash Z^0_j$ is a countable union of embedded $C^1$ hypersurfaces and $\mathscr H^{n-1}((Z_j \backslash Z^0_j)\cap K) < \infty$ for any compact set $K\subseteq\Omega$. In particular, if $u$ is non-constant then
%	%	\[
%	%		\{ x \in \Omega : \nabla u(x) = 0 \text{ and } \Hess u(x) = 0 \}
%	%	\]
%	%	is of Hausdorff dimension $\leq n-2$ and the set
%	%	\[
%	%		\{ x \in \Omega : \nabla u(x) = 0 \text{ and } \Hess u(x) \neq 0 \}
%	%	\]
%	%	is contained in a countable union of embedded $C^1$ hypersurfaces and has locally finite $(n-1)$-dimensional Hausdorff measure.
%\end{lemma}
%
%\begin{proof}
%	Let us set
%	\[
%	Z_i = \{ x \in \Omega : \partial_i u(x) = 0 \} \qquad \text{for } \, i = 1,\dots,n
%	\]
%	and
%	\[
%	Z = \bigcap_{i=1}^n Z_i \equiv \{ x \in \Omega : \nabla u(x) = 0 \} \, .
%	\]
%	The claim of the lemma is that $\mathscr H^{n-1}(Z\cap K) < \infty$ for all compact sets $K\subseteq\Omega$. %so it is sufficient to show that $\mathcal H^{n-1}(Z_i\cap K) < \infty$ for all compact $K\subseteq\Omega$ and $1\leq i\leq n$.
%	For $i=1,\dots,n$ the function $w_i = \partial_i u \in C^2(\Omega)$ satisfies the linear elliptic equation
%	\[
%	\diver(A(\nabla u)\nabla w_i) + f'(u) w_i = 0 \qquad \text{in } \, \Omega \, .
%	\]
%	Since $u$ is non-constant, there exists some $j\in\{1,\dots,n\}$ such that $w_j\not\equiv 0$ in $\Omega$. As $Z \subseteq Z_j$, it suffices to show that $\mathscr H^{n-1}(Z_j\cap K) < \infty$ for all compact $K\subseteq\Omega$. The assumptions on $a$, $f$ and $u$ ensure that the coefficient matrix function $\Omega \ni x \mapsto A(\nabla u(x))$ is locally uniformly elliptic and Lipschitz continuous, and that $\Omega \ni x \mapsto f'(u(x))$ is continuous and locally bounded. Therefore, the unique continuation principle holds and so $w_j$ has finite order of vanishing at every point of $Z_j$. By Theorem (1.7) from \cite{HS89}, for each point $x\in Z_j$ there exist $\rho_0(x)>0$ and $C=C(n,x)>0$ such that $B_{\rho_0(x)}(x) \subseteq \Omega$ and
%	\[
%	\mathscr{H}^{n-1}(B_\rho(x) \cap Z_j) \leq C \rho^{n-1} \qquad \text{for all } \, 0 < \rho \leq \rho_0(x) \, .
%	\]
%	%	and
%	%	\[
%	%		\dim(B_\rho(x_0) \cap Z^0_j) \leq n-2 \, .
%	%	\]
%	%	Therefore, $Z^0_j$ is of Hausdorff dimension less or equal than $n-2$ and the set $Z_j\backslash Z^0_j$, which is a countable union of embedded hypersurfaces, has intersection of finite $(n-1)$-dimensional measure with any compact subset of $\Omega$.
%	By covering the compact set $Z_j \cap K$ with a finite number of balls $B_{\rho_0(x_1)}(x_1),\dots,B_{\rho_0(x_k)}(x_k)$ centered at points $x_1,\dots,x_k$ of $Z_j\cap K$ we deduce $\mathscr{H}^{n-1}(Z_j\cap K) < \infty$, as desired.
%\end{proof}



%We have the following ``localization'' on subdomains $\Omega_0$ as in Lemma \ref{Om_finper} of the geometric Poincaré formula of Theorem \ref{teo_Poincare}.
%
%\begin{proposition}\label{prop_Poincare_loc}
%	Let $\Omega\subset\R^n$ be a domain with $C^3$ boundary, let $a$ satisfy \eqref{assu_A_superstrong} and $f \in C^1(\R)$. Let $u \in C^3(\Omega) \cap C^{2,\alpha}_\loc(\overline\Omega)$ be a stable solution to
%	\[
%	\left\{
%	\begin{array}{l@{\qquad}l}
%		\Delta_a u + f(u) = 0 & \text{in } \Omega \\[0.3cm]
%		u, \partial_\eta u \text{ locally constant} & \text{on } \partial\Omega .
%	\end{array}
%	\right.
%	\]
%	Let $\Omega_0\subseteq\Omega$ be an open set such that $\partial\Omega_0 = \Gamma_0 \cup \Gamma_1$ for some closed sets $\Gamma_0$, $\Gamma_1$ satisfying
%	\[
%		\Gamma_0 \subseteq \{ x \in \Omega : \nabla u(x) = 0 \} \, , \qquad \Gamma_1 \subseteq \partial\Omega \, .
%	\]
%	Let $X$ be a Killing vector field in $\overline{\Omega}_0$ and suppose that $w = \langle \nabla u, X \rangle$ satisfies $w \ge 0$, $w \not\equiv 0$ in $\Omega_0$. Then, for each $\varphi\in\lip_c(\R^n)$ it holds
%	\begin{align*}\label{piufinalissima_teo_3}
%		\int_{\Omega_0} |\nabla u|^2 \langle A(\nabla u) \nabla \varphi, \nabla \varphi \rangle & \geq \int_{\Omega_0} \varphi^2 \left\{\lambda_1 \left|\nabla^\top u_\nu\right|^2 +\lambda_2\left[u_\nu^2|\II|^2 +\Ric(\nabla u, \nabla u)\right]\right\} \\
%		& \phantom{=\;} + \int_{\Omega_0} w^2 \langle A(\nabla u) \nabla \left(\frac{\varphi|\nabla u|}{w}\right), \nabla \left(\frac{\varphi|\nabla u|}{w}\right) \rangle \, .
%	\end{align*}
%\end{proposition}
%
%\begin{proof}
%	Let $\varphi\in \lip_c(\R^n)$ and $\eps>0$ be given. %We apply the Gauss-Green theorem on $\Omega_0$ to the Lipschitz vector field
%%	\[
%%		V = \frac{\varphi^2|\nabla u|^2}{w+\eps} A(\nabla u)\nabla w - \frac{1}{2} \frac{\varphi^2 w}{w+\eps} A(\nabla u)\nabla|\nabla u|^2
%%	\]
%	For each sufficiently small $\tau,\delta>0$ let $\psi\in\lip_c(\Omega)$ be the cut-off function as in \tcr{the proof of Proposition 27}. By Lemma \ref{Om_finper} the set $\Omega_0$ has locally finite perimeter, so we apply the Gauss-Green theorem to the Lipschitz vector field
%	\[
%		V = \frac{\psi\varphi^2|\nabla u|^2}{w+\eps} A(\nabla u)\nabla w
%	\]
%	to obtain
%	\[
%		\int_{\Omega_0} \diver\, V = - \int_{\partial\Omega_0} \langle V,\eta\rangle = 0
%	\]
%	where the second equality holds because $V$ vanishes everywhere on $\partial\Omega_0$, as $|\nabla u|=0$ on $\Gamma_0$ and $\psi=0$ on $\Gamma_1$. Since $w$ satisfies
%	\[
%		\diver(A(\nabla u)\nabla w) + f'(u)w = 0 \qquad \text{weakly in } \, \Omega
%	\]
%	we have
%	\[
%		\diver\,V = - \frac{\psi\varphi^2|\nabla u|^2}{w+\eps} f'(u)w + \langle A(\nabla u)\nabla w,\nabla\left(\frac{\psi\varphi^2|\nabla u|^2}{w+\eps}\right)\rangle
%	\]
%	and so
%	\[
%		\int_{\Omega_0} \psi \varphi^2 \frac{w}{w+\eps} f'(u)|\nabla u|^2 = \int_{\Omega_0} \langle A(\nabla u)\nabla w,\nabla\left(\frac{\psi\varphi^2|\nabla u|^2}{w+\eps}\right)\rangle \, .
%	\]
%	Letting $\delta\to0$ and then $\tau\to0$ as in the proof of \tcr{Proposition 27} we obtain
%	\[
%		\int_{\Omega_0} \frac{\varphi^2 w}{w+\eps} f'(u)|\nabla u|^2 = \int_{\Omega_0} \langle A(\nabla u)\nabla w,\nabla\left(\frac{\varphi^2|\nabla u|^2}{w+\eps}\right)\rangle + \int_{\Gamma_1} \frac{\varphi^2|\nabla u|^2}{w+\eps} \langle A(\nabla u)\nabla w,\eta\rangle \, .
%	\]
%%	that we rewrite as
%%	\begin{align*}
%%		\int_{\Omega_0} \langle A(\nabla u)\nabla(\varphi|\nabla u|),\nabla(\varphi|\nabla u|)\rangle & = \int_{\Omega_0} \varphi^2 \frac{w}{w+\eps} f'(u)|\nabla u|^2 - \int_{\Gamma_0} \frac{\varphi^2|\nabla u|^2}{w+\eps} \langle A(\nabla u)\nabla w,\eta\rangle \\
%%		& \phantom{=\;} + \int_{\Omega_0} (w+\eps)^2 \langle A(\nabla u)\nabla\left(\frac{\varphi|\nabla u|}{w+\eps}\right),\nabla\left(\frac{\varphi|\nabla u|}{w+\eps}\right)\rangle \, .
%%	\end{align*}
%	We also apply the Gauss-Green theorem to the vector field
%	\[
%		W = \frac{1}{2} \frac{\varphi^2 w}{w+\eps} A(\nabla u)\nabla|\nabla u|^2
%	\]
%	to obtain
%	\[
%		\int_{\Omega_0} \diver\, W = -\int_{\partial\Omega_0} \langle W,\eta\rangle = - \int_{\Gamma_1} \langle W,\eta\rangle
%	\]
%	where the second equality holds because $\nabla u=0$ on $\Gamma_0$. By the Bochner formula \eqref{bochner_betterfornablau} and using that $\Delta_a u = -f(u)$ we have
%	\begin{align*}
%		\diver\, W & = \frac{\varphi^2 w}{w+\eps} |\nabla u|\diver(A(\nabla u)\nabla|\nabla u|) + \langle A(\nabla u)\nabla|\nabla u|, \nabla\left(\frac{\varphi^2|\nabla u|w}{w+\eps}\right)\rangle \\
%		& = \frac{\varphi^2 w}{w+\eps} \left\{ \lambda_1|\nabla^\top|\nabla u||^2 + \lambda_2 [|\nabla u|^2|\II|^2 + \Ric(\nabla u,\nabla u)] \right\} - \varphi^2 \frac{w}{w+\eps} f'(u)|\nabla u|^2 \\
%		& \phantom{=\;} + \langle A(\nabla u)\nabla|\nabla u|, \nabla\left(\frac{\varphi^2|\nabla u|w}{w+\eps}\right)\rangle
%	\end{align*}
%	in $\Omega_0$, so we obtain
%	\begin{align*}
%		\int_{\Omega_0} \frac{\varphi^2 w}{w+\eps} f'(u)|\nabla u|^2 & = \int_{\Omega_0} \frac{\varphi^2 w}{w+\eps} \left\{ \lambda_1|\nabla^\top|\nabla u||^2 + \lambda_2 [|\nabla u|^2|\II|^2 + \Ric(\nabla u,\nabla u)] \right\} \\
%		& \phantom{=\;} + \int_{\Omega_0} \langle A(\nabla u)\nabla|\nabla u|, \nabla\left(\frac{\varphi^2|\nabla u|w}{w+\eps}\right)\rangle \\
%		& \phantom{=\;} + \frac{1}{2} \int_{\Gamma_1} \varphi^2 \frac{w}{w+\eps} \langle A(\nabla u)\nabla|\nabla u|^2,\eta\rangle
%	\end{align*}
%	and putting everything together we obtain
%	\begin{align*}
%		0 & = \int_{\Omega_0} \frac{\varphi^2 w}{w+\eps} \left\{ \lambda_1|\nabla^\top|\nabla u||^2 + \lambda_2 [|\nabla u|^2|\II|^2 + \Ric(\nabla u,\nabla u)] \right\} \\
%		& \phantom{=\;} + \int_{\Omega_0} \langle A(\nabla u)\nabla|\nabla u|, \nabla\left(\frac{\varphi^2|\nabla u|w}{w+\eps}\right)\rangle - \langle A(\nabla u)\nabla w,\nabla\left(\frac{\varphi^2|\nabla u|^2}{w+\eps}\right)\rangle \\
%		& \phantom{=\;} + \int_{\Gamma_1} \frac{\varphi^2}{w+\eps} \langle w A(\nabla u)\nabla\frac{|\nabla u|^2}{2} - |\nabla u|^2 A(\nabla u)\nabla w,\eta\rangle \, .
%	\end{align*}
%	The integral over $\Gamma_1$ vanishes by Lemma \ref{lem_bordo}. Straightforward computations yield
%	\begin{align*}
%		\langle A(\nabla u)\nabla|\nabla u|, \nabla\left(\frac{\varphi^2|\nabla u|w}{w+\eps}\right)\rangle & = \varphi^2 |\nabla u| \langle A(\nabla u)\nabla|\nabla u|,\nabla\frac{w}{w+\eps}\rangle \\
%		& \phantom{=\;} + \frac{w}{w+\eps} \langle A(\nabla u)\nabla|\nabla u|,\nabla(\varphi^2|\nabla u|)\rangle \\
%		& = - \frac{\eps}{2} \frac{\varphi^2}{(w+\eps)^2} \langle A(\nabla u)\nabla|\nabla u|^2,\nabla w\rangle \\
%		& \phantom{=\;} + \frac{w}{w+\eps} \langle A(\nabla u)\nabla(\varphi|\nabla u|),\nabla(\varphi|\nabla u|)\rangle \\
%		& \phantom{=\;} - \frac{w}{w+\eps} |\nabla u|^2 \langle A(\nabla u)\nabla\varphi,\nabla\varphi\rangle
%	\end{align*}
%	and
%	\begin{align*}
%		\langle A(\nabla u)\nabla w,\nabla\left(\frac{\varphi^2|\nabla u|^2}{w+\eps}\right)\rangle & = \varphi|\nabla u|\langle A(\nabla u)\nabla w,\nabla\left(\frac{\varphi|\nabla u|}{w+\eps}\right)\rangle \\
%		& \phantom{=\;} + \frac{\varphi|\nabla u|}{w+\eps} \langle A(\nabla u)\nabla w,\nabla(\varphi|\nabla u|)\rangle \\
%		& = -(w+\eps)^2 \langle A(\nabla u)\nabla\left(\frac{\varphi|\nabla u|}{w+\eps}\right),\nabla\left(\frac{\varphi|\nabla u|}{w+\eps}\right)\rangle \\
%		& \phantom{=\;} + \langle A(\nabla u)\nabla(\varphi|\nabla u|),\nabla(\varphi|\nabla u|)\rangle
%	\end{align*}
%	so we arrive at
%	\begin{align*}
%		\int_{\Omega_0} \frac{w|\nabla u|^2}{w+\eps} \langle A(\nabla u)\nabla\varphi,\nabla\varphi\rangle & = \int_{\Omega_0} \frac{\varphi^2 w}{w+\eps} \left\{ \lambda_1|\nabla^\top|\nabla u||^2 + \lambda_2 [|\nabla u|^2|\II|^2 + \Ric(\nabla u,\nabla u)] \right\} \\
%		& \phantom{=\;} + \int_{\Omega_0} (w+\eps)^2 \langle A(\nabla u)\nabla\left(\frac{\varphi|\nabla u|}{w+\eps}\right),\nabla\left(\frac{\varphi|\nabla u|}{w+\eps}\right)\rangle \\
%		& \phantom{=\;} - \int_{\Omega_0} \frac{\eps}{2} \frac{\varphi^2}{(w+\eps)^2} \langle A(\nabla u)\nabla|\nabla u|^2,\nabla w\rangle \\
%		& \phantom{=\;} - \int_{\Omega_0} \frac{\eps}{w+\eps} \langle A(\nabla u)\nabla(\varphi|\nabla u|),\nabla(\varphi|\nabla u|)\rangle
%	\end{align*}
%	and then letting $\eps\to0$ and repeating the arguments developed in Theorem  \ref{teo_Poincare} we obtain the desired conclusion.
%\end{proof}

%\begin{proposition}
%	Let $\Omega\subseteq\R^n$ be an open set, $f : \Omega \to \R$ a non-constant real analytic function such that the set
%	\[
%		\Gamma = \{x \in \Omega : \nabla f(x) = 0\}
%	\]
%	satisfies $\overline{\Gamma} \cap \partial\Omega = \emptyset$. Let $\Sigma$ be the $(n-1)$-dimensional stratum of $\Gamma$. If $\Hess f \neq 0$ everywhere on $\overline{\Sigma}$, then $\Sigma = \overline{\Sigma} = \Gamma$.
%\end{proposition}
%
%\begin{proof}
%	$\Sigma$ is a disjoint union of embedded real analytic submanifolds of dimension $n-1$ in $\Omega$. Since $f$ is constant on each of these hypersurfaces, the restriction to $\Sigma$ of the Hessian of $f$ is given by
%	\[
%		\Hess f = (\Delta f) \, \nu_\flat \otimes \nu_\flat
%	\]
%	where $\nu$ is the normal vector field to $\Sigma$. In particular, we have $(\Delta f)^2 = |\Hess f|^2$ on $\Sigma$, hence also on $\overline{\Sigma}$ by continuity. Since $\Hess f \neq 0$ on $\overline{\Sigma}$, we also have $\Delta f \neq 0$ on $\overline{\Sigma}$. Then, we infer that $\Delta f$ is a simple eigenvalue for $\Hess f$ on the whole $\overline{\Sigma}$ and, by continuity, also on a neighbourhood $U\subseteq\Omega$ of $\overline{\Sigma}$. Let $v : U \to \R^n$ be a continuous vector field of unit eigenvectors of $\Hess f$ corresponding to the eigenvalue $\Delta f$. By analyticity of $\Hess f$, $\Delta f$ and simplicity of the eigenvalue, the vector field $v$ is analytic in $U$. Also, on each connected component of $\Sigma$ we have either $\nu \equiv v$ or $\nu \equiv -v$.
%\end{proof}

We are ready to prove our main results of this section.

\begin{lemma} \label{lem_Du0}
	Let $\Omega\subseteq\R^n$ be a $C^1$ domain, let $a$ satisfy \eqref{assu_A_superstrong} and let $f \in C^1(\R)$. Assume that $u \in C^3(\Omega)\cap C^1(\overline{\Omega})$ is a non-constant, stable solution to
	\[
		\Delta_a u + f(u) = 0 \qquad \text{in } \Omega.
	\]
	%with moderate energy growth.
	%
	%
%	Assume that there exists a sequence $\{\phi_j\}$
%	
%	satisfying
%	\begin{equation} \label{Du_growth}
%		\int_{\Omega\cap B_R} |\nabla u|^2 \lambda_{\max}(|\nabla u|) = O(R^2\log R) \qquad \text{as } \, R \to \infty \, .
%	\end{equation}
	Let $\Omega_0\subseteq\Omega$ be a connected open subset such that $\nabla u = 0$ everywhere on $\partial\Omega_0$. If $u$ has moderate energy growth in $\Omega_0$, then $\Omega_0 \equiv \Omega$ and there exists $X\in\mathbb{S}^{n-1}$ such that $\langle\nabla u,X\rangle > 0$ in $\Omega$. In particular, $\critu = \emptyset$.
\end{lemma}

\begin{proof}
	By unique continuation, see the beginning of this section, since $u$ is non-constant in $\Omega$ then it must be non-constant in $\Omega_0$. Each partial derivative $w_j = \partial_j u$ satisfies \eqref{eq_partialder} and $w_j \equiv 0$ on $\partial\Omega_0$. Since $w_j^2 \leq |\nabla u|^2$ and $u$ has moderate energy growth in $\Omega_0$, by Lemma \ref{Lemma_w-} we deduce that $w_j$ either vanishes identically in $\Omega_0$ or it has a constant strict sign in $\Omega_0$. Since $u$ is non-constant in $\Omega_0$, for some index $j\in\{1,\dots,n\}$ we have $w_j \not \equiv 0$ and thus either $w_j>0$ or $w_j<0$ everywhere in $\Omega_0$. Then, taking either $X = e_j$ or $X = -e_j$ we have $\langle\nabla u,X\rangle > 0$ everywhere in $\Omega_0$. We now show that $\Omega_0 = \Omega$. Suppose, by contradiction, that $\Omega_0\subsetneq\Omega$. Then there exists $x\in\Omega\cap\partial\Omega_0\neq\emptyset$. The stability of $u$ guarantees the existence of $0<v\in C^{1,\alpha}_\loc(\Omega)$ satisfying
	\[
		\diver(A(\nabla u)\nabla v) + f'(u)v = 0 \qquad \text{in } \, \Omega \, .
	\]
	By Remark \ref{rem-Lemma_w-} we have $w_j = \lambda v$ in $\Omega_0$ for some $\lambda\in\R\backslash\{0\}$. By continuity, $w_j(x) = \lambda v(x) \neq 0$, contradiction.
\end{proof}

\begin{theorem}\label{teo_twopieces}
%	Let $\Omega\subseteq\R^n$ be a $C^2$ domain whose boundary consists of exactly two connected components $\partial_1\Omega$, $\partial_2\Omega$ such that, for $i=1,2$, there exists a unit vector $X_i\in\mathbb{S}^{n-1}$ with
%	\[
%		\langle\eta,X_i\rangle \geq 0 \qquad \text{on } \, \partial_i\Omega \, .
%	\]
	Let $\Omega\subseteq\R^n$ be a $C^2$ domain with disconnected boundary. Let $a$ satisfy \eqref{assu_A_superstrong}, let $f \in C^1(\R)$ and let $u \in C^3(\Omega) \cap C^2(\overline\Omega)$ be a non-constant, stable solution of moderate energy growth to
	\[
	\left\{
	\begin{array}{l@{\qquad}l}
		\Delta_a u + f(u) = 0 & \text{in } \Omega \\[0.3cm]
		u, \partial_\eta u \text{ locally constant} & \text{on } \partial\Omega.
%		u = b_1, \; \partial_\eta u = c_1 & \text{on } \partial_1\Omega \\[0.3cm]
%		u = b_2, \; \partial_\eta u = c_2 & \text{on } \partial_2 \Omega
	\end{array}
	\right.
	\]
	%for some $b_1,b_2,c_1,c_2\in\R$.
	Suppose that, for $j=1,2$, there exists $X_j \in \mathbb{S}^{n-1}$ such that
	\[
	\langle \eta, X_j \rangle \ge 0, \quad \langle \eta,X_j \rangle \not\equiv 0 \qquad \text{on } \, \partial_j \Omega, 
	\]
	and also assume that $\partial_\eta u = 0$ on the (possibly empty) set $\partial\Omega \backslash (\partial_1\Omega \cup \partial_2\Omega)$. If $\critu\neq\emptyset$, then $\nabla u \neq 0$ everywhere on $\partial_1\Omega \cup \partial_2\Omega$. %and $\critu$ is a closed subsets of $\R^n$ disjoint from $\partial\Omega$. Moreover, if \tcr{there exists a closed set $\Gamma\subseteq\critu$ such that $\Omega\backslash\Gamma$ is disconnected}  %$\Omega\backslash\critu$ is disconnected.
	Moreover, if there exists $\Gamma\subseteq\critu$ such that $\Omega\backslash\Gamma$ is disconnected then $\partial_1\Omega$ and $\partial_2\Omega$ are parallel hyperplanes, $\partial\Omega = \partial_1\Omega \cup \partial_2\Omega$, $u$ is 1D and $\critu$ is a hyperplane parallel to each $\partial_j\Omega$.
%	\begin{equation} \label{sg-cyl}
%		\int_{\Omega \cap B_R} |\nabla u|^2 \lambda_{\max}(|\nabla u|) = O(R^2 \log R) \qquad \text{as } \, R \to \infty \, ,
%	\end{equation}
%	\tcb{the existence of a $C^3$ hypersurface $\Gamma \subset\critu$ properly embedded in $\Omega$. If $u$ has moderate energy growth, then $c_1$ and $c_2$ are both nonzero, $\partial_1\Omega$ and $\partial_2\Omega$ are parallel hyperplanes, $u$ is 1D and $\Gamma \equiv \crit(u)$ is a hyperplane parallel to each $\partial_j\Omega$.}
\end{theorem}

\begin{remark}\label{rem_keyR2}
	In dimension $n=2$ the condition $\langle \eta, X_j \rangle \not \equiv 0$ in $\partial_j\Omega$ can be omitted. Indeed, if $\langle \eta, X_j \rangle \equiv 0$ then $X_j$ is tangent to $\partial_j\Omega$, hence $\partial_j\Omega$ is a line. One can therefore rotate $X_j$ to find $\widetilde X_j \in \mathbb{S}^{1}$ such that $\langle \eta, \widetilde X_j \rangle >0$ on  $\partial_j\Omega$.
\end{remark}

\begin{proof}
	%	As a preparation for the complete proof of the theorem, we prove the following preliminary claim: If there exists a connected component $\Omega_0$ of $\Omega\backslash\critu$ such that $\partial\Omega_0 \subseteq \critu$, then $\Omega_0$ is either a slab or a half-space, that is, there exist a unit vector $v\in\mathbb{S}^{n-1}$ and an open interval $I\subseteq\R$ such that
	%	\begin{equation} \label{Om-slab}
		%		\Omega_0 = \{ x \in \R^n : \langle x,v\rangle \in I \}
		%	\end{equation}
	%	and there exists a strictly monotone $u_0 : I \to \R$ such that
	%	\begin{equation} \label{u0-slab}
		%		u(x) = u_0(\langle x,v\rangle) \qquad \text{for all } \, x \in I \, .
		%	\end{equation}
	%	We prove the claim. Suppose that $\Omega_0$ is a connected component of $\Omega\backslash\critu$ such that $\partial\Omega_0 \subseteq \critu$. In particular, for each $j=1,\dots,n$ the function $w_j = \partial_j u$ satisfies $w_j \equiv 0$ on $\partial\Omega_0$ and
	%	\[
	%		\diver(A(\nabla u)\nabla w_j) + f'(u)w_j = 0 \qquad \text{in } \, \Omega \, .
	%	\]
	%	Since $w_j^2 \leq |\nabla u|^2$ and \eqref{sg-cyl} holds, we can apply Lemma \ref{Lemma_w-} deduce that each function $w_j$ either vanishes identically on $\Omega_0$ or has a constant strict sign in $\Omega_0$. Suppose, by contradiction, that $w_j \equiv 0$ on $\Omega_0$ for $j=1,\dots,n$. By our assumptions on $a$ and $u$ the matrix function $\Omega \ni x \mapsto A(\nabla u(x))$ is locally Lipschitz continuous and elliptic, so the unique continuation principle implies that $w_j \equiv 0$ on $\Omega$ for $j=1,\dots,n$, that is, $\nabla u\equiv0$ in $\Omega$, so $u$ must be constant in $\Omega$, contradiction (since $u=0$ on $\partial\Omega$ but $u$ also attains a positive maximum in $\Omega$). Therefore, there exists $i\in\{1,\dots,n\}$ such that either $w_i>0$ everywhere in $\Omega_0$ or $w_i<0$ everywhere in $\Omega_0$, that is, $u$ is strictly monotone in the direction of some parallel vector field on $\Omega_0$. Then, repeating the argument of the proof of Theorem \ref{teo_splitting} we have that $\Omega_0$ and $u$ split as claimed. \tcr{[Per questo argomento è sufficiente che sia $\critu\subseteq\{\nabla u = 0\}$ e che valgano Picone+Poincaré su aperti di perimetro localmente finito tali che $\nabla u\equiv0$ sulla frontiera.]} %closure of $\Omega_0$ in $\Omega$ is a proper subset of $\Omega$. Let $\Omega_1$ be a connected component of $\Omega\backslash\overline{\Omega}_0$. Then $\Omega_1\backslash A$ is a union of connected components of $\Omega\backslash A$, so $\partial\Omega_1 \subseteq A$. Since $A\subseteq\Omega$ is properly embedded, we have $\Omega_0\subsetneq\Omega$ and $\Omega$ is connected, we have $\partial\Omega_0\cap\Omega \neq \emptyset$. Hence, $\partial\Omega\subsetneq\partial\Omega_0$ and we have Since $\Omega_0 \cup \partial\Omega_0 = \overline{\Omega}_0 \subseteq \overline{\Omega} = \Omega \cup \partial\Omega$,
	
	%For $i=1,2$, let $c_i$ denote the constant value of $\partial_\eta u$ on $\partial_i\Omega$. We first show that $c_1$ and $c_2$ are both nonzero if $\critu\neq\emptyset$; from this fact it will follow that $\nabla u\neq 0$ everywhere on $\partial\Omega$, implying that $\critu$ is a closed subset of $\overline{\Omega}$, hence of $\R^n$, disjoint from $\partial\Omega$. Suppose by contradiction that the claim is false. Then either $c_1=c_2=0$ or $c_i = 0$ for exactly one index $i\in\{1,2\}$.
	We first show that if $\critu\neq\emptyset$ then $\nabla u\neq 0$ everywhere on $\partial_1\Omega \cup \partial_2\Omega$. Suppose by contradiction that this is false. Then, denoting by $c_j$ the constant value of $\partial_\eta u$ on $\partial_j\Omega$, we either have $c_j = 0$ for each $j$ or, up to switching index 1 and 2, $c_1 \neq 0$ and $c_j = 0$ for $j \ge 2$.
	
	\textbf{Case 1.} Suppose that $c_j = 0$ for each $j$. Then $\nabla u = 0$ everywhere on $\partial\Omega$. Since $u$ is non-constant in $\Omega$, by Lemma \ref{lem_Du0} applied with $\Omega_0 = \Omega$ we get $\critu=\emptyset$, contradiction.
	
%	there exists $X\in\mathbb{S}^{n-1}$ such that $\langle\nabla u,X\rangle > 0$ everywhere in $\Omega$. 
%	In particular, for each $j=1,\dots,n$ the function $w_j = \partial_j u$ satisfies
%	\[
%		\diver(A(\nabla u)\nabla w_j) + f'(u)w_j = 0 \qquad \text{in } \, \Omega
%	\]
%	and $w_j \equiv 0$ on $\partial\Omega$. Since $w_j^2 \leq |\nabla u|^2$ and \eqref{sg-cyl} holds, by Lemma \ref{Lemma_w-} to deduce that each function $w_j$ either vanishes identically on $\Omega_0$ or has a constant strict sign in $\Omega_0$. Since $u$ is non-constant, there exists some index $j\in\{1,\dots,n\}$ such that $w_j\not\equiv 0$, so that either $w_j>0$ everywhere in $\Omega$ or $w_j<0$ everywhere in $\Omega$.
%	But then $\nabla u$ never vanishes in $\Omega$, so $\critu=\emptyset$, contradiction.
	
	\textbf{Case 2.} Suppose that $c_1\neq 0$ and $c_j=0$ for $j \ge 2$. By our assumptions on $\partial_1\Omega$, the inner product
	\[
		\langle\nabla u,X_1\rangle = \langle\nabla u,\eta\rangle \langle\eta,X_1\rangle = c_1 \langle\eta, X_1\rangle
	\]
	does not change sign and does not vanish identically on $\partial_1\Omega$, while $\langle\nabla u,X_1\rangle = 0$ everywhere on $\partial\Omega\backslash\partial_1\Omega$. Up to replacing $X_1$ with $-X_1$, we can ensure that
	\[
	\langle \nabla u, X_1 \rangle \ge 0, \quad \langle \nabla u,X_1 \rangle \not\equiv 0 \qquad \text{on } \, \partial \Omega. 
	\]
	We can therefore apply Theorem \ref{teo_monotonicity} to deduce that $\langle\nabla u,X_1\rangle > 0$ everywhere in $\Omega$, yielding again $\critu=\emptyset$ and therefore the desired contradiction.
	 
	Now, suppose that there exists $\Gamma \subset \critu$ disconnecting $\Omega$. Then $\Omega\setminus\overline{\Gamma}$ is also disconnected. Indeed, since $\Omega\setminus\Gamma$ is disconnected there exist open sets $U_1,U_2\subseteq\R^n$ such that
	\[
		\Omega\backslash\Gamma \subseteq U_1\cup U_2 \, , \qquad U_1\cap U_2\cap(\Omega\backslash\Gamma) = \emptyset \, , \qquad U_i \cap (\Omega\backslash\Gamma) \neq \emptyset \quad \text{for } i = 1,2.
	\]
	Since $\Omega\backslash\overline{\Gamma} \subseteq \Omega\backslash\Gamma$ we also have
	\[
		\Omega\backslash\overline{\Gamma} \subseteq U_1\cup U_2 \, , \qquad U_1\cap U_2\cap(\Omega\backslash\overline{\Gamma}) = \emptyset
	\]
	Moreover, since $(\Omega\backslash\Gamma)\backslash(\Omega\backslash\overline{\Gamma}) = \Omega\cap(\overline{\Gamma}\backslash\Gamma) \subseteq \critu$ and $\critu$ has empty interior by non-constancy of $u$ and unique continuation, 
%	none of the sets $U_i$ can have intersection with $\Omega\backslash\Gamma$ entirely contained in the difference set $(\Omega\backslash\Gamma)\backslash(\Omega\backslash\overline{\Gamma})$, that is, 
	we also have
	\[
		\qquad U_i \cap (\Omega\backslash\overline{\Gamma}) \neq \emptyset \quad \text{for } i = 1,2.
	\]
	This shows that $\Omega\backslash\overline{\Gamma}$ is disconnected, as claimed. Since $\nabla u=0$ everywhere on $\overline{\Gamma}$ by continuity of $|\nabla u|$ in $\overline{\Omega}$ and, on the other hand, $\nabla u\neq0$ everywhere on $\partial_1\Omega\cup\partial_2\Omega$, we have for $j=1,2$ that each $\partial_j\Omega$ is disjoint from the closed set $\overline{\Gamma}$ and therefore entirely contained in the boundary of some connected component of $\Omega\backslash\overline{\Gamma}$. On the other hand, each component of $\Omega\backslash\overline{\Gamma}$ has boundary contained in $\partial\Omega \cup \overline{\Gamma}$, and $\partial_1\Omega \cup \partial_2\Omega$ is precisely the subset of $\partial\Omega \cup \overline{\Gamma}$ where $\nabla u\neq 0$. By Lemma \ref{lem_Du0}, $\nabla u$ cannot vanish identically on the boundary of a component of $\Omega\backslash\overline{\Gamma}$, otherwise we would deduce $\critu=\emptyset$, contradiction. Summarizing,  we conclude that $\Omega\backslash\overline{\Gamma}$ has exactly two connected components, call them $\Omega_1$ and $\Omega_2$, and without loss of generality we can assume $\partial_j\Omega\subseteq\partial\Omega_j$ for $j=1,2$.
	%Since $\partial\Omega$ is disjoint from $\critu$, hence from $\Gamma$, we have that each connected component of $\partial\Omega$ is entirely contained in the boundary of some connected component of the disconnected open set $\Omega\backslash\tcr{\Gamma}$. On the other hand, each component of $\Omega\backslash\tcr{\Gamma}$ has boundary contained in $\partial\Omega\cup\critu$. None of them can have boundary entirely contained in $\critu$, otherwise by Lemma \ref{lem_Du0} we would deduce $\critu = \emptyset$, contradiction. Hence, $\Omega\backslash\tcr{\Gamma}$ has exactly two components, $\Omega_1$ and $\Omega_2$, and without loss of generality we can assume $\partial_i\Omega\subseteq\partial\Omega_i$ for $i=1,2$.
	%We claim that they are actually contained in the boundaries of different components of $\Omega\setminus\critu$. Indeed, if this were false then there would exist another connected component $\Omega_0$ of $\Omega\setminus\critu$ whose boundary is entirely contained in $\critu$. Applying Lemma \ref{lem_Du0} to $\Omega_0$ we would deduce $\critu = \emptyset$, contradiction. reach a contradiction. Whence, up to renaming we can assume $\partial_j \Omega \subseteq \partial \Omega_j$.
	
%	since $\Gamma \subseteq \critu$ is properly embedded in $\Omega$ then $\Gamma$ is properly embedded also in $\R^n$. Thus, by the Jordan separation theorem (see \cite{Lima}) $\R^n \backslash \Gamma$ has two connected components $U_1,U_2$. By restriction, $\Omega \backslash \Gamma$ has two components $\Omega_j = \Omega \cap U_j$, as well. Each of them has boundary (as a subset of $\R^n$) contained in $\critu \cup \partial\Omega = \critu \cup \partial_1\Omega \cup \partial_2\Omega$. We claim that the sets $\partial_1\Omega$ and $\partial_2\Omega$ are contained in different components of $\R^n\setminus\Gamma$. Indeed, if this were false then $\Omega_j = U_j$ for some $j$, say $j=1$, but then $\partial\Omega_1 = \Gamma \subseteq\critu$. Applying Lemma \ref{lem_Du0} with the choice $\Omega_0 = \Omega_1$ we reach a contradiction. Whence, up to renaming we can assume $\partial_j \Omega \subseteq \partial \Omega_j$.
	
%	Having established $c_i\neq0$ for $i=1,2$, we have $\nabla u\neq 0$ everywhere on $\partial\Omega$. Therefore, \tcb{since $\Gamma \subseteq \critu$ is properly embedded in $\Omega$ then $\Gamma$ is properly embedded also in $\R^n$. Thus, by the Jordan separation theorem (see \cite{Lima}) $\R^n \backslash \Gamma$ has two connected components $U_1,U_2$. By restriction, $\Omega \backslash \Gamma$ has two components $\Omega_j = \Omega \cap U_j$, as well. Each of them has boundary (as a subset of $\R^n$) contained in $\critu \cup \partial\Omega = \critu \cup \partial_1\Omega \cup \partial_2\Omega$. We claim that the sets $\partial_1\Omega$ and $\partial_2\Omega$ are contained in different components of $\R^n\setminus\Gamma$. Indeed, if this were false then $\Omega_j = U_j$ for some $j$, say $j=1$, but then $\partial\Omega_1 = \Gamma \subseteq\critu$. Applying Lemma \ref{lem_Du0} with the choice $\Omega_0 = \Omega_1$ we reach a contradiction. Whence, up to renaming we can assume $\partial_j \Omega \subseteq \partial \Omega_j$.}
	
	For $j=1,2$, the inner product $\langle\nabla u,X_j\rangle = c_j\langle\eta,X_j\rangle$ does not change sign and does not vanish identically on $\partial_j\Omega$, while $\langle\nabla u,X_j\rangle = 0$ everywhere else on $\partial\Omega_j$. Up to replacing $X_j$ with $-X_j$ we can assume that
	\[
	\langle \nabla u, X_j \rangle \ge 0, \quad \langle \nabla u,X_j \rangle \not\equiv 0 \qquad \text{on } \, \partial \Omega_j, 
	\]
	so indeed $\langle\nabla u,X_j\rangle >0$ in $\Omega_j$ because of  Theorem \ref{teo_monotonicity}. By Lemma \ref{Om_finper}, each $\Omega_j$ has locally finite perimeter in $\Omega$. We then apply Theorem \ref{teo_splitting} to conclude that each $\Omega_j$ is a slab in $\R^n$ bounded by two parallel hyperplanes, and in each of them $u$ is strictly monotone in the direction perpendicular to the boundary. Since $\Omega_1 \cap \Omega_2 = \emptyset$, the hyperplanes bounding $\Omega_1$ and $\Omega_2$ must all be parallel. In particular, $\partial_1\Omega$ and $\partial_2\Omega$ are parallel hyperplanes and the set $\Omega\cap\overline{\Gamma} = \Omega \backslash (\Omega_1\cup\Omega_2)$, having empty interior, must be a hyperplane parallel to both $\partial_1\Omega$ and $\partial_2\Omega$. In particular, $\Omega$ is the open slab bounded between $\partial_1\Omega$ and $\partial_2\Omega$, so $\partial\Omega$ has no other component besides them. Lastly, since $\Omega\cap\overline{\Gamma}\subseteq\critu$ and $\nabla u\neq 0$ everywhere in $\Omega_1\cup\Omega_2$ we have $\critu = \Omega\cap\overline{\Gamma}$, so $\critu$ is a hyperplane parallel to both $\partial_j\Omega$ as claimed.
\end{proof}
	
	
	
	
	
	
	
	
	
	
	 
%	Each connected component of $\Omega\backslash \tcb{\Gamma}$ is an open set whose boundary (as a subset of $\R^n$) is contained in $\critu \cup \partial\Omega = \critu \cup \partial_1\Omega \cup \partial_2\Omega$. We claim that for $i=1,2$ the set $\partial_i\Omega$ is entirely contained in the boundary of a unique connected component $\Omega_i$ of $\tcb{\Omega\backslash \Gamma}$ and is disjoint from the closure of any other component of $\tcb{\Omega\backslash\Gamma}$. Note that we are not claiming $\Omega_1 \neq \Omega_2$; in fact, this will be shown later in the proof.
%	
%	To prove the claim in case $i=1$, let $x_0\in\partial_1\Omega$ be fixed and let $x\in\partial_1\Omega$ be any other point. Let $\sigma : [0,1] \to \partial_1\Omega$ be a curve such that $\sigma(0) = x_0$ and $\sigma(1) = x$. Consider a variation $V : [0,\eps] \times [0,1] \to \overline{\Omega}$ of $\sigma$ such that
%	\[
%		V(t,0) = x_0 \, , \qquad V(t,1) = x \, , \qquad V(t,s) \in \Omega\backslash\critu
%	\]
%	for all $t\in(0,\eps]$ and $s\in(0,1)$, which exists since $\Gamma$ is closed and disjoint from $\partial\Omega$. Then the image $V((0,\eps]\times(0,1))$ is entirely contained in some connected component of $\tcb{\Omega\backslash\Gamma}$. By arbitrariness of $x$, this shows that $\partial_1\Omega$ is entirely contained in the closure of some connected component $\Omega_1$ of $\tcb{\Omega\backslash\Gamma}$. Since each point $x\in\partial_1\Omega$ has a neighbourhood whose intersection with $\Omega$ is diffeomorphic to a Euclidean half ball, hence connected, we also see that $x$ cannot be in the closure of more than one connected component of $\tcb{\Omega\backslash\Gamma}$, proving uniqueness of $\Omega_1$.
%	
%	%	For $i=1,2$, let $\Omega_i$ be the connected component of $\Omega\backslash\critu$ such that $\critu_i\subseteq\partial\Omega_i$. We claim that $\Omega_1\neq\Omega_2$. Suppose, by contradiction, that this is false, so that $\Omega_1$ and $\Omega_2$ are the same connected component of $\Omega\backslash\critu$. Since $\Omega\backslash\critu$ is not connected, there exists another component $\Omega_0$ of $\Omega\backslash\critu$ different from $\Omega_1\equiv\Omega_2$. By what we observed above, the boundary of $\Omega_0$ does not meet $\critu_1$ or $\critu_2$, so it must be $\partial\Omega_0\subseteq\critu$. From the preliminary claim we have that $\Omega_0$ and $u$ must be as in \eqref{Om-slab}-\eqref{u0-slab} for some unit vector $v\in\mathbb{S}^{n-1}$, open interval $I\subseteq\R$ and some $u_0 : I \to \R$. Up to a rotation of coordinates, we can assume that $v = \mathbf{e}_1$, so that $\partial_j u \equiv 0$ in $\Omega_0$ for $j=2,\dots,n$. Since $u$ is constant on each connected component of $\partial\Omega$ but it is not constant in $\Omega$, we have that $\critu_1$ and $\critu_2$ are both hyperplanes orthogonal to $v$. \tcr{[Giustificare bene.]} Then $\Omega$, being connected, must be the slab contained between $\critu_1$ and $\critu_2$. Since $\Omega_0\subseteq\Omega$, the hyperplanes $\critu_1$ and $\critu_2$ must lie on opposite sides with respect to $\Omega_0$, but then we have $\Omega_1\neq\Omega_2$, contradiction.
%	
%	%For $i=1,2$, let $\Omega_i$ be the connected component of $\Omega\backslash\critu$ such that $\critu_i\subseteq\partial\Omega_i$.
%	We now claim that $\Omega_1\cup\Omega_2 = \tcb{\Omega\backslash\Gamma}$. Suppose, by contradiction, that this is false, that is, there exists a component $\Omega_0$ of $\tcb{\Omega\backslash\Gamma}$ whose boundary contains neither $\partial_1\Omega$ nor $\partial_2\Omega$. By what we observed above it must be $\partial\Omega_0\subseteq\critu$, so we have $\nabla u = 0$ on $\partial\Omega_0$. Since $\Omega_0\subsetneq\Omega$, by Lemma \ref{lem_Du0} we have $u$ constant in $\Omega_0$. \tcb{But then $\Omega_0 \subseteq \critu$, which is impossible since at any point of $\critu$ the partial derivatives $\partial_j u$ have finite order of vanishing unless they are constant on the entire $\Omega$, see the argument at the end of Lemma \ref{Om_finper} \tcr{DO WE PUT IT IN A SEPARATE REMARK AT THE BEGINNING OF SECTION 5?}.} 
%%	For each $j=1,\dots,n$ the function $w_j = \partial_j u$ satisfies
%%	\[
%%		\diver(A(\nabla u)\nabla w_j) + f'(u)w_j = 0 \qquad \text{in } \, \Omega_0
%%	\]
%%	and
%%	\[
%%		w_j \equiv 0 \qquad \text{on } \, \partial\Omega_0 \, .
%%	\]
%%	Similarly to how we did at the beginning of the proof, by Lemma \ref{Lemma_w-} we deduce that each function $w_j$ either vanishes identically on $\Omega_0$ or has a constant strict sign in $\Omega_0$. Suppose that $w_j \equiv 0$ on $\Omega_0$ for $j=1,\dots,n$. By our assumptions on $a$ and $u$ the matrix function $\Omega \ni x \mapsto A(\nabla u(x))$ is locally Lipschitz continuous and elliptic, so the unique continuation principle implies that $w_j \equiv 0$ on $\Omega$ for $j=1,\dots,n$, that is, $\nabla u\equiv0$ in $\Omega$, so $u$ must be constant in $\Omega$ and we reach a contradiction (since $u=0$ on $\partial\Omega$ but $u$ attains a positive maximum in $\Omega$). Therefore, there exists $j\in\{1,\dots,n\}$ such that $w_j\neq 0$ in $\Omega_0$. By stability of $u$, we also have the existence of $0<v\in C^2(\Omega)$ satisfying
%%	\[
%%		\diver(A(\nabla u)\nabla v) + f'(u)v = 0 \qquad \text{in } \, \Omega \, .
%%	\]
%%	By Remark \ref{rem-Lemma_w-} we have that $w_j$ is a nonzeroconstant multiple of $v$ in $\Omega_0$. Since $v>0$ on $\partial\Omega_0\subseteq\critu\subseteq\Omega$, by continuity of $v$ and $w_j$ in $\Omega$ we must also have $w_j\neq 0$ everywhere on $\partial\Omega_0$, and so we reach again a contradiction.
%	
%	So, we established that $\Omega_1\cup\Omega_2 = \tcb{\Omega\backslash\Gamma}$. Since $\tcb{\Omega\backslash\Gamma}$ is disconnected, it must be $\Omega_1\neq\Omega_2$, so $\partial\Omega_i \subseteq \partial_i\Omega \cup \critu$ for $i=1,2$. We also know that $c_i\neq 0$ for $i=1,2$. By our assumptions on the structure of $\partial_i\Omega$, for $i=1,2$ the inner product $\langle\nabla u,X_i\rangle = c_i\langle\eta,X_i\rangle$ has a constant weak sign on $\partial_i\Omega$ and is not everywhere vanishing on $\partial_i\Omega$, while $\langle\nabla u,X_i\rangle = 0$ on $\Gamma$. Up to replacing $X_i$ with $-X_i$ we can assume that
%	\[
%		\langle\nabla u,X_i\rangle \geq 0 \, , \quad \langle\nabla u,X_i\rangle \not\equiv 0 \qquad \text{on } \, \partial\Omega_i
%	\]
%	for $i=1,2$, \tcb{so indeed $\langle\nabla u,X_i\rangle >0$ on $\Omega_i$ because of Theorem \ref{teo_monotonicity}.} We then apply Theorem \ref{teo_splitting} to both $\Omega_1$ and $\Omega_2$ to conclude that each of these sets is a slab in $\R^n$ bounded by two parallel hyperplanes, and in each of them $u$ is strictly monotone in the direction perpendicular to the boundary. Since $\Omega_1 \cap \Omega_2 = \emptyset$, the hyperplanes bounding $\Omega_1$ and $\Omega_2$ must all be parallel. In particular, $\partial_1\Omega$ and $\partial_2\Omega$ are parallel hyperplanes and \tcb{$\Gamma \equiv \critu = \Omega \backslash (\Omega_1\cup\Omega_2)$}	 is a hyperplane parallel to both $\partial_1\Omega$ and $\partial_2\Omega$.
%	
%	%	So, we established that $\Omega_1\neq\Omega_2$. In particular, $\partial\Omega_i \subseteq \critu_i \cup \critu$ for $i=1,2$. We now show that $\Omega_1$ and $\Omega_2$ are slabs and that $u$ is one-dimensional in each of these regions. We prove the claim for $\Omega_1$, the argument being the same for $\Omega_2$. If $c_1=0$, then $\nabla u\equiv 0$ on $\partial\Omega_1$ and by the same reasoning used to prove the preliminary claim we conclude that $\Omega_1$ is a slab and that $u$ is one-dimensional. So, let us suppose that $c_1\neq 0$. By our assumptions on the structure of $\critu_1$, the inner product
%	%	\[
%	%		\langle\nabla u,X_1\rangle = \langle\nabla u,\eta\rangle \langle\eta,X_1\rangle = c_1 \langle\eta, X_1\rangle
%	%	\]
%	%	has a constant weak sign on $\critu_1$ and is not vanishing everywhere on $\critu_1$, while clearly $\langle\nabla u,X_1\rangle = 0$ on $\critu$. Up to replacing $X_1$ with $-X_1$, we can assume that
%	%	\[
%	%		\langle\nabla u,X_1\rangle \geq 0 \, , \quad \langle\nabla u,X_1\rangle \not\equiv 0 \qquad \text{on } \, \partial\Omega_1 \, .
%	%	\]
%	%	We then apply Theorem \ref{teo_splitting} to conclude that $\Omega_1$ is a slab and that $u$ is one-dimensional and strictly monotone in the direction perpendicular to the boundary. \tcr{[Per questo passaggio occorre che Picone+Poincaré valgano su domini soddisfacenti le ipotesi dell'insieme $\Omega_0$ nel Lemma \ref{Om_finper}.]}
%	
%	%	We obtained that each connected component of $\Omega$ is a slab bounded between a pair of parallel hyperplanes. Since these components are pairwise disjoint, the bounding hyperplanes must be all parallel. In particular, $\critu_1$ and $\critu_2$ are parallel and $\critu$ is a finite union of hyperplanes parallel to both $\critu_1$ and $\critu_2$. \tcr{[Nota: le ipotesi di regolarità su $\critu$ necessarie per applicare il teorema di spezzamento sono quelle da imporre \emph{sulle componenti di $\partial\Omega$ dove $\nabla u = 0$} affinché sia valido il Teorema \ref{teo_splitting}; potrebbe essere meno di $C^3$? Inoltre, anche qui non si usa il fatto che $\critu$ sia l'insieme dei massimi: se $\critu\subseteq\{\nabla u = 0\}$ allora su ogni componente connessa di $\critu$ vale $u=\text{const.}$ e $\nabla u = 0$.]}
%\end{proof}



\section{Proof of the main theorems for capillary graphs}\label{sec_proof}

\begin{proof}[Proof of Theorem \ref{teo_main_MC_manifolds}]
	Under assumptions (A), (B) or (C) we apply items $(ii)$, $(i)$ or $(v)$, respectively, in Theorem \ref{teo_goodcutoff_MC} for the choice $a(t) = 1/\sqrt{1+t^2}$ to infer, in view of \eqref{eq_quadraticgrowth}, that $u$ has moderate energy growth. On the other hand, under assumption (D) we apply Theorem \ref{prop_grad} to deduce that $|\nabla u|\in L^\infty(\Omega)$ and then Proposition \ref{prop_moderate_boundgradient} to conclude, again, that $u$ has moderate energy growth. Condition \eqref{condi_Hcj} can be rewritten as $\langle \nabla u, X \rangle \ge 0$ and $\langle \nabla u, X \rangle \not \equiv 0$ on $\partial \Omega$, whence $\langle\nabla u,X\rangle>0$ in $\Omega$ by Theorem \ref{teo_monotonicity}. The conclusion follows by Theorem \ref{teo_splitting}.
\end{proof}


\begin{proof}[Proof of Theorem \ref{teo_CMC_intro_R2}]
	First, observe that by regularity theory $u$ is analytic in $\Omega$. Under any of the assumptions (i)-(iv) we have that $|\nabla u| \in L^\infty(\partial\Omega)$, since $|\nabla u| = |\partial_\eta u|$ is locally constant on $\partial\Omega$ and the subset $\partial_\star \Omega \subseteq \partial\Omega$ where $\partial_\eta u\neq 0$ has finitely many connected components. Hence, $|\nabla u| \in L^\infty(\Omega)$ by Theorem \ref{prop_grad}. Clearly, any subset $\Omega \subseteq \R^2$ satisfies $|\Omega \cap B_R| \le CR^2$, so $u$ has moderate energy growth by Proposition \ref{prop_moderate_boundgradient}.
%	We claim that $u$ has moderate energy growth in each of the following cases:
%	\begin{itemize}
%		\item[-] Under the assumptions in $(i)$ or $(ii)$;
%		\item[-] Under the assumptions in $(iii)$, if $b_1 = b_2$;
%		\item[-] Under the assumptions in $(iii)$, if $b_1 \neq b_2$ and \eqref{eq_areagrowth_intro} holds.
%	\end{itemize}
%	Indeed, the claim follows from a direct application of Theorem \ref{teo_goodcutoff_MC}.} %with the improvements in Remark \ref{rem_very_useful} to deal with $(ii)$ and $(iii)$.
	%We also observe that in case $(i)$ the gradient of $u$ does not vanish identically on $\partial \Omega$, since otherwise one could apply Lemma \ref{lem_Du0} with $\Omega_0 \equiv \Omega$ to deduce the existence of $X \in \mathbb{S}^1$ such that $\langle \nabla u, X \rangle >0$ in $\Omega$. Then, by Theorem \ref{teo_splitting} one would have that $\Omega = (0,\infty) \times \R$ and that $u$ is a half-plane, is 1D, therefore either a piece of non-constant plane or cylinder, none of which has vanishing normal derivative on the whole boundary.
	
	$(i)$. Since $\partial\Omega$ is connected we have that $\partial_\eta u = c_1$ is constant on $\partial\Omega$. Moreover, by Remark \ref{rem_keyR2} we can assume that $\langle \eta, v \rangle \ge 0$ and $\not \equiv 0$ on $\partial \Omega$. We claim that $c_1\neq 0$. In this case, by Theorem \ref{teo_monotonicity} we obtain $\langle \nabla u, X \rangle > 0$ in $\Omega$ with $X = v$. A direct application of Theorem \ref{teo_splitting} yields that $\Omega = (0,\infty) \times \R$ and that $u$ is 1D. Hence, necessarily $H=0$ and $\Omega$, $\GG_u$ are half-planes. To show the claim, if it were $c_1 = 0$ then $\nabla u$ would vanish identically on $\partial\Omega$ and we could apply Lemma \ref{lem_Du0} with $\Omega_0 \equiv \Omega$ to deduce the existence of $X \in \mathbb{S}^1$ such that $\langle \nabla u, X \rangle >0$ in $\Omega$. Then, again, by Theorem \ref{teo_splitting} we would have that $\Omega = (0,\infty) \times \R$ and that $u$ is a non-constant affine function, contradicting the vanishing of $\nabla u$ on $\partial\Omega$.
	
	$(ii)$. By assumption $c_1 = \partial_\eta u \neq 0$, $c_j = 0$ for each $j\neq 1$. By Remark \ref{rem_keyR2}, up to changing the sign of $v$ and rotating it we can assume $\langle\eta,v\rangle > 0$ on $\partial_1\Omega$. Then, assumption \eqref{condi_Hcj} in Theorem \ref{teo_main_MC_manifolds} is satisfied for $X = v$ and the conclusion follows by a direct application of that theorem (note that, in particular, the a priori disconnected set $\partial\Omega$ has exactly two connected components and $\partial_\eta u = c_2 = 0$ on $\partial_2\Omega$; this excludes the possibility that $u$ be affine, forcing $H\neq 0$ and $\GG_u$ to be a piece of cylinder). 
%	On the other hand, if $\langle\eta,v\rangle\equiv0$ on $\partial_1\Omega$ then the embedded complete curve $\partial_1\Omega \subseteq \R^2$ is everywhere tangent to the parallel vector field $v$, so $\partial_1\Omega$ is a line and there exists a different vector $w\in\mathbb{S}^1$ such that $\langle\eta,w\rangle > 0$ everywhere on $\partial_1\Omega$. Then, again, the assumption \eqref{condi_Hcj} in Theorem \ref{teo_main_MC_manifolds} is satisfied for $X=w$ and we conclude as desired.
	
	$(iii)$. If $\langle\eta,v\rangle\not\equiv 0$ on $\partial_\star\Omega$ then we apply Theorem \ref{teo_main_MC_manifolds} for $X$ equal to either $v$ or $-v$ to deduce that $\Omega =(0,T)\times\R$ and $u$ is 1D and strictly monotone, hence $\GG_u$ is a piece of a cylinder or of a half-plane. On the other hand, if $\langle\eta,v\rangle\equiv 0$ on $\partial_\star\Omega$ then $\partial_1\Omega$ and $\partial_2\Omega$ are parallel lines in $\R^2$, both tangent to $v$, and then the connected set $\Omega$ must be contained in the strip bounded by them. In particular $\eta$ points in opposite directions on $\partial_1\Omega$ and $\partial_2\Omega$, so there exists $w\in\mathbb{S}^1$ such that $\langle\eta,w\rangle>0$ on $\partial_1\Omega$ and $\langle\eta,w\rangle<0$ on $\partial_2\Omega$ and again we reach the desired conclusion by applying Theorem \ref{teo_main_MC_manifolds} with $X$ equal to either $w$ or $-w$.
	
	$(iv)$. %Let $\partial_\star \Omega = \partial\Omega \cap \{\partial_\eta u \neq 0\}$. If $u$ is constant on $\partial_\star\Omega$, then $u$ has moderate energy growth in $\Omega$ by Theorem \ref{teo_goodcutoff_MC} applied for the choice $\Omega_0 = \Omega$ and the conclusion follows from Theorem \ref{teo_twopieces}. On the other hand, if $u$ is not constant on $\partial_\star \Omega$ then necessarily $\partial_\star \Omega$ consists of exactly two components $\partial_1\Omega$ and $\partial_2\Omega$. Each of these two sets is entirely contained in the boundary of some connected component of $\Omega\backslash\Gamma$; conversely, no connected component of $\Omega\backslash\Gamma$ can be such that $\nabla u = 0$ everywhere on its boundary, otherwise $u$ would be of moderate energy growth on that component by Theorem \ref{teo_goodcutoff_MC} and then we would reach a contradiction by Lemma \ref{lem_Du0}. Therefore, $\Omega\backslash\Gamma$ consists of exactly two connected components $\Omega_1$ and $\Omega_2$ satisfying $\partial_j\Omega \subseteq \partial\Omega_j$ for $j=1,2$. Now, for each $j\in\{1,2\}$ we have that $u$ is constant on $\overline{\Omega}_j \cap \partial_\star\Omega \equiv \partial_j\Omega$, so by Theorem \ref{teo_goodcutoff_MC} we deduce that $u$ has moderate energy growth in each of them. Then, we can conclude as in the proof of Theorem \ref{teo_twopieces}.
	Since $u$ has moderate energy growth, we apply Theorem \ref{teo_twopieces} (together with Remark \ref{rem_keyR2}) to deduce that $\Omega$ is a slab $(0,T) \times \R$, $\Gamma$ is a straight line and $u$ is 1D. Since the graph $\GG_u$ is a surface of constant mean curvature and $\partial_\eta u\neq0$ on $\partial \Omega$, $\GG_u$ must be a piece of cylinder.
	
	We now prove that if $\inf_\Omega u < b_j$ for each $j\geq 3$ and $\Gamma$ has an accumulation point in $\overline{\Omega}$, then indeed $\Gamma$ disconnects $\Omega$. %Clearly $\Gamma\subseteq\critu$ and by assumption $\Gamma\neq\emptyset$. Hence, by Theorem \ref{teo_twopieces} we have $c_i\neq0$ for $i=1,2$, $\critu$ is a closed subset of $\overline{\Omega}$ disjoint from $\Omega$ and so the same is true of $\Gamma$.
		%Notice that $c_1$ and $c_2$ must both be nonzero, since otherwise the same argument outlined in the first part of the proof of Theorem \ref{teo_twopieces} would yield $\critu = \emptyset$, contradicting the assumption that $\Gamma\neq\emptyset$. Having established $c_1c_2\neq0$, we have $\nabla u\neq 0$ everywhere on $\partial\Omega$ and so $\critu$ is a closed subset of $\overline{\Omega}$ disjoint from $\Omega$. Therefore, the same is true of $\Gamma$.}
	By the \L ojasiewicz structure theorem (see \cite{KP}), the analytic set $\Gamma$ %$\critu$
	can be decomposed as
	\[
		\Gamma = \Gamma_0\sqcup\Gamma_1,
	\]
	where $\Gamma_0$ is a finite set and $\Gamma_1$ is a finite union of real analytic curves embedded in $\Omega$. The assumption that $\Gamma$ has an accumulation point grants that $\Gamma_1\neq\emptyset$. Since $u$ attains an interior minimum, by the strong maximum principle it must be $H < 0$ and so we have $|\nabla \di u(x)|\neq 0$ for each $x \in \Omega$. We can thus apply \cite[Corollary 3.4]{BCM} to deduce that the closure of each component $\gamma$ of $\Gamma_1$ still belongs to $\Gamma_1$, and therefore $\gamma$ is a real analytic curve properly embedded in $\Omega$. We claim that $\gamma$ is properly embedded in the entire $M$. Otherwise, there would exist a sequence $\{p_j\} \subset \gamma$ with $p_j \to p \in \partial \Omega$, and since $u \in C^1(\overline\Omega)$ we would get $u(p) = \inf_\Omega u$ and $\nabla u(p) = 0$. The strict inequality $\inf_\Omega u < b_j$ for $j \ge 3$ guarantees that this is not possible. To conclude, by the smooth Jordan-Brouwer separation theorem (see \cite{Lima}) $\gamma$ would disconnect $\R^2$, hence also $\Omega$.
	%\tcr{REMOVABLE!!!!: and the assumption $\sup_\Omega u>\max\{b_1,b_2\}$ implies that $\critu\cap\partial\Omega=\emptyset$.} \tcr{THIS IS REMOVABLE: By the smooth Jordan-Brouwer separation theorem (see \cite{Lima}), each connected component $\Gamma$ of $\Gamma_1$ splits $\R^2$ into two connected components. In particular, $\Omega \backslash \Gamma$ is disconnected.}
%
%\tcr{Previous final argument, to remove if the blue part is checked}\\
%Let us consider the connected components of $\Omega\backslash\critu$. If the boundary of one such component (say $\Omega_0$) were entirely a subset of $\critu$, then $u$ would be constant and equal to $\sup_\Omega u$ on $\partial\Omega_0$ and by Theorem \ref{teo_goodcutoff_MC} $u$ would have moderate energy growth. Since $\nabla u$ would vanish on $\partial\Omega_0$ in this case, Lemma \ref{lem_Du0} applied in $\Omega_0$ would yield a contradiction. We therefore deduce that the boundary of each component of $\Omega\backslash\critu$ either contains $\partial_1\Omega$ or $\partial_2\Omega$. This implies that $\Omega\backslash\critu$ has exactly two connected components, say $\Omega_1$ and $\Omega_2$ and, for $i=1,2$, $\partial\Omega_i=\partial_i\Omega\cup\critu$. We can now apply Theorem \ref{teo_main_MC_manifolds} to $\Omega_1$ and $\Omega_2$ separately to deduce that  $\Omega_j =(0,T_j)\times\R$ and $u$ is 1D. In particular, $\critu = \Gamma_1$ is connected and totally geodesic (hence, a straight line) and from $\partial_\eta u\neq0$ on $\partial_1\Omega \cup \partial_2 \Omega$  the graph of $u$ must be a piece of cylinder.
\end{proof}

\begin{proof}[Proof of Theorem \ref{teo_CMC_intro_R3}]
	Again, by regularity theory $u$ is analytic in $\Omega$. As in the proof of Theorem \ref{teo_CMC_intro_R2}, under any of the assumptions $(i)-(iv)$ we have that $|\nabla u| \in L^\infty(\Omega)$ and then $u$ has moderate energy growth by Proposition \ref{prop_moderate_boundgradient} and assumption \eqref{eq_quadraticgrowth}.
	
	$(i)$. Assume by contradiction that $\partial \Omega$ is connected. Then, proceeding verbatim as in $(i)$ of Theorem \ref{teo_CMC_intro_R2} above, we deduce that $\Omega = (0,\infty) \times \R^2$. This in turn would imply $|\Omega\cap B_R|\sim\frac23\pi R^3$ as $R\to \infty$, a contradiction.	
%	
%	
%	Let $\partial_1\Omega$ be mildly $v$-transverse and let us assume by contradiction that $\partial\Omega$ has only one connected component. Since $u$ is constant on $\partial\Omega$, by Theorem \ref{teo_goodcutoff_MC} it has moderate energy growth and by Lemma \ref{lem_Du0} $\nabla u$ cannot vanish identically on $\partial\Omega$, forcing $c_1\neq0$. Theorem \ref{teo_main_MC_manifolds} now implies that $\Omega=(0,\infty)\times\R^2$, 
	
	$(ii)$-$(iii)$-$(iv)$. We proceed as in cases $(ii)$, $(iii)$ and $(iv)$ of Theorem \ref{teo_CMC_intro_R2}.	
%	\tcr{Again, by Theorem \ref{teo_goodcutoff_MC} $u$ has moderate energy growth and again by Lemma \ref{lem_Du0} we must have $c_1\neq0$.} Now the claim follows by applying Theorem \ref{teo_main_MC_manifolds} and observing that the half-hyperplane case can be ruled out because of the vanishing of $c_2$.\par
	Concerning case $(iv)$, we only have to prove that $\Gamma$ disconnects $\Omega$ in case $\inf_\Omega u < b_j$ for any $j\geq 3$ and $\haus^2(\Gamma) > 0$. As observed in the proof of $(iv)$ of Theorem \ref{teo_CMC_intro_R2}, by the strong maximum principle it must be $H < 0$ and thus $\nabla \di u \neq0$ everywhere in $\Omega$. Applying \L ojasiewicz's structure theorem together with \cite[Corollary 3.4]{BCM}, from $\haus^{2}(\Gamma)>0$ we deduce that the top stratum $\Gamma_2$ of $\Gamma$ is non-empty, and that any of its components is a complete real analytic surface without boundary, properly embedded in $\Omega$. The conclusion then follows again by the smooth Jordan-Brower theorem.
\end{proof}


\begin{proof}[Proof of Theorem \ref{teo_CMC_intro_R2_gen}]
	In cases $(i)$, $(ii)$ and $(iii)$ one can always deduce that $u$ has moderate energy growth by suitable application of either Theorem \ref{teo_goodcutoff_MC} or Proposition \ref{prop_moderate_boundgradient} together with Theorem \ref{prop_grad}, according to which condition is assumed among $(A)$, $(B)$, $(C)$ or $(D)$. Then, the proof proceeds as that of cases $(i)$, $(ii)$ and $(iii)$ of Theorem \ref{teo_CMC_intro_R2}. Case $(iv)$ requires a bit of care.
	
	$(iv)$. First, note that each connected component of $\Omega\setminus\Gamma$ has locally finite perimeter in $\Omega$ by Lemma \ref{Om_finper}. If $(D)$ holds then $|\nabla u|\in L^\infty(\Omega)$ by Theorem \eqref{prop_grad} and $u$ has moderate energy growth by Proposition \ref{prop_moderate_boundgradient}, so the desired conclusion follows from Theorem \ref{teo_twopieces} together with Remark \ref{rem_keyR2}. Suppose, instead, that $(A)$ or $(C)$ holds. Let $\partial_\star \Omega = \partial\Omega \cap \{\partial_\eta u \neq 0\}$. If $u$ is constant on $\partial_\star\Omega$, then $u$ has moderate energy growth in $\Omega$ by either item $(ii)$ or $(v)$ in Theorem \ref{teo_goodcutoff_MC} and the conclusion follows from Theorem \ref{teo_twopieces}. On the other hand, suppose that $u$ is not constant on $\partial_\star \Omega$. Each of the connected sets $\partial_1\Omega$ and $\partial_2\Omega$ is entirely contained in the boundary of some connected component of $\Omega\backslash\Gamma$; conversely, no connected component of $\Omega\backslash\Gamma$ can be such that $\nabla u = 0$ everywhere on its boundary, otherwise $u$ would be of moderate energy growth on that component by Theorem \ref{teo_goodcutoff_MC} (note that in this scenario the connected component $\Omega_0$ under discussion would satisfy $\overline{\Omega}_0 \cap \partial_\star\Omega = \emptyset$, so any item in Theorem \ref{teo_goodcutoff_MC} would be applicable) and then we would reach a contradiction by Lemma \ref{lem_Du0}. Therefore, $\Omega\backslash\Gamma$ consists of exactly two connected components $\Omega_1$ and $\Omega_2$ satisfying $\partial_j\Omega \subseteq \partial\Omega_j$ for $j=1,2$. Now, for each $j\in\{1,2\}$ we have that $u$ is constant on $\overline{\Omega}_j \cap \partial_\star\Omega \equiv \partial_j\Omega$, so by item $(ii)$ of Theorem \ref{teo_goodcutoff_MC} we deduce that $u$ has moderate energy growth in each of them. From this point on, we can conclude again as in the proof of Theorem \ref{teo_twopieces}.
	
	The only claim left to prove is that $\Gamma$ disconnects $\Omega$ whenever $u_\ast \doteq \inf_\Omega u > b_j$ for all $j\geq 3$, $f$ is analytic in a neighbourhood of $u_\ast$, $f(u_\ast) \neq 0$ and $\Gamma$ has an accumulation point in $\overline{\Omega}$. To this aim, we need only show that the top stratum $\Gamma_1$ of $\Gamma$ is a finite union of real analytic curves properly embedded in $\Omega$, and then the conclusion will follow by Jordan-Brouwer theorem as in the proof of Theorem \ref{teo_CMC_intro_R2}. The assumption $f(u_\ast) \neq 0$, the analyticity of $f$ on $I_\eps \doteq (u_\ast-\eps, u_\ast+\eps)$ and the fact that $\Gamma$ has an accumulation point enable us to use \cite[Corollary 3.4]{BCM} to deduce that the top stratum $\Gamma_1$ of $\Gamma$ is non-empty and given by a finite union of real analytic curves properly embedded in $\Omega \cap u^{-1}(I_\eps)$. Each of these curves is also properly embedded in $\Omega$. Indeed, if at least one of them were not properly embedded in $\Omega$, parametrizing it by $\gamma : \R \to \Omega$ there would exist a sequence of points $t_j \to \infty$ such that $x_j \doteq \gamma(t_j) \to x \in \Omega$. Since $\gamma(\R)$ is properly embedded in $\Omega \cap u^{-1}(I_\eps)$, we would have $u(x) = u_\ast-\eps$, contradicting the fact that $u(x_j)= u_\ast$ and the continuity of $u$. This concludes the proof of the claim.
\end{proof}


%\input{Monotonicity}
%
%\input{CollectionNotes}

\section*{Appendix}

We here prove a form of the divergence theorem which is suited for our purposes. The main point requiring some care is that the subset $\Omega_0$ considered below just has locally finite perimeter in $\Omega$ and not in $\overline\Omega$, while the vector field $W$ does not have compact support in $\Omega$. 


\begin{lemma} \label{lem_divergence}
	Let $M$ be a complete Riemannian manifold, $\Omega\subseteq M$ a $C^1$ domain with interior normal vector $\eta$ and let $\Omega_0\subseteq\Omega$ be an open set with locally finite perimeter in $\Omega$. Let $W$ be a compactly supported vector field in $\overline\Omega$, continuous in $\overline{\Omega}$ and locally Lipschitz in $\Omega$, satisfying $W=0$ on $\partial\Omega_0\cap\Omega$. If $h\in L^1(\Omega)$ is such that $\diver\,W \geq h$ weakly in $\Omega$, then
	\begin{equation} \label{GGloc}
		\int_{\overline{\Omega}_0\cap\partial\Omega} \langle W,\eta\rangle + \int_{\Omega_0} h \leq 0 .
	\end{equation}
\end{lemma}

\begin{proof}
	Let $K = \{x\in M : \dist(x, {\rm spt}\, W) < 1\}$. Let
	\[
		\Psi \, : \, \R \times \partial\Omega \to M, \qquad \Psi(t,x) = \exp_x(t\eta_x) \, .
	\]
	Since $\partial\Omega$ is $C^1$ regular and $K$ is relatively compact and open in $M$, there exists $\tau_0\in(0,1]$ such that the restriction of $\Psi$ to the open set
	\[
	E := (-\tau_0,\tau_0) \times (K\cap\partial\Omega)
	\]
	is injective and realizes a $C^1$ diffeomorphism onto $\Psi(E)$. From our choice of $\tau_0$, we have that the distance $r$ from $\partial \Omega$ is $C^1$ regular in the set
	\[
		U := K \cap \{x\in\Omega : 0 < r(x) < \tau_0\}
	\]
	and
	\begin{equation} \label{rPsiX}
		U\subseteq\Psi(E) \qquad \text{and} \qquad r = \pi_\R \circ (\Psi_{|E})^{-1} \quad \text{in } \, U
	\end{equation}
	where $\pi_\R : \R \times \partial\Omega \to \R$ denotes the projection onto the first factor. For $\tau \in (0, \tau_0/2)$, consider the cut-off function $\psi\in\lip(M)$ given by 
	\begin{equation} \label{def_psiX}
		\psi(x) = \left\{ \begin{array}{ll}
			0 & \quad \text{if } \, r(x) < \tau \\[0.1cm]
			\dfrac{r(x)-\tau}{\tau} & \quad \text{if } \, r(x) \in [\tau, 2\tau] \\[0.3cm]
			1 & \quad \text{if } \, r(x) > 2\tau.
		\end{array}\right.
	\end{equation}
	Applying the divergence theorem to the vector field $\psi W$, which is compactly supported in $\Omega$ and globally Lipschitz, we get
	\begin{equation} \label{div_X}
		0 = \int_{\Omega_0} \diver (\psi W) = \int_{\Omega_0} \langle W, \nabla \psi \rangle  + \int_{\Omega_0} \psi \, \diver W \geq \int_{\Omega_0} \langle W, \nabla \psi \rangle  + \int_{\Omega_0} \psi h \, .
	\end{equation}
	Using the definition of $\psi$ and the coarea formula,
	\begin{equation} \label{GGloc1}
		\int_{\Omega_0} \langle W, \nabla\psi\rangle = \frac{1}{\tau} \int_\tau^{2\tau} \int_{\Omega_0\cap\{r=t\}}  \langle W, \nabla r\rangle \, \di t \, .
	\end{equation}
	We claim that when passing to limits as $\tau\to0$ it holds
	\begin{align}
		\label{GGloc2a}
		\int_{\Omega_0} \psi h & \to \int_{\Omega_0} h \\
		\label{GGloc2b}
		\frac{1}{\tau} \int_\tau^{2\tau} \int_{\Omega_0\cap\{r=t\}} 
		\langle W, \nabla r\rangle \, \di t & \to \int_{\overline{\Omega}_0\cap\partial\Omega} \langle W,\eta\rangle	
	\end{align}
	so that \eqref{GGloc} follows from \eqref{div_X} and \eqref{GGloc1}. Clearly \eqref{GGloc2a} holds by the Lebesgue convergence theorem since $h\in L^1(\Omega)$, so we focus on justifying \eqref{GGloc2b}, which requires a bit of care. Fix $b>0$ and consider
	\[
	F_b = K \cap \partial \Omega \cap \{|W| > b\}, \qquad F^b = K \cap \partial \Omega \cap \{|W| \leq b\} \, .
	\] 
	For $t \in (0, \tau_0)$, define also $F_b(t) = \Psi(t, F_b)$ and $F^b(t) = \Psi(t, F^b)$, which are disjoint sets being $\Psi(t,\cdot)$ a diffeomorphism. Then, for each $t\in(0,\tau_0)$
	\begin{equation} \label{eq_int}
		\disp \int_{\Omega_0\cap\{r=t\}} \langle W, \nabla r\rangle = \int_{\Omega_0 \cap F_b(t)} \langle W, \nabla r\rangle + \int_{\Omega_0\cap F^b(t)} \langle W, \nabla r\rangle .
	\end{equation}
	Since the function $|W|$ is continuous and compactly supported in $\overline{\Omega}$, it is also uniformly continuous. Denoting by $\sigma : \R^+ \to \R^+$ its modulus of continuity, we have then
	\begin{equation} \label{in_sigma}
		|W| \geq b - \sigma(t) \quad \text{on } \, F_b(t) \qquad \text{and} \qquad |W| \leq b + \sigma(t) \quad \text{on } \, F^b(t)
	\end{equation}
	for each $t\in(0,\tau_0)$. Fix $t_b\in(0,\tau_0)$ small enough so that $\sigma(t_b) < b$. Then, for each $t\in(0,t_b)$ we have $\sigma(t) < b$, hence $|W|>0$ on $F_b(t)$ by \eqref{in_sigma} and from this together with the assumption that $W=0$ on $\partial\Omega_0 \cap \Omega$ we infer that $F_b(t) \cap \partial\Omega_0 = \emptyset$. In particular, the flow $\Psi$ starting from points of $\overline{\Omega}_0 \cap F_b$ does not hit $\partial\Omega_0 \cap \Omega$ before time $t_b$ and we have
	\[
	\Omega_0 \cap F_b(t) \equiv \Psi(t,\overline{\Omega}_0 \cap F_b) \qquad \text{for all } \, t \in (0,t_b) \, .
	\] 
	The family of hypersurfaces $t \mapsto \Omega_0 \cap F_b(t)$ is thus continuous in the $C^1$ topology and converges to $\overline\Omega_0 \cap F_b$ as $t\to0$. Therefore, we have
	\[
	\lim_{t\to0} \int_{\Omega_0 \cap F_b(t)} \langle W, \nabla r \rangle = \int_{\overline{\Omega}_0\cap F_b} \langle W, \eta\rangle
	\]
	and from the integral mean value theorem we obtain
	\begin{equation}\label{eq_int2}
		\lim_{\tau\to0} \frac{1}{\tau} \int_\tau^{2\tau} \int_{\Omega_0\cap F_b(t)} \langle W, \nabla r \rangle \, \di t = \int_{\overline\Omega_0\cap F_b} \langle W, \eta\rangle \, .
	\end{equation}
	On the other hand, for each $t\in(0,\tau_0)$ we have
	\[
	\left| \int_{\Omega_0\cap F^b(t)} \langle W, \nabla r\rangle \right| \le (b+\sigma(t)) \haus^{n-1}(F^b(t))
	\]
	by \eqref{in_sigma}, and since $\haus^{n-1}(F^b(t)) \to \haus^{n-1}(F^b) \leq \haus^{n-1}(K\cap\partial\Omega) < \infty$ as $t \to 0$ we get
	\begin{equation}\label{eq_int1}
		\limsup_{\tau\to0} \left| \frac{1}{\tau} \int_\tau^{2\tau} \int_{\Omega_0\cap F^b(t)} \langle W, \nabla r \rangle \, \di t \right| \le C b
	\end{equation}
	with $C>0$ independent of $b$. Then, from \eqref{eq_int}, \eqref{eq_int2} and \eqref{eq_int1} we obtain
	\[
		\limsup_{\tau\to0} \left| \frac{1}{\tau} \int_\tau^{2\tau} \int_{\Omega_0\cap\{r=t\}} \langle W,\nabla\psi\rangle \, \di t - \int_{\overline{\Omega}_0\cap F_b} \langle W,\eta\rangle \right| \leq Cb
	\]
	for each $b>0$. Then \eqref{GGloc2b} follows by letting $b\to0$ and noting that
	\[
		\lim_{b\to0} \int_{\overline{\Omega}_0\cap F_b} \langle W,\eta\rangle = \int_{\overline{\Omega}_0\cap \partial\Omega} \langle W,\eta\rangle
	\]
	by Lebesgue's convergence theorem.
\end{proof}

\begin{theorem} \label{thm_prop_grad}
	Let $(M^n,\metric)$ be a complete Riemannian manifold and $\Omega\subseteq M$ a domain. Let $f\in C^1(\R)$ and let $u\in C^3(\Omega) \cap C^1(\overline{\Omega})$ be a solution to
	\[
		\MM[u] + f(u) = 0 \qquad \text{in } \, \Omega \, .
	\]
	Suppose that
	\[
		\sup_\Omega |f(u)| < \infty \, , \qquad f'(u) \leq 0 \, , \qquad \inf_\Omega u > -\infty
	\]
	and that
	\[
		\Ric \geq 0 \quad \text{in } \, \Omega , \qquad \Sec \geq - \kappa^2 \quad \text{in } \, M
	\]
	for some $\kappa\in\R^+$. Then
	\[
		\sup_\Omega |\nabla u| \leq \max\left\{ \sqrt{n/2}, \, \sup_{\partial\Omega} |\nabla u|\right\} .
	\]
\end{theorem}

\begin{proof}
%	Let $W = \sqrt{1+|\nabla u|^2}$. For any $C>0$ we consider the function $z_C = W e^{-Cu} \leq W$. We shall prove that
%	\[
%		\sup_\Omega z_C \leq \max\left\{ A, \sup_{\partial\Omega} W \right\} 
%	\]
%	for some constant $A>1$ not depending on $C$. Since $z_C \to W$ pointwise in $\Omega$ as $C\to0$, this will yield \eqref{supDu} for $C_1 = \sqrt{A^2-1}$.
	
	Without loss of generality, we can assume $u\geq 0$ in $\Omega$ and choose $C_0>0$ so that
	\[
		|f(u)| \leq C_0 \qquad \text{in } \, \Omega \, .
	\]
	Let $F : \Omega \to \R\times M$ be the graph map given by $F(x) = (u(x),x)$ for each $x\in\Omega$ and denote by $g$ the Riemannian metric on $\Omega$ obtained by pulling back via $F$ the ambient product metric on $\R\times M$. We have
	\[
		g = \metric + \di u \otimes \di u \, .
	\]
	In components, if we write
	\[
		\metric = \sigma_{ij} \, \di x^i \otimes \di x^j, \qquad (\sigma^{ij}) = (\sigma_{ij})^{-1}, \qquad \di u = u_i \, \di x^i, \qquad u^i = \sigma^{ij} u_j
	\]
	then the components $g_{ij}$ of the metric $g$ and of the inverse matrix $(g^{ij}) = (g_{ij})^{-1}$ are
	\begin{equation} \label{gij}
		g_{ij} = \sigma_{ij} + u_i u_j \qquad \text{and} \qquad g^{ij} = \sigma^{ij} - \frac{u^i u^j}{W^2}
	\end{equation}
	where $W = \sqrt{1+|\nabla u|^2}$. We denote by $\nabla^g$, $\|\cdot\|$, $\diver_g$ and $\Delta_g$ the gradient, tensor norm, divergence and Laplace-Beltrami operator associated to $g$, respectively. For any $\varphi\in C^2(\Omega)$ we have (see, for instance, formula (25) in \cite{cmr})
	\begin{equation} \label{Lapg}
		\Delta_g \varphi = g^{ij} \varphi_{ij} - \varphi_i u^i \frac{H}{W}
	\end{equation}
	where $\varphi_i$ and $\varphi_{ij}$ are the components of $\di\varphi$ and $\nabla\di\varphi$, respectively, and $H$ is the (non-normalized) mean curvature of the graph of $u$ in $\R\times M$ in the direction of the upward normal vector $\mathbf{n}$. From the Jacobi equation we have
	\[
		\Delta_g \frac{1}{W} + (\|\II\|^2 + \overline{\Ric}(\mathbf{n},\mathbf{n})) \frac{1}{W} + g(\nabla^g H,\nabla^g u) = 0, \qquad \Delta_g u = \frac{H}{W}
	\]
	where $\II$ is the second fundamental form of the graph of $u$ in $\R\times M$ and $\overline{\Ric}$ is the Ricci tensor of $\R\times M$. %$\mathbf{n}$ is the upward pointing normal vector to $\Sigma$ in $\R\times M$ and $H$ is the (non-normalized) mean curvature of $\Sigma$ in the direction of $\mathbf{n}$.
	Since
	\[
		H = \MM[u] = -f(u), \qquad g(\nabla^g H,\nabla^g u) = -f'(u)\|\nabla^g u\|^2 \geq 0 \qquad \text{and} \qquad \overline{\Ric} \geq 0,
	\]
	we have
	\begin{equation} \label{DeltaWu}
		\Delta_g \frac{1}{W} + \frac{\|\II\|^2}{W} \leq 0, \qquad \Delta_g u = - \frac{f(u)}{W} \, .
	\end{equation}
	Note that the first inequality implies
	\begin{equation} \label{DeltaW2}
		\Delta_g W \geq \|\II\|^2 W + 2 \frac{\|\nabla^g W\|^2}{W} \, .
	\end{equation}
	
	Let $x_0\in\Omega$ be fixed. Let $C,\eps,\delta>0$ be given, with $C$ and $\delta$ satisfying
	\begin{equation} \label{deltaC}
		\delta < e^{-Cu(x_0)} \, ,
	\end{equation}
	and define
	\[
		\eta = e^{-Cu-\eps r^2} - \delta, \qquad z = W \eta, \qquad \Omega_0 = \{ x \in \Omega : \eta(x) > 0 \}
	\]
	where $r(x) = \dist(x_0,x)$ denotes the Riemannian distance function in $M$ from $x_0$. Since $\eta(x_0) = e^{-Cu(x_0)} - \delta > 0$, we clearly have $x_0 \in \Omega_0 \neq \emptyset$. Since $u\geq 0$, we also have
	\begin{equation} \label{Omega0}
		\Omega_0 \subseteq \{ e^{-\eps r^2} > \delta \} = \{ \eps r^2 < \log(1/\delta) \}
	\end{equation}
	(note that $\log(1/\delta) > 0$ since $\delta<1$ by the constraint \eqref{deltaC} and the assumption $u\geq 0$), hence $\Omega_0$ is relatively compact. By continuity of $z$ in $\overline{\Omega}$, there exists $x_1\in\overline{\Omega}_0$ such that
	\[
		z(x_1) = \max_{\overline{\Omega}_0} z > 0 \, .
	\]
	
	If $x_1\in\partial\Omega$ then we have
	\begin{equation} \label{West1}
		W(x_0) e^{-Cu(x_0)} = z(x_0) \leq z(x_1) \leq W(x_1) \leq \sup_{\partial\Omega} W
	\end{equation}
	where we used that $z \leq W$ since $u\geq 0$. %If $x_1\in\Omega$ we shall prove that $z(x_0) \leq A$ for some constant $A\geq 1$ independent from $C,\eps,\delta$ or $x_0$, yielding again
%%	If $x_1\in\Omega$ we shall prove that $z$ must satisfy $z(x_1) \leq A$ for some constant $A>1$ independent from $C,\eps,\delta$ or $x_0$, yielding again
%	\[
%		W(x_0) e^{-Cu(x_0)} \leq A \, .
%	\]
%	In both cases, when letting $C\to0$, we will deduce
%	\[
%		W(x_0) \leq \max\left\{A,\sup_{\partial\Omega} W\right\}
%	\]
%	that is
%	\[
%		|\nabla u|(x_0) \leq \max\left\{ \sqrt{A^2-1},\sup_{\partial\Omega} |\nabla u| \right\}.
%	\]
	Let us suppose that $x_1\in\Omega$ instead. A direct computation yields
	\[
		W^2 \diver_g(W^{-2} \nabla^g z) = \eta \Delta_g W + W\Delta_g\eta - 2 \eta \frac{\|\nabla^g W\|^2}{W} \qquad \text{in } \, \Omega
	\]
	from which, using also \eqref{DeltaW2}, we infer
	\begin{equation} \label{divWz}
		W^2 \diver_g(W^{-2} \nabla^g z) \geq \left(\|\II\|^2 + \frac{\Delta_g \eta}{\eta} \right) z \qquad \text{in } \, \Omega_0 \, .
	\end{equation}
	By direct computation,
	\begin{equation} \label{Lapeta}
		\Delta_g \eta = (\eta+\delta)(-C\Delta_g u - \eps\Delta_g r^2 + \|C\nabla^g u + \eps\nabla^g r^2\|^2) \qquad \text{in } \, \Omega \, .
	\end{equation}
	By \eqref{Lapg}, \eqref{gij}, the Hessian comparison theorem and $|H| = |f(u)| \leq C_0$ we estimate
	\begin{equation} \label{Lapr2}
		\Delta_g r^2 = g^{ij} (r^2)_{ij} - \frac{H}{W} \langle \nabla u,\nabla r^2\rangle \leq 2n \kappa r \coth(\kappa r) + 2C_0 r \leq C_2(1+r)
	\end{equation}
	where $C_2 = C_2(n,\kappa) > 0$. By \eqref{gij} we have
	\[
		\|\nabla^g u\|^2 = g^{ij} u_i u_j = \frac{|\nabla u|^2}{W^2} = 1 - \frac{1}{W^2}, \qquad \|\nabla^g r\|^2 = g^{ij} r_i r_j \leq \sigma^{ij} r_i r_j = |\nabla r|^2 \leq 1
	\]
	hence, by Young's inequality,
	\begin{equation} \label{CeYoung}
		\|C\nabla^g u + \eps\nabla^g r^2\|^2 \geq \frac{1}{2} C^2\|\nabla^g u\|^2 - \eps^2 \|\nabla^g r^2\|^2 \geq \frac{C^2}{2} (1-W^{-2}) - 4\eps^2 r^2 \, .
	\end{equation}
	From \eqref{Lapeta}, \eqref{Lapr2}, \eqref{CeYoung} and the second in \eqref{DeltaWu} we get
	\[
		\frac{\Delta_g\eta}{\eta} \geq \left(1+\frac{\delta}{\eta}\right)\left(\frac{Cf(u)}{W} + \frac{C^2}{2} \left(1-\frac{1}{W^2}\right) - C_2\eps(1+r) - 4 \eps^2 r^2 \right) \qquad \text{in } \, \Omega_0
	\]
	and therefore, by \eqref{divWz},
	\begin{align*}
		W^2 \diver_g(W^{-2} \nabla^g z) & \geq \left(\|\II\|^2 + \left(1+\frac{\delta}{\eta}\right) \frac{Cf(u)}{W}\right) z \\
		& \phantom{=\;} + \left(1+\frac{\delta}{\eta}\right) \left( \frac{C^2}{2} \left(1-\frac{1}{W^2}\right) - C_2 \eps (1+r) - 4 \eps^2 r^2 \right) z \qquad \text{in } \, \Omega_0 \, .
	\end{align*}
	Since we assumed $x_1\in\Omega_0$ to be an interior maximum point for $z$, we have
	\[
		W^2 \diver_g(W^{-2}\nabla^g z) \leq 0 \qquad \text{at } \, x_1
	\]
	by ellipticity, hence
	\begin{equation} \label{max_bon}
		0 \geq \|\II\|^2 + \left(1+\frac{\delta}{\eta}\right) \frac{Cf(u)}{W} + \left(1+\frac{\delta}{\eta}\right) \left( \frac{C^2}{2} \left(1-\frac{1}{W^2}\right) - C_2 \eps (1+r) - 4 \eps^2 r^2 \right)
	\end{equation}
	at $x_1$. We proceed to suitably bound the various terms on the right hand side of \eqref{max_bon} from below. We first consider the quantity
	\[
		(\ast) \doteq \|\II\|^2 + \left(1+\frac{\delta}{\eta}\right) \frac{Cf(u)}{W} = \|\II\|^2 + \frac{Cf(u)}{W} + \frac{\delta C f(u)}{z} \, .
	\]
	Clearly $(\ast) \geq 0$ if $f(u) \geq 0$. On the other hand, if $f(u) < 0$ we can estimate
	\[
		\|\II\|^2 \geq \frac{(\mathrm{tr}_g \II)^2}{n} = \frac{H^2}{n} = \frac{f(u)^2}{n} \qquad \text{and} \qquad \frac{Cf(u)}{W} \geq \frac{Cf(u)}{z} \, ,
	\]
	where in the second inequality we used the fact that $Cf(u) < 0$ and $z \leq W$, whence
	\[
		(\ast) \geq \frac{f(u)^2}{n} + (1+\delta) \frac{Cf(u)}{z} \geq - \frac{(1+\delta)^2 n C^2}{4z^2} \qquad \text{if } \, f(u) < 0
	\]
	where we used Young's inequality to estimate
	\[
		\frac{f(u)^2}{n} + \frac{(1+\delta)^2 n C^2}{4z^2} \geq - (1+\delta) \frac{Cf(u)}{z} \, .
	\]
	Therefore, since $\delta/\eta>0$ in $\Omega_0$, we have
	\[
		(\ast) \geq \left\{
			\begin{array}{l@{\quad}l}
				0 & \text{if } \, f(u) \geq 0 \\
				- \left(1+\dfrac{\delta}{\eta}\right)\dfrac{(1+\delta)^2 n C^2}{4 z^2} & \text{if } \, f(u) < 0 \, .
			\end{array}
		\right.
	\]
	Since $z\leq W$ we also have
	\[
		\frac{C^2}{2} \left(1-\frac{1}{W^2}\right) \geq \frac{C^2}{2} \left(1-\frac{1}{z^2}\right) .
	\]
	As observed in \eqref{Omega0}, we have $\eps r^2 < \log(1/\delta)$ in $\Omega_0$ and therefore
	\[
		C_2\eps(1+r) + 4 \eps^2 r^2 \leq C_2\left(\eps + \sqrt{\eps \log(1/\delta)}\right) + 4\eps\log(1/\delta) \qquad \text{in } \, \Omega_0 \, .
	\]
	Then, from \eqref{max_bon} we deduce
	\begin{equation} \label{max_bon2}
		\frac{C^2}{2} \left(1-\frac{1}{z^2}\right) - C_2\left(\eps + \sqrt{\eps \log(1/\delta)}\right) - 4\eps\log(1/\delta) \leq \frac{(1+\delta)^2 n C^2}{4 z^2} \qquad \text{at } \, x_1,
	\end{equation}
	that is,
	\[
		1 - \frac{C_\delta}{z(x_1)^2} \leq \frac{2\sqrt{\eps}}{C^2} \left(C_2\sqrt{\log(1/\delta)} + \sqrt{\eps}(C_2+4\log(1/\delta))\right)
	\]
	where $C_\delta = 1 + (1+\delta)^2n/2$. Since $z(x_0) \leq z(x_1)$, we further have
	\begin{equation} \label{max_bon3}
		1 - \frac{C_\delta}{z(x_0)^2} \leq \frac{2\sqrt{\eps}}{C^2} \left(C_2\sqrt{\log(1/\delta)} + \sqrt{\eps}(C_2+4\log(1/\delta))\right) .
	\end{equation}
	(Note that if $f(u) \geq 0$ then \eqref{max_bon2} holds true, more strongly, with $(1+\delta)^2 n C^2 / 4z^2$ replaced by $0$, so \eqref{max_bon3} still holds with $C_\delta$ replaced by $1$.) We emphasize that \eqref{max_bon3} holds for \emph{any} choice of parameters $C,\eps,\delta>0$ provided that \eqref{deltaC} is satisfied and $z$ attains its global maximum at an interior point of $\Omega_0$, so in view of \eqref{West1} we have
	\[
		z(x_0) \leq \max\left\{ A(C,\eps,\delta), \sup_{\partial\Omega} W \right\} \, .
	\]
	where
	\[
		A(C,\eps,\delta) = \sup\left\{ t > 0 : 1 - \frac{C_\delta}{t^2} \leq \frac{2\sqrt{\eps}}{C^2} \left(C_2\sqrt{\log(1/\delta)} + \sqrt{\eps}(C_2+4\log(1/\delta))\right) \right\} .
	\]
	For a fixed $C>0$, letting first $\eps\to0$ and then $\delta\to0$ we get $A(C,\eps,\delta) \to A = \sqrt{1+n/2}$, hence 
	\[
		W(x_0) e^{-Cu(x_0)} = z(x_0) \leq \max\left\{A,\sup_{\partial\Omega} W\right\} .
	\]
	Passing to the limit as $C\to0$ we obtain
	\[
		W(x_0) \leq \max\left\{A,\sup_{\partial\Omega} W\right\},
	\]
	that is,
	\[
		|\nabla u|(x_0) \leq \max\left\{ \sqrt{A^2-1},\,\sup_{\partial\Omega} |\nabla u| \right\} = \max\left\{ \sqrt{n/2},\,\sup_{\partial\Omega} |\nabla u| \right\}
	\]
	as desired.
\end{proof}




%
%
%
%
%
%
%
%\vspace{0.5cm}
%%
%\noindent \textbf{Acknowledgements: } The second author is indebted to Jorge Herbert de Lira for a stimulating discussion about Killing fields on Riemannian manifolds.
%The first and the third
%authors were supported by the
%ERC grant EPSILON ({\it Elliptic Pde's and Symmetry of Interfaces and Layers
%for Odd Nonlinearities}).






\bibliographystyle{plain}

%\bibliography{bibliosemilinear}

\begin{thebibliography}{10}

\bibitem{abbm}
V. Agostiniani, C. Bernardini, S. Borghini and L. Mazzieri, 
\newblock On the Serrin problem for ring-shaped domains: the $n$-dimesional case
\newblock Available at arXiv:2505.04458.

\bibitem{abm}
V. Agostiniani, S. Borghini and L. Mazzieri, 
\newblock On the Serrin problem for ring-shaped domains.
\newblock {\em J. Eur. Math. Soc. (JEMS)} 27 (2025), no. 7, 2705--2749.

% 
%\bibitem{albertiambrosiocabre}
%G.~Alberti, L.~Ambrosio, and X.~Cabr{\'e}.
%\newblock On a long-standing conjecture of {E}. {D}e {G}iorgi: symmetry in 3{D}
%  for general nonlinearities and a local minimality property.
%\newblock {\em Acta Appl. Math.}, 65(1-3):9--33, 2001.
%\newblock Special issue dedicated to Antonio Avantaggiati on the occasion of
%  his 70th birthday.

%\bibitem{ambrosiocabre}
%L.~Ambrosio and X.~Cabr{\'e}.
%\newblock Entire solutions of semilinear elliptic equations in {$ \R^3$} and a
%  conjecture of {D}e {G}iorgi.
%\newblock {\em J. Amer. Math. Soc.}, 13(4):725--739 (electronic), 2000.

\bibitem{ambrosiofuscopallara}
L. Ambrosio, N. Fusco and D. Pallara.
\newblock {\em Functions of Bounded Variation and Free Discontinuity Problems}.
\newblock Oxford Science Publications, 2000.

%\bibitem{bercaffnire2}
%H.~Berestycki, L.~Caffarelli, and L.~Nirenberg.
%\newblock Further qualitative properties for elliptic equations in unbounded
%  domains.
%\newblock {\em Ann. Scuola Norm. Sup. Pisa Cl. Sci. (4)}, 25(1-2):69--94
%  (1998), 1997.
%\newblock Dedicated to Ennio De Giorgi.

\bibitem{bcn}
H.~Berestycki, L.A. Caffarelli, and L.~Nirenberg.
\newblock Monotonicity for elliptic equations in unbounded {L}ipschitz domains.
\newblock {\em Comm. Pure Appl. Math.}, 50(11):1089--1111, 1997.

%\bibitem{berenire}
%H.~Berestycki and L.~Nirenberg.
%\newblock On the method of moving planes and the sliding method.
%\newblock {\em Bol. Soc. Brasil. Mat. (N.S.)}, 22(1):1--37, 1991.

%\bibitem{bernirevara}
%H.~Berestycki, L.~Nirenberg, and S.R.S.~Varadhan.
%\newblock The principal eigenvalue and maximum principle for second-order
%  elliptic operators in general domains.
%\newblock {\em Comm. Pure Appl. Math.}, 47(1):47--92, 1994.

\bibitem{BCM}
S.~Borghini, P. T.~Chru\'sciel, and L.~Mazzieri.
\newblock On the Uniqueness of Schwarzschild--de Sitter Spacetime.
\newblock {\em Arch. Rational Mech. Anal.}, 247 no.~2, Paper No. 22, 35 pp., 2023.

%\bibitem{brooks}
%R.~Brooks.
%\newblock A relation between growth and the spectrum of the {L}aplacian.
%\newblock {\em Math. Z.}, 178(4):501--508, 1981.

%\bibitem{calabi_vol}
%E.~Calabi.
%\newblock On manifolds with non-negative {R}icci curvature {II}.
%\newblock {\em Notices Amer. Math. Soc.}, 22:A205, 1975.

%\bibitem{cheegergromoll}
%J.~Cheeger and D.~Gromoll.
%\newblock The splitting theorem for manifolds of nonnegative {R}icci curvature.
%\newblock {\em J. Diff. Geom.}, 6:119--128, 1971/72.

%\bibitem{chengyau75}
%S.Y. Cheng and S.T. Yau.
%\newblock Differential equations on {R}iemannian manifolds and their geometric
%  applications.
%\newblock {\em Comm. Pure Appl. Math.}, 28(3):333--354, 1975.

%\bibitem{chengyau77}
%S.Y. Cheng and S.T. Yau.
%\newblock Hypersurfaces with constant scalar curvature.
%\newblock {\em Math. Ann.}, 225:195--204, 1977.

\bibitem{choe} 
J. Choe, 
\newblock On the analytic reflection of a minimal surface. 
\newblock {\em Pacific J. Math.} 157 (1993), no. 1, 29--36.

\bibitem{cmmr} 
G. Colombo, M. Magliaro, L. Mari and M. Rigoli, 
\newblock Bernstein and half-space properties for minimal graphs under Ricci lower bounds.
\newblock {\em Int. Math. Res. Not. IMRN} 2022, no. 23, 18256--18290.

\bibitem{cmr}
G. Colombo, L. Mari and M. Rigoli, 
\newblock A splitting theorem for capillary graphs under Ricci lower bounds.
\newblock {\em J. Funct. Anal.} 281 (2021), no. 8, Paper No. 109136, 50 pp.


\bibitem{cui} 
Q. Cui, 
\newblock A note on an overdetermined problem for minimal surface equation.
\newblock {\em Proc. Amer. Math. Soc.} 150 (2022), no. 12, 5403--5409.


\bibitem{duc_eels}
D.M. Duc and J. Eells,
\newblock 
Regularity of exponentially harmonic functions.
\newblock {\em Internat. J. Math.} 2 (1991), no. 4, 395--408.

%\bibitem{dupfar}
%L.~Dupaigne and A.~Farina.
%\newblock Stable solutions of {$- \Delta u = f (u)$} in {$\R^N$}.
%\newblock {\em J. Eur. Math. Soc.}, 12(4):855--882, 2010.

\bibitem{eels_lemaire}
J. Eells and L. Lemaire, 
\newblock Some properties of exponentially harmonic maps.
\newblock {\em Partial differential equations, Part 1,2} (Warsaw, 1990), 129--136.
Banach Center Publ., 27, Part 1, 2
Polish Academy of Sciences, Institute of Mathematics, Warsaw, 1992

%\bibitem{FTh}
%A.~Farina.
%\newblock Finite-energy solutions, quantization effects and {L}iouville-type
%  results for a variant of the {G}inzburg-{L}andau systems in {${\R}^K$}.
%\newblock {\em Differential Integral Equations}, 11(6):875--893, 1998.

%\bibitem{FAR-H}
%A.~Farina.
%\newblock Propri\'et\'es qualitatives de solutions d'\'equations et syst\`emes
%  d'\'equations non-lin\'eaires.
%\newblock 2002.
%\newblock Habilitation \`a diriger des recherches, Paris VI.

\bibitem{FarMarVal}
A.~Farina, L.~Mari, and E.~Valdinoci.
\newblock Splitting theorems, symmetry results and overdetermined problems for 
Riemannian manifolds.
\newblock {\em Communications in Partial Differential Equations}, 
38(10):1818--1862, 2013.


\bibitem{FSV}
A.~Farina, B.~Sciunzi, and E.~Valdinoci.
\newblock Bernstein and {D}e {G}iorgi type problems: new results via a
  geometric approach.
\newblock {\em Ann. Sc. Norm. Super. Pisa Cl. Sci. (5)}, 7(4):741--791, 2008.

%\bibitem{FarSirVal}
%A.~Farina, Y.~Sire, and E.~Valdinoci.
%\newblock Stable solutions of elliptic equations on {R}iemannian manifolds.
%\newblock {\em J. Geom. Anal.}, 2012.

%\bibitem{le}
%\tcr{NOT CITED YET} A.~Farina and E.~Valdinoci.
%\newblock The state of the art for a conjecture of {D}e {G}iorgi and related
%  problems.
%\newblock In {\em Recent progress on reaction-diffusion systems and viscosity
%  solutions}, pages 74--96. World Sci. Publ., Hackensack, NJ, 2009.

\bibitem{Arma}
A.~Farina and E.~Valdinoci.
\newblock Flattening results for elliptic {PDE}s in unbounded domains with
  applications to overdetermined problems.
\newblock {\em Arch. Ration. Mech. Anal.}, 195(3):1025--1058, 2010.

\bibitem{Arma2}
A.~Farina and E.~Valdinoci.
\newblock Correction of a technical point and some remarks about the paper appeared in arch. ration. mech. anal. 195(2010), no. 3, 1025-1058.
\newblock {\em Arch. Ration. Mech. Anal.},
DOI: 10.1007/s00205-012-0568-6,
{\tt http://www.springerlink.com/content/p1m68k1l856667w7/}

\bibitem{federer}
H. Federer.
\newblock {\em Geometric Measure Theory}.
\newblock Springer-Verlag, 1969.

\bibitem{finn} R. Finn. Equilibrium capillary surfaces. Vol. 284. \emph{Grundlehren der mathematischen Wissenschaften.} Springer-Verlag, New York, 1986, pp. xvi+245.


\bibitem{fischercolbrieschoen}
D.~Fischer-Colbrie and R.~Schoen.
\newblock The structure of complete stable minimal surfaces in $3$-manifolds of
  non negative scalar curvature.
\newblock {\em Comm. Pure Appl. Math.}, XXXIII:199--211, 1980.

\bibitem{gilbargtrudinger}
D.~Gilbarg and N.~Trudinger.
\newblock {\em Elliptic {P}artial {D}ifferential {E}quations of {S}econd
  {O}rder}.
\newblock third ed., Springer-Verlag, 1998.

%\bibitem{DG}
%\tcr{NOT CITED YET} E.~De Giorgi.
%\newblock Convergence problems for functionals and operators.
%\newblock In {\em Proceedings of the {I}nternational {M}eeting on {R}ecent
%  {M}ethods in {N}onlinear {A}nalysis ({R}ome, 1978)}, pages 131--188, Bologna,
%  1979. Pitagora.



%\bibitem{grigoryan}
%A.~Grigor'yan.
%\newblock Analytic and geometric background of recurrence and non-explosion of
%  the {B}rownian motion on {R}iemannian manifolds.
%\newblock {\em Bull. Amer. Math. Soc.}, 36:135--249, 1999.

\bibitem{HS89}
R. Hardt and L. Simon.
\newblock Nodal sets for solutions of elliptic equations.
\newblock {\em J. Differential Geom.}, 30:505--522, 1989.

%\bibitem{hamilton4man}
%R.S. Hamilton.
%\newblock Four-manifolds with positive curvature operator.
%\newblock {\em J. Differential Geom.}, 24(2):153--179, 1986.

%\bibitem{Jimbo}
%S.~Jimbo.
%\newblock On a semilinear diffusion equation on a {R}iemannian manifold and its
%  stable equilibrium solutions.
%\newblock {\em Proc. Japan Acad. Ser. A Math. Sci.}, 60(10):349--352, 1984.

\bibitem{karp}
L. Karp. 
\newblock 
Subharmonic functions on real and complex manifolds. 
\newblock {\em Math. Z.} 179, 535--554, 1982.

\bibitem{KP}
S.~G. Krantz and H.~R. Parks.
\newblock A primer of real analytic functions, second edition.
\newblock {\em Birkh\"auser Advanced Texts: Basler Lehrb\"ucher, Birkh\"auser Boston, Boston, MA, 2002.}

%\bibitem{Pli}
%P.~Li.
%\newblock Lecture notes on geometric analysis.
%\newblock {\em Lecture Notes Series No. 6, Research Institute of Mathematics,
%  Global Analysis Research Center, Seoul National University, Korea (1993)}.

\bibitem{lian_sicbaldi} Y. Lian, P. Sicbaldi, \emph{Rigidity results for the capillary overdetermined problem.} Available at arXiv:2503.14215. 


\bibitem{Lima}
E.~L. Lima.
\newblock The Jordan-Brouwer separation theorem for smooth hypersurfaces.
\newblock {\em Amer. Math. Monthly}, 95(1): 39--42, 1988.

%\bibitem{liu}
%G.~Liu.
%\newblock 3-manifolds with nonnegative {R}icci curvature.
%\newblock {\em available at arXiv:1108.1888}.

\bibitem{marcellini}
P. Marcellini, 
\newblock On the definition and the lower semicontinuity of certain quasiconvex integrals.
\newblock {\em Ann. Inst. H. Poincar\'e Anal. Non Lin\'eaire} 3 (1986), no. 5, 391--409.

%\bibitem{mastroliarigoli}
%P.~Mastrolia and M.~Rigoli.
%\newblock Diffusion-type operators, {L}iouville theorems and gradient estimates
%  on complete manifolds.
%\newblock {\em Nonlinear Anal.}, 72(9-10):3767--3785, 2010.

%\bibitem{mckean}
%H.P. McKean.
%\newblock An upper bound to the spectrum of {$\Delta$} on a manifold of
%  negative curvature.
%\newblock {\em J. Diff. Geom.}, 4:359--366, 1970.

\bibitem{mosspie}
W.F. Moss and J.~Piepenbrink.
\newblock Positive solutions of elliptic equations.
\newblock {\em Pac. J. Math.}, 75:219--226, 1978.

%\bibitem{omori}
%H.~Omori.
%\newblock Isometric immersions of {R}iemannian manifolds.
%\newblock {\em J. Math. Soc. Japan}, 19:205--214, 1967.

%\bibitem{petersen}
%P.~Petersen.
%\newblock {\em Riemannian {G}eometry}.
%\newblock Springer-Verlag, 1997.

%\bibitem{prsmemoirs}
%S.~Pigola, M.~Rigoli, and A.G. Setti.
%\newblock Maximum principles on {R}iemannian manifolds and applications.
%\newblock {\em Mem. Amer. Math. Soc.}, 174:no.822, 2005.

\bibitem{prs}
S.~Pigola, M.~Rigoli, and A.G. Setti.
\newblock {\em Vanishing and finiteness results in {G}eometric {A}nalisis. {A}
  generalization of the {B}ochner technique}, volume 266 of {\em Progress in
  Math.}
\newblock Birk\"auser, 2008.

%\bibitem{pigolaveronelliIJM}
%S.~Pigola and G.~Veronelli.
%\newblock Remarks on {$L^p$}-vanishing results in {G}eometric {A}nalysis.
%\newblock {\em Internat. J. Math.}, 23(1):18pp., 2012.

\bibitem{DPV} S. Dipierro, A. Pinamonti and E. Valdinoci. Rigidity results for elliptic boundary value problems: stable solutions for quasilinear equations with Neumann or Robin boundary conditions. \emph{Int. Math. Res. Not. IMRN} 2020, no. 5, 1366--1384.

%\bibitem{protterweinberger}
%M.H. Protter and H.F. Weinberger.
%\newblock {\em Maximum principles in differential equations}.
%\newblock (Corrected reprint of the 1967 original), Springer-Verlag, 1984.

\bibitem{pucci_serrin} P. Pucci and J. Serrin, 
\newblock The maximum principle. 
\newblock Progress in Nonlinear Differential Equations and their Applications, 73, Birkh\"auser Verlag, Basel, 2007, x+235 pp.

\bibitem{reichel_1}
W. Reichel, 
\newblock Radial symmetry for an electrostatic, a capillarity and some fully nonlinear overdetermined problems on exterior domains.
\newblock {\em Z. Anal. Anwendungen} 
15 (1996), no. 3, 619--635.

%\bibitem{reichel_arma}
%\tcr{NOT YET CITED} W. Reichel, 
%\newblock Radial symmetry for elliptic boundary-value problems on exterior domains.
%\newblock {\em Arch. Rational Mech. Anal.} 
%137 (1997), no. 4, 381--394.

%\bibitem{rigolisetti}
%M.~Rigoli and A.G.~Setti.
%\newblock Liouville type theorems for $\varphi$-subharmonic functions.
%\newblock{\em Rev. Mat. Iberoam.}, 17(3):471--520, 2001.

%\bibitem{ros_sic}
%A.~Ros and P.~Sicbaldi.
%\newblock Geometry and topology of some overdetermined elliptic problems.
%\newblock {\em Preprint}, 2012.

%\bibitem{Saloff-Coste}
%L.~Saloff-Coste.
%\newblock {\em Aspects of Sobolev-Type Inequalities}.
%\newblock London Mathematical Society Lecture Notes Series, 2001.

%\bibitem{sattinger}
%D.H. Sattinger.
%\newblock Monotone methods in nonlinear elliptic and parabolic boundary value
%  problems.
%\newblock {\em Indiana Univ. Math. J.}, 21:979--1000, 1971/72.

\bibitem{serrin}
J. Serrin, 
\newblock A symmetry problem in potential theory.
\newblock {\em Arch. Rational Mech. Anal.} 
43 (1971), 304–318.

\bibitem{serrin2}
J. Serrin, 
\newblock The problem of Dirichlet for quasilinear elliptic differential equations with many independent variables.
\newblock {\em Philos. Trans. Roy. Soc. London Ser. A}
264 (1969), 413-496.


\bibitem{sibner_sibner}
L.M. Sibner and R.J. Sibner, 
\newblock A non-linear Hodge-de-Rham theorem.
\newblock {\em Acta Math.} 125 (1970), 57--73.

\bibitem{simon}
L. Simon, 
\newblock Singular Sets and Asymptotics in Geometric Analysis, 
\newblock Lipschitz Lectures. Institut f\"ur Angewandte Mathematik, http://math.stanford.edu/$\sim$lms/lipschitz/lipschitz.pdf,
Bonn, 2007.

\bibitem{sirakov} 
B. Sirakov, 
\newblock Symmetry for exterior elliptic problems and two conjectures in potential theory.
\newblock {\em Ann. Inst. H. Poincar\'e C Anal. Non Lin\'eaire} 
18 (2001), no. 2, 135--156.

%\bibitem{sternbergzumbrun2}
%P.~Sternberg and K.~Zumbrun.
%\newblock Connectivity of phase boundaries in strictly convex domains.
%\newblock {\em Arch. Rational Mech. Anal.}, 141(4):375--400, 1998.

%\bibitem{sternbergzumbrun1}
%P.~Sternberg and K.~Zumbrun.
%\newblock A {P}oincar\'e inequality with applications to volume-constrained
%  area-minimizing surfaces.
%\newblock {\em J. Reine Angew. Math.}, 503:63--85, 1998.


\bibitem{sternbergzumbrun1}
P.~Sternberg and K.~Zumbrun.
\newblock Connectivity of phase boundaries in strictly convex domains. 
\newblock {\em Arch. Rational Mech. Anal.}, 141 (1998), no. 4, 375--400.


%\bibitem{traizet}
%M. Traizet, 
%\newblock Classification of the solutions to an overdetermined elliptic problem in
%the plane.
%\newblock {\em Geom. Funct. Anal.} 24 (2014), no. 2, 690--720.

%\bibitem{walter}
%W.~Walter.
%\newblock A theorem on elliptic differential inequalities with an application
%  to gradient bounds.
%\newblock {\em Math. Z.}, 200(2):293--299, 1989.

%\bibitem{yau2}
%S.T. Yau.
%\newblock Some function-theoretic properties of complete {R}iemannian manifold
%  and their applications to geometry.
%\newblock {\em Indiana Univ. Math. J.}, 25(7):659--670, 1976.

\bibitem{zhikov}
V.V. Zhikov, 
\newblock Averaging of functionals of the calculus of variations and elasticity theory.
\newblock {\em Izv. Akad. Nauk SSSR Ser. Mat.} 50 (1986), no. 4, 675--710, 877.


\end{thebibliography}
%
%\end{document}
%



\end{document}




